\input{preamble}

% OK, commencer ici.
%
\begin{document}

\title{Espaces algébriques formels}


\maketitle

\phantomsection
\label{section-phantom}

\tableofcontents

\section{Introduction}
\label{section-introduction}

\noindent
Les schémas formels ont été introduits dans \cite{EGA}. Une version plus générale
des schémas formels a été introduite dans \cite{McQuillan}, et une autre
dans \cite{Yasuda}. Les espaces algébriques formels ont été introduits dans \cite{Kn}.
On trouvera des résultats connexes, ainsi que bien d'autres, dans
\cite{Abbes} et \cite{Fujiwara-Kato}.
Ce chapitre introduit la notion d'espace algébrique formel
que nous utiliserons. Notre définition est assez générale pour que la plupart
des classes de schémas ou d'espaces formels considérées dans la littérature en soient des sous-catégories pleines.

\medskip\noindent
Bien que nous comparions certaines de ces théories alternatives
à la nôtre, nous ne donnons pas toujours tous les détails lorsqu'ils ne sont pas nécessaires
au développement logique de la théorie.

\medskip\noindent
Outre l'introduction des espaces algébriques formels, nous établissons aussi quelques
propriétés très élémentaires et étudions plusieurs types de morphismes.










\section{Schémas formels à la EGA}
\label{section-formal-schemes-EGA}

\noindent
Dans cette section, nous rappelons la construction des schémas formels de \cite{EGA}.
Cette notion, bien que très utile en géométrie algébrique,
n'est pas toujours celle qu'il convient de considérer. Il serait peut-être plus juste de dire
que la mise en place de la théorie comporte un certain nombre de choix, alors que, suivant
le but poursuivi, d'autres choix pourraient être préférables. On trouve en effet dans la littérature
de nombreuses théories voisines
adaptées au problème que les auteurs souhaitent étudier. Néanmoins, l'un
des grands avantages de la théorie esquissée ici est qu'elle permet de travailler
avec des objets géométriques bien déterminés.

\medskip\noindent
Avant de commencer, signalons une difficulté concernant la condition de faisceau
pour les faisceaux d'anneaux topologiques ou, plus généralement, d'espaces
topologiques. En effet, les grandes catégories suivantes :
\begin{enumerate}
\item la catégorie des espaces topologiques,
\item la catégorie des groupes topologiques,
\item la catégorie des anneaux topologiques,
\item la catégorie des modules topologiques sur un anneau topologique donné,
\end{enumerate}
munies de leurs foncteurs d'oubli naturels vers $\textit{Sets}$, ne sont pas
des exemples de types de structures algébriques au sens de
Sheaves, Section \ref{sheaves-section-algebraic-structures}.
Nous ne pouvons donc pas leur appliquer sans précaution les méthodes développées dans ce
chapitre. En revanche, chacune des catégories
énumérées ci-dessus admet des limites et des égalisateurs, et le foncteur d'oubli
vers les ensembles, les groupes, les anneaux ou les modules commute à ceux-ci
(voir Topology, Lemmas \ref{topology-lemma-limits},
\ref{topology-lemma-topological-group-limits},
\ref{topology-lemma-topological-ring-limits}, et
\ref{topology-lemma-topological-module-limits}).
Nous pouvons donc définir la notion de
faisceau comme dans Sheaves, Definition
\ref{sheaves-definition-sheaf-values-in-category},
et le préfaisceau sous-jacent d'ensembles, de groupes, d'anneaux ou de modules
est un faisceau. La différence essentielle est que, pour un recouvrement ouvert
$U = \bigcup_{i \in I} U_i$, le diagramme
$$
\xymatrix{
\mathcal{F}(U) \ar[r]
&
\prod\nolimits_{i\in I}
\mathcal{F}(U_i)
\ar@<1ex>[r] \ar@<-1ex>[r]
&
\prod\nolimits_{(i_0, i_1) \in I \times I}
\mathcal{F}(U_{i_0} \cap U_{i_1})
}
$$
doit être un diagramme égalisateur dans la catégorie des espaces
topologiques, des groupes topologiques, des anneaux topologiques ou des modules topologiques,
c'est-à-dire que la première application identifie
$\mathcal{F}(U)$ à un sous-espace de $\prod_{i \in I} \mathcal{F}(U_i)$,
ce dernier étant muni de la topologie produit.

\medskip\noindent
La fibre $\mathcal{F}_x$ d'un faisceau $\mathcal{F}$
d'espaces topologiques, de groupes topologiques, d'anneaux topologiques ou de
modules topologiques en un point $x \in X$ est définie comme la limite inductive sur les
voisinages ouverts
$$
\mathcal{F}_x = \colim_{x\in U} \mathcal{F}(U)
$$
dans la catégorie correspondante. Cela revient à prendre
la limite inductive au niveau des ensembles, des groupes, des anneaux ou des modules
(voir Topology, Lemmas \ref{topology-lemma-colimits},
\ref{topology-lemma-topological-group-colimits},
\ref{topology-lemma-topological-ring-colimits}, et
\ref{topology-lemma-topological-module-colimits}),
mais cette limite est munie d'une topologie. Attention :
la topologie obtenue dépend de la catégorie dans laquelle on travaille ; voir
Examples, Section \ref{examples-section-colimit-topology}.
On peut associer un faisceau aux préfaisceaux d'espaces topologiques,
de groupes topologiques, d'anneaux topologiques ou de modules topologiques,
et la prise des fibres commute à cette opération ; voir
Remark \ref{remark-sheafification-of-presheaves-in-top}.

\medskip\noindent
Soit $f : X \to Y$ une application continue d'espaces topologiques.
Il existe un foncteur $f_*$ de la catégorie des faisceaux d'espaces
topologiques, de groupes topologiques, d'anneaux topologiques ou de modules topologiques
vers la catégorie correspondante des faisceaux sur $Y$, défini en posant, comme d'habitude,
$f_*\mathcal{F}(V) = \mathcal{F}(f^{-1}V)$.
(Nous remettons à plus tard l'étude de l'image inverse dans ce cadre.)
Nous définissons la notion de $f$-application $\xi : \mathcal{G} \to \mathcal{F}$
entre un faisceau d'espaces topologiques $\mathcal{G}$ sur $Y$ et
un faisceau d'espaces topologiques $\mathcal{F}$ sur $X$ exactement de la
même manière que dans Sheaves, Definition \ref{sheaves-definition-f-map},
avec la condition supplémentaire que
$\xi_V : \mathcal{G}(V) \to \mathcal{F}(f^{-1}V)$ soit continue
pour tout ouvert $V \subset Y$. On a
$$
\{f\text{-applications de }\mathcal{G}\text{ vers }\mathcal{F}\} =
\Mor_{\Sh(Y, \textit{Top})}(\mathcal{G}, f_*\mathcal{F})
$$
comme dans Sheaves, Lemma \ref{sheaves-lemma-f-map}. Il en va de même pour les
faisceaux de groupes topologiques, d'anneaux topologiques et de modules topologiques. Enfin,
soit $\xi : \mathcal{G} \to \mathcal{F}$ une $f$-application comme ci-dessus.
Étant donné $x \in X$, d'image $y = f(x)$, il existe une application
continue
$$
\xi_x : \mathcal{G}_y \longrightarrow \mathcal{F}_x
$$
entre les fibres, définie exactement de la même manière que dans la discussion
qui suit Sheaves, Definition \ref{sheaves-definition-composition-f-maps}.

\medskip\noindent
La discussion précédente permet de définir une catégorie $LTRS$ des
« espaces localement annelés topologiques ». Un objet est un couple
$(X, \mathcal{O}_X)$ formé d'un espace topologique
$X$ et d'un faisceau d'anneaux topologiques $\mathcal{O}_X$ dont les fibres
$\mathcal{O}_{X, x}$ sont des anneaux locaux (lorsqu'on oublie la topologie).
Un morphisme $(X, \mathcal{O}_X) \to (Y, \mathcal{O}_Y)$ dans
$LTRS$ est un couple $(f, f^\sharp)$, où $f : X \to Y$ est une application
continue d'espaces topologiques et $f^\sharp : \mathcal{O}_Y \to \mathcal{O}_X$
est une $f$-application telle que, pour tout $x \in X$, l'application induite
$$
f^\sharp_x : \mathcal{O}_{Y, f(x)} \longrightarrow \mathcal{O}_{X, x}
$$
soit un homomorphisme local d'anneaux locaux (en oubliant les topologies).
La composition se définit exactement comme celle des
morphismes d'espaces localement annelés.

\medskip\noindent
Supposons maintenant que l'espace topologique $X$ possède une base formée
d'ouverts quasi-compacts. Étant donné un faisceau $\mathcal{F}$ d'ensembles, de groupes,
d'anneaux ou de modules sur un anneau, on peut munir $\mathcal{F}$
d'une structure de faisceau d'espaces topologiques, de groupes topologiques,
d'anneaux topologiques ou de modules topologiques.
En effet, si $U \subset X$ est un ouvert quasi-compact,
nous munissons $\mathcal{F}(U)$ de la topologie discrète. Si $U \subset X$
est quelconque, nous choisissons un recouvrement ouvert $U = \bigcup_{i \in I} U_i$
par des ouverts quasi-compacts et munissons $\mathcal{F}(U)$
de la topologie induite par $\prod_{i \in I} \mathcal{F}(U_i)$
(comme l'impose la discussion précédente).
Le lecteur vérifiera (nous omettons les détails) que l'on obtient ainsi un faisceau d'espaces topologiques,
de groupes topologiques, d'anneaux topologiques ou de modules topologiques.
Nous dirons qu'un faisceau d'espaces topologiques, de groupes topologiques,
d'anneaux topologiques ou de modules topologiques est
{\it pseudo-discret} si la topologie de $\mathcal{F}(U)$ est
discrète pour tout ouvert quasi-compact $U \subset X$. Alors
la construction précédente est adjointe au foncteur d'oubli
et induit une équivalence entre la catégorie des faisceaux
d'ensembles et celle des faisceaux pseudo-discrets d'espaces topologiques
(et de même pour les groupes, les anneaux et les modules).

\medskip\noindent
Grothendieck et Dieudonné définissent d'abord les schémas formels affines.
Ceux-ci correspondent aux anneaux topologiques admissibles $A$ ; voir
More on Algebra, Definition \ref{more-algebra-definition-topological-ring}.
Plus précisément, étant donné $A$, on considère un système fondamental $I_\lambda$ d'idéaux
de définition de l'anneau $A$. (Dans tout anneau topologique admissible,
la famille de tous les idéaux de définition forme un système fondamental.)
Pour tout $\lambda$, on peut considérer le
schéma $\Spec(A/I_\lambda)$. Si $I_\lambda \subset I_\mu$, le
morphisme induit
$$
\Spec(A/I_\mu) \to \Spec(A/I_\lambda)
$$
est un épaississement, car $I_\mu^n \subset I_\lambda$ pour un certain $n$.
On peut aussi le voir en remarquant que chacune des
applications
$$
\Spec(A/I_\lambda) \to \Spec(A)
$$
est un homéomorphisme sur l'ensemble des idéaux premiers ouverts de $A$.
Ceci motive la définition
$$
\text{Spf}(A) = \{\text{idéaux premiers ouverts }\mathfrak p \subset A\}
$$
muni de la topologie induite par $\Spec(A)$. Pour tout $\lambda$,
nous pouvons considérer le faisceau structural $\mathcal{O}_{\Spec(A/I_\lambda)}$
comme un faisceau sur $\text{Spf}(A)$. Soit $\mathcal{O}_\lambda$ le
faisceau pseudo-discret d'anneaux topologiques correspondant ; voir ci-dessus.
On pose alors
$$
\mathcal{O}_{\text{Spf}(A)} = \lim \mathcal{O}_\lambda
$$
où la limite est prise dans la catégorie des faisceaux d'anneaux topologiques.
Le couple $(\text{Spf}(A), \mathcal{O}_{\text{Spf}(A)})$ est appelé le
{\it spectre formel} de $A$.

\medskip\noindent
Il convient alors de vérifier plusieurs points. Premièrement,
les fibres $\mathcal{O}_{\text{Spf}(A), x}$ sont des anneaux locaux
(lorsqu'on oublie la topologie). Deuxièmement, étant donné
$f \in A$, pour l'ouvert correspondant $D(f) \cap \text{Spf}(A)$,
on a
$$
\Gamma(D(f) \cap \text{Spf}(A), \mathcal{O}_{\text{Spf}(A)})
= A_{\{f\}} = \lim (A/I_\lambda)_f
$$
comme anneaux topologiques, où $I_\lambda$ est un système fondamental d'idéaux
de définition comme ci-dessus. De plus, l'anneau $A_{\{f\}}$ est lui aussi admissible, et
$(\text{Spf}(A_f), \mathcal{O}_{\text{Spf}(A_{\{f\}})})$
est isomorphe à
$(D(f) \cap \text{Spf}(A),
\mathcal{O}_{\text{Spf}(A)}|_{D(f) \cap \text{Spf}(A)})$.
Enfin, pour une paire d'anneaux topologiques admissibles $A, B$,
on a
\begin{equation}
\label{equation-morphisms-affine-formal-schemes}
\Mor_{LTRS}((\text{Spf}(B), \mathcal{O}_{\text{Spf}(B)}),
(\text{Spf}(A), \mathcal{O}_{\text{Spf}(A)}))
= \Hom_{cont}(A, B)
\end{equation}
où $LTRS$ est la catégorie des « espaces localement annelés topologiques »
définie ci-dessus.

\medskip\noindent
Cela étant, on définit dans \cite{EGA} un {\it schéma formel} comme un couple
$(\mathfrak X, \mathcal{O}_\mathfrak X)$ où $\mathfrak X$
est un espace topologique et $\mathcal{O}_\mathfrak X$ un faisceau
d'anneaux topologiques tel que tout point possède un voisinage ouvert
isomorphe (dans $LTRS$) à un schéma formel affine.
Un {\it morphisme de schémas formels}
$f : (\mathfrak X, \mathcal{O}_\mathfrak X) \to
(\mathfrak Y, \mathcal{O}_\mathfrak Y)$
est un morphisme de la catégorie $LTRS$.

\medskip\noindent
Soit $A$ un anneau muni de la topologie discrète. Alors $A$ est
admissible et le schéma formel $\text{Spf}(A)$ est égal à
$\Spec(A)$. Le faisceau structural $\mathcal{O}_{\text{Spf}(A)}$
est le faisceau pseudo-discret d'anneaux topologiques associé
à $\mathcal{O}_{\Spec(A)}$ ; autrement dit, son faisceau
d'anneaux sous-jacent est égal à $\mathcal{O}_{\Spec(A)}$, et l'anneau
$\mathcal{O}_{\text{Spf}(A)}(U) = \mathcal{O}_{\Spec(A)}(U)$
sur un ouvert quasi-compact $U$ est muni de la topologie discrète,
mais pas en général. On peut donc associer à tout schéma affine
un schéma formel affine. De la même manière, on peut partir
d'un schéma quelconque $(X, \mathcal{O}_X)$ et lui associer
$(X, \mathcal{O}'_X)$, où $\mathcal{O}'_X$ est le
faisceau pseudo-discret d'anneaux topologiques dont le faisceau
d'anneaux sous-jacent est $\mathcal{O}_X$. Cette construction est
compatible aux morphismes et définit un foncteur
\begin{equation}
\label{equation-compare-schemes-formal-schemes}
\textit{Schemes} \longrightarrow \textit{Formal Schemes}
\end{equation}
Il résulte immédiatement de
(\ref{equation-morphisms-affine-formal-schemes})
que ce foncteur est pleinement fidèle.

\medskip\noindent
Soit $\mathfrak X$ un schéma formel. Définissons la {\it taille}
du schéma formel par la formule
$\text{size}(\mathfrak X) = \max(\aleph_0, \kappa_1, \kappa_2)$,
où $\kappa_1$ est le cardinal de l'ensemble des ouverts formels affines de
$\mathfrak X$ et $\kappa_2$ le suprémum des cardinaux
des $\mathcal{O}_\mathfrak X(\mathfrak U)$, où
$\mathfrak U \subset \mathfrak X$ est un tel ouvert formel affine.

\begin{lemma}
\label{lemma-fully-faithful}
Choisissons une catégorie de schémas $\Sch_\alpha$
comme dans Sets, Lemma \ref{sets-lemma-construct-category}.
À tout schéma formel $\mathfrak X$, associons
$$
h_\mathfrak X : (\Sch_\alpha)^{opp} \longrightarrow \textit{Sets},\quad
h_\mathfrak X(S) = \Mor_{\textit{Formal Schemes}}(S, \mathfrak X)
$$
son foncteur des points. Alors on a
$$
\Mor_{\textit{Formal Schemes}}(\mathfrak X, \mathfrak Y) =
\Mor_{\textit{PSh}(\Sch_\alpha)}(h_\mathfrak X, h_\mathfrak Y)
$$
pourvu que la taille de $\mathfrak X$ ne soit pas trop grande.
\end{lemma}

\begin{proof}
Observons d'abord que $h_\mathfrak X$ vérifie la condition de faisceau pour
la topologie de Zariski, quel que soit le schéma formel $\mathfrak X$ (voir
Schemes, Definition \ref{schemes-definition-representable-by-open-immersions}).
Cela découle du caractère local des morphismes dans la catégorie
des schémas formels. En outre, pour une immersion ouverte
$\mathfrak V \to \mathfrak W$ de schémas formels,
la transformation de foncteurs correspondante $h_\mathfrak V \to h_\mathfrak W$
est injective et représentable par des immersions ouvertes (voir
Schemes, Definition \ref{schemes-definition-representable-by-open-immersions}).
Choisissons un recouvrement ouvert $\mathfrak X = \bigcup \mathfrak U_i$
d'un schéma formel par des schémas formels affines $\mathfrak U_i$.
La famille de foncteurs
$h_{\mathfrak U_i}$ recouvre alors $h_\mathfrak X$ (voir
Schemes, Definition \ref{schemes-definition-representable-by-open-immersions}).
Enfin, remarquons que
$$
h_{\mathfrak U_i} \times_{h_\mathfrak X} h_{\mathfrak U_j} =
h_{\mathfrak U_i \cap \mathfrak U_j}
$$
Par conséquent, la donnée d'une application $h_\mathfrak X \to h_\mathfrak Y$
équivaut à celle d'une famille d'applications
$h_{\mathfrak U_i} \to h_\mathfrak Y$ qui coïncident sur les intersections.
Nous pouvons donc ramener la bijectivité (resp. l'injectivité) de l'application
du lemme à la bijectivité (resp. l'injectivité) pour les couples
$(\mathfrak U_i, \mathfrak Y)$
et à l'injectivité (resp. aucune condition)
pour $(\mathfrak U_i \cap \mathfrak U_j, \mathfrak Y)$.
On se ramène ainsi au cas où $\mathfrak X$ est un
schéma formel affine. Écrivons $\mathfrak X = \text{Spf}(A)$
pour un anneau topologique admissible $A$. Choisissons aussi un
système fondamental d'idéaux de définition $I_\lambda \subset A$.

\medskip\noindent
Nous pouvons aussi nous localiser sur $\mathfrak Y$.
Supposons en effet que $\mathfrak V \subset \mathfrak Y$ soit un
sous-schéma formel ouvert, et soit $\varphi : h_\mathfrak X \to h_\mathfrak Y$.
Alors
$$
h_\mathfrak V \times_{h_\mathfrak Y, \varphi} h_\mathfrak X \to h_\mathfrak X
$$
est représentable par des immersions ouvertes. Après changement de base à
$\Spec(A/I_\lambda)$ pour tout $\lambda$, on obtient un sous-schéma ouvert
$U_\lambda \subset \Spec(A/I_\lambda)$. Toutefois, si
$I_\lambda \subset I_\mu$, le morphisme $\Spec(A/I_\lambda) \to \Spec(A/I_\mu)$
transforme par image inverse $U_\mu$ en $U_\lambda$. Ceux-ci se recollent donc pour donner
un sous-schéma formel ouvert $\mathfrak U \subset \mathfrak X$.
Un argument immédiat (omis) montre que
$$
h_\mathfrak U = h_\mathfrak V \times_{h_\mathfrak Y} h_\mathfrak X
$$
On voit ainsi qu'un recouvrement ouvert
$\mathfrak Y = \bigcup \mathfrak V_j$ et une transformation
de foncteurs $\varphi :  h_\mathfrak X \to h_\mathfrak Y$
déterminent un recouvrement ouvert correspondant de $\mathfrak X$.
Comme $\mathfrak X$ est affine, on peut raffiner ce recouvrement en un
recouvrement ouvert fini
$\mathfrak X = \mathfrak U_1 \cup \ldots \cup \mathfrak U_n$
par des sous-schémas formels affines. Autrement dit, pour tout $i$, il
existe un $j$ et une application $\varphi_i : h_{\mathfrak U_i} \to h_{\mathfrak V_j}$
tels que le diagramme
$$
\xymatrix{
h_{\mathfrak U_i} \ar[r]_{\varphi_i} \ar[d] & h_{\mathfrak V_j} \ar[d] \\
h_{\mathfrak X} \ar[r]^\varphi & h_\mathfrak Y
}
$$
soit commutatif. Quelques arguments supplémentaires (que nous omettons) montrent alors
qu'il suffit d'établir la bijectivité du lemme lorsque
$\mathfrak X$ et $\mathfrak Y$ sont tous deux des schémas formels affines.

\medskip\noindent
Supposons que $\mathfrak X$ et $\mathfrak Y$ soient des schémas formels affines.
Écrivons $\mathfrak X = \text{Spf}(A)$ et $\mathfrak Y = \text{Spf}(B)$.
Soit $\varphi : h_\mathfrak X \to h_\mathfrak Y$ une transformation
de foncteurs. Soit $I_\lambda \subset A$ un système fondamental
d'idéaux de définition. Le morphisme d'inclusion canonique
$i_\lambda : \Spec(A/I_\lambda) \to \mathfrak X$ est envoyé sur un morphisme
$\varphi(i_\lambda) : \Spec(A/I_\lambda) \to \mathfrak Y$.
D'après (\ref{equation-morphisms-affine-formal-schemes}), celui-ci correspond
à une application continue $\chi_\lambda : B \to A/I_\lambda$.
Comme $\varphi$ est une transformation de foncteurs, il s'ensuit
que, pour $I_\lambda \subset I_\mu$, le composé
$B \to A/I_\lambda \to A/I_\mu$ est égal à $\chi_\mu$.
Autrement dit, on obtient un homomorphisme d'anneaux
$$
\chi = \lim \chi_\lambda : B \longrightarrow \lim A/I_\lambda = A
$$
C'est un homomorphisme continu, car l'image réciproque
de $I_\lambda$ est ouverte pour tout $\lambda$ (puisque $A/I_\lambda$ est muni de la topologie
discrète et que $\chi_\lambda$ est continue). On obtient donc
un morphisme $\text{Spf}(\chi) : \mathfrak X \to \mathfrak Y$ par
(\ref{equation-morphisms-affine-formal-schemes}).
Nous omettons la vérification que cette construction est l'inverse
de l'application du lemme dans ce cas.

\medskip\noindent
Remarques ensemblistes. Pour que cela fonctionne sur la catégorie donnée
de schémas $\Sch_\alpha$, il suffit de s'assurer que tous les
schémas utilisés dans la démonstration ci-dessus sont isomorphes à des objets de $\Sch_\alpha$.
En fait, une analyse attentive montre qu'il suffit que les
schémas $\Spec(A/I_\lambda)$ intervenant ci-dessus soient isomorphes à des
objets de $\Sch_\alpha$. Pour cela, il suffit certainement de supposer que
la taille de $\mathfrak X$ est au plus celle
d'un schéma appartenant à $\Sch_\alpha$.
\end{proof}

\begin{lemma}
\label{lemma-formal-scheme-sheaf-fppf}
\begin{slogan}
Les schémas formels sont des faisceaux fpqc
\end{slogan}
Soit $\mathfrak X$ un schéma formel. Le foncteur des points
$h_\mathfrak X$ (voir Lemma \ref{lemma-fully-faithful})
vérifie la condition de faisceau pour les recouvrements fpqc.
\end{lemma}

\begin{proof}
Topologies, Lemma \ref{topologies-lemma-sheaf-property-fpqc}
nous ramène au cas d'un recouvrement de Zariski et d'un recouvrement
$\{\Spec(S) \to \Spec(R)\}$ où $R \to S$ est fidèlement plat.
Nous avons observé dans la démonstration de Lemma \ref{lemma-fully-faithful} 
que $h_\mathfrak X$ vérifie la condition de faisceau pour les recouvrements de Zariski.

\medskip\noindent
Supposons que $R \to S$ soit un homomorphisme d'anneaux fidèlement plat.
Notons $\pi : \Spec(S) \to \Spec(R)$ le
morphisme de schémas correspondant. Il est surjectif et plat.
Soit $f : \Spec(S) \to \mathfrak X$ un morphisme
tel que $f \circ \text{pr}_1 = f \circ \text{pr}_2$
comme applications $\Spec(S \otimes_R S) \to \mathfrak X$.
D'après Descent, Lemma \ref{descent-lemma-equiv-fibre-product},
l'application $f$ entre les ensembles sous-jacents
est de la forme $f = g \circ \pi$ pour une certaine
application (ensembliste) $g : \Spec(R) \to \mathfrak X$.
D'après Morphisms, Lemma \ref{morphisms-lemma-fpqc-quotient-topology},
et puisque $f$ est continue, on voit que $g$
est continue.

\medskip\noindent
Choisissons $y \in \Spec(R)$. Prenons un sous-schéma formel ouvert affine
$\mathfrak U \subset \mathfrak X$ contenant $g(y)$.
Écrivons $\mathfrak U = \text{Spf}(A)$ pour un anneau topologique
admissible $A$. D'après ce qui précède, on peut choisir $r \in R$ tel que
$y \in D(r) \subset g^{-1}(\mathfrak U)$.
La restriction de $f$ à $\pi^{-1}(D(r))$, à valeurs dans $\mathfrak U$,
correspond à une application continue d'anneaux $A \to S_r$ d'après
(\ref{equation-morphisms-affine-formal-schemes}). Les deux homomorphismes d'anneaux induits
$A \to S_r \otimes_{R_r} S_r = (S \otimes_R S)_r$ sont égaux
par l'hypothèse sur $f$.
Remarquons que $R_r \to S_r$ est fidèlement plat.
D'après Descent, Lemma \ref{descent-lemma-ff-exact}, l'égalisateur des
deux flèches $S_r \to S_r \otimes_{R_r} S_r$ est $R_r$.
On en déduit que $A \to S_r$ se factorise de manière unique par une application $A \to R_r$
qui est également continue, car elle a le même noyau (ouvert) que
l'application $A \to S_r$. Cette application donne à son tour un morphisme $D(r) \to \mathfrak U$ par
(\ref{equation-morphisms-affine-formal-schemes}).

\medskip\noindent
Qu'avons-nous démontré jusqu'ici ? Pour tout $y \in \Spec(R)$,
il existe un ouvert affine standard
$y \in D(r) \subset \Spec(R)$ tel que le morphisme
$f|_{\pi^{-1}(D(r))} : \pi^{-1}(D(r)) \to \mathfrak X$ se factorise de manière unique
par un certain morphisme $D(r) \to \mathfrak X$. Nous omettons la
vérification que ces morphismes se recollent pour donner le
morphisme recherché $\Spec(R) \to \mathfrak X$.
\end{proof}

\begin{remark}[Variante de McQuillan]
\label{remark-mcquillan}
Il existe une variante de la construction des schémas formels due à
McQuillan ; voir \cite{McQuillan}.
Il propose d'affaiblir légèrement la condition d'admissibilité.
Rappelons en effet qu'un anneau topologique admissible est un anneau topologique $A$
complet (et séparé selon nos conventions),
linéairement topologisé, qui possède un
idéal de définition, c'est-à-dire un
idéal ouvert $I$ tel que tout voisinage de $0$ contienne $I^n$
pour un certain $n \geq 1$.
McQuillan travaille avec ce que nous appellerons des anneaux topologiques
{\it faiblement admissibles}. Un anneau topologique faiblement admissible $A$ est un anneau topologique
complet (et séparé selon nos conventions),
linéairement topologisé, qui possède un
{\it idéal faible de définition}, c'est-à-dire un idéal ouvert $I$ tel que,
pour tout $f \in I$, on ait
$f^n \to 0$ lorsque $n \to \infty$. Comme dans le cas admissible,
si $I$ est un idéal faible de définition et si $J \subset A$ est un
idéal ouvert, alors $I \cap J$ est un idéal faible de définition.
Les idéaux faibles de définition forment donc un système fondamental de
voisinages ouverts de $0$, et
l'on peut suivre essentiellement la même démarche que ci-dessus
pour définir une catégorie plus grande de schémas formels fondée
sur cette notion. Les analogues de Lemmas \ref{lemma-fully-faithful} et
\ref{lemma-formal-scheme-sheaf-fppf}
restent valables dans ce cadre (avec la même démonstration).
\end{remark}

\begin{remark}[Faisceau associé à un préfaisceau d'espaces topologiques]
\label{remark-sheafification-of-presheaves-in-top}
\begin{reference}
\cite{Gray}
\end{reference}
Dans cette remarque, nous étudions brièvement la construction du faisceau associé à un préfaisceau
d'espaces topologiques. Les mêmes arguments s'appliquent exactement aux
préfaisceaux de groupes abéliens topologiques, d'anneaux topologiques et de
modules topologiques (sur un anneau topologique donné). Pour traiter
la question dans toute sa généralité, travaillons sur un site
$\mathcal{C}$. Le lecteur qui s'intéresse au cas des (pré)faisceaux
sur un espace topologique $X$ doit voir les objets de $\mathcal{C}$
comme les ouverts de $X$, les morphismes de $\mathcal{C}$ comme les inclusions
d'ouverts, et les recouvrements de $\mathcal{C}$ comme les recouvrements de $X$ ; voir
Sites, Example \ref{sites-example-site-topological}.
Notons $\Sh(\mathcal{C}, \textit{Top})$ la catégorie des faisceaux
d'espaces topologiques sur $\mathcal{C}$ et
$\textit{PSh}(\mathcal{C}, \textit{Top})$ celle des préfaisceaux
d'espaces topologiques sur $\mathcal{C}$.
Soit $\mathcal{F}$ un préfaisceau d'espaces topologiques sur $\mathcal{C}$.
Le faisceau associé $\mathcal{F}^\#$ doit vérifier la formule
$$
\Mor_{\textit{PSh}(\mathcal{C}, \textit{Top})}(\mathcal{F}, \mathcal{G})
=
\Mor_{\Sh(\mathcal{C}, \textit{Top})}(\mathcal{F}^\#, \mathcal{G})
$$
fonctoriellement en $\mathcal{G}$ dans $\Sh(\mathcal{C}, \textit{Top})$.
Autrement dit, nous cherchons à construire l'adjoint à gauche
du foncteur d'inclusion
$\Sh(\mathcal{C}, \textit{Top}) \to \textit{PSh}(\mathcal{C}, \textit{Top})$.
Affirmons d'abord que $\Sh(\mathcal{C}, \textit{Top})$ admet des limites
et que le foncteur d'inclusion commute à celles-ci.
En effet, étant donnés une catégorie $\mathcal{I}$ et un foncteur
$i \mapsto \mathcal{G}_i$ à valeurs dans $\Sh(\mathcal{C}, \textit{Top})$,
nous définissons simplement
$$
(\lim \mathcal{G}_i)(U) = \lim \mathcal{G}_i(U)
$$
où la limite est prise dans la catégorie des espaces topologiques
(Topology, Lemma \ref{topology-lemma-limits}). Cela définit un faisceau,
car les limites commutent aux limites
(Categories, Lemma \ref{categories-lemma-colimits-commute}),
et en particulier aux produits et aux égalisateurs (qui sont les
opérations utilisées dans l'axiome des faisceaux). Enfin, un morphisme
de préfaisceaux $\mathcal{F} \to \lim \mathcal{G}_i$ est
manifestement la même chose qu'un système compatible de morphismes
$\mathcal{F} \to \mathcal{G}_i$. Autrement dit, l'objet
$\lim \mathcal{G}_i$ est la limite dans la catégorie
des préfaisceaux d'espaces topologiques, et a fortiori dans la
catégorie des faisceaux d'espaces topologiques.
Notre seconde affirmation est que tout morphisme de préfaisceaux
$\mathcal{F} \to \mathcal{G}$, où $\mathcal{G}$ est un objet de
$\Sh(\mathcal{C}, \textit{Top})$, se factorise par un sous-faisceau
$\mathcal{G}' \subset \mathcal{G}$ de taille bornée.
Nous définissons ici la {\it taille} $|\mathcal{H}|$
d'un faisceau d'espaces topologiques $\mathcal{H}$ comme le cardinal
$\sup_{U \in \Ob(\mathcal{C})} |\mathcal{H}(U)|$.
Pour démontrer l'affirmation, posons
$$
\mathcal{G}'(U) =
\left\{
\quad
s \in \mathcal{G}(U)
\quad \middle| \quad
\begin{matrix}
\text{il existe un recouvrement }\{U_i \to U\}_{i \in I} \\
\text{tel que }
s|_{U_i} \in \Im(\mathcal{F}(U_i) \to \mathcal{G}(U_i))
\end{matrix}
\quad
\right\}
$$
Nous munissons $\mathcal{G}'(U)$ de la topologie induite.
Alors $\mathcal{G}'$ est un faisceau d'espaces topologiques (détails omis),
et $\mathcal{G}' \to \mathcal{G}$ est un morphisme par lequel
l'application donnée $\mathcal{F} \to \mathcal{G}$ se factorise. De plus,
la taille de $\mathcal{G}'$ est bornée par un cardinal
$\kappa$ qui ne dépend que de $\mathcal{C}$ et du préfaisceau $\mathcal{F}$
(indication : utiliser le fait que, selon nos conventions, les recouvrements de $\mathcal{C}$
forment un ensemble). En rassemblant ces observations, on voit
que les hypothèses de Categories, Theorem
\ref{categories-theorem-adjoint-functor}
sont vérifiées, et l'on obtient le foncteur faisceau associé comme adjoint à gauche
du foncteur d'inclusion des faisceaux dans les préfaisceaux.
Enfin, soit $p$ un point du
site $\mathcal{C}$ défini par un foncteur $u : \mathcal{C} \to \textit{Sets}$ ;
voir Sites, Definition \ref{sites-definition-point}.
Pour un espace topologique $M$, le préfaisceau défini par la règle
$$
U \mapsto \text{Map}(u(U), M) = \prod\nolimits_{x \in u(U)} M
$$
et muni de la topologie produit est un faisceau d'espaces topologiques.
Le même argument que celui donné dans la démonstration de
Sites, Lemma \ref{sites-lemma-point-pushforward-sheaf}, montre donc que
$\mathcal{F}_p = \mathcal{F}^\#_p$ ; autrement dit, le foncteur faisceau associé
commute à la prise de la fibre en un point.
\end{remark}




\section{Conventions et notation}
\label{section-conventions}

\noindent
Les conventions adoptées désormais seront semblables à celles de
Properties of Spaces, Section \ref{spaces-properties-section-conventions}.
Ainsi, nous supposerons toujours que tous les schémas appartiennent
à un grand site fppf $\Sch_{fppf}$. De plus, tout anneau $A$ considéré aura la
propriété que $\Spec(A)$ est (isomorphe à) un objet de ce grand site.
Pour les anneaux topologiques $A$, nous supposons seulement que tous leurs quotients discrets ont
cette propriété (mais nous faisons habituellement des hypothèses plus fortes ; comparer à
Remark \ref{remark-set-theoretic}).

\medskip\noindent
Soient $S$ un schéma et $X$ un « espace » sur $S$, c'est-à-dire un faisceau sur
$(\Sch/S)_{fppf}$. Dans ce chapitre, nous noterons $X \times_S X$ le
produit de $X$ par lui-même dans la catégorie des faisceaux sur $(\Sch/S)_{fppf}$,
et non $X \times X$. De plus, si $X$ et $Y$ sont des « espaces »,
nous dirons « soit $f : X \to Y$ un morphisme » pour indiquer que $f$ est une
transformation naturelle de foncteurs, c'est-à-dire une application de faisceaux sur
$(\Sch/S)_{fppf}$. De même, si $U$ est un schéma sur $S$ et si
$X$ est un « espace » sur $S$, nous dirons
« soit $f : U \to X$ un morphisme » ou
« soit $g : X \to U$ un morphisme » pour indiquer que $f$ ou $g$
est une application de faisceaux $h_U \to X$ ou $X \to h_U$, où $h_U$ est défini dans
Categories, Example \ref{categories-example-hom-functor}.






\section{Anneaux et modules topologiques}
\label{section-topological-module}

\noindent
Cette section prolonge
More on Algebra, Section \ref{more-algebra-section-topological-ring}.
Soient $R$ un anneau topologique et $M$ un $R$-module linéairement topologisé.
Lorsque nous dirons « {\it soit $M_\lambda$ un système fondamental de
sous-modules ouverts} », nous entendrons que tout $M_\lambda$ est un sous-module ouvert
et que tout voisinage de $0$ contient l'un des $M_\lambda$.
Autrement dit, cela signifie que les $M_\lambda$ forment un système fondamental
de voisinages de $0$ dans $M$, constitué de sous-modules.
De même, si $R$ est un anneau linéairement topologisé, nous dirons
« {\it soit $I_\lambda$ un système fondamental d'idéaux ouverts} »
pour signifier que les $I_\lambda$ forment un système fondamental
de voisinages de $0$ dans $R$, constitué d'idéaux.

\begin{example}
\label{example-what-does-it-mean}
Soient $R$ un anneau linéairement topologisé et $M$ un
$R$-module linéairement topologisé. Soient $I_\lambda$ un système fondamental
d'idéaux ouverts de $R$ et $M_\mu$ un système fondamental de
sous-modules ouverts de $M$. La continuité de $+ : M \times M \to M$
est automatique, et la continuité de $R \times M \to M$ signifie que
$$
\forall f, x, \mu\ \exists \lambda, \nu,\ (f + I_\lambda)(x + M_\nu)
\subset fx + M_\mu
$$
Puisque $fM_\nu + I_\lambda M_\nu \subset M_\mu$ dès que
$M_\nu \subset M_\mu$, on voit que la condition équivaut à
$$
\forall x, \mu\ \exists \lambda\ I_\lambda x \subset M_\mu
$$
Cependant, il ne faut pas nécessairement qu'étant donné $\mu$, il existe $\lambda$
tel que $I_\lambda M \subset M_\mu$. Considérons par exemple
$R = k[[t]]$ muni de la topologie $t$-adique et
$M = \bigoplus_{n \in \mathbf{N}} R$, avec le système fondamental de
sous-modules ouverts donné par
$$
M_m = \bigoplus\nolimits_{n \in \mathbf{N}} t^{nm}R
$$
Comme tout $x \in M$ n'a qu'un nombre fini de coordonnées non nulles, on voit
qu'étant donnés $m$ et $x$, il existe $k$ tel que $t^k x \in M_m$.
Ainsi, $M$ est un $R$-module linéairement topologisé, mais il n'est pas vrai
qu'étant donné $m$, il existe $k$ tel que $t^kM \subset M_m$.
En revanche, si $R \to S$ est un homomorphisme continu d'anneaux
linéairement topologisés, l'énoncé correspondant est vrai :
pour tout idéal ouvert $J \subset S$, il existe un idéal ouvert
$I \subset R$ tel que $IS \subset J$ (comme le lecteur le déduira aisément
de la continuité de l'homomorphisme $R \to S$).
\end{example}

\begin{lemma}
\label{lemma-closed}
Soit $R$ un anneau topologique. Soient $M$ un $R$-module linéairement
topologisé et $M_\lambda$, $\lambda \in \Lambda$, un système fondamental
de sous-modules ouverts. Soit $N \subset M$ un sous-module.
L'adhérence de $N$ est $\bigcap_{\lambda \in \Lambda} (N + M_\lambda)$.
\end{lemma}

\begin{proof}
Comme chaque $N + M_\lambda$ est ouvert, il est également fermé. Leur
intersection est donc fermée. Si $x \in M$ n'appartient pas à l'adhérence de $N$,
alors $(x + M_\lambda) \cap N = 0$ pour un certain $\lambda$. Ainsi,
$x \not \in N + M_\lambda$. Cela démontre le lemme.
\end{proof}

\noindent
Sauf mention contraire, nous munissons les sous-modules et les modules quotients
de la topologie induite. Soient $M$ un module linéairement topologisé
sur un anneau topologique $R$, et $0 \to N \to M \to Q \to 0$
une suite exacte courte de $R$-modules. Si les $M_\lambda$ forment un
système fondamental de sous-modules ouverts de $M$, alors
les $N \cap M_\lambda$ forment un système fondamental de sous-modules ouverts de $N$.
Si $\pi : M \to Q$ est l'application quotient, les $\pi(M_\lambda)$ forment un
système fondamental de sous-modules ouverts de $Q$. En particulier, ces topologies
induites sont des topologies linéaires.

\begin{lemma}
\label{lemma-closure}
Soit $R$ un anneau topologique. Soit $M$ un $R$-module linéairement
topologisé. Soit $N \subset M$ un sous-module. Alors :
\begin{enumerate}
\item $0 \to N^\wedge \to M^\wedge \to (M/N)^\wedge$ est exacte, et
\item $N^\wedge$ est l'adhérence de l'image de $N \to M^\wedge$.
\end{enumerate}
\end{lemma}

\begin{proof}
Soient $M_\lambda$, $\lambda \in \Lambda$, un système fondamental de
sous-modules ouverts. Alors les $N \cap M_\lambda$ forment un système fondamental
de sous-modules ouverts de $N$, et les $M_\lambda + N/N$ un système fondamental
de sous-modules ouverts de $M/N$. On voit donc que (1) découle de
l'exactitude des suites
$$
0 \to N/N \cap M_\lambda \to M/M_\lambda \to M/(M_\lambda + N) \to 0
$$
et du fait que la prise de limites commute aux limites. La seconde
assertion en résulte, ainsi que du fait que $N \to N^\wedge$
est d'image dense et que le noyau de $M^\wedge \to (M/N)^\wedge$ est fermé.
\end{proof}

\begin{lemma}
\label{lemma-quotient-by-closed}
Soit $R$ un anneau topologique. Soit $M$ un $R$-module complet et linéairement
topologisé. Soit $N \subset M$ un sous-module fermé. Si $M$ possède un
système fondamental dénombrable de voisinages de $0$, alors
$M/N$ est complet et l'application $M \to M/N$ est ouverte.
\end{lemma}

\begin{proof}
Soient $M_n$, $n \in \mathbf{N}$, un système fondamental de sous-modules ouverts de $M$.
Nous pouvons supposer $M_{n + 1} \subset M_n$
pour tout $n$. Les $(M_n + N)/N$ forment un système fondamental dans $M/N$.
Il faut donc montrer que $M/N = \lim M/(M_n + N)$. Considérons
les suites exactes courtes
$$
0 \to N/N \cap M_n \to M/M_n \to M/(M_n + N) \to 0
$$
Comme les applications de transition du système $\{N/N\cap M_n\}$ sont surjectives,
on voit que $M = \lim M/M_n$ (puisque $M$ est complet) se surjecte sur
$\lim M/(M_n + N)$ d'après
Algebra, Lemma \ref{algebra-lemma-ML-exact-sequence}.
Puisque $N$ est fermé, le noyau de $M \to \lim M/(M_n + N)$
est $N$ (voir Lemma \ref{lemma-closed}). Enfin, $M \to M/N$
est ouverte par définition de la topologie quotient.
\end{proof}

\begin{lemma}
\label{lemma-ses}
\begin{reference}
\cite[Theorem 8.1]{Ma}
\end{reference}
Soit $R$ un anneau topologique. Soit $M$ un $R$-module linéairement
topologisé. Soit $N \subset M$ un sous-module. Supposons que $M$ possède un
système fondamental dénombrable de voisinages de $0$. Alors :
\begin{enumerate}
\item $0 \to N^\wedge \to M^\wedge \to (M/N)^\wedge \to 0$ est exacte,
\item $N^\wedge$ est l'adhérence de l'image de $N \to M^\wedge$,
\item $M^\wedge \to (M/N)^\wedge$ est ouverte.
\end{enumerate}
\end{lemma}

\begin{proof}
La suite $0 \to N^\wedge \to M^\wedge \to (M/N)^\wedge$ est exacte,
et l'assertion (2) résulte de Lemma \ref{lemma-closure}.
On obtient ainsi une application canonique $c : M^\wedge/N^\wedge \to (M/N)^\wedge$.
Le module $M^\wedge/N^\wedge$ est complet, et
$M^\wedge \to M^\wedge/N^\wedge$ est ouverte d'après
Lemma \ref{lemma-quotient-by-closed}.
La propriété universelle du complété donne une application canonique
$b : (M/N)^\wedge \to M^\wedge/N^\wedge$.
Alors $b$ et $c$ sont inverses l'une de l'autre, puisqu'ils le sont sur un sous-ensemble dense.
\end{proof}

\begin{lemma}
\label{lemma-completion-adic-star}
Soit $R$ un anneau topologique. Soit $M$ un $R$-module topologique.
Soit $I \subset R$ un idéal de type fini. Supposons que $M$
possède un sous-module ouvert dont la topologie est $I$-adique. Alors $M^\wedge$
possède un sous-module ouvert dont la topologie est $I$-adique, et l'on a
$M^\wedge/I^n M^\wedge = M/I^nM$ pour tout $n \geq 1$.
\end{lemma}

\begin{proof}
Soit $M' \subset M$ un sous-module ouvert dont la topologie est $I$-adique.
Alors $\{I^nM'\}_{n \geq 1}$ est un système fondamental de sous-modules ouverts
de $M$. Ainsi, $M^\wedge = \lim M/I^nM'$ contient
$(M')^\wedge = \lim M'/I^nM'$ comme sous-module ouvert,
et la topologie sur $(M')^\wedge$ est
$I$-adique d'après Algebra, Lemma \ref{algebra-lemma-hathat-finitely-generated}.
Comme $I$ est de type fini, $I^n$ est de type fini ;
soit $f_1, \ldots, f_r$ un système de générateurs. Observons que la surjection
$(f_1, \ldots, f_r) : M^{\oplus r} \to I^n M$ est continue
et ouverte d'après la description précédente de la topologie sur $M$.
En appliquant Lemma \ref{lemma-ses} à cette surjection et à la
suite exacte courte $0 \to I^nM \to M \to M/I^nM \to 0$,
on conclut que
$$
(f_1, \ldots, f_r) :
(M^\wedge)^{\oplus r} \longrightarrow M^\wedge
$$
se surjecte sur le noyau de la surjection $M^\wedge \to M/I^nM$.
Comme $f_1, \ldots, f_r$ engendrent $I^n$, cela conclut la démonstration.
\end{proof}

\begin{definition}
\label{definition-toplogy-tensor-product}
\begin{reference}
Comparer à \cite[0, Section 7.7]{EGA}
\end{reference}
Soit $R$ un anneau topologique. Soient $M$ et $N$ des $R$-modules
linéairement topologisés. Le {\it produit tensoriel} de $M$ et $N$
est le produit tensoriel (usuel) $M \otimes_R N$, muni
de la topologie linéaire définie en déclarant que les
$$
\Im(M_\mu \otimes_R N + M \otimes_R N_\nu \longrightarrow M \otimes_R N)
$$
forment un système fondamental de sous-modules ouverts, lorsque
$M_\mu \subset M$ et $N_\nu \subset N$ parcourent des systèmes fondamentaux
de sous-modules ouverts de $M$ et de $N$.
Le {\it produit tensoriel complété}
$$
M \widehat{\otimes}_R N =
\lim M \otimes_R N/(M_\mu \otimes_R N + M \otimes_R N_\nu) =
\lim M/M_\mu \otimes_R N/N_\nu
$$
est le complété du produit tensoriel.
\end{definition}

\noindent
Observons que la topologie de $R$ n'intervient pas dans la construction
du produit tensoriel ni du produit tensoriel complété.
Si $R \to A$ et $R \to B$ sont des homomorphismes continus d'anneaux
linéairement topologisés, la construction précédente
fournit un produit tensoriel $A \otimes_R B$ et un produit
tensoriel complété $A \widehat{\otimes}_R B$.

\medskip\noindent
Consignons ici les notions introduites dans Remark \ref{remark-mcquillan}.

\begin{definition}
\label{definition-weakly-admissible}
Soit $A$ un anneau linéairement topologisé.
\begin{enumerate}
\item Un élément $f \in A$ est dit {\it topologiquement nilpotent}
si $f^n \to 0$ lorsque $n \to \infty$.
\item Un {\it idéal faible de définition} de $A$ est un idéal ouvert
$I \subset A$ entièrement constitué d'éléments topologiquement nilpotents.
\item On dit que $A$ est {\it faiblement préadmissible} si $A$ possède un idéal faible
de définition.
\item On dit que $A$ est {\it faiblement admissible} si $A$ est faiblement préadmissible
et complet\footnote{Selon nos conventions, cela inclut la séparation.}.
\end{enumerate}
\end{definition}

\noindent
Étant donnés un idéal faible de définition $I$ dans un anneau linéairement topologisé
$A$ et un idéal ouvert $J$, l'intersection $I \cap J$ est un
idéal faible de définition. Par conséquent, s'il existe un idéal faible de définition,
il existe un système fondamental d'idéaux ouverts
constitué d'idéaux faibles de définition. En particulier,
si $A$ est un anneau topologique faiblement admissible, alors
$A = \lim A/I_\lambda$, où $\{I_\lambda\}$ est un système fondamental
d'idéaux faibles de définition.

\begin{lemma}
\label{lemma-weakly-admissible-henselian}
Soit $A$ un anneau topologique faiblement admissible. Soit $I \subset A$
un idéal faible de définition. Alors $(A, I)$ est un couple hensélien.
\end{lemma}

\begin{proof}
Soit $A \to A'$ un homomorphisme d'anneaux étale et soit $\sigma : A' \to A/I$
un homomorphisme de $A$-algèbres. D'après More on Algebra, Lemma
\ref{more-algebra-lemma-characterize-henselian-pair}, il suffit
de relever $\sigma$ en un homomorphisme de $A$-algèbres $A' \to A$.
Pour cela, comme $A$ est complet, il suffit de trouver,
pour tout idéal ouvert $J \subset I$, un unique homomorphisme de $A$-algèbres $A' \to A/J$
relevant $\sigma$. Comme $I$ est un idéal faible de définition,
l'idéal $I/J$ est localement nilpotent. On conclut par
More on Algebra, Lemma \ref{more-algebra-lemma-locally-nilpotent-henselian}.
\end{proof}

\begin{lemma}
\label{lemma-topologically-nilpotent}
Soit $B$ un anneau linéairement topologisé. L'ensemble des éléments topologiquement nilpotents
de $B$ est un idéal radical fermé de $B$.
Soit $\varphi : A \to B$ un homomorphisme continu d'anneaux linéairement topologisés.
\begin{enumerate}
\item Si $f \in A$ est topologiquement nilpotent, alors $\varphi(f)$ est
topologiquement nilpotent.
\item Si $I \subset A$ est constitué d'éléments topologiquement nilpotents,
alors l'adhérence de $\varphi(I)B$ est constituée d'éléments topologiquement
nilpotents.
\end{enumerate}
\end{lemma}

\begin{proof}
Soit $\mathfrak b \subset B$ l'ensemble des éléments topologiquement nilpotents.
Nous omettons la démonstration du fait que $\mathfrak b$ est un idéal radical
(c'est un bon exercice sur les définitions). Soit $g$ un élément de l'adhérence
de $\mathfrak b$. Nous voulons montrer que $g$ est topologiquement nilpotent.
Soit $J \subset B$ un idéal ouvert. Il faut montrer que
$g^e \in J$ pour un certain $e \geq 1$. On a $g \in \mathfrak b + J$
d'après Lemma \ref{lemma-closed}. Ainsi, $g = f + h$
pour certains $f \in \mathfrak b$ et $h \in J$. Choisissons $m \geq 1$ tel que
$f^m \in J$. Alors $g^{m + 1} \in J$, comme voulu.

\medskip\noindent
Soit $\varphi : A \to B$ comme dans l'énoncé du lemme.
L'assertion (1) est claire, et l'assertion (2) résulte de celle-ci et
du fait que $\mathfrak b$ est un idéal fermé.
\end{proof}

\begin{lemma}
\label{lemma-closure-image-ideal}
Soit $A \to B$ un homomorphisme continu d'anneaux linéairement topologisés.
Soit $I \subset A$ un idéal. L'adhérence de $IB$
est le noyau de $B \to B \widehat{\otimes}_A A/I$.
\end{lemma}

\begin{proof}
Soient les $J_\mu$ un système fondamental d'idéaux ouverts de $B$.
L'adhérence de $IB$ est $\bigcap (IB + J_\lambda)$ d'après Lemma \ref{lemma-closed}.
Soient les $I_\mu$ un système fondamental d'idéaux ouverts de $A$.
Alors
$$
B \widehat{\otimes}_A A/I = \lim (B/J_\lambda \otimes_A A/(I_\mu + I)) =
\lim B/(J_\lambda + I_\mu B + I B)
$$
Comme $A \to B$ est continu, pour tout $\lambda$, il
existe $\mu$ tel que $I_\mu B \subset J_\lambda$ ; voir la discussion de
Example \ref{example-what-does-it-mean}. La limite peut donc
s'écrire $\lim B/(J_\lambda + IB)$, et le résultat est clair.
\end{proof}

\begin{lemma}
\label{lemma-completed-tensor-product}
Soient $B \to A$ et $B \to C$ des homomorphismes continus d'anneaux
linéairement topologisés.
\begin{enumerate}
\item Si $A$ et $C$ sont faiblement préadmissibles, alors
$A \widehat{\otimes}_B C$ est faiblement admissible.
\item Si $A$ et $C$ sont préadmissibles, alors
$A \widehat{\otimes}_B C$ est admissible.
\item Si $A$ et $C$ possèdent un système fondamental dénombrable d'idéaux
ouverts, alors $A \widehat{\otimes}_B C$ possède un système fondamental dénombrable
d'idéaux ouverts.
\item Si $A$ et $C$ sont préadiques et possèdent des idéaux de définition
de type fini, alors $A \widehat{\otimes}_B C$ est adique et possède
un idéal de définition de type fini.
\item Si $A$ et $C$ sont des anneaux noethériens préadiques et si
$B/\mathfrak b \to A/\mathfrak a$ est de type fini,
où $\mathfrak a \subset A$ et $\mathfrak b \subset B$
sont les idéaux des éléments topologiquement nilpotents, alors
$A \widehat{\otimes}_B C$ est noethérien adique.
\end{enumerate}
\end{lemma}

\begin{proof}
Soient $I_\lambda \subset A$, $\lambda \in \Lambda$, et
$J_\mu \subset C$, $\mu \in M$,
des systèmes fondamentaux d'idéaux ouverts ; alors, par définition,
$$
A \widehat{\otimes}_B C =
\lim_{\lambda, \mu} A/I_\lambda \otimes_B C/J_\mu
$$
muni de la topologie limite. Un système fondamental d'idéaux ouverts
est donc donné par les noyaux $K_{\lambda, \mu}$ des applications
$A \widehat{\otimes}_B C \to A/I_\lambda \otimes_B C/J_\mu$.
Remarquons que $K_{\lambda, \mu}$ est l'adhérence de l'idéal
$I_\lambda(A \widehat{\otimes}_B C) + J_\mu(A \widehat{\otimes}_B C)$.
Enfin, nous avons un homomorphisme d'anneaux
$\tau : A \otimes_B C \to A \widehat{\otimes}_B C$ d'image dense.

\medskip\noindent
Démonstration de (1). Si $I_\lambda$ et $J_\mu$ sont constitués d'éléments topologiquement
nilpotents, il en va de même de $K_{\lambda, \mu}$ d'après
Lemma \ref{lemma-topologically-nilpotent}. Ainsi, 
$A \widehat{\otimes}_B C$ est faiblement admissible par définition.

\medskip\noindent
Démonstration de (2). Supposons que, pour certains $\lambda_0$ et $\mu_0$, les idéaux
$I = I_{\lambda_0} \subset A$ et $J_{\mu_0} \subset C$ soient des idéaux de
définition. Alors, pour tout $\lambda$, il existe $n$ tel que
$I^n \subset I_\lambda$. Pour tout $\mu$, il existe $m$ tel que
$J^m \subset J_\mu$. Par conséquent,
$$
\left(I(A \widehat{\otimes}_B C) + J(A \widehat{\otimes}_B C)\right)^{n + m}
\subset
I_\lambda(A \widehat{\otimes}_B C) + J_\mu(A \widehat{\otimes}_B C)
$$
Il s'ensuit que l'idéal ouvert $K = K_{\lambda_0, \mu_0}$
vérifie $K^{n + m} \subset K_{\lambda, \mu}$. Ainsi, $K$
est un idéal de définition de $A \widehat{\otimes}_B C$,
et $A \widehat{\otimes}_B C$ est admissible par définition.

\medskip\noindent
Démonstration de (3). Si $\Lambda$ et $M$ sont dénombrables, il en va de même de
$\Lambda \times M$.

\medskip\noindent
Démonstration de (4). Supposons $\Lambda = \mathbf{N}$ et $M = \mathbf{N}$,
et soient $I \subset A$ et $J \subset C$ des idéaux de type fini
tels que $I_n = I^n$ et $J_n = J^n$. Alors
$$
I(A \widehat{\otimes}_B C) + J(A \widehat{\otimes}_B C)
$$
est un idéal de type fini, et l'on voit facilement que
$A \widehat{\otimes}_B C$ est le complété de
$A \otimes_B C$ pour cet idéal. L'assertion (4)
résulte donc de Algebra, Lemma \ref{algebra-lemma-hathat-finitely-generated}.

\medskip\noindent
Démonstration de (5).
Soit $\mathfrak c \subset C$ l'idéal des éléments topologiquement nilpotents.
Comme $A$ et $C$ sont noethériens adiques, on voit que
$\mathfrak a$ et $\mathfrak c$ sont des idéaux de définition (détails omis).
D'après (4), nous savons déjà que
$A \widehat{\otimes}_B C$ est adique et que
$\mathfrak a(A \widehat{\otimes}_B C) + \mathfrak c(A \widehat{\otimes}_B C)$
est un idéal de définition de type fini. Puisque
$$
A \widehat{\otimes}_B C /
\left(\mathfrak a(A \widehat{\otimes}_B C) +
\mathfrak c(A \widehat{\otimes}_B C)\right)
=
A/\mathfrak a \otimes_{B/\mathfrak b} C/\mathfrak c
$$
est noethérien, car c'est une algèbre de type fini sur l'anneau noethérien
$C/\mathfrak c$, on conclut par
Algebra, Lemma \ref{algebra-lemma-completion-Noetherian}.
\end{proof}









\section{Homomorphismes d'anneaux tendus}
\label{section-taut-ring-maps}

\noindent
Il est commode de donner un nom à la propriété suivante
des homomorphismes continus entre anneaux linéairement topologisés.

\begin{definition}
\label{definition-taut}
Soit $\varphi : A \to B$ un homomorphisme continu d'anneaux linéairement topologisés.
On dit que $\varphi$ est {\it tendu}\footnote{Cette terminologie n'est pas standard.
La définition s'étend aux modules : on dit qu'un
$A$-module linéairement topologisé $M$ est $A$-tendu si, pour tout idéal ouvert $I \subset A$, l'adhérence
de $IM$ dans $M$ est ouverte et si ces adhérences forment un système fondamental de
voisinages de $0$ dans $M$.}
si, pour tout idéal ouvert $I \subset A$, l'adhérence de l'idéal $\varphi(I)B$
est ouverte et si ces adhérences forment un système fondamental d'idéaux ouverts.
\end{definition}

\noindent
Si $\varphi : A \to B$ est un homomorphisme continu d'anneaux linéairement topologisés
et si les $I_\lambda$ forment un système fondamental d'idéaux ouverts de $A$, alors $\varphi$
est tendu si et seulement si
les adhérences des $I_\lambda B$ sont ouvertes et forment un système fondamental
d'idéaux ouverts de $A$.

\begin{lemma}
\label{lemma-taut-weakly-admissible}
Soit $\varphi : A \to B$ un homomorphisme continu d'anneaux topologiques
faiblement admissibles. Les conditions suivantes sont équivalentes :
\begin{enumerate}
\item $\varphi$ est tendu ;
\item pour tout idéal faible de définition $I \subset A$, l'adhérence de
$\varphi(I)B$ est un idéal faible de définition de $B$, et ces idéaux forment un
système fondamental d'idéaux faibles de définition de $B$.
\end{enumerate}
\end{lemma}

\begin{proof}
Les remarques qui suivent Definition \ref{definition-taut} montrent que
(2) implique (1). Réciproquement, supposons $\varphi$ tendu. Si $I \subset A$
est un idéal faible de définition, alors l'adhérence de $\varphi(I)B$
est ouverte par définition du caractère tendu et est constituée d'éléments topologiquement
nilpotents d'après Lemma \ref{lemma-topologically-nilpotent}.
L'adhérence de $\varphi(I)B$ est donc un idéal faible de définition.
En outre, par définition du caractère tendu, ces idéaux forment un
système fondamental d'idéaux ouverts ; la condition (2) est donc vérifiée.
\end{proof}

\begin{lemma}
\label{lemma-completion-taut}
Soit $A$ un anneau linéairement topologisé. L'homomorphisme $A \to A^\wedge$
de $A$ vers son complété est tendu.
\end{lemma}

\begin{proof}
Soient les $I_\lambda$ un système fondamental d'idéaux ouverts de $A$.
Rappelons que $A^\wedge = \lim A/I_\lambda$, muni de la topologie limite,
ce qui signifie que les noyaux $J_\lambda = \Ker(A^\wedge \to A/I_\lambda)$
forment un système fondamental d'idéaux ouverts de $A^\wedge$.
Puisque $J_\lambda$ est l'adhérence de $I_\lambda A^\wedge$
(comparer à Lemma \ref{lemma-closure-image-ideal}), on conclut.
\end{proof}

\begin{lemma}
\label{lemma-composition-taut}
Soient $A \to B$ et $B \to C$ des homomorphismes continus d'anneaux
linéairement topologisés. Si $A \to B$ et $B \to C$ sont tendus, alors
$A \to C$ est tendu.
\end{lemma}

\begin{proof}
Omis. Indication : si $I \subset A$ est un idéal et si $J$ est l'adhérence
de $IB$, alors l'adhérence de $JC$ est égale à celle de $IC$.
\end{proof}

\begin{lemma}
\label{lemma-permanence-taut}
Soient $A \to B$ et $B \to C$ des homomorphismes continus d'anneaux
linéairement topologisés. Si $A \to C$ est tendu, alors
$B \to C$ est tendu.
\end{lemma}

\begin{proof}
Soit $J \subset B$ un idéal ouvert d'image réciproque $I \subset A$.
Alors l'adhérence de $JC$ contient celle de $IC$. Cette
adhérence est donc ouverte, puisque $A \to C$ est tendu. Soient les $I_\lambda$ un système fondamental
d'idéaux ouverts de $A$. Soit $K_\lambda$ l'adhérence de
$I_\lambda C$. Puisque $A \to C$ est tendu, ils forment un système fondamental
d'idéaux ouverts de $C$. Notons $J_\lambda \subset B$ l'image réciproque
de $K_\lambda$. Alors l'adhérence de $J_\lambda C$ est $K_\lambda$.
On voit donc que les adhérences des idéaux $JC$, lorsque $J$ parcourt
les idéaux ouverts de $B$, forment un système fondamental d'idéaux ouverts de $C$.
\end{proof}

\begin{lemma}
\label{lemma-base-change-taut}
Soient $A \to B$ et $A \to C$ des homomorphismes continus d'anneaux
linéairement topologisés. Si $A \to B$ est tendu, alors
$C \to B \widehat{\otimes}_A C$ est tendu.
\end{lemma}

\begin{proof}
Soit $K \subset C$ un idéal ouvert. Choisissons un idéal ouvert $I \subset A$
dont l'image dans $C$ est contenue dans $J$. Par hypothèse, l'adhérence $J$
de $IB$ est ouverte. Comme $A \to B$ est tendu, on voit que
$B \widehat{\otimes}_A C$ est la limite des anneaux $B/J \otimes_{A/I} C/K$
pour tous les choix de $K$ et de $I$, c'est-à-dire que
les idéaux $J(B \widehat{\otimes}_A C) + K(B \widehat{\otimes}_A C)$
forment un système fondamental d'idéaux ouverts. Or,
comme $B \to B \widehat{\otimes}_A C$ est continu, on voit
que $J$ s'envoie dans l'adhérence de $K(B \widehat{\otimes}_A C)$
(puisque $I$ s'envoie dans $K$). Cette adhérence est donc égale à
$J(B \widehat{\otimes}_A C) + K(B \widehat{\otimes}_A C)$,
ce qui achève la démonstration.
\end{proof}

\begin{lemma}
\label{lemma-taut-ascent-countable}
Soit $\varphi : A \to B$ un homomorphisme continu d'anneaux
linéairement topologisés. Si $\varphi$ est tendu et si $A$
possède un système fondamental dénombrable d'idéaux ouverts, alors
$B$ possède un système fondamental dénombrable d'idéaux ouverts.
\end{lemma}

\begin{proof}
Cela résulte immédiatement des définitions.
\end{proof}

\begin{lemma}
\label{lemma-taut-ascent-weakly-admissible}
Soit $\varphi : A \to B$ un homomorphisme continu d'anneaux
linéairement topologisés. Si $\varphi$ est tendu et si $A$
est faiblement préadmissible, alors $B$ est faiblement préadmissible.
\end{lemma}

\begin{proof}
Soit $I \subset A$ un idéal faible de définition. Alors l'adhérence
$J$ de $IB$ est ouverte et constituée d'éléments topologiquement nilpotents
d'après Lemma \ref{lemma-topologically-nilpotent}. Par conséquent, $J$ est un idéal faible
de définition de $B$.
\end{proof}

\begin{lemma}
\label{lemma-taut-ascent-admissible}
Soit $\varphi : A \to B$ un homomorphisme continu d'anneaux
linéairement topologisés. Si $\varphi$ est tendu et si $A$
est préadmissible, alors $B$ est préadmissible.
\end{lemma}

\begin{proof}
Soit $I \subset A$ un idéal de définition. Soient
les $I_\lambda \subset A$ un système fondamental d'idéaux ouverts.
Alors l'adhérence $J$ de $IB$ est ouverte, et les
adhérences $J_\lambda$ des $I_\lambda B$ sont ouvertes et forment un
système fondamental d'idéaux ouverts de $B$.
Pour tout $\lambda$, il existe $n$ tel que $I^n \subset I_\lambda$.
Observons que $J^n$ est contenu dans l'adhérence de $I^nB$.
Ainsi, $J^n \subset J_\lambda$, et l'on conclut que $J$ est un
idéal de définition.
\end{proof}

\begin{lemma}
\label{lemma-dense-image-surjective}
Soit $\varphi : A \to B$ un homomorphisme continu d'anneaux
linéairement topologisés. Supposons que :
\begin{enumerate}
\item $\varphi$ soit tendu et d'image dense ;
\item $A$ soit complet et possède un système fondamental dénombrable
d'idéaux ouverts ; et
\item $B$ soit séparé.
\end{enumerate}
Alors $\varphi$ est surjectif et ouvert, $B$ est complet et $B = A/K$ pour
un certain idéal fermé $K \subset A$.
\end{lemma}

\begin{proof}
En combinant le lemme de l'application ouverte (More on Algebra, Lemma
\ref{more-algebra-lemma-open-mapping})
et le caractère tendu de $\varphi$, 
on voit que l'application $\varphi$ est ouverte. Comme l'image de $\varphi$
est dense, $\varphi$ est surjectif. Le noyau $K$
de $\varphi$ est fermé, puisque $\varphi$ est continu.
Il s'ensuit que $B = A/K$ est complet ; voir par exemple
Lemma \ref{lemma-quotient-by-closed}.
\end{proof}







\section{Homomorphismes d'anneaux adiques}
\label{section-adic-ring-maps}

\noindent
Posons la définition suivante.

\begin{definition}
\label{definition-adic-homomorphism}
Soient $A$ et $B$ des anneaux topologiques préadiques. Un homomorphisme d'anneaux
$\varphi : A \to B$ est dit {\it adique}\footnote{Cette terminologie n'est peut-être pas standard.}
s'il existe un idéal de définition $I \subset A$ tel que
la topologie de $B$ soit la topologie $I$-adique.
\end{definition}

\noindent
Si $\varphi : A \to B$ est un homomorphisme adique d'anneaux préadiques, alors
$\varphi$ est continu et la topologie de $B$ est la topologie $I$-adique
pour tout idéal de définition $I$ de $A$.

\begin{lemma}
\label{lemma-composition-adic}
Soient $A \to B$ et $B \to C$ des homomorphismes continus d'anneaux
préadiques. Si $A \to B$ et $B \to C$ sont adiques, alors
$A \to C$ est adique.
\end{lemma}

\begin{proof}
Omis.
\end{proof}

\begin{lemma}
\label{lemma-permanence-adic}
Soient $A \to B$ et $B \to C$ des homomorphismes continus d'anneaux
préadiques. Si $A \to C$ est adique, alors
$B \to C$ est adique.
\end{lemma}

\begin{proof}
Choisissons un idéal de définition $I$ de $A$. Comme $A \to C$ est adique,
on voit que $IC$ est un idéal de définition de $C$.
Comme $B \to C$ est continu, on peut trouver un
idéal de définition $J \subset B$ dont l'image est contenue dans $IC$.
Comme $A \to B$ est continu, l'image réciproque $I' \subset I$
de $J$ dans $I$ est aussi un idéal de définition de $A$.
Ainsi, $I'C \subset JC \subset IC$ est compris entre
deux idéaux de définition et est donc lui-même un idéal de définition.
\end{proof}

\begin{lemma}
\label{lemma-adic-taut}
Soit $\varphi : A \to B$ un homomorphisme continu entre anneaux
topologiques préadiques. Si $\varphi$ est adique, alors $\varphi$ est tendu.
\end{lemma}

\begin{proof}
Cela résulte immédiatement des définitions.
\end{proof}

\noindent
Le lemme suivant affirme deux choses :
\begin{enumerate}
\item la propriété d'être adique monte le long des homomorphismes tendus
d'anneaux complets linéairement topologisés ; et
\item les propriétés « $\varphi$ est tendu » et « $\varphi$ est adique »
sont équivalentes pour les homomorphismes continus $\varphi : A \to B$ entre anneaux adiques
si $A$ possède un idéal de définition de type fini.
\end{enumerate}
En raison de (2), on peut dire que le « caractère tendu » généralise le « caractère adique »
aux homomorphismes continus entre anneaux linéairement topologisés quelconques.
Voir aussi Section \ref{section-adic}.

\begin{lemma}
\label{lemma-taut-is-adic}
Soit $\varphi : A \to B$ un homomorphisme continu d'anneaux linéairement topologisés.
Si $\varphi$ est tendu, si $A$ est préadique et possède un idéal de définition
de type fini, et si $B$ est complet, alors $B$ est adique et possède un idéal de définition
de type fini, et l'homomorphisme d'anneaux $\varphi$ est adique.
\end{lemma}

\begin{proof}
Choisissons un idéal de définition de type fini $I$ de $A$.
Soit $J_n$ l'adhérence de $\varphi(I^n)B$ dans $B$.
Comme $B$ est complet, on a $B = \lim B/J_n$.
Soit $B' = \lim B/I^nB$ le complété $I$-adique de $B$.
D'après Algebra, Lemma \ref{algebra-lemma-hathat-finitely-generated},
la topologie $I$-adique sur $B'$ est complète, et
$B'/I^nB' = B/I^nB$. Ainsi, l'homomorphisme d'anneaux $B' \to B$ est continu
et d'image dense, puisque $B' \to B/I^nB \to B/J_n$ est surjectif
pour tout $n$. Enfin, l'homomorphisme $B' \to B$ est tendu,
car $(I^nB')B = I^nB$ et $A \to B$ est tendu.
D'après Lemma \ref{lemma-dense-image-surjective}, $B' \to B$ est ouvert
et surjectif. La topologie de $B$ est donc la topologie $I$-adique,
ce qui achève la démonstration.
\end{proof}








\section{Anneaux faiblement adiques}
\label{section-weakly-adic}

\noindent
Nous conseillons au lecteur de sauter cette section. La notion suivante est une généralisation
naturelle des anneaux adiques.

\begin{definition}
\label{definition-weakly-adic}
\begin{reference}
\cite[Definition 8.3.8]{Gabber-Ramero}
\end{reference}
Soit $A$ un anneau linéairement topologisé.
\begin{enumerate}
\item On dit que $A$ est {\it faiblement préadique}\footnote{Dans \cite{Gabber-Ramero}, les
auteurs disent que $A$ est {\it $c$-adique}.} s'il existe un idéal
$I \subset A$ tel que l'adhérence de $I^n$ soit ouverte pour tout $n \geq 0$
et que ces adhérences forment un système fondamental d'idéaux ouverts.
\item On dit que $A$ est {\it faiblement adique} si $A$ est faiblement préadique
et complet\footnote{Selon nos conventions, cela inclut la séparation.}.
\end{enumerate}
\end{definition}

\noindent
Pour les anneaux complets linéairement topologisés, on a les implications suivantes :
$$
\xymatrix{
\text{adique + noethérien} \ar@{=>}[d] \\
\text{adique + idéal de définition de type fini} \ar@{=>}[d] \\
\text{adique} \ar@{=>}[d] \\
\text{faiblement adique} \ar@{=>}[d] \\
\text{admissible + à base dénombrable} \ar@{=>}[d] \ar@{=>}[r] &
\text{admissible} \ar@{=>}[d] \\
\text{faiblement admissible + à base dénombrable} \ar@{=>}[r] &
\text{faiblement admissible}
}
$$
où « à base dénombrable » signifie que notre anneau topologique possède un système fondamental
dénombrable d'idéaux ouverts. On dispose d'un diagramme d'implications analogue
pour les anneaux linéairement topologisés non complets (c'est-à-dire en utilisant les notions
de préadique, faiblement préadique, préadmissible et faiblement préadmissible).
Contrairement à ce qui se passe pour les anneaux préadiques, le complété d'un anneau
faiblement préadique est faiblement adique, comme le montre le lemme suivant
qui caractérise les anneaux faiblement préadiques.

\begin{lemma}
\label{lemma-weakly-adic}
Soit $A$ un anneau linéairement topologisé. Les conditions suivantes sont équivalentes :
\begin{enumerate}
\item $A$ est faiblement préadique ;
\item il existe un homomorphisme continu tendu d'anneaux $A' \to A$,
où $A'$ est un anneau topologique préadique ;
\item $A$ est préadmissible et il existe un idéal de définition $I$
tel que l'adhérence de $I^n$ soit ouverte pour tout $n \geq 1$ ; et
\item $A$ est préadmissible et, pour tout idéal de définition $I$,
l'adhérence de $I^n$ est ouverte pour tout $n \geq 1$.
\end{enumerate}
Le complété d'un anneau faiblement préadique est faiblement adique.
Si $A$ est faiblement adique, alors $A$ est admissible et possède un système fondamental
dénombrable d'idéaux ouverts.
\end{lemma}

\begin{proof}
Supposons (1). Choisissons un idéal $I$ tel que l'adhérence
de $I^n$ soit ouverte pour tout $n$ et que ces adhérences forment
un système fondamental d'idéaux ouverts. Notons $A' = A$ muni de
la topologie $I$-adique. Alors $A' \to A$ est tendu par définition,
et la condition (2) est vérifiée.

\medskip\noindent
Supposons (2). Soit $I' \subset A'$ un idéal de définition.
Notons $I$ l'adhérence de $I'A$. Le caractère tendu de $A' \to A$ signifie
que les adhérences $I_n$ des $(I')^nA$ sont ouvertes et forment un système fondamental
d'idéaux ouverts. Ainsi, $I = I_1$ est ouvert, et les adhérences de $I^n$
sont égales à $I_n$ ; elles sont donc ouvertes et forment un système fondamental
d'idéaux ouverts. En particulier, $I$ est un idéal de définition
tel que l'adhérence de $I^n$ soit ouverte pour tout $n$. La condition (3) est donc vérifiée.

\medskip\noindent
Si $I \subset A$ est comme dans (3), alors $I$ est un idéal comme dans
Definition \ref{definition-weakly-adic}, et la condition (1) est vérifiée.
De plus, si $I' \subset A$ est un autre idéal de définition quelconque, alors
$I'$ est ouvert (voir More on Algebra, Definition
\ref{more-algebra-definition-topological-ring})
et contient donc $I^n$ pour un certain $n \geq 1$.
Ainsi, $(I')^m$ contient $I^{nm}$ pour tout $m \geq 1$, et l'on
conclut que les adhérences de $(I')^m$ sont ouvertes pour tout $m$.
On voit ainsi que (3) implique (4).
L'implication (4) $\Rightarrow$ (3) est triviale.

\medskip\noindent
Soit $A$ faiblement préadique. Choisissons $A' \to A$ comme dans (2).
D'après Lemmas \ref{lemma-completion-taut} et
\ref{lemma-composition-taut},
le composé $A' \to A^\wedge$ est tendu. Ainsi, $A^\wedge$
est faiblement préadique par l'équivalence de (2) et (1).
Comme le complété d'un anneau linéairement topologisé $A$ est complet
(More on Algebra, Section \ref{more-algebra-section-topological-ring}),
on voit que $A^\wedge$ est faiblement adique.

\medskip\noindent
Soit $A$ faiblement adique. Alors $A$ est complet
et préadmissible par (1) $\Rightarrow$ (3) ; par conséquent, $A$
est admissible. Bien entendu, par définition, $A$ possède un système fondamental
dénombrable d'idéaux ouverts.
\end{proof}

\noindent
Nous donnons deux critères qui garantissent qu'un anneau faiblement adique
est adique et possède un idéal de définition de type fini.

\begin{lemma}
\label{lemma-weakly-adic-finite-generation-bis}
Soit $A$ un anneau complet linéairement topologisé. Soit $I \subset A$ un idéal
de type fini tel que l'adhérence de $I^n$ soit ouverte pour tout
$n \geq 0$ et que ces adhérences forment un système fondamental d'idéaux ouverts.
Alors $A$ est adique et possède un idéal de définition de type fini.
\end{lemma}

\begin{proof}
Notons $A'$ l'anneau $A$ muni de la topologie $I$-adique.
Les hypothèses nous disent que $A' \to A$ est tendu.
On conclut par Lemma \ref{lemma-taut-is-adic} (plus précisément, ce
lemme nous dit également que $I$ est un idéal de définition).
\end{proof}

\begin{lemma}
\label{lemma-weakly-adic-finite-generation}
Soit $A$ un anneau topologique faiblement adique. Soit $I$ un
idéal de définition tel que $I/I_2$ soit un module de type fini,
où $I_2$ est l'adhérence de $I^2$.
Alors $A$ est adique et possède un idéal de définition de type fini.
\end{lemma}

\begin{proof}
Nous utilisons la caractérisation de Lemma \ref{lemma-weakly-adic}
sans autre mention.
Choisissons $f_1, \ldots, f_r \in I$ dont les images engendrent $I/I_2$.
Posons $I' = (f_1, \ldots, f_r)$. On a $I' + I_2 = I$.
Alors $I_2$ est l'adhérence de $I^2 = (I' + I_2)^2 \subset I' + I_3$,
où $I_3$ est l'adhérence de $I^3$. Par conséquent, $I' + I_3 = I$.
En poursuivant ainsi, on voit que $I' + I_n = I$ pour tout $n \geq 2$,
où $I_n$ est l'adhérence de $I^n$.
Autrement dit, l'adhérence de $I'$ dans $A$ est $I$.
Par conséquent, l'adhérence de $(I')^n$ est $I_n$. Les adhérences de $(I')^n$
forment donc un système fondamental d'idéaux ouverts de $A$.
On conclut par Lemma \ref{lemma-weakly-adic-finite-generation-bis}.
\end{proof}

\noindent
Une propriété essentielle de la condition « faiblement préadique » est qu'elle se transmet au but
par les homomorphismes d'anneaux tendus entre anneaux linéairement topologisés.

\begin{lemma}
\label{lemma-taut-ascent-weakly-adic}
Soit $\varphi : A \to B$ un homomorphisme continu d'anneaux
linéairement topologisés. Si $\varphi$ est tendu et si $A$
est faiblement préadique, alors $B$ est faiblement préadique.
\end{lemma}

\begin{proof}
Soit $I \subset A$ un idéal tel que l'adhérence $I_n$ de $I^n$
soit ouverte et que ces adhérences définissent un système fondamental d'idéaux ouverts.
Alors l'adhérence de $I^nB$ est égale à celle de $I_nB$.
Comme $\varphi$ est tendu, ces adhérences sont ouvertes et forment un système fondamental
d'idéaux ouverts de $B$. Ainsi, $B$ est faiblement préadique.
\end{proof}

\begin{lemma}
\label{lemma-completed-tensor-product-weakly-adic}
Soient $B \to A$ et $B \to C$ des homomorphismes continus d'anneaux
linéairement topologisés. Si $A$ et $C$ sont faiblement préadiques, alors
$A \widehat{\otimes}_B C$ est faiblement adique.
\end{lemma}

\begin{proof}
Nous utiliserons la caractérisation de Lemma \ref{lemma-weakly-adic}
sans autre mention. D'après Lemma \ref{lemma-completed-tensor-product},
nous savons que $A \widehat{\otimes}_B C$ est admissible.
De plus, la démonstration de ce lemme montre que
l'adhérence $K \subset A \widehat{\otimes}_B C$
est un idéal de définition lorsque $I \subset A$ et $J \subset C$
le sont, cette adhérence étant celle de $I(A \widehat{\otimes}_B C) + J(A \widehat{\otimes}_B C)$.
Il suffit alors de montrer que
l'adhérence de $K^n$ est ouverte pour tout $n \geq 1$.
Puisque l'idéal $K^n$ contient
$I^n(A \widehat{\otimes}_B C) + J^n(A \widehat{\otimes}_B C)$,
puisque l'adhérence de $I^n$ dans $A$ est ouverte et
puisque l'adhérence de $J^n$ dans $C$ est ouverte, on
voit que l'adhérence de $K^n$ est ouverte dans $A \widehat{\otimes}_B C$.
\end{proof}






\section{Descente des propriétés}
\label{section-taut-descent}

\noindent
Dans cette section, nous considérons la situation suivante :
\begin{enumerate}
\item $\varphi : A \to B$ est un homomorphisme continu d'anneaux
topologiques linéairement topologisés ;
\item $\varphi$ est tendu ; et
\item pour tout idéal ouvert $I \subset A$, si $J \subset B$
désigne l'adhérence de $IB$, alors l'homomorphisme $A/I \to B/J$
est fidèlement plat.
\end{enumerate}
Nous allons montrer que, dans cette situation, les propriétés de $B$ sont héritées
par $A$.

\begin{lemma}
\label{lemma-taut-descent-countable}
Dans la situation ci-dessus, si $B$ possède un système fondamental dénombrable
d'idéaux ouverts, alors $A$ possède un système fondamental dénombrable d'idéaux ouverts.
\end{lemma}

\begin{proof}
Choisissons un système fondamental $B \supset J_1 \supset J_2 \supset \ldots$
d'idéaux ouverts. Puisque $\varphi$ est tendu, pour tout $n$, on peut trouver un
idéal ouvert $I_n$ tel que $J_n \supset I_nB$. Nous affirmons que les
$I_n$ forment un système fondamental d'idéaux ouverts de $A$.
En effet, supposons que $I \subset A$ soit ouvert. Comme $\varphi$ est tendu,
l'adhérence de $IB$ est ouverte et contient donc $J_n$ pour un certain $n$
assez grand. Ainsi, $I_nB \subset IB$. Soit $J$ l'adhérence de
$IB$ dans $B$. Puisque $A/I \to B/J$ est fidèlement plat, il est injectif.
Par conséquent, comme $I_n \to A/I \to B/J$ est nul puisque $I_nB \subset IB \subset J$,
on conclut que $I_n \to A/I$ est nul. Ainsi, $I_n \subset I$,
ce qui achève la démonstration.
\end{proof}

\begin{lemma}
\label{lemma-taut-descent-weakly-admissible}
Dans la situation ci-dessus, si $B$ est faiblement préadmissible, alors
$A$ est faiblement préadmissible.
\end{lemma}

\begin{proof}
Soit $J \subset B$ un idéal faible de définition. Soit $I \subset A$
un idéal ouvert tel que $IB \subset J$. Pour montrer que $I$ est un idéal faible
de définition, il faut montrer que tout $f \in I$ est topologiquement nilpotent.
Soit $I' \subset A$ un idéal ouvert. Notons $J' \subset B$ l'adhérence
de $I'B$. Alors $A/I' \to B/J'$ est fidèlement plat, donc injectif.
Ainsi, pour montrer que $f^n \in I'$, il suffit de montrer
que $\varphi(f)^n \in J'$. C'est le cas pour $n \gg 0$, puisque
$\varphi(f) \in J$, que l'idéal $J$ est un idéal faible de définition de $B$
et que $J'$ est ouvert dans $B$.
\end{proof}

\begin{lemma}
\label{lemma-taut-descent-admissible}
Dans la situation ci-dessus, si $B$ est préadmissible, alors $A$
est préadmissible.
\end{lemma}

\begin{proof}
Soit $J \subset B$ un idéal faible de définition. Soit $I \subset A$
un idéal ouvert tel que $IB \subset J$. Soit $I' \subset A$ un idéal ouvert.
Pour montrer que $I$ est un idéal de définition, il faut montrer que
$I^n \subset I'$ pour $n \gg 0$. Notons $J' \subset B$ l'adhérence
de $I'B$. Alors $A/I' \to B/J'$ est fidèlement plat, donc injectif.
Ainsi, pour montrer que $I^n \subset I'$, il suffit de montrer
que $\varphi(I)^n \subset J'$. C'est le cas pour $n \gg 0$, puisque
$\varphi(I) \subset J$, que l'idéal $J$ est un idéal de définition de $B$
et que $J'$ est ouvert dans $B$.
\end{proof}

\begin{lemma}
\label{lemma-taut-descent-weakly-adic}
Dans la situation ci-dessus, si $B$ est faiblement préadique, alors $A$
est faiblement préadique.
\end{lemma}

\begin{proof}
Nous utiliserons sans autre mention la caractérisation des anneaux faiblement préadiques donnée dans
Lemma \ref{lemma-weakly-adic}.
D'après Lemma \ref{lemma-taut-descent-admissible},
l'anneau topologique $A$ est préadmissible.
Soit $I \subset A$ un idéal de définition.
Fixons $n \geq 1$. Pour démontrer le lemme, il faut
montrer que l'adhérence de $I^n$ est ouverte.
Soient les $I_\lambda \subset A$ un système fondamental d'idéaux ouverts.
Notons $J \subset B$, resp.\ $J_\lambda \subset B$,
l'adhérence de $IB$, resp.\ celle de $I_\lambda B$.
Comme $B$ est faiblement préadique, l'adhérence de $J^n$ est ouverte.
Il existe donc $\lambda$ tel que
$$
J_\lambda \subset \bigcap\nolimits_\mu (J^n + J_\mu)
$$
car le membre de droite est l'adhérence de $J^n$ d'après
Lemma \ref{lemma-closed}.
Cela signifie que l'image de $J_\lambda$ dans $B/J_\mu$ est contenue dans
l'image de $J^n$ dans $B/J_\mu$. Observons que l'image de $J^n$ dans
$B/J_\mu$ est égale à celle de $I^nB$ dans $B/J_\mu$ (puisque tout élément
de $J$ est congru à un élément de $IB$ modulo $J_\mu$).
Comme $A/I_\mu \to B/J_\mu$ est fidèlement plat et que
$I_\lambda B \subset J_\lambda$,
on conclut que l'image de $I_\lambda$ dans $A/I_\mu$
est contenue dans l'image de $I^n$.
On en déduit que $I_\lambda$ est contenu dans l'adhérence
de $I^n$, ce qui achève la démonstration.
\end{proof}

\begin{lemma}
\label{lemma-taut-descent-adic-star}
Dans la situation ci-dessus, si $B$ est adique et possède un idéal de définition
de type fini, et si $A$ est complet, alors $A$ est adique et
possède un idéal de définition de type fini.
\end{lemma}

\begin{proof}
Nous savons déjà que $A$ est faiblement adique, et a fortiori admissible, d'après
Lemma \ref{lemma-taut-descent-weakly-adic}
(et Lemma \ref{lemma-weakly-adic} pour voir que les anneaux adiques sont faiblement adiques).
Soit $I \subset A$ un idéal de définition.
Soit $J \subset B$ un idéal de définition de type fini.
Comme l'adhérence de $IB$ est ouverte, on peut trouver $n > 0$ tel que
$J^n$ soit contenu dans l'adhérence de $IB$. Après avoir remplacé $J$
par $J^n$, on peut donc supposer que $J$ est un idéal de définition
de type fini contenu dans l'adhérence de $IB$. D'après Lemma \ref{lemma-closed},
cela implique en particulier que
$$
J \subset IB + J^2
$$
Considérons le $A$-module de type fini $M = (J + IB)/IB$.
L'équation affichée montre que $JM = M$.
D'après Lemma \ref{lemma-weakly-admissible-henselian}, par exemple,
on voit que $J$ est contenu dans le radical de Jacobson de $B$.
Le lemme de Nakayama, plus précisément la partie (2) de
Algebra, Lemma \ref{algebra-lemma-NAK}, donne donc $M = 0$.
Ainsi, $J \subset IB$.

\medskip\noindent
Comme $J$ est de type fini, on peut trouver un idéal de type fini
$I' \subset I$ tel que $J \subset I'B$.
Comme $A \to B$ est continu, que $J \subset B$ est ouvert et que
$I$ est un idéal de définition, on peut trouver $n > 0$ tel que
$I^nB \subset J$. Soit $J_{n + 1} \subset B$ l'adhérence de
$I^{n + 1}B$.
On a
$$
I^n \cdot (B/J_{n + 1}) \subset J \cdot (B/J_{n + 1}) \subset
I' \cdot (B/J_{n + 1})
$$
Puisque $A/I^{n + 1} \to B/J_{n + 1}$ est fidèlement plat, cela
implique $I^n \cdot (A/I^{n + 1}) \subset I' \cdot (A/I^{n + 1})$,
ce qui signifie à son tour que
$$
I^n \subset I' + I^{n + 1}
$$
Cela implique $I^n \subset I'  + I^{n + k}$ pour tout $k \geq 1$,
ce qui implique à son tour $I^{nm} \subset (I')^m + I^{nm + k}$
pour tous $k, m \geq 1$. Il s'ensuit que l'adhérence de
$(I')^m$ contient $I^{nm}$. Comme l'adhérence de $I^{nm}$ est ouverte,
puisque $A$ est faiblement adique,
on conclut que l'adhérence de $(I')^m$ est ouverte pour tout $m$.
Comme ces adhérences forment un système fondamental d'idéaux ouverts
de $A$ (puisqu'il en est ainsi des adhérences de $I^n$),
on conclut par Lemma \ref{lemma-weakly-adic-finite-generation-bis}.
\end{proof}









\section{Espaces algébriques formels affines}
\label{section-affine-formal-algebraic-spaces}

\noindent
Dans cette section, nous introduisons les espaces algébriques formels affines.
Ceux-ci sont en fait les mêmes objets que ceux appelés schémas
formels affines dans \cite{BVGD}. Nous les appellerons cependant
espaces algébriques formels affines, afin d'éviter toute confusion avec
la notion de schéma formel affine définie dans \cite{EGA}.

\medskip\noindent
Rappelons qu'un épaississement de schémas est une immersion
fermée qui induit une surjection sur les espaces topologiques
sous-jacents ; voir More on Morphisms, Definition
\ref{more-morphisms-definition-thickening}.

\begin{definition}
\label{definition-affine-formal-algebraic-space}
Soit $S$ un schéma. On dit qu'un faisceau $X$ sur $(\Sch/S)_{fppf}$ est un
{\it espace algébrique formel affine} s'il existe :
\begin{enumerate}
\item un ensemble ordonné filtrant $\Lambda$ ;
\item un système $(X_\lambda, f_{\lambda \mu})$ indexé par $\Lambda$
dans $(\Sch/S)_{fppf}$ tel que :
\begin{enumerate}
\item chaque $X_\lambda$ soit affine ;
\item chaque $f_{\lambda \mu} : X_\lambda \to X_\mu$ soit un épaississement.
\end{enumerate}
\end{enumerate}
et tels que
$$
X \cong \colim_{\lambda \in \Lambda} X_\lambda
$$
comme faisceaux fppf, et que $X$ satisfasse une condition ensembliste
(voir Remark \ref{remark-set-theoretic}). Un
{\it morphisme d'espaces algébriques formels affines}
sur $S$ est une application de faisceaux.
\end{definition}

\noindent
Observons que le système $(X_\lambda, f_{\lambda \mu})$ ne fait pas
partie des données. Supposons que $U$ soit un schéma quasi-compact sur $S$.
Comme les applications de transition sont des monomorphismes, on a
$$
X(U) = \colim X_\lambda(U)
$$
d'après Sites, Lemma \ref{sites-lemma-directed-colimits-sections}.
Ainsi, la construction du faisceau fppf inhérente à la limite inductive de la
définition se réduit à la construction du faisceau de Zariski, laquelle ne change
rien pour les schémas quasi-compacts.

\begin{lemma}
\label{lemma-diagonal-affine-formal-algebraic-space}
Soit $S$ un schéma. Si $X$ est un espace algébrique formel affine sur
$S$, alors le morphisme diagonal $\Delta : X \to X \times_S X$
est représentable et est une immersion fermée.
\end{lemma}

\begin{proof}
Supposons donnés $U \to X$ et $V \to X$, où $U, V$ sont des schémas sur $S$.
Montrons que $U \times_X V$ est représentable. Écrivons $X = \colim X_\lambda$
comme dans Definition \ref{definition-affine-formal-algebraic-space}.
La discussion précédente montre que, localement pour Zariski sur $U$ et $V$, les morphismes
se factorisent par un certain $X_\lambda$. Dans ce cas,
$U \times_X V = U \times_{X_\lambda} V$, qui est un schéma.
La diagonale est donc représentable ; voir
Spaces, Lemma \ref{spaces-lemma-representable-diagonal}.
Étant donné $(a, b) : W \to X \times_S X$, où $W$ est un schéma sur $S$,
considérons l'application $X \times_{\Delta, X \times_S X, (a, b)} W \to W$.
Comme précédemment, localement sur $W$, les morphismes $a$ et $b$ se factorisent par
le schéma affine $X_\lambda$ pour un certain $\lambda$ ; on obtient alors
le morphisme
$X_\lambda
\times_{\Delta_\lambda, X_\lambda \times_S X_\lambda, (a, b)} W \to W$.
Il est obtenu par changement de base à partir de
$\Delta_\lambda : X_\lambda \to X_\lambda \times_S X_\lambda$,
qui est une immersion fermée puisque $X_\lambda \to S$ est séparé
(car $X_\lambda$ est affine).
Ainsi, $X \to X \times_S X$ est une immersion fermée.
\end{proof}

\noindent
Un morphisme de schémas $X \to X'$ est un épaississement s'il est une immersion fermée
et induit une surjection sur les ensembles de points sous-jacents ; voir
(More on Morphisms, Definition
\ref{more-morphisms-definition-thickening}).
Par conséquent, la propriété d'être un épaississement est conservée par tout
changement de base et est locale pour la topologie fpqc sur le but ; voir
Spaces, Section \ref{spaces-section-lists}.
Ainsi, Spaces, Definition \ref{spaces-definition-relative-representable-property},
s'applique à la propriété « être un épaississement », et nous savons ce que signifie, pour une
transformation représentable $F \to G$ de
préfaisceaux sur $(\Sch/S)_{fppf}$, être un épaississement.
Observons que cela est compatible avec notre définition
(More on Morphisms of Spaces, Definition
\ref{spaces-more-morphisms-definition-thickening})
des épaississements lorsque $F$ et $G$ sont des espaces algébriques.

\begin{lemma}
\label{lemma-covering-by-thickenings}
Soient $X_\lambda, \lambda \in \Lambda$ et $X = \colim X_\lambda$
comme dans Definition \ref{definition-affine-formal-algebraic-space}.
Alors $X_\lambda \to X$ est représentable et est un épaississement.
\end{lemma}

\begin{proof}
L'énoncé a un sens d'après la discussion de
Spaces, Section \ref{spaces-section-representable} et
\ref{spaces-section-representable-properties}.
D'après Lemma \ref{lemma-diagonal-affine-formal-algebraic-space},
les morphismes $X_\lambda \to X$ sont représentables.
Étant donné $U \to X$, où $U$ est un schéma,
la discussion qui suit
Definition \ref{definition-affine-formal-algebraic-space}
montre que, localement pour Zariski sur $U$, le
morphisme se factorise par un certain $X_\mu$ avec $\lambda \leq \mu$.
Dans ce cas, $U \times_X X_\lambda = U \times_{X_\mu} X_\lambda$,
de sorte que $U \times_X X_\lambda \to U$ est obtenu par changement de base à partir de
l'épaississement $X_\lambda \to X_\mu$.
\end{proof}

\begin{lemma}
\label{lemma-factor-through-thickening}
Soient $X_\lambda, \lambda \in \Lambda$ et $X = \colim X_\lambda$
comme dans Definition \ref{definition-affine-formal-algebraic-space}.
Si $Y$ est un espace algébrique quasi-compact sur $S$, tout
morphisme $Y \to X$ se factorise par un $X_\lambda$.
\end{lemma}

\begin{proof}
Choisissons un schéma affine $V$ et un morphisme étale surjectif
$V \to Y$. Le composé $V \to Y \to X$ se factorise par
$X_\lambda$ pour un certain $\lambda$, d'après la discussion qui suit
Definition \ref{definition-affine-formal-algebraic-space}.
Comme $V \to Y$ est une surjection de faisceaux, on conclut.
\end{proof}

\begin{lemma}
\label{lemma-characterize-affine-formal-algebraic-space}
Soit $S$ un schéma. Soit $X$ un faisceau sur $(\Sch/S)_{fppf}$.
Alors $X$ est un espace algébrique formel affine si et seulement si
les conditions suivantes sont satisfaites :
\begin{enumerate}
\item tout morphisme $U \to X$, où $U$ est un schéma affine sur $S$,
se factorise par un morphisme $T \to X$ qui est représentable et est un
épaississement, avec $T$ un schéma affine sur $S$ ; et
\item une condition ensembliste comme dans Remark \ref{remark-set-theoretic}.
\end{enumerate}
\end{lemma}

\begin{proof}
Il résulte de Lemmas \ref{lemma-covering-by-thickenings} et
\ref{lemma-factor-through-thickening} qu'un espace algébrique formel affine
satisfait (1) et (2). Pour démontrer la réciproque, nous pouvons
supposer que $X$ n'est pas vide.
Soit $\Lambda$ la catégorie des morphismes représentables $T \to X$ qui sont des
épaississements, où $T$ est un schéma affine sur $S$. Cette catégorie
est filtrante. Puisque $X$ n'est pas vide, $\Lambda$ contient au moins un
objet. Si $T \to X$ et $T' \to X$ appartiennent à $\Lambda$, on peut
factoriser $T \amalg T' \to X$ par $T'' \to X$ dans $\Lambda$. Entre
deux objets quelconques de $\Lambda$, il existe une flèche unique ou aucune. Ainsi,
$\Lambda$ est un ensemble ordonné filtrant et, par hypothèse,
$X = \colim_{T \to X\text{ dans }\Lambda} T$. Pour achever la démonstration,
il faut montrer que toute flèche $T \to T'$ de $\Lambda$ est un épaississement.
C'est vrai parce que $T' \to X$ est un monomorphisme de faisceaux, de sorte que
$T = T \times_{T'} T' = T \times_X T'$ ; par conséquent, le morphisme
$T \to T'$ est égal à la projection $T \times_X T' \to T'$, qui est
un épaississement puisque $T \to X$ en est un.
\end{proof}

\noindent
Pour un espace algébrique formel affine général $X$, rien ne garantit
que $X$ possède assez de fonctions pour séparer les points, par exemple.
Voir Examples, Section \ref{examples-section-affine-formal-algebraic-space}.
Pour caractériser ceux qui en possèdent assez, nous proposons le lemme suivant.

\begin{lemma}
\label{lemma-mcquillan-affine-formal-algebraic-space}
Soit $S$ un schéma. Soit $X$ un faisceau fppf sur $(\Sch/S)_{fppf}$
qui satisfait la condition ensembliste de
Remark \ref{remark-set-theoretic}.
Les conditions suivantes sont équivalentes :
\begin{enumerate}
\item il existe un anneau topologique faiblement admissible $A$ sur $S$
(voir Remark \ref{remark-mcquillan}) tel que
$X = \colim_{I \subset A\text{ idéal faible de définition}} \Spec(A/I)$ ;
\item $X$ est un espace algébrique formel affine et
il existe une $S$-algèbre $A$ et une application $X \to \Spec(A)$
telles que, pour une immersion fermée $T \to X$ avec $T$ schéma affine,
le composé $T \to \Spec(A)$ soit une immersion fermée ;
\item $X$ est un espace algébrique formel affine et
il existe une $S$-algèbre $A$ et une application $X \to \Spec(A)$
telles que, pour une immersion fermée $T \to X$ avec $T$ schéma,
le composé $T \to \Spec(A)$ soit une immersion fermée ;
\item $X$ est un espace algébrique formel affine et, pour un choix de
$X = \colim X_\lambda$ comme dans
Definition \ref{definition-affine-formal-algebraic-space},
les projections $\lim \Gamma(X_\lambda, \mathcal{O}_{X_\lambda})
\to \Gamma(X_\lambda, \mathcal{O}_{X_\lambda})$ sont surjectives ;
\item $X$ est un espace algébrique formel affine et, pour tout choix
de $X = \colim X_\lambda$ comme dans
Definition \ref{definition-affine-formal-algebraic-space},
les projections $\lim \Gamma(X_\lambda, \mathcal{O}_{X_\lambda})
\to \Gamma(X_\lambda, \mathcal{O}_{X_\lambda})$ sont surjectives.
\end{enumerate}
De plus, l'anneau topologique faiblement admissible est
$A = \lim \Gamma(X_\lambda, \mathcal{O}_{X_\lambda})$,
muni de sa topologie limite, et les idéaux faibles de définition
classifient exactement les morphismes $T \to X$ qui sont représentables
et sont des épaississements.
\end{lemma}

\begin{proof}
Il est clair que (5) implique (4).

\medskip\noindent
Supposons (4) pour $X = \colim X_\lambda$ comme dans
Definition \ref{definition-affine-formal-algebraic-space}.
Posons $A = \lim \Gamma(X_\lambda, \mathcal{O}_{X_\lambda})$.
Soit $T \to X$ une immersion fermée, avec $T$ un schéma
(remarquons que $T \to X$ est représentable d'après
Lemma \ref{lemma-diagonal-affine-formal-algebraic-space}).
Puisque $X_\lambda \to X$ est un épaississement, il en va de même de
$X_\lambda \times_X T \to T$. D'autre part,
$X_\lambda \times_X T \to X_\lambda$ est une immersion fermée ;
$X_\lambda \times_X T$ est donc affine. Par conséquent, $T$ est affine
d'après Limits, Proposition \ref{limits-proposition-affine}.
Alors $T \to X$ se factorise par $X_\lambda$ pour un certain $\lambda$,
d'après Lemma \ref{lemma-factor-through-thickening}.
Ainsi, $A \to \Gamma(X_\lambda, \mathcal{O}) \to \Gamma(T, \mathcal{O})$
est surjectif. La condition (3) est donc vérifiée.

\medskip\noindent
Il est clair que (3) implique (2).

\medskip\noindent
Supposons (2) pour $A$ et $X \to \Spec(A)$. Écrivons $X = \colim X_\lambda$
comme dans Definition \ref{definition-affine-formal-algebraic-space}.
Alors $A_\lambda = \Gamma(X_\lambda, \mathcal{O})$ est un quotient
de $A$ par l'hypothèse (2). Ainsi, $A^\wedge = \lim A_\lambda$
est un anneau topologique complet ; voir la discussion de
More on Algebra, Section \ref{more-algebra-section-topological-ring}.
Les applications $A^\wedge \to A_\lambda$ sont surjectives, puisque $A \to A_\lambda$ l'est.
Nous affirmons que, pour tout $\lambda$, le noyau $I_\lambda \subset A^\wedge$ de
$A^\wedge \to A_\lambda$ est un idéal faible de définition.
En effet, il est ouvert par définition de la topologie limite.
Si $f \in I_\lambda$, alors, pour tout $\mu \in \Lambda$,
l'image de $f$ dans $A_\mu$ est nulle dans tous les corps résiduels
des points de $X_\mu$. C'est donc un élément nilpotent
de $A_\mu$. Par conséquent, une puissance $f^n \in I_\mu$. Ainsi, $f^n \to 0$
lorsque $n \to 0$. L'anneau $A^\wedge$ est donc faiblement admissible.
Enfin, supposons que $I \subset A^\wedge$ soit un idéal faible
de définition. Alors $I \subset A^\wedge$ est ouvert ; il existe donc
$\lambda$ tel que $I \supset I_\lambda$. On obtient ainsi un morphisme
$\Spec(A^\wedge/I) \to \Spec(A_\lambda) \to X$.
Il s'ensuit que $X = \colim \Spec(A^\wedge/I)$, où
la limite inductive porte maintenant sur tous les idéaux faibles de définition.
La condition (1) est donc vérifiée.

\medskip\noindent
Supposons (1). Dans ce cas, il est clair que $X$ est un espace algébrique
formel affine. Soit $X = \colim X_\lambda$ une présentation quelconque comme dans
Definition \ref{definition-affine-formal-algebraic-space}.
Pour tout $\lambda$, on peut trouver un idéal faible de définition
$I \subset A$ tel que $X_\lambda \to X$ se factorise par
$\Spec(A/I) \to X$ ; voir Lemma \ref{lemma-factor-through-thickening}.
Alors $X_\lambda = \Spec(A/I_\lambda)$ avec $I \subset I_\lambda$.
Réciproquement, pour tout idéal faible de définition $I \subset A$,
le morphisme $\Spec(A/I) \to X$ se factorise par $X_\lambda$
pour un certain $\lambda$, c'est-à-dire que $I_\lambda \subset I$.
Il s'ensuit que chaque $I_\lambda$ est un idéal faible de définition
et qu'ils forment une partie cofinale de l'ensemble des idéaux faibles
de définition. Ainsi, $A = \lim A/I = \lim A/I_\lambda$,
et l'on voit que (5) est vraie et, de plus, que
$A = \lim \Gamma(X_\lambda, \mathcal{O}_{X_\lambda})$.
\end{proof}

\noindent
Ce lemme permet de poser la définition suivante.

\begin{definition}
\label{definition-types-affine-formal-algebraic-space}
Soit $S$ un schéma. Soit $X$ un espace algébrique formel affine sur $S$.
On dit que $X$ est {\it de McQuillan} si $X$ satisfait les conditions équivalentes
de Lemma \ref{lemma-mcquillan-affine-formal-algebraic-space}. Soit $A$
l'anneau topologique faiblement admissible associé à $X$. On dit que :
\begin{enumerate}
\item $X$ est {\it classique} si $X$ est de McQuillan et si $A$ est admissible
(More on Algebra, Definition \ref{more-algebra-definition-topological-ring}) ;
\item $X$ est {\it faiblement adique} si $X$ est de McQuillan et si $A$ est faiblement adique
(Definition \ref{definition-weakly-adic}) ;
\item $X$ est {\it adique} si $X$ est de McQuillan et si $A$ est adique
(More on Algebra, Definition \ref{more-algebra-definition-topological-ring}) ;
\item $X$ est {\it adique*} si $X$ est de McQuillan, si $A$ est adique et si $A$
possède un idéal de définition de type fini ; et
\item $X$ est {\it noethérien} si $X$ est de McQuillan et si $A$ est
à la fois noethérien et adique.
\end{enumerate}
\end{definition}

\noindent
Dans \cite{Fujiwara-Kato}, les auteurs emploient la terminologie « de type idéal fini »
pour exprimer qu'un anneau topologique adique $A$ contient un idéal de définition
de type fini. Étant donné un espace algébrique formel affine $X$,
voici les implications entre les notions introduites dans la définition :
$$
\xymatrix{
X\text{ noethérien} \ar@{=>}[r] &
X\text{ adique*} \ar@{=>}[r] &
X\text{ adique} \ar@{=>}[lld] \\
X\text{ faiblement adique} \ar@{=>}[r] &
X\text{ classique} \ar@{=>}[r] &
X\text{ de McQuillan}
}
$$
Voir la discussion de Section \ref{section-weakly-adic} et, pour un énoncé précis,
Lemma \ref{lemma-implications-between-types}.

\begin{remark}
\label{remark-compare-with-affine-formal-schemes}
Les espaces algébriques formels affines classiques correspondent aux
schémas formels affines considérés dans EGA (\cite{EGA}). Pour l'expliquer,
supposons que notre schéma de base soit $\Spec(\mathbf{Z})$. Soit
$\mathfrak X = \text{Spf}(A)$ un schéma formel affine.
Soit $h_\mathfrak X$ son foncteur des points comme dans
Lemma \ref{lemma-fully-faithful}.
Alors $h_\mathfrak X = \colim h_{\Spec(A/I)}$, où la limite inductive
porte sur la famille des idéaux de définition de l'anneau topologique
admissible $A$. Cela résulte de
(\ref{equation-morphisms-affine-formal-schemes})
après évaluation sur les schémas affines, et il suffit de vérifier l'égalité
sur les schémas affines puisque les deux membres sont des faisceaux fppf ; voir
Lemma \ref{lemma-formal-scheme-sheaf-fppf}.
Ainsi, $h_\mathfrak X$ est un espace algébrique formel affine.
En fait, c'est un espace algébrique formel affine classique
d'après Definition \ref{definition-types-affine-formal-algebraic-space}.
Lemma \ref{lemma-fully-faithful} nous dit donc que
la catégorie des schémas formels affines est équivalente à celle
des espaces algébriques formels affines classiques.
\end{remark}

\noindent
Le lien avec les schémas formels affines étant établi,
il est naturel de poser la définition suivante.

\begin{definition}
\label{definition-affine-formal-spectrum}
Soit $S$ un schéma. Soit $A$ un anneau topologique faiblement admissible sur
$S$ ; voir Definition \ref{definition-weakly-admissible}\footnote{Voir
More on Algebra, Definition
\ref{more-algebra-definition-topological-ring}
pour le cas classique, et Remark \ref{remark-mcquillan}
pour une discussion des différences.}.
Le {\it spectre formel} de $A$ est l'espace algébrique formel affine
$$
\text{Spf}(A) = \colim \Spec(A/I)
$$
où la limite inductive est prise sur l'ensemble des idéaux faibles de définition de $A$
et dans la catégorie $\Sh((\Sch/S)_{fppf})$.
\end{definition}

\noindent
Un tel spectre formel est de McQuillan par construction et, réciproquement,
tout espace algébrique formel affine de McQuillan est isomorphe à un
spectre formel. Précisons que notre théorie contient des espaces algébriques
formels affines qui ne sont le spectre formel d'aucun
anneau topologique faiblement admissible.
Suivant \cite{Yasuda}, nous pourrions introduire les $S$-pro-anneaux
comme les pro-objets de la catégorie des $S$-algèbres ; voir Categories,
Remark \ref{categories-remark-pro-category}. Tout
espace algébrique formel affine sur $S$ serait alors le spectre formel d'un tel
$S$-pro-anneau. Nous ne le ferons pas et travaillerons directement avec les
espaces algébriques formels affines correspondants.

\medskip\noindent
La construction du spectre formel est fonctorielle. Pour l'expliquer,
soit $\varphi : B \to A$ un homomorphisme continu d'anneaux topologiques
faiblement admissibles sur $S$. Alors
$$
\text{Spf}(\varphi) : \text{Spf}(B) \to \text{Spf}(A)
$$
est l'unique morphisme d'espaces algébriques formels affines tel que
les diagrammes
$$
\xymatrix{
\Spec(B/J) \ar[d] \ar[r] & \Spec(A/I) \ar[d] \\
\text{Spf}(B) \ar[r] & \text{Spf}(A)
}
$$
soient commutatifs pour tous les idéaux faibles de définition $I \subset A$ et $J \subset B$
tels que $\varphi(I) \subset J$. Comme la continuité de $\varphi$
implique que, pour tout idéal faible de définition $J \subset B$,
il existe un idéal faible de définition $I \subset A$ ayant la propriété
requise, on voit que les commutativités demandées déterminent
et définissent uniquement $\text{Spf}(\varphi)$.

\begin{lemma}
\label{lemma-morphism-between-formal-spectra}
Soit $S$ un schéma. Soient $A$, $B$ des anneaux topologiques faiblement admissibles
sur $S$. Tout morphisme $f : \text{Spf}(B) \to \text{Spf}(A)$
d'espaces algébriques formels affines sur $S$
est égal à $\text{Spf}(f^\sharp)$ pour un unique homomorphisme continu
de $S$-algèbres $f^\sharp : A \to B$.
\end{lemma}

\begin{proof}
Soit $f : \text{Spf}(B) \to \text{Spf}(A)$ comme dans le lemme.
Soit $J \subset B$ un idéal faible de définition. D'après
Lemma \ref{lemma-factor-through-thickening},
il existe un idéal faible de définition $I \subset A$ tel que
$\Spec(B/J) \to \text{Spf}(B) \to \text{Spf}(A)$
se factorise par $\Spec(A/I)$. D'après
Schemes, Lemma \ref{schemes-lemma-morphism-into-affine},
on obtient un homomorphisme de $S$-algèbres $A/I \to B/J$.
Ces homomorphismes sont compatibles lorsque $J$ varie et définissent
l'homomorphisme $f^\sharp : A \to B$. Celui-ci est continu, car,
pour tout idéal faible de définition $J \subset B$, il existe un
idéal faible de définition $I \subset A$ tel que
$f^\sharp(I) \subset J$. L'égalité $f = \text{Spf}(f^\sharp)$
résulte du choix des homomorphismes d'anneaux $A/I \to B/J$ qui constituent $f^\sharp$.
\end{proof}

\begin{lemma}
\label{lemma-presentation-representable}
Soit $S$ un schéma. Soit $f : X \to Y$ une application
de préfaisceaux sur $(\Sch/S)_{fppf}$. Si $X$ est un espace algébrique formel
affine et si $f$ est représentable par des espaces algébriques et localement quasi-fini,
alors $f$ est représentable (par des schémas).
\end{lemma}

\begin{proof}
Soient $T$ un schéma sur $S$ et $T \to Y$ une application. Il faut montrer que
l'espace algébrique $X \times_Y T$ est un schéma. Écrivons $X = \colim X_\lambda$
comme dans Definition
\ref{definition-affine-formal-algebraic-space}.
Soit $W \subset X \times_Y T$
un sous-espace ouvert quasi-compact. La restriction de la projection
$X \times_Y T \to X$ à $W$ se factorise par $X_\lambda$ pour un certain $\lambda$.
Alors
$$
W \to X_\lambda \times_S T
$$
est un monomorphisme (donc séparé) et est localement quasi-fini (car
$W \to X \times_Y T \to T$ est localement quasi-fini par l'hypothèse
sur $X \to Y$ ; voir Morphisms of Spaces, Lemma
\ref{spaces-morphisms-lemma-permanence-quasi-finite}).
Par conséquent, $W$ est un schéma d'après
Morphisms of Spaces, Proposition
\ref{spaces-morphisms-proposition-locally-quasi-finite-separated-over-scheme}.
Ainsi, $X \times_Y T$ est un schéma d'après
Properties of Spaces, Lemma \ref{spaces-properties-lemma-subscheme}.
\end{proof}






\section{Espaces algébriques formels affines à indexation dénombrable}
\label{section-countably-indexed}

\noindent
Ce sont les espaces algébriques formels affines du lemme suivant.

\begin{lemma}
\label{lemma-countable-affine-formal-algebraic-space}
Soient $S$ un schéma et $X$ un espace algébrique formel affine sur $S$.
Les assertions suivantes sont équivalentes :
\begin{enumerate}
\item il existe un système $X_1 \to X_2 \to X_3 \to \ldots$
d'épaississements de schémas affines sur $S$ tel que $X = \colim X_n$,
\item il existe une écriture $X = \colim X_\lambda$ comme dans la
Définition \ref{definition-affine-formal-algebraic-space}
telle que $\Lambda$ soit dénombrable.
\end{enumerate}
\end{lemma}

\begin{proof}
Ceci résulte de l'observation qu'un ensemble ordonné filtrant dénombrable
possède une partie cofinale isomorphe à $(\mathbf{N}, \geq)$.
Voir la démonstration de Algebra, Lemma \ref{algebra-lemma-ML-limit-nonempty}.
\end{proof}

\begin{definition}
\label{definition-countable}
Soient $S$ un schéma et $X$ un espace algébrique formel affine sur $S$.
On dit que $X$ est {\it à indexation dénombrable} si les conditions équivalentes du
Lemme \ref{lemma-countable-affine-formal-algebraic-space} sont satisfaites.
\end{definition}

\noindent
Dans la terminologie de \cite{BVGD}, cela s'exprime en disant que
$X$ est un $\aleph_0$-ind-schéma.

\begin{lemma}
\label{lemma-implications-between-types}
Soit $X$ un espace algébrique formel affine sur un schéma $S$.
\begin{enumerate}
\item Si $X$ est noethérien, alors $X$ est adique*.
\item Si $X$ est adique*, alors $X$ est adique.
\item Si $X$ est adique, alors $X$ est faiblement adique.
\item Si $X$ est faiblement adique, alors $X$ est classique.
\item Si $X$ est faiblement adique, alors $X$ est à indexation dénombrable.
\item Si $X$ est à indexation dénombrable, alors $X$ est de McQuillan.
\end{enumerate}
\end{lemma}

\begin{proof}
Les assertions (1), (2), (3) et (4) résultent de l'écriture $X = \text{Spf}(A)$,
où $A$ est un anneau faiblement admissible (donc complet) linéairement
topologisé, et des implications entre les différents
types de tels anneaux discutées dans la Section \ref{section-weakly-adic}.

\medskip\noindent
Démonstration de (5). Par définition, il existe un anneau topologique faiblement adique $A$
tel que $X = \colim \Spec(A/I)$, où la colimite porte sur les idéaux
de définition de $A$. Comme $A$ est faiblement adique, il existe en particulier
un système fondamental dénombrable d'idéaux ouverts $I_\lambda$, voir la
Définition \ref{definition-weakly-adic}.
Alors $X = \colim \Spec(A/I_n)$ par définition de $\text{Spf}(A)$.
Ainsi, $X$ est à indexation dénombrable.

\medskip\noindent
Démonstration de (6). Écrivons $X = \colim X_n$ pour un système
$X_1 \to X_2 \to X_3 \to \ldots$ d'épaississements de schémas
affines sur $S$. Alors
$$
A = \lim \Gamma(X_n, \mathcal{O}_{X_n})
$$
l'homomorphisme canonique vers chacun des $\Gamma(X_n, \mathcal{O}_{X_n})$ est surjectif, car les applications de transition
sont surjectives puisque les morphismes $X_n \to X_{n + 1}$ sont des immersions
fermées. Par conséquent, $X$ est de McQuillan.
\end{proof}

\begin{lemma}
\label{lemma-countably-indexed}
Soient $S$ un schéma et $X$ un préfaisceau sur $(\Sch/S)_{fppf}$.
Les assertions suivantes sont équivalentes :
\begin{enumerate}
\item $X$ est un espace algébrique formel affine à indexation dénombrable,
\item $X = \text{Spf}(A)$, où $A$ est une $S$-algèbre topologique faiblement
admissible qui possède un système fondamental dénombrable de voisinages de $0$,
\item $X = \text{Spf}(A)$, où $A$ est une $S$-algèbre topologique faiblement
admissible qui possède un système fondamental
$A \supset I_1 \supset I_2 \supset I_3 \supset \ldots$
d'idéaux faibles de définition,
\item $X = \text{Spf}(A)$, où $A$ est une $S$-algèbre topologique complète
ayant pour système fondamental de voisinages ouverts de $0$ une
suite dénombrable $A \supset I_1 \supset I_2 \supset I_3 \supset \ldots$
d'idéaux tels que $I_n/I_{n + 1}$ soit localement nilpotent, et
\item $X = \text{Spf}(A)$, où $A = \lim B/J_n$, muni de la topologie limite,
et où $B \supset J_1 \supset J_2 \supset J_3 \supset \ldots$ est une
suite d'idéaux d'une $S$-algèbre $B$ telle que $J_n/J_{n + 1}$ soit
localement nilpotent.
\end{enumerate}
\end{lemma}

\begin{proof}
Supposons (1). D'après le Lemme \ref{lemma-implications-between-types},
on peut écrire $X = \text{Spf}(A)$, où $A$ est une $S$-algèbre topologique
faiblement admissible. Pour toute présentation $X = \colim X_n$ comme dans la
partie (1) du Lemme \ref{lemma-countable-affine-formal-algebraic-space},
on voit que $A = \lim A_n$, avec $X_n = \Spec(A_n)$ et
$A_n = A/I_n$ pour un idéal faible de définition $I_n \subset A$.
Ceci résulte de la dernière assertion du
Lemme \ref{lemma-mcquillan-affine-formal-algebraic-space},
qui implique en outre que $\{I_n\}$ est un système fondamental
de voisinages ouverts de $0$. Nous avons donc une suite
$$
A \supset I_1 \supset I_2 \supset I_3 \supset \ldots
$$
d'idéaux faibles de définition telle que $A = \lim A/I_n$. On voit ainsi
que la condition (1) implique chacune des conditions (2) -- (5).

\medskip\noindent
Supposons (5). Notons d'abord que la topologie limite sur
$A = \lim B/J_n$ est complète et linéairement topologisée, voir
More on Algebra, Section \ref{more-algebra-section-topological-ring}.
Si $f \in A$ s'envoie sur zéro dans $B/J_1$, une certaine puissance s'envoie sur zéro
dans $B/J_2$, puisque son image dans $J_1/J_2$ est nilpotente, puis une nouvelle
puissance s'envoie sur zéro dans $J_2/J_3$, et ainsi de suite. On voit ainsi que
l'idéal ouvert $\Ker(A \to B/J_1)$ est un idéal faible de définition.
Ainsi, $A$ est faiblement admissible. On voit donc que (5) implique (2).

\medskip\noindent
Il est clair que (4) est un cas particulier de (5), en prenant $B = A$.
Il est clair que (3) est un cas particulier de (2).

\medskip\noindent
Supposons que $A$ soit comme en (2). Soit $E_n$ un système fondamental dénombrable
de voisinages de $0$ dans $A$. Comme $A$ est un anneau topologique faiblement
admissible, on peut trouver des idéaux ouverts $I_n \subset E_n$.
On peut aussi choisir un idéal faible de définition $J \subset A$.
Les $J \cap I_n$ forment alors un système fondamental d'idéaux faibles de définition
de $A$, et l'on obtient
$X = \text{Spf}(A) = \colim \Spec(A/(J \cap I_n))$,
ce qui montre que $X$ est un espace algébrique formel affine à indexation dénombrable.
\end{proof}

\begin{lemma}
\label{lemma-characterize-noetherian-affine}
Soient $S$ un schéma et $X$ un espace algébrique formel affine.
Les assertions suivantes sont équivalentes :
\begin{enumerate}
\item $X$ est noethérien,
\item $X$ est adique* et, pour toute immersion fermée $T \to X$ où $T$ est un schéma,
$T$ est noethérien,
\item $X$ est adique* et, pour un certain choix de $X = \colim X_\lambda$ comme dans la
Définition \ref{definition-affine-formal-algebraic-space},
les schémas $X_\lambda$ sont noethériens, et
\item $X$ est faiblement adique et, pour un certain choix de $X = \colim X_\lambda$
comme dans la Définition \ref{definition-affine-formal-algebraic-space},
les schémas $X_\lambda$ sont noethériens.
\end{enumerate}
\end{lemma}

\begin{proof}
Supposons $X$ noethérien. Alors $X = \text{Spf}(A)$, où
$A$ est un anneau adique noethérien.
Soit $T \to X$ une immersion fermée, où $T$ est un schéma.
D'après le Lemme \ref{lemma-mcquillan-affine-formal-algebraic-space},
on voit que $T$ est affine et que $T \to \Spec(A)$ est une immersion
fermée. Comme $A$ est noethérien, on voit que $T$ est noethérien.
On voit ainsi que (1) $\Rightarrow$ (2).

\medskip\noindent
Les implications (2) $\Rightarrow$ (3) et (2) $\Rightarrow$ (4) sont
immédiates (voir le Lemme \ref{lemma-implications-between-types}).

\medskip\noindent
Pour démontrer (3) $\Rightarrow$ (1), écrivons $X = \text{Spf}(A)$
pour un anneau adique $A$ muni d'un idéal de définition $I$ de type fini.
On sait également que les anneaux $A/I_\lambda$ sont noethériens
pour un certain système fondamental d'idéaux ouverts $I_\lambda$.
Comme $I$ est ouvert, on peut trouver un $\lambda$ tel que $I_\lambda \subset I$.
Alors $A/I$ est noethérien, et l'on conclut que $A$ est noethérien par
Algebra, Lemma \ref{algebra-lemma-completion-Noetherian}.

\medskip\noindent
Pour démontrer (4) $\Rightarrow$ (3), écrivons $X = \text{Spf}(A)$
pour un anneau faiblement adique $A$. Alors $A$ est admissible et
possède un idéal de définition $I$, tandis que l'adhérence $I_2$ de $I^2$
est ouverte, voir le Lemme \ref{lemma-weakly-adic}.
On sait également que les anneaux $A/I_\lambda$ sont noethériens
pour un certain système fondamental d'idéaux ouverts $I_\lambda$.
Choisissons un $\lambda$ tel que $I_\lambda \subset I_2$.
Alors $A/I_2$ est noethérien en tant que quotient de $A/I_\lambda$.
Par conséquent, $I/I_2$ est un $A$-module de type fini.
Ainsi, $A$ est un anneau adique possédant un idéal de définition
de type fini par le Lemme \ref{lemma-weakly-adic-finite-generation}.
Donc $X$ est adique* et (3) est satisfaite.
\end{proof}





\section{Espaces algébriques formels}
\label{section-formal-algebraic-spaces}

\noindent
Nous dérogeons ici à notre habitude d'introduire les notions nouvelles d'abord
pour les anneaux, puis pour les schémas et enfin pour les espaces algébriques, en
introduisant les espaces algébriques formels sans avoir d'abord introduit les
schémas formels. L'idée générale sera qu'un espace algébrique formel
est un faisceau pour la topologie fppf qui, localement pour la topologie étale, est un
schéma formel affine au sens de \cite{BVGD}.
On trouvera des résultats connexes dans \cite{Yasuda}.

\medskip\noindent
Dans la définition d'un espace algébrique formel, nous allons
emprunter de la terminologie à
Bootstrap, Sections
\ref{bootstrap-section-morphism-representable-by-spaces} et
\ref{bootstrap-section-representable-by-spaces-properties}.

\begin{definition}
\label{definition-formal-algebraic-space}
Soit $S$ un schéma. On dit qu'un faisceau $X$ sur $(\Sch/S)_{fppf}$ est un
{\it espace algébrique formel} s'il existe une famille de morphismes
$\{X_i \to X\}_{i \in I}$ de faisceaux telle que
\begin{enumerate}
\item $X_i$ soit un espace algébrique formel affine,
\item $X_i \to X$ soit représentable par des espaces algébriques et étale,
\item $\coprod X_i \to X$ soit surjectif en tant que morphisme de faisceaux,
\end{enumerate}
et que $X$ satisfasse une condition ensembliste
(voir la Remarque \ref{remark-set-theoretic}). Un
{\it morphisme d'espaces algébriques formels}
sur $S$ est un morphisme de faisceaux.
\end{definition}

\noindent
Discussion. Contrôle de cohérence : un espace algébrique formel affine est
un espace algébrique formel. Dans la situation de la définition,
les morphismes $X_i \to X$ sont représentables (par des schémas), voir le
Lemme \ref{lemma-presentation-representable}.
D'après Bootstrap, Lemma
\ref{bootstrap-lemma-surjective-flat-locally-finite-presentation},
au lieu de demander que $\coprod X_i \to X$ soit
surjectif en tant que morphisme de faisceaux, on pourrait demander qu'il soit
surjectif (ce qui a un sens puisqu'il est représentable).

\medskip\noindent
Notre notion d'espace algébrique formel est {\bf très générale}.
En fait, même les espaces algébriques formels affines définis ci-dessus
sont des objets très pathologiques.

\begin{lemma}
\label{lemma-diagonal-formal-algebraic-space}
Soit $S$ un schéma. Si $X$ est un espace algébrique formel sur
$S$, alors le morphisme diagonal $\Delta : X \to X \times_S X$
est représentable, est un monomorphisme, est localement quasi-fini,
est localement de type fini et est séparé.
\end{lemma}

\begin{proof}
Supposons donnés $U \to X$ et $V \to X$, avec $U, V$ des schémas sur $S$.
Alors $U \times_X V$ est un faisceau. Choisissons $\{X_i \to X\}$ comme dans la
Définition \ref{definition-formal-algebraic-space}.
Pour tout $i$, le morphisme
$$
(U \times_X X_i) \times_{X_i} (V \times_X X_i)
= (U \times_X V) \times_X X_i \to U \times_X V
$$
est représentable et étale, car c'est un changement de base de $X_i \to X$,
et sa source est un schéma (utiliser les
Lemmes \ref{lemma-diagonal-affine-formal-algebraic-space} et
\ref{lemma-presentation-representable}). Ces morphismes forment une famille surjective ;
ainsi, $U \times_X V$ est un espace algébrique par
Bootstrap, Theorem \ref{bootstrap-theorem-final-bootstrap}.
Le morphisme $U \times_X V \to U \times_S V$ est un monomorphisme.
Il est également localement quasi-fini, car, en le précomposant par
le morphisme affiché ci-dessus, on obtient la composée
$$
(U \times_X X_i) \times_{X_i} (V \times_X X_i)
\to (U \times_X X_i) \times_S (V \times_X X_i)
\to U \times_S V
$$
qui est localement quasi-finie en tant que composée d'une immersion
fermée (Lemme \ref{lemma-diagonal-affine-formal-algebraic-space})
et d'un morphisme étale, voir
Descent on Spaces, Lemma
\ref{spaces-descent-lemma-locally-quasi-finite-etale-local-source}.
On conclut donc que $U \times_X V$ est un schéma par
Morphisms of Spaces, Proposition
\ref{spaces-morphisms-proposition-locally-quasi-finite-separated-over-scheme}.
Ainsi, $\Delta$ est représentable, voir
Spaces, Lemma \ref{spaces-lemma-representable-diagonal}.

\medskip\noindent
En fait, puisque nous avons montré ci-dessus que les morphismes de schémas
$U \times_X V \to U \times_S V$ sont toujours des monomorphismes et sont
localement quasi-finis, on conclut que $\Delta : X \to X \times_S X$
est un monomorphisme et est localement quasi-fini, voir
Spaces, Lemma \ref{spaces-lemma-transformation-diagonal-properties}.
On peut ensuite appliquer le principe de
Spaces, Lemma
\ref{spaces-lemma-representable-transformations-property-implication}
pour voir que $\Delta$ est séparé et localement de type fini.
En effet, un monomorphisme de schémas est séparé
(Schemes, Lemma \ref{schemes-lemma-monomorphism-separated})
et un morphisme de schémas localement quasi-fini est
localement de type fini
(ceci résulte de la définition donnée dans
Morphisms, Section \ref{morphisms-section-quasi-finite}).
\end{proof}

\begin{lemma}
\label{lemma-space-to-formal-space}
Soit $S$ un schéma. Soit $f : X \to Y$ un morphisme d'un
espace algébrique sur $S$ vers un espace algébrique formel sur $S$.
Alors $f$ est représentable par des espaces algébriques.
\end{lemma}

\begin{proof}
Soit $Z \to Y$ un morphisme, où $Z$ est un schéma sur $S$.
Il faut montrer que $X \times_Y Z$ est un espace algébrique.
Choisissons un schéma $U$ et un morphisme étale surjectif $U \to X$.
Alors $U \times_Y Z \to X \times_Y Z$ est représentable, surjectif et étale
(Spaces, Lemma
\ref{spaces-lemma-base-change-representable-transformations-property}),
et $U \times_Y Z$ est un schéma d'après le
Lemme \ref{lemma-diagonal-formal-algebraic-space}.
Le résultat en découle par
Bootstrap, Theorem \ref{bootstrap-theorem-final-bootstrap}.
\end{proof}

\begin{remark}
\label{remark-compare-with-formal-schemes}
À des questions ensemblistes près, la catégorie des schémas formels au sens d'EGA
(voir la Section \ref{section-formal-schemes-EGA}) est équivalente à une sous-catégorie
pleine de la catégorie des espaces algébriques formels. Pour l'expliquer,
supposons que notre schéma de base soit $\Spec(\mathbf{Z})$. D'après le
Lemme \ref{lemma-formal-scheme-sheaf-fppf}, le foncteur des points
$h_\mathfrak X$ associé à un schéma formel $\mathfrak X$ est un faisceau
pour la topologie fppf. D'après le Lemme \ref{lemma-fully-faithful},
le foncteur $\mathfrak X \mapsto h_\mathfrak X$ est un plongement pleinement fidèle
de la catégorie des schémas formels dans la catégorie des
faisceaux fppf. Étant donné un schéma formel $\mathfrak X$, choisissons un recouvrement ouvert
$\mathfrak X = \bigcup \mathfrak X_i$, où les $\mathfrak X_i$ sont des
schémas formels affines. Alors $h_{\mathfrak X_i}$
est un espace algébrique formel affine d'après la
Remarque \ref{remark-compare-with-affine-formal-schemes}.
Les morphismes $h_{\mathfrak X_i} \to h_\mathfrak X$ sont représentables
et sont des immersions ouvertes. Ainsi, $\{h_{\mathfrak X_i} \to h_\mathfrak X\}$
est une famille comme dans la Définition \ref{definition-formal-algebraic-space},
et l'on voit que $h_\mathfrak X$ est un espace algébrique formel.
\end{remark}

\begin{remark}
\label{remark-set-theoretic}
Soient $S$ un schéma et $(\Sch/S)_{fppf}$ un grand site fppf comme dans
Topologies, Definition \ref{topologies-definition-big-small-fppf}.
Comme condition ensembliste sur $X$ dans les
Définitions \ref{definition-affine-formal-algebraic-space} et
\ref{definition-formal-algebraic-space}, nous prenons la suivante :
il existe des objets $U, R$ de $(\Sch/S)_{fppf}$, un
morphisme $U \to X$ qui soit une surjection de faisceaux fppf, et
un morphisme $R \to U \times_X U$ qui soit une surjection de faisceaux fppf.
Autrement dit, nous demandons que notre faisceau soit le conoyau de
deux morphismes entre faisceaux représentables.
Voici quelques observations qui impliquent que cette notion se comporte
de façon raisonnable :
\begin{enumerate}
\item Supposons $X = \colim_{\lambda \in \Lambda} X_\lambda$
et que le système satisfasse les conditions (1) et (2) de la
Définition \ref{definition-affine-formal-algebraic-space}. Alors
$U = \coprod_{\lambda \in \Lambda} X_\lambda \to X$ est une surjection
de faisceaux fppf. De plus, $U \times_X U$ est un sous-schéma fermé
de $U \times_S U$ d'après le Lemme \ref{lemma-diagonal-affine-formal-algebraic-space}.
Ainsi, si $U$ est représentable par un objet de $(\Sch/S)_{fppf}$,
alors $U \times_S U$ l'est aussi (voir Sets, Lemma \ref{sets-lemma-what-is-in-it}),
et la condition ensembliste est satisfaite. C'est toujours le cas
si $\Lambda$ est dénombrable, voir Sets, Lemma \ref{sets-lemma-what-is-in-it}.
\item Contrôle de cohérence. Soit $\{X_i \to X\}_{i \in I}$ comme dans la
Définition \ref{definition-formal-algebraic-space}
(avec la condition ensembliste formulée ci-dessus),
et supposons que chaque $X_i$ soit en fait un schéma affine.
Alors $X$ est un espace algébrique. En effet, si l'on choisit un
grand site fppf $(\Sch'/S)_{fppf}$ plus grand tel que $U' = \coprod X_i$
et $R' = \coprod X_i \times_X X_j$ soient représentables par des objets
de ce site, alors $X' = U'/R'$ sera un objet de la catégorie
des espaces algébriques pour ce choix. Une application de
Spaces, Lemma \ref{spaces-lemma-fully-faithful} montre alors que
$X$ est un espace algébrique pour $(\Sch/S)_{fppf}$.
\item Soit $\{X_i \to X\}_{i \in I}$ une famille de morphismes de faisceaux
satisfaisant les conditions (1), (2), (3) de la
Définition \ref{definition-formal-algebraic-space}.
Pour chaque $i$, on peut choisir $U_i \in \Ob((\Sch/S)_{fppf})$
et $U_i \to X_i$, qui soit une surjection de faisceaux.
Ainsi, si $I$ n'est pas trop grand (s'il est par exemple dénombrable),
$U = \coprod U_i \to X$ est une surjection de faisceaux et
$U$ est représentable par un objet de $(\Sch/S)_{fppf}$.
Pour obtenir un $R \in \Ob((\Sch/S)_{fppf})$ muni d'une surjection vers $U \times_X U$,
il suffit de supposer que la diagonale $\Delta : X \to X \times_S X$ n'est pas
trop pathologique ; c'est par exemple toujours possible si la diagonale de $X$ est
quasi-compacte, c'est-à-dire si $X$ est quasi-séparé.
\end{enumerate}
\end{remark}







\section{La réduction}
\label{section-reduction}

\noindent
Tout espace algébrique formel possède un espace algébrique réduit
sous-jacent, comme le montre le lemme suivant.

\begin{lemma}
\label{lemma-reduction-formal-algebraic-space}
Soit $S$ un schéma. Soit $X$ un espace algébrique formel sur $S$.
Il existe un espace algébrique réduit $X_{red}$ et un morphisme représentable
$X_{red} \to X$ qui est un épaississement. Tout morphisme $U \to X$,
où $U$ est un espace algébrique réduit, se factorise de manière unique par $X_{red}$.
\end{lemma}

\begin{proof}
Supposons d'abord que $X$ soit un espace algébrique formel affine.
Écrivons $X = \colim X_\lambda$ comme dans la
Définition \ref{definition-affine-formal-algebraic-space}.
Puisque les morphismes de transition sont des épaississements, les schémas
affines $X_\lambda$ ont tous des réductions isomorphes $X_{red}$.
Le morphisme $X_{red} \to X$ est représentable et est un épaississement
par le Lemme \ref{lemma-covering-by-thickenings} et par le fait que
les composés d'épaississements sont des épaississements. Nous omettons la
vérification de la propriété
universelle (utiliser Schemes, Definition
\ref{schemes-definition-reduced-induced-scheme},
Schemes, Lemma \ref{schemes-lemma-map-into-reduction},
Properties of Spaces, Definition
\ref{spaces-properties-definition-reduced-induced-space}, et
Properties of Spaces, Lemma \ref{spaces-properties-lemma-map-into-reduction}).

\medskip\noindent
Soient $X$ et $\{X_i \to X\}_{i \in I}$ comme dans la
Définition \ref{definition-formal-algebraic-space}.
Pour tout $i$, soit $X_{i, red} \to X_i$ la réduction construite
ci-dessus. Pour $i, j \in I$, la projection
$X_{i, red} \times_X X_j \to X_{i, red}$ est un morphisme étale (par hypothèse)
de schémas (par le Lemme \ref{lemma-presentation-representable}).
Ainsi $X_{i, red} \times_X X_j$ est réduit (voir
Descent, Lemma \ref{descent-lemma-reduced-local-smooth}).
La projection $X_{i, red} \times_X X_j \to X_j$ se factorise donc
par $X_{j, red}$ en vertu de la propriété universelle. Nous en déduisons que
$$
R_{ij} = X_{i, red} \times_X X_j = X_{i, red} \times_X X_{j, red} =
X_i \times_X X_{j, red}
$$
car les morphismes $X_{i, red} \to X_i$ sont des injections de faisceaux.
Posons $U = \coprod X_{i, red}$, posons
$R = \coprod R_{ij}$, et notons $s, t : R \to U$ les deux
projections. En tant que faisceau, $R = U \times_X U$, et $s$ et $t$
sont étales. Alors $(t, s) : R \to U$ définit une relation d'équivalence
étale d'après les observations précédentes. Ainsi $X_{red} = U/R$ est un
espace algébrique par Spaces, Theorem \ref{spaces-theorem-presentation}.
Par construction, le diagramme
$$
\xymatrix{
\coprod X_{i, red} \ar[r] \ar[d] & \coprod X_i \ar[d] \\
X_{red} \ar[r] & X
}
$$
est cartésien. Puisque la flèche verticale de droite est étale surjective
et que la flèche horizontale supérieure est représentable et est un épaississement,
nous concluons que $X_{red} \to X$ est représentable par
Bootstrap, Lemma \ref{bootstrap-lemma-after-fppf-sep-lqf}
(pour vérifier les hypothèses du lemme, utiliser qu'un morphisme étale
surjectif est surjectif, plat et localement de présentation
finie, et que les épaississements sont séparés et localement quasi-finis).
Nous pouvons ensuite appliquer Spaces, Lemma
\ref{spaces-lemma-descent-representable-transformations-property}
pour conclure que $X_{red} \to X$ est un épaississement
(utiliser qu'être un épaississement équivaut à être
une immersion fermée surjective).

\medskip\noindent
Enfin, supposons que $U \to X$ soit un morphisme où
$U$ est un espace algébrique réduit sur $S$. Alors chaque $X_i \times_X U$
est étale sur $U$, donc réduit (d'après notre définition des
espaces algébriques réduits dans Properties of Spaces, Section
\ref{spaces-properties-section-types-properties}).
Alors $X_i \times_X U \to X_i$ se factorise par $X_{i, red}$.
Ainsi $U \to X$ se factorise par $X_{red}$, puisque
$\{X_i \times_X U \to U\}$ est un recouvrement étale.
\end{proof}

\begin{example}
\label{example-reduction-affine-formal-spectrum}
Soit $A$ un anneau topologique faiblement admissible. Dans ce cas, on a
$$
\text{Spf}(A)_{red} = \Spec(A/\mathfrak a)
$$
où $\mathfrak a \subset A$ est l'idéal des éléments topologiquement
nilpotents. En effet, $\mathfrak a$ est un idéal radical
(Lemme \ref{lemma-topologically-nilpotent})
qui est ouvert puisque $A$ est faiblement admissible.
\end{example}

\begin{lemma}
\label{lemma-reduction-smooth}
Soit $S$ un schéma. Soit $f : X \to Y$ un morphisme d'espaces
algébriques formels sur $S$ qui est représentable par des
espaces algébriques et lisse (par exemple étale).
Alors $X_{red} = X \times_Y Y_{red}$.
\end{lemma}

\begin{proof}
(Le cas étale résulte directement de la construction de
l'espace algébrique réduit sous-jacent dans la démonstration du
Lemme \ref{lemma-reduction-formal-algebraic-space}.)
Supposons $f$ lisse. Observons que $X \times_Y Y_{red} \to Y_{red}$
est un morphisme lisse d'espaces algébriques. Ainsi $X \times_Y Y_{red}$
est un espace algébrique réduit par Descent on Spaces, Lemma
\ref{spaces-descent-lemma-reduced-local-smooth}.
La propriété universelle de la réduction montre alors que le morphisme canonique
$X_{red} \to X \times_Y Y_{red}$ est un isomorphisme.
\end{proof}

\begin{lemma}
\label{lemma-reduction-surjective}
Soit $S$ un schéma. Soit $f : X \to Y$ un morphisme d'espaces
algébriques formels sur $S$ qui est représentable par des
espaces algébriques. Alors $f$ est surjectif au sens de
Bootstrap, Definition \ref{bootstrap-definition-property-transformation}
si et seulement si $f_{red} : X_{red} \to Y_{red}$ est un
morphisme surjectif d'espaces algébriques.
\end{lemma}

\begin{proof}
Omis.
\end{proof}






\section{Limites inductives d'espaces algébriques le long d'épaississements}
\label{section-global-colimits}

\noindent
Un type particulier d'espace algébrique formel est celui qui peut s'écrire
globalement comme une limite inductive cofiltrante d'espaces algébriques le long
d'épaississements, ainsi que dans le lemme suivant. Nous verrons plus loin
(à la Section \ref{section-quasi-compact-quasi-separated})
que tout espace algébrique formel quasi-compact et quasi-séparé
est une telle limite inductive globale.

\begin{lemma}
\label{lemma-colimit-is-formal}
Soit $S$ un schéma. Supposons donnés un ensemble ordonné filtrant
$\Lambda$ et un système d'espaces algébriques $(X_\lambda, f_{\lambda \mu})$
sur $\Lambda$, où chaque $f_{\lambda \mu} : X_\lambda \to X_\mu$ est un
épaississement. Alors $X = \colim_{\lambda \in \Lambda} X_\lambda$
est un espace algébrique formel sur $S$.
\end{lemma}

\begin{proof}
Puisque nous prenons la limite inductive dans la catégorie des faisceaux fppf,
nous voyons que $X$ est un faisceau. Choisissons et fixons $\lambda \in \Lambda$. Choisissons un
recouvrement étale $\{X_{i, \lambda} \to X_\lambda\}$ où $X_i$ est un schéma
affine sur $S$, voir Properties of Spaces, Lemma
\ref{spaces-properties-lemma-cover-by-union-affines}.
Pour chaque $\mu \geq \lambda$, il existe un diagramme cartésien
$$
\xymatrix{
X_{i, \lambda} \ar[r] \ar[d] & X_{i, \mu} \ar[d] \\
X_\lambda \ar[r] & X_\mu
}
$$
dont les flèches verticales sont étales, voir
More on Morphisms of Spaces, Theorem
\ref{spaces-more-morphisms-theorem-topological-invariance}
(on utilise aussi qu'un épaississement est une immersion fermée surjective qui
satisfait aux conditions du théorème). De plus, ces diagrammes sont
uniques à isomorphisme unique près, et donc
$X_{i, \mu} = X_\mu \times_{X_{\mu'}} X_{i, \mu'}$ pour
$\mu' \geq \mu$. Les morphismes $X_{i, \mu} \to X_{i, \mu'}$
sont des épaississements comme changements de base d'épaississements. Chaque $X_{i, \mu}$
est un schéma affine par Limits of Spaces, Proposition
\ref{spaces-limits-proposition-affine} et par le fait que
$X_{i, \lambda}$ est affine.
Posons $X_i = \colim_{\mu \geq \lambda} X_{i, \mu}$. Alors $X_i$ est
un espace algébrique formel affine. Le morphisme $X_i \to X$
est étale car, étant donnés un schéma affine $U$ et un morphisme $U \to X$,
celui-ci se factorise par $X_\mu$ pour un certain $\mu \geq \lambda$ (détails omis).
Nous voyons ainsi que $X$ est un espace algébrique formel.
\end{proof}

\noindent
Soit $S$ un schéma. Soit $X$ un espace algébrique formel sur $S$.
Comment démontrer ou vérifier que $X$ est une limite inductive globale comme dans le
Lemme \ref{lemma-colimit-is-formal} ? Pour ce faire, nous cherchons des morphismes
$i : Z \to X$ où $Z$ est un espace algébrique sur $S$ et où $i$ est
surjectif et une immersion fermée ; autrement dit, $i$ est un épaississement.
Cela a un sens puisque $i$ est représentable par des espaces algébriques
(Lemme \ref{lemma-space-to-formal-space}) et que nous pouvons appliquer
Bootstrap, Definition \ref{bootstrap-definition-property-transformation}
comme précédemment.

\begin{example}
\label{example-david-hansen}
Soit $(A, \mathfrak m, \kappa)$ un anneau de valuation qui est
complet pour la topologie $(\pi)$-adique pour un certain $\pi \in \mathfrak m$ non nul.
Supposons aussi que $\mathfrak m$ ne soit pas de type fini.
Un exemple est $A = \mathcal{O}_{\mathbf{C}_p}$ et $\pi = p$,
où $\mathcal{O}_{\mathbf{C}_p}$ est l'anneau des entiers
du corps des nombres complexes $p$-adiques $\mathbf{C}_p$
(qui est le complété de la clôture algébrique de
$\mathbf{Q}_p$). Un autre exemple est
$$
A =
\left\{
\sum\nolimits_{\alpha \in \mathbf{Q},\ \alpha \geq 0} a_\alpha t^\alpha
\middle|
\begin{matrix}
a_\alpha \in \kappa \text{ et, pour tout }n\text{, il n'y a qu'un} \\
\text{nombre fini de coefficients non nuls }a_\alpha
\text{ tels que }\alpha \leq n
\end{matrix}
\right\}
$$
et $\pi = t$. Alors $X = \text{Spf}(A)$ est un espace algébrique formel
affine, et $\Spec(\kappa) \to X$ est un épaississement qui correspond
à l'idéal faible de définition $\mathfrak m \subset A$,
lequel n'est toutefois pas un idéal de définition.
\end{example}

\begin{remark}[Idéaux faibles de définition]
\label{remark-weak-ideals-of-definition}
Soit $\mathfrak X$ un schéma formel au sens de McQuillan, voir
Remarque \ref{remark-mcquillan}. Un {\it idéal faible de définition}
de $\mathfrak X$ est un faisceau d'idéaux
$\mathcal{I} \subset \mathcal{O}_\mathfrak X$ tel que,
pour tout sous-schéma formel ouvert affine $\mathfrak U \subset \mathfrak X$,
l'idéal
$\mathcal{I}(\mathfrak U) \subset \mathcal{O}_\mathfrak X(\mathfrak U)$
soit un idéal faible de définition de l'anneau topologique faiblement admissible
$\mathcal{O}_\mathfrak X(\mathfrak U)$.
Il suffit de vérifier la condition sur les membres d'un recouvrement ouvert affine.
Il existe une correspondance biunivoque
$$
\{\text{idéaux faibles de définition de }\mathfrak X\}
\leftrightarrow
\{\text{épaississements }i : Z \to h_\mathfrak X\text{ comme ci-dessus}\}
$$
Cette correspondance associe à $\mathcal{I}$ le schéma
$Z = (\mathfrak X, \mathcal{O}_\mathfrak X/\mathcal{I})$
muni du morphisme évident vers $\mathfrak X$.
Un {\it système fondamental d'idéaux faibles de définition}
est une famille d'idéaux faibles de définition
$\mathcal{I}_\lambda$ telle que, sur tout sous-schéma formel ouvert
affine $\mathfrak U \subset \mathfrak X$, les
idéaux
$$
I_\lambda = \mathcal{I}_\lambda(\mathfrak U) \subset
A = \Gamma(\mathfrak U, \mathcal{O}_\mathfrak X)
$$
forment un système fondamental d'idéaux faibles de définition de
l'anneau topologique faiblement admissible $A$. Il suffit de le vérifier
sur les membres d'un recouvrement ouvert affine. Nous concluons que
l'espace algébrique formel $h_\mathfrak X$ associé au
schéma formel de McQuillan $\mathfrak X$ est une limite inductive de schémas comme
dans le Lemme \ref{lemma-colimit-is-formal} si et seulement s'il
existe un système fondamental d'idéaux faibles de définition
de $\mathfrak X$.
\end{remark}

\begin{remark}[Idéaux de définition]
\label{remark-ideals-of-definition}
Soit $\mathfrak X$ un schéma formel \`a la EGA.
Un {\it idéal de définition} de $\mathfrak X$ est un faisceau d'idéaux
$\mathcal{I} \subset \mathcal{O}_\mathfrak X$ tel que,
pour tout sous-schéma formel ouvert affine $\mathfrak U \subset \mathfrak X$,
l'idéal
$\mathcal{I}(\mathfrak U) \subset \mathcal{O}_\mathfrak X(\mathfrak U)$
soit un idéal de définition de l'anneau topologique admissible
$\mathcal{O}_\mathfrak X(\mathfrak U)$.
Il suffit de vérifier la condition sur les membres d'un recouvrement ouvert affine.
Nous n'obtenons {\bf pas} la même correspondance entre idéaux de définition
et épaississements $Z \to h_\mathfrak X$ que dans la
Remarque \ref{remark-weak-ideals-of-definition} ; un exemple
est donné dans l'Exemple \ref{example-david-hansen}.
Un {\it système fondamental d'idéaux de définition}
est une famille d'idéaux de définition
$\mathcal{I}_\lambda$ telle que, sur tout sous-schéma formel ouvert
affine $\mathfrak U \subset \mathfrak X$, les
idéaux
$$
I_\lambda = \mathcal{I}_\lambda(\mathfrak U) \subset
A = \Gamma(\mathfrak U, \mathcal{O}_\mathfrak X)
$$
forment un système fondamental d'idéaux de définition de
l'anneau topologique admissible $A$. Il suffit de le vérifier
sur les membres d'un recouvrement ouvert affine. Supposons que $\mathfrak X$
soit quasi-compact et que $\{\mathcal{I}_\lambda\}_{\lambda \in \Lambda}$
soit un système fondamental d'idéaux faibles de définition.
Si $A$ est un anneau topologique admissible, alors tous les
idéaux ouverts suffisamment petits sont des idéaux de définition
(en effet, tout idéal ouvert contenu dans un idéal de définition
est un idéal de définition). Puisqu'il suffit donc de vérifier
la propriété sur les membres, en nombre fini, d'un recouvrement ouvert affine,
nous voyons que $\mathcal{I}_\lambda$ est un idéal de définition
pour $\lambda$ assez grand. D'après la discussion de la
Remarque \ref{remark-weak-ideals-of-definition}, nous concluons que
l'espace algébrique formel $h_\mathfrak X$ associé au
schéma formel quasi-compact $\mathfrak X$ \`a la EGA
est une limite inductive de schémas comme dans le Lemme \ref{lemma-colimit-is-formal}
si et seulement s'il existe un système fondamental d'idéaux de définition
de $\mathfrak X$.
\end{remark}






\section{Complété le long d'une partie fermée}
\label{section-completion}

\noindent
Notre notion d'espace algébrique formel se prête bien à la formation du
complété le long d'une partie fermée.

\begin{lemma}
\label{lemma-completion-affine-is-affine-formal-algebraic-space}
Soit $S$ un schéma. Soit $X$ un schéma affine sur $S$.
Soit $T \subset |X|$ une partie fermée. Alors le foncteur
$$
(\Sch/S)_{fppf} \longrightarrow \textit{Sets},\quad
U \longmapsto \{f : U \to X \mid f(|U|) \subset T\}
$$
est un espace algébrique formel affine de McQuillan.
\end{lemma}

\begin{proof}
Écrivons $X = \Spec(A)$ ; la partie $T$ correspond à l'idéal radical $I \subset A$.
Soit $U = \Spec(B)$ un schéma affine sur $S$ et soit
$f : U \to X$ un élément de $F(U)$. Alors $f$ correspond à un
homomorphisme d'anneaux $\varphi : A \to B$ tel que tout idéal premier de $B$ contienne
$\varphi(I) B$. Ainsi, tout élément de $\varphi(I)$ est nilpotent dans $B$ ; voir
Algebra, Lemma \ref{algebra-lemma-Zariski-topology}.
En posant $J = \Ker(\varphi)$, nous concluons que $I/J$ est un idéal localement
nilpotent de $A/J$. De manière équivalente, $V(J) = V(I) = T$.
Autrement dit, le foncteur du lemme est égal à
$\colim \Spec(A/J)$, où la limite inductive porte sur la
famille des idéaux $J$ tels que $V(J) = T$.
Notre foncteur est donc un espace algébrique formel affine. Il est de McQuillan
(Definition \ref{definition-types-affine-formal-algebraic-space})
car les applications $A \to A/J$ sont surjectives
et, par conséquent, $A^\wedge = \lim A/J \to A/J$ est surjective ; voir le
Lemma \ref{lemma-mcquillan-affine-formal-algebraic-space}.
\end{proof}

\begin{lemma}
\label{lemma-completion-is-formal-algebraic-space}
Soit $S$ un schéma. Soit $X$ un espace algébrique sur $S$.
Soit $T \subset |X|$ une partie fermée. Alors le foncteur
$$
(\Sch/S)_{fppf} \longrightarrow \textit{Sets},\quad
U \longmapsto \{f : U \to X \mid f(|U|) \subset T\}
$$
est un espace algébrique formel.
\end{lemma}

\begin{proof}
Notons $F$ le foncteur. Soit $\{U_i \to U\}$ un recouvrement fppf.
Alors $\coprod |U_i| \to |U|$ est surjectif. Puisque $X$ est un
faisceau fppf, il s'ensuit que $F$ est un faisceau fppf.

\medskip\noindent
Soit $\{g_i : X_i \to X\}$ un recouvrement étale tel que $X_i$ soit affine
pour tout $i$ ; voir Properties of Spaces, Lemma
\ref{spaces-properties-lemma-cover-by-union-affines}.
Les morphismes $F \times_X X_i \to F$ sont étales
(voir Spaces, Lemma
\ref{spaces-lemma-base-change-representable-transformations-property})
et l'application $\coprod F \times_X X_i \to F$ est une surjection de faisceaux.
Il suffit donc de démontrer que $F \times_X X_i$ est un espace algébrique
formel affine. Un point à valeurs dans $U$ de $F \times_X X_i$ est un
morphisme $U \to X_i$ dont l'image est contenue dans la partie fermée
$g_i^{-1}(T) \subset |X_i|$. L'assertion résulte donc du
Lemma \ref{lemma-completion-affine-is-affine-formal-algebraic-space}.
\end{proof}

\begin{definition}
\label{definition-completion}
Soit $S$ un schéma. Soit $X$ un espace algébrique sur $S$.
Soit $T \subset |X|$ une partie fermée. L'espace algébrique formel
du Lemme \ref{lemma-completion-is-formal-algebraic-space}
est appelé le {\it complété de $X$ le long de $T$}.
\end{definition}

\noindent
Dans \cite[Chapter I, Section 10.8]{EGA}, la notation $X_{/T}$
sert à désigner le complété, et nous l'emploierons aussi
à l'occasion. Soit $f : X \to X'$ un
morphisme d'espaces algébriques sur un schéma $S$. Supposons
que $T \subset |X|$ et $T' \subset |X'|$ soient des parties fermées
telles que $|f|(T) \subset T'$. Il est alors clair que
$f$ définit un morphisme d'espaces algébriques formels
$$
X_{/T} \longrightarrow X'_{/T'}
$$
entre les complétés.

\begin{lemma}
\label{lemma-map-completions-representable}
Soit $S$ un schéma. Soit $f : X' \to X$ un morphisme
d'espaces algébriques sur $S$. Soit $T \subset |X|$
une partie fermée et posons $T' = |f|^{-1}(T) \subset |X'|$.
Alors
$$
\xymatrix{
X'_{/T'} \ar[r] \ar[d] & X' \ar[d]^f \\
X_{/T} \ar[r] & X
}
$$
est un diagramme cartésien de faisceaux. En particulier, le morphisme
$X'_{/T'} \to X_{/T}$ est représentable par des espaces algébriques.
\end{lemma}

\begin{proof}
En effet, supposons que $Y \to X$ soit un morphisme d'un schéma vers $X$ tel
que $|Y|$ s'envoie dans $T$. Alors $Y \times_X X' \to X$ est un morphisme
d'espaces algébriques tel que $|Y \times_X X'|$ s'envoie dans $T'$. Le
foncteur $Y \times_{X_{/T}} X'_{/T'}$ est donc représenté par $Y \times_X X'$,
ce qui démontre le lemme.
\end{proof}

\begin{lemma}
\label{lemma-reduction-completion}
Soit $S$ un schéma. Soit $X$ un espace algébrique sur $S$.
Soit $T \subset |X|$ une partie fermée. La réduction $(X_{/T})_{red}$
du complété $X_{/T}$ de $X$ le long de $T$ est
le sous-espace fermé induit réduit $Z$ de $X$ correspondant à $T$.
\end{lemma}

\begin{proof}
Cela résulte du Lemme \ref{lemma-reduction-formal-algebraic-space},
Properties of Spaces, Definition
\ref{spaces-properties-definition-reduced-induced-space}
(qui utilise Properties of Spaces, Lemma
\ref{spaces-properties-lemma-reduced-closed-subspace} pour construire $Z$),
et de la définition de $X_{/T}$ que
$Z$ et $(X_{/T})_{red}$ sont des espaces algébriques réduits
caractérisés par la même propriété d'applications :
un morphisme $g : Y \to X$ dont la source est un espace algébrique réduit
se factorise par l'un ou l'autre si et seulement si $|Y|$ s'envoie dans $T \subset |X|$.
\end{proof}

\begin{lemma}
\label{lemma-affine-formal-completion-types}
Soit $S$ un schéma. Soit $X = \Spec(A)$ un schéma affine sur $S$.
Soit $T \subset X$ une partie fermée. Soit $X_{/T}$ le
complété formel de $X$ le long de $T$.
\begin{enumerate}
\item Si $X \setminus T$ est quasi-compact, c'est-à-dire si $T$ est constructible,
alors $X_{/T}$ est adique*.
\item Si $T = V(I)$ pour un idéal de type fini $I \subset A$,
alors $X_{/T} = \text{Spf}(A^\wedge)$, où $A^\wedge$ est le
complété $I$-adique de $A$.
\item Si $X$ est noethérien, alors $X_{/T}$ est noethérien.
\end{enumerate}
\end{lemma}

\begin{proof}
Si (1) est vérifié, Algebra, Lemma \ref{algebra-lemma-qc-open} fournit
un idéal $I \subset A$ comme en (2). Si (3) est vérifié, nous pouvons
également trouver un idéal $I \subset A$ comme en (2). De plus, les complétés
d'anneaux noethériens sont noethériens d'après
Algebra, Lemma \ref{algebra-lemma-completion-Noetherian-Noetherian}.
Il suffit donc de démontrer (2).

\medskip\noindent
Preuve de (2).
Soit $I = (f_1, \ldots, f_r) \subset A$ un idéal définissant $T$.
Si $Z = \Spec(B)$ est un schéma affine et si $g : Z \to X$ est
un morphisme tel que $g(Z) \subset T$ (au sens ensembliste), alors
$g^\sharp(f_i)$ est nilpotent dans $B$ pour tout $i$. Ainsi,
$I^n$ s'envoie sur zéro dans $B$ pour un certain $n$. Nous voyons donc que
$X_{/T} = \colim \Spec(A/I^n) = \text{Spf}(A^\wedge)$.
\end{proof}

\noindent
Le lemme suivant est dû à Ofer Gabber.

\begin{lemma}
\label{lemma-completion-countably-indexed}
\begin{reference}
Courriel d'Ofer Gabber du 11 septembre 2014.
\end{reference}
Soit $S$ un schéma. Soit $X = \Spec(A)$ un schéma affine sur $S$.
Soit $T \subset X$ un sous-schéma fermé.
\begin{enumerate}
\item Si le complété formel $X_{/T}$ est à indexation dénombrable
et s'il existe une famille dénombrable $f_1, f_2, f_3, \ldots \in A$ telle que
$T = V(f_1, f_2, f_3, \ldots)$, alors $X_{/T}$ est adique*.
\item La conclusion de (1) est fausse si l'on omet l'hypothèse selon laquelle
$T$ peut être défini par une famille dénombrable de fonctions sur $X$.
\end{enumerate}
\end{lemma}

\begin{proof}
L'hypothèse que $X_{/T}$ est à indexation dénombrable signifie qu'il existe une
suite d'idéaux
$$
A \supset J_1 \supset J_2 \supset J_3 \supset \ldots
$$
vérifiant $V(J_n) = T$ et telle que, pour tout idéal $J \subset A$ avec $V(J) = T$,
il existe un $n$ tel que $J \supset J_n$.

\medskip\noindent
Pour construire un exemple pour (2), soit $\omega_1$ le premier ordinal non
dénombrable. Soit $k$ un corps et soit
$A$ la $k$-algèbre engendrée par les $x_\alpha$, $\alpha \in \omega_1$,
et les $y_{\alpha \beta}$, où $\alpha \in \beta \in \omega_1$,
soumis aux relations $x_\alpha = y_{\alpha \beta} x_\beta$.
Soit $T = V(x_\alpha)$. Soit $J_n = (x_\alpha^n)$.
Si $J \subset A$ est un idéal tel que
$V(J) = T$, alors $x_\alpha^{n_\alpha} \in J$ pour un certain $n_\alpha \geq 1$.
L'un des ensembles $\{\alpha \mid n_\alpha = n\}$ doit être non borné dans
$\omega_1$. Les relations impliquent alors que $J_n \subset J$.

\medskip\noindent
Pour établir (2), il suffit maintenant de montrer que $A^\wedge = \lim A/J_n$
n'est complet pour la topologie définie par aucun idéal de type fini.
Pour $\gamma \in \omega_1$, soit $A_\gamma$ le quotient de $A$
par l'idéal engendré par les $x_\alpha$, $\alpha \in \gamma$, et les
$y_{\alpha \beta}$, $\alpha \in \gamma$. Comme $A/J_1$ est réduit,
tout élément topologiquement nilpotent $f$ de $\lim A/J_n$ appartient à
$J_1^\wedge = \lim J_1/J_n$. Cela signifie que $f$ est une série infinie
ne faisant intervenir qu'un nombre dénombrable de générateurs. Ainsi, $f$ s'annule dans
$A_\gamma^\wedge = \lim A_\gamma/J_nA_\gamma$ pour un certain $\gamma$.
Notons que $A^\wedge \to A_\gamma^\wedge$ est continu et ouvert d'après le
Lemma \ref{lemma-ses}.
Si la topologie de $A^\wedge$ était $I$-adique pour un idéal de type fini
$I \subset A^\wedge$, alors $I$ s'enverrait sur zéro dans un certain
$A_\gamma^\wedge$. Cela signifierait que $A_\gamma^\wedge$ est discret,
ce qui n'est pas le cas, puisqu'il existe une application surjective, continue et ouverte
(d'après le Lemme \ref{lemma-ses})
$A_\gamma^\wedge \to k[[t]]$ donnée par
$x_\alpha \mapsto t$, $y_{\alpha \beta} \mapsto 1$ lorsque
$\gamma = \alpha$ ou $\gamma \in \alpha$.

\medskip\noindent
Avant de démontrer (1), établissons d'abord l'assertion suivante : si $I \subset A^\wedge$ est
un idéal de type fini dont l'adhérence $\bar I$ est ouverte, alors $I = \bar I$.
Puisque $V(J_n^2) = T$, il existe un $m$ tel que $J_n^2 \supset J_m$.
En passant à une sous-suite, nous pouvons donc supposer que $J_n^2 \supset J_{n + 1}$ pour tout
$n$. Posons $J_n^\wedge = \lim_{k \geq n} J_n/J_k \subset A^\wedge$.
Comme l'adhérence $\bar I = \bigcap (I + J_n^\wedge)$
(Lemme \ref{lemma-closed}) est ouverte, il existe un $m$ tel que
$I + J_n^\wedge \supset J_m^\wedge$ pour tout $n \geq m$. Fixons un tel $m$.
Nous avons
$$
J_{n - 1}^\wedge I + J_{n + 1}^\wedge \supset
J_{n - 1}^\wedge (I + J_{n + 1}^\wedge) \supset
J_{n - 1}^\wedge J_m^\wedge
$$
pour tout $n \geq m + 1$. En effet, la première inclusion est triviale et la
seconde a été établie ci-dessus. Puisque $J_{n - 1}J_m \supset J_{n - 1}^2 \supset J_n$,
ces inclusions montrent que l'image de $J_n$ dans $A^\wedge$
est contenue dans l'idéal $J_{n - 1}^\wedge I + J_{n + 1}^\wedge$.
Comme cet idéal est ouvert, nous concluons que
$$
J_{n - 1}^\wedge I + J_{n + 1}^\wedge \supset J_n^\wedge.
$$
Écrivons $I = (g_1, \ldots, g_t)$. Prenons $f \in J_{m + 1}^\wedge$.
À l'aide de la dernière inclusion affichée, valable pour tout $n \geq m + 1$,
nous pouvons écrire, par récurrence sur $c \geq 0$,
$$
f = \sum f_{i, c} g_i \mod J_{m + 1+ c}^\wedge
$$
avec $f_{i, c} \in J_m^\wedge$ et
$f_{i, c} \equiv f_{i, c - 1} \bmod J_{m + c}^\wedge$.
Il s'ensuit que $IJ_m^\wedge \supset J_{m + 1}^\wedge$.
En combinant ceci avec $I + J_{m + 1}^\wedge \supset J_m^\wedge$,
nous concluons que $I$ est ouvert.

\medskip\noindent
Preuve de (1). Supposons $T = V(f_1, f_2, f_3, \ldots)$.
Soit $I_m \subset A^\wedge$ l'idéal engendré par $f_1, \ldots, f_m$.
Distinguons deux cas.

\medskip\noindent
Cas I : pour un certain $m$, l'adhérence de $I_m$ est ouverte.
Alors $I_m$ est ouvert d'après le résultat du paragraphe précédent.
Pour tout $n$, nous avons $(J_n)^2 \supset J_{n+1}$ par construction ; par conséquent,
l'adhérence de $(J_n^\wedge)^2$ contient $J_{n+1}^\wedge$
et est donc ouverte. En prenant $n$ assez grand, il s'ensuit que l'adhérence
du produit de deux idéaux ouverts quelconques de $A^\wedge$ est ouverte.
Démontrons que $I_m^k$ est ouvert pour $k \ge 1$, par récurrence sur $k$.
Le cas $k = 1$ est précisément notre hypothèse sur $m$ dans le cas I.
Pour $k > 1$, supposons $I_m^{k - 1}$ ouvert. Alors
$I_m^k = I_m^{k - 1} \cdot I_m$ est le produit de deux idéaux ouverts
et son adhérence est donc ouverte. Mais comme $I_m^k$
est de type fini, il s'ensuit que $I_m^k$
est ouvert d'après le paragraphe précédent (appliqué à $I = I_m^k$) ;
nous pouvons donc poursuivre la récurrence sur $k$.
Comme tout élément de $I_m$ est topologiquement nilpotent, nous concluons
que $I_m$ est un idéal de définition, ce qui prouve que $A^\wedge$
est adique et possède un idéal de définition de type fini, c'est-à-dire que
$X_{/T}$ est adique*.

\medskip\noindent
Cas II. Pour tout $m$, l'adhérence $\bar I_m$ de $I_m$ n'est pas ouverte.
La topologie de $A^\wedge/\bar I_m$ n'est alors pas discrète. Cela signifie
que nous pouvons choisir $\phi(m) \geq m$ tel que
$$
\Im(J_{\phi(m)} \to A/(f_1, \ldots, f_m)) \not =
\Im(J_{\phi(m) + 1} \to A/(f_1, \ldots, f_m))
$$
Pour le voir, nous avons utilisé l'égalité
$A^\wedge/(\bar I_m + J_n^\wedge) = A/((f_1, \ldots, f_m) + J_n)$.
Choisissons des exposants $e_i > 0$ tels que $f_i^{e_i} \in J_{\phi(m) + 1}$
pour $0 < m < i$. Posons $J = (f_1^{e_1}, f_2^{e_2}, f_3^{e_3}, \ldots)$.
Alors $V(J) = T$. Nous affirmons que $J \not \supset J_n$ pour tout $n$,
ce qui est une contradiction et montre que le cas II ne se produit pas.
En effet, l'image de $J$ dans $A/(f_1, \ldots, f_m)$ est contenue
dans l'image de $J_{\phi(m) + 1}$, laquelle est strictement contenue dans
l'image de $J_m$.
\end{proof}
















\section{Produits fibrés}
\label{section-fibre-products}

\noindent
Section obligée sur les produits fibrés d'espaces algébriques formels.

\begin{lemma}
\label{lemma-etale-covering-by-formal-algebraic-spaces}
Soit $S$ un schéma. Soit $\{X_i \to X\}_{i \in I}$ une famille d'applications
de faisceaux sur $(\Sch/S)_{fppf}$. Supposons (a) que $X_i$ soit un
espace algébrique formel sur $S$, (b) que $X_i \to X$ soit représentable
par des espaces algébriques et étale, et (c) que $\coprod X_i \to X$
soit une surjection de faisceaux. Alors $X$ est un espace algébrique formel
sur $S$.
\end{lemma}

\begin{proof}
Pour chaque $i$, choisissons $\{X_{ij} \to X_i\}_{j \in J_i}$ comme dans la
Definition \ref{definition-formal-algebraic-space}.
Alors $\{X_{ij} \to X\}_{i \in I, j \in J_i}$ est une famille
comme dans la Définition \ref{definition-formal-algebraic-space}
pour $X$.
\end{proof}

\begin{lemma}
\label{lemma-fibre-products-general}
Soit $S$ un schéma. Soient $X, Y$ des espaces algébriques formels sur $S$
et soit $Z$ un faisceau dont la diagonale est représentable par des
espaces algébriques. Soient $X \to Z$ et $Y \to Z$ des applications de faisceaux.
Alors $X \times_Z Y$ est un espace algébrique formel.
\end{lemma}

\begin{proof}
Choisissons $\{X_i \to X\}$ et $\{Y_j \to Y\}$ comme dans la
Definition \ref{definition-formal-algebraic-space}.
Alors $\{X_i \times_Z Y_j \to X \times_Z Y\}$ est une famille
d'applications représentables par des espaces algébriques et étales.
Le Lemme \ref{lemma-etale-covering-by-formal-algebraic-spaces}
montre donc qu'il suffit d'établir que $X \times_Z Y$ est un espace algébrique
formel lorsque $X$ et $Y$ sont des espaces algébriques formels affines.

\medskip\noindent
Supposons que $X$ et $Y$ soient des espaces algébriques formels affines.
Écrivons $X = \colim X_\lambda$ et $Y = \colim Y_\mu$ comme
dans la Définition \ref{definition-affine-formal-algebraic-space}.
Alors $X \times_Z Y = \colim X_\lambda \times_Z Y_\mu$.
Chaque $X_\lambda \times_Z Y_\mu$ est un espace algébrique.
Pour $\lambda \leq \lambda'$ et $\mu \leq \mu'$, le morphisme
$$
X_\lambda \times_Z Y_\mu \to
X_\lambda \times_Z Y_{\mu'} \to
X_{\lambda'} \times_Z Y_{\mu'}
$$
est un épaississement, car c'est un composé de changements de base d'épaississements.
Nous concluons donc en appliquant le Lemme \ref{lemma-colimit-is-formal}.
\end{proof}

\begin{lemma}
\label{lemma-fibre-products}
Soit $S$ un schéma. La catégorie des espaces algébriques formels sur $S$
possède des produits fibrés.
\end{lemma}

\begin{proof}
C'est un cas particulier du Lemme \ref{lemma-fibre-products-general},
car les espaces algébriques formels ont des diagonales représentables ; voir le
Lemme \ref{lemma-diagonal-formal-algebraic-space}.
\end{proof}

\begin{lemma}
\label{lemma-reduction-fibre-products}
Soit $S$ un schéma. Soient $X \to Z$ et $Y \to Z$ des morphismes
d'espaces algébriques formels sur $S$. Alors
$(X \times_Z Y)_{red} = (X_{red} \times_{Z_{red}} Y_{red})_{red}$.
\end{lemma}

\begin{proof}
Cela résulte de la propriété universelle de la réduction
du Lemme \ref{lemma-reduction-formal-algebraic-space}.
\end{proof}

\noindent
Nous avons déjà démontré le lemme suivant (sans savoir que
les produits fibrés existent).

\begin{lemma}
\label{lemma-diagonal-morphism-formal-algebraic-spaces}
Soit $S$ un schéma. Soit $f : X \to Y$ un morphisme d'espaces algébriques formels
sur $S$. Le morphisme diagonal $\Delta : X \to X \times_Y X$
est représentable (par des schémas), est un monomorphisme, est localement quasi-fini,
localement de type fini et séparé.
\end{lemma}

\begin{proof}
Soit $T$ un schéma et soit $T \to X \times_Y X$ un morphisme.
Alors
$$
T \times_{(X \times_Y X)} X = T \times_{(X \times_S X)} X
$$
Le résultat découle donc immédiatement du
Lemme \ref{lemma-diagonal-formal-algebraic-space}.
\end{proof}






\section{Axiomes de séparation pour les espaces algébriques formels}
\label{section-separation}

\noindent
Cette section porte sur les conditions de séparation « absolues » des espaces
algébriques formels. Nous étudierons plus loin les conditions de séparation
des morphismes d'espaces algébriques formels.

\begin{lemma}
\label{lemma-characterize-quasi-separated}
Soit $S$ un schéma. Soit $X$ un espace algébrique formel sur $S$.
Les conditions suivantes sont équivalentes :
\begin{enumerate}
\item la réduction de $X$
(Lemme \ref{lemma-reduction-formal-algebraic-space}) est un
espace algébrique quasi-séparé ;
\item pour $U \to X$, $V \to X$, où $U$, $V$ sont des schémas quasi-compacts,
le produit fibré $U \times_X V$ est quasi-compact ;
\item pour $U \to X$, $V \to X$, où $U$, $V$ sont affines,
le produit fibré $U \times_X V$ est quasi-compact.
\end{enumerate}
\end{lemma}

\begin{proof}
Observons que $U \times_X V$ est un schéma d'après le
Lemme \ref{lemma-diagonal-formal-algebraic-space}.
Soient $U_{red}, V_{red}, X_{red}$ les réductions de $U, V, X$.
Alors
$$
U_{red} \times_{X_{red}} V_{red} = U_{red} \times_X V_{red} \to U \times_X V
$$
est un épaississement de schémas. L'équivalence de (1) et (2) en découle
clairement, compte tenu du lemme analogue pour les espaces algébriques ; voir
Properties of Spaces, Lemma
\ref{spaces-properties-lemma-characterize-quasi-separated}.
Nous omettons la démonstration de l'équivalence de (2) et (3).
\end{proof}

\begin{lemma}
\label{lemma-characterize-separated}
Soit $S$ un schéma. Soit $X$ un espace algébrique formel sur $S$.
Les conditions suivantes sont équivalentes :
\begin{enumerate}
\item la réduction de $X$
(Lemme \ref{lemma-reduction-formal-algebraic-space}) est un espace
algébrique séparé ;
\item pour $U \to X$, $V \to X$, où $U$, $V$ sont affines,
le produit fibré $U \times_X V$ est affine et l'application
$$
\mathcal{O}(U) \otimes_\mathbf{Z} \mathcal{O}(V)
\longrightarrow
\mathcal{O}(U \times_X V)
$$
est surjective.
\end{enumerate}
\end{lemma}

\begin{proof}
Si (2) est vérifiée, alors $X_{red}$ est un espace algébrique séparé : on applique
Properties of Spaces, Lemma
\ref{spaces-properties-lemma-characterize-quasi-separated}
aux morphismes $U \to X_{red}$ et $V \to X_{red}$,
avec $U, V$ affines, et on utilise que $U \times_{X_{red}} V = U \times_X V$.

\medskip\noindent
Supposons (1). Soient $U \to X$ et $V \to X$ comme en (2).
Observons que $U \times_X V$ est un schéma d'après le
Lemme \ref{lemma-diagonal-formal-algebraic-space}.
Soient $U_{red}, V_{red}, X_{red}$ les réductions de $U, V, X$.
Alors
$$
U_{red} \times_{X_{red}} V_{red} = U_{red} \times_X V_{red} \to U \times_X V
$$
est un épaississement de schémas. Il s'ensuit que
$(U \times_X V)_{red} = (U_{red} \times_{X_{red}} V_{red})_{red}$.
En particulier, nous voyons que $(U \times_X V)_{red}$ est un schéma affine
et que l'application
$$
\mathcal{O}(U) \otimes_\mathbf{Z} \mathcal{O}(V)
\longrightarrow
\mathcal{O}((U \times_X V)_{red})
$$
est surjective ; voir Properties of Spaces, Lemma
\ref{spaces-properties-lemma-characterize-quasi-separated}.
Alors $U \times_X V$ est affine d'après
Limits of Spaces, Proposition \ref{spaces-limits-proposition-affine}.
D'autre part, le morphisme $U \times_X V \to U \times V$
de schémas affines est le composé
$$
U \times_X V = X \times_{(X \times_S X)} (U \times_S V)
\to U \times_S V \to U \times V
$$
Le premier morphisme est un monomorphisme localement de type fini
(Lemma \ref{lemma-diagonal-formal-algebraic-space}).
Le second morphisme est une immersion
(Schemes, Lemma \ref{schemes-lemma-fibre-product-after-map}).
Le composé est donc un monomorphisme localement de type fini.
D'autre part, le composé est entier, car l'application entre les schémas
affines réduits sous-jacents est une immersion fermée
d'après ce qui précède, donc universellement fermée (utiliser
Morphisms, Lemma \ref{morphisms-lemma-integral-universally-closed}).
Ainsi, l'homomorphisme d'anneaux
$$
\mathcal{O}(U) \otimes_\mathbf{Z} \mathcal{O}(V)
\longrightarrow
\mathcal{O}(U \times_X V)
$$
est un épimorphisme entier de type fini,
donc fini, donc surjectif (utiliser
Morphisms, Lemma \ref{morphisms-lemma-finite-integral}
et
Algebra, Lemma \ref{algebra-lemma-finite-epimorphism-surjective}).
\end{proof}

\begin{definition}
\label{definition-separated}
Soit $S$ un schéma. Soit $X$ un espace algébrique formel sur $S$.
On dit que
\begin{enumerate}
\item $X$ est {\it quasi-séparé} si les conditions équivalentes du
Lemme \ref{lemma-characterize-quasi-separated} sont satisfaites ;
\item $X$ est {\it séparé} si les conditions équivalentes du
Lemme \ref{lemma-characterize-separated} sont satisfaites.
\end{enumerate}
\end{definition}

\noindent
Le lemme suivant implique en particulier que le produit tensoriel
complété d'anneaux topologiques faiblement admissibles est
un anneau topologique faiblement admissible.

\begin{lemma}
\label{lemma-fibre-product-affines-over-separated}
Soit $S$ un schéma. Soient $X \to Z$ et $Y \to Z$ des morphismes
d'espaces algébriques formels sur $S$. Supposons $Z$ séparé.
\begin{enumerate}
\item Si $X$ et $Y$ sont des espaces algébriques formels affines, alors
$X \times_Z Y$ l'est aussi ;
\item si $X$ et $Y$ sont des espaces algébriques formels affines de McQuillan, alors
$X \times_Z Y$ l'est aussi ;
\item si $X$, $Y$ et $Z$ sont des espaces algébriques formels affines de McQuillan
correspondant aux $S$-algèbres topologiques faiblement admissibles
$A$, $B$ et $C$, alors $X \times_Z Y$ correspond à
$A \widehat{\otimes}_C B$.
\end{enumerate}
\end{lemma}

\begin{proof}
Écrivons $X = \colim X_\lambda$ et $Y = \colim Y_\mu$ comme
dans la Définition \ref{definition-affine-formal-algebraic-space}.
Alors $X \times_Z Y = \colim X_\lambda \times_Z Y_\mu$.
Puisque $Z$ est séparé, les produits fibrés sont affines ;
nous voyons donc que (1) est vérifiée. Supposons que $X$ et $Y$ correspondent
aux $S$-algèbres topologiques faiblement admissibles $A$ et $B$,
et que $X_\lambda = \Spec(A/I_\lambda)$ et $Y_\mu = \Spec(B/J_\mu)$.
Alors
$$
X_\lambda \times_Z Y_\mu \to
X_\lambda \times Y_\mu \to \Spec(A \otimes B)
$$
est une immersion fermée. L'une des conditions du
Lemme \ref{lemma-mcquillan-affine-formal-algebraic-space}
est donc satisfaite et nous concluons que $X \times_Z Y$ est de McQuillan.
Si de plus $Z$ est de McQuillan et correspond à $C$, alors
$$
X_\lambda \times_Z Y_\mu = \Spec(A/I_\lambda \otimes_C B/J_\mu)
$$
nous voyons donc que l'anneau topologique faiblement admissible
correspondant à $X \times_Z Y$ est le produit tensoriel complété
(voir Définition \ref{definition-toplogy-tensor-product}).
\end{proof}

\begin{lemma}
\label{lemma-separated-from-separated}
Soit $S$ un schéma. Soit $X$ un espace algébrique formel sur $S$.
Soit $U \to X$ un morphisme, où $U$ est un espace algébrique
séparé sur $S$. Alors $U \to X$ est séparé.
\end{lemma}

\begin{proof}
L'énoncé a un sens, car $U \to X$ est représentable par des
espaces algébriques (Lemme \ref{lemma-space-to-formal-space}).
Soit $T$ un schéma et soit $T \to X$ un morphisme. Nous devons montrer
que $U \times_X T \to T$ est séparé. Puisque $U \times_X T \to U \times_S T$
est un monomorphisme, il suffit de montrer que $U \times_S T \to T$
est séparé. Cela résulte du fait qu'il s'agit du changement de base de $U \to S$.
Nous avons utilisé dans l'argument précédent :
Morphisms of Spaces, Lemmas
\ref{spaces-morphisms-lemma-base-change-separated},
\ref{spaces-morphisms-lemma-composition-separated},
\ref{spaces-morphisms-lemma-monomorphism-separated}, et
\ref{spaces-morphisms-lemma-separated-implies-morphism-separated}.
\end{proof}




\section{Espaces algébriques formels quasi-compacts}
\label{section-quasi-compact}

\noindent
Voici la caractérisation des espaces algébriques formels
quasi-compacts.

\begin{lemma}
\label{lemma-characterize-quasi-compact}
Soit $S$ un schéma. Soit $X$ un espace algébrique formel sur $S$.
Les conditions suivantes sont équivalentes :
\begin{enumerate}
\item la réduction de $X$
(Lemme \ref{lemma-reduction-formal-algebraic-space}) est un espace
algébrique quasi-compact ;
\item on peut trouver $\{X_i \to X\}_{i \in I}$ comme dans la
Définition \ref{definition-formal-algebraic-space}, avec $I$ fini ;
\item il existe un morphisme $Y \to X$ représentable par des espaces
algébriques, étale et surjectif, tel que
$Y$ soit un espace algébrique formel affine.
\end{enumerate}
\end{lemma}

\begin{proof}
Omis.
\end{proof}

\begin{definition}
\label{definition-quasi-compact}
Soit $S$ un schéma. Soit $X$ un espace algébrique formel sur $S$.
On dit que $X$ est {\it quasi-compact} si les conditions équivalentes du
Lemme \ref{lemma-characterize-quasi-compact} sont satisfaites.
\end{definition}

\begin{lemma}
\label{lemma-characterize-quasi-compact-morphism}
Soit $S$ un schéma. Soit $f : X \to Y$ un morphisme d'espaces algébriques
formels sur $S$. Les conditions suivantes sont équivalentes :
\begin{enumerate}
\item l'application induite $f_{red} : X_{red} \to Y_{red}$ entre les réductions
(Lemme \ref{lemma-reduction-formal-algebraic-space}) est un morphisme
quasi-compact d'espaces algébriques ;
\item pour tout schéma quasi-compact $T$ et tout morphisme $T \to Y$,
le produit fibré $X \times_Y T$ est un espace algébrique formel
quasi-compact ;
\item pour tout schéma affine $T$ et tout morphisme $T \to Y$,
le produit fibré $X \times_Y T$ est un espace algébrique formel
quasi-compact ;
\item il existe un recouvrement $\{Y_j \to Y\}$ comme dans la
Définition \ref{definition-formal-algebraic-space}
tel que chaque $X \times_Y Y_j$ soit un espace algébrique formel quasi-compact.
\end{enumerate}
\end{lemma}

\begin{proof}
Omis.
\end{proof}

\begin{definition}
\label{definition-quasi-compact-morphism}
Soit $S$ un schéma. Soit $f : X \to Y$ un morphisme
d'espaces algébriques formels sur $S$.
On dit que $f$ est {\it quasi-compact} si les conditions équivalentes du
Lemme \ref{lemma-characterize-quasi-compact-morphism} sont satisfaites.
\end{definition}

\noindent
Cette notion coïncide avec celle qui existe déjà lorsque le morphisme
est représentable par des espaces algébriques (et, en particulier, lorsqu'il est
représentable).

\begin{lemma}
\label{lemma-quasi-compact-representable}
Soit $S$ un schéma. Soit $f : X \to Y$ un morphisme d'espaces algébriques
formels sur $S$, représentable par des espaces algébriques.
Alors $f$ est quasi-compact au sens de la
Définition \ref{definition-quasi-compact-morphism}
si et seulement si $f$ est quasi-compact au sens de
Bootstrap, Definition \ref{bootstrap-definition-property-transformation}.
\end{lemma}

\begin{proof}
Cela résulte immédiatement des définitions et du
Lemme \ref{lemma-characterize-quasi-compact-morphism}.
\end{proof}





\section{Espaces algébriques formels quasi-compacts et quasi-séparés}
\label{section-quasi-compact-quasi-separated}

\noindent
Le résultat suivant est dû à Yasuda ; voir
\cite[Proposition 3.32]{Yasuda}.

\begin{lemma}
\label{lemma-structure-quasi-compact-quasi-separated}
\begin{reference}
\cite[Proposition 3.32]{Yasuda}
\end{reference}
Soit $S$ un schéma. Soit $X$ un espace algébrique formel quasi-compact et
quasi-séparé sur $S$. Alors $X = \colim X_\lambda$
pour un système d'espaces algébriques $(X_\lambda, f_{\lambda \mu})$
indexé par un ensemble ordonné filtrant $\Lambda$, où chaque
$f_{\lambda \mu} : X_\lambda \to X_\mu$ est un épaississement.
\end{lemma}

\begin{proof}
D'après le Lemme \ref{lemma-characterize-quasi-compact}, nous pouvons choisir un
espace algébrique formel affine $Y$ et un morphisme étale surjectif
représentable $Y \to X$. Écrivons $Y = \colim Y_\lambda$ comme dans la
Définition \ref{definition-affine-formal-algebraic-space}.

\medskip\noindent
Prenons $\lambda \in \Lambda$. Alors $Y_\lambda \times_X Y$ est un schéma d'après le
Lemme \ref{lemma-presentation-representable}. La réduction
(Lemma \ref{lemma-reduction-formal-algebraic-space})
de $Y_\lambda \times_X Y$ est égale à la réduction de
$Y_{red} \times_{X_{red}} Y_{red}$, laquelle est quasi-compact puisque $X$
est quasi-séparé et que $Y_{red}$ est affine.
Par conséquent, $Y_\lambda \times_X Y$ est un schéma quasi-compact.
Il existe donc un $\mu \geq \lambda$ tel que
$\text{pr}_2 : Y_\lambda \times_X Y \to Y$ se factorise
par $Y_\mu$ ; voir le Lemme \ref{lemma-factor-through-thickening}.
Soit $Z_\lambda$ l'image schématique du morphisme
$\text{pr}_2 : Y_\lambda \times_X Y \to Y_\mu$.
Celle-ci est indépendante du choix de $\mu$, et nous pouvons donc
considérer $Z_\lambda \subset Y$ comme l'image schématique
du morphisme $\text{pr}_2 : Y_\lambda \times_X Y \to Y$.
Observons que $Z_\lambda$ est aussi égale à l'image schématique
du morphisme $\text{pr}_1 : Y \times_X Y_\lambda \to Y$, puisque
celui-ci est isomorphe au morphisme ayant servi à définir $Z_\lambda$.
Nous affirmons que $Z_\lambda \times_X Y = Y \times_X Z_\lambda$ comme sous-foncteurs
de $Y \times_X Y$. En effet, puisque $Y \to X$ est étale, nous voyons que
$Z_\lambda \times_X Y$ est l'image schématique du morphisme
$$
\text{pr}_{13} = \text{pr}_1 \times \text{id}_Y :
Y \times_X Y_\lambda \times_X Y \longrightarrow Y \times_X Y
$$
d'après Morphisms of Spaces, Lemma
\ref{spaces-morphisms-lemma-quasi-compact-scheme-theoretic-image}.
De même, $Y \times_X Z_\lambda$ est l'image schématique
du morphisme
$$
\text{pr}_{13} = \text{id}_Y \times \text{pr}_2 : 
Y \times_X Y_\lambda \times_X Y \longrightarrow Y \times_X Y
$$
L'affirmation en résulte. Alors
$R_\lambda = Z_\lambda \times_X Y = Y \times_X Z_\lambda$,
muni du morphisme $R_\lambda \to Z_\lambda \times_S Z_\lambda$,
définit une relation d'équivalence étale. Nous obtenons ainsi un espace
algébrique $X_\lambda = Z_\lambda/R_\lambda$. Par construction, le diagramme
$$
\xymatrix{
Z_\lambda \ar[r] \ar[d] & Y \ar[d] \\
X_\lambda \ar[r] & X
}
$$
est cartésien (car $X$ est le conoyau des deux projections
$R = Y \times_X Y \to Y$, car $Z_\lambda \subset Y$ est $R$-invariant,
et car $R_\lambda$ est la restriction de $R$ à $Z_\lambda$).
Ainsi, $X_\lambda \to X$ est représentable et est une immersion fermée ; voir
Spaces, Lemma
\ref{spaces-lemma-morphism-sheaves-with-P-effective-descent-etale}.
D'autre part, puisque $Y_\lambda \subset Z_\lambda$, nous voyons que
$(X_\lambda)_{red} = X_{red}$ ; autrement dit, $X_\lambda \to X$
est un épaississement. Enfin, nous affirmons que
$$
X = \colim X_\lambda
$$
Nous avons $Y \times_X X_\lambda = Z_\lambda \supset Y_\lambda$. Tout
morphisme $T \to X$, où $T$ est un schéma sur $S$, se relève localement pour la topologie étale
en un morphisme vers $Y$, lequel se relève localement pour la topologie étale en un morphisme
vers un certain $Y_\lambda$. Ainsi, $T \to X$ se relève localement pour la topologie étale sur
$T$ en un morphisme vers $X_\lambda$. Ceci achève la démonstration.
\end{proof}

\begin{remark}
\label{remark-structure-quasi-compact-quasi-separated}
Dans cette remarque, nous traduisons l'énoncé et la démonstration du
Lemme \ref{lemma-structure-quasi-compact-quasi-separated}
dans le langage des schémas formels \`a la EGA.
La Remarque \ref{remark-ideals-of-definition} montre
que le lemme se traduit comme suit :
\begin{itemize}
\item[(*)] Tout schéma formel quasi-compact et quasi-séparé
possède un système fondamental d'idéaux de définition.
\end{itemize}
Pour le démontrer, nous utilisons d'abord le principe de récurrence (reformulé pour les
schémas formels quasi-compacts et quasi-séparés) de
Cohomology of Schemes, Lemma \ref{coherent-lemma-induction-principle}
pour nous ramener à la situation suivante :
$\mathfrak X = \mathfrak U \cup \mathfrak V$
où $\mathfrak U$, $\mathfrak V$ sont des sous-schémas formels ouverts,
où $\mathfrak V$ est affine, et où le résultat est vrai pour $\mathfrak U$,
$\mathfrak V$ et $\mathfrak U \cap \mathfrak V$. Choisissons des idéaux
de définition quelconques $\mathcal{I} \subset \mathcal{O}_\mathfrak U$
et $\mathcal{J} \subset \mathcal{O}_\mathfrak V$.
Par hypothèse, nous disposons d'un système fondamental d'idéaux
de définition sur $\mathfrak U$ et $\mathfrak V$ ; comme
$\mathfrak U \cap \mathfrak V$ est quasi-compact, nous pouvons trouver
des idéaux de définition $\mathcal{I}' \subset \mathcal{I}$
et $\mathcal{J}' \subset \mathcal{J}$
tels que
$$
\mathcal{I}'|_{\mathfrak U \cap \mathfrak V} \subset
\mathcal{J}|_{\mathfrak U \cap \mathfrak V}
\quad\text{et}\quad
\mathcal{J}'|_{\mathfrak U \cap \mathfrak V} \subset
\mathcal{I}|_{\mathfrak U \cap \mathfrak V}
$$
Soient $U \to U' \to \mathfrak U$ et $V \to V' \to \mathfrak V$ les
immersions fermées déterminées par les idéaux de définition
$\mathcal{I}' \subset \mathcal{I} \subset \mathcal{O}_\mathfrak U$
et
$\mathcal{J}' \subset \mathcal{J} \subset \mathcal{O}_\mathfrak V$.
Notons $\mathfrak U \cap V$ le sous-schéma ouvert de $V$ dont
l'espace topologique sous-jacent est celui de $\mathfrak U \cap \mathfrak V$.
Par notre choix de $\mathcal{I}'$, il existe une factorisation
$\mathfrak U \cap V \to U'$.
Nous définissons de même $U \cap \mathfrak V$, qui se factorise par $V'$.
Considérons alors
$$
Z_U = \text{image schématique de }
U \amalg (\mathfrak U \cap V) \longrightarrow U'
$$
et
$$
Z_V = \text{image schématique de }
(U \cap \mathfrak V) \amalg V \longrightarrow V'
$$
Comme la formation de l'image schématique d'un morphisme quasi-compact
commute à la restriction aux ouverts (Morphisms, Lemma
\ref{morphisms-lemma-quasi-compact-scheme-theoretic-image})
on voit que $Z_U \cap \mathfrak V = \mathfrak U \cap Z_V$.
Ainsi, $Z_U$ et $Z_V$ se recollent en un schéma $Z$ muni
d'un morphisme $Z \to \mathfrak X$. Comme dans la discussion de la
Remark \ref{remark-weak-ideals-of-definition}
nous voyons que $Z$ correspond à un idéal faible
de définition $\mathcal{I}_Z \subset \mathcal{O}_\mathfrak X$.
Notons que $Z_U \subset U'$ et que
$Z_V \subset V'$. La famille de tous les $\mathcal{I}_Z$
ainsi construits forme donc un système fondamental d'idéaux faibles
de définition. Une sous-famille fournit dès lors un système fondamental d'idéaux
de définition ; voir la Remarque \ref{remark-ideals-of-definition}.
\end{remark}

\begin{lemma}
\label{lemma-characterize-affine}
Soit $S$ un schéma. Soit $X$ un espace algébrique formel sur $S$.
Alors $X$ est un espace algébrique formel affine si et seulement si
sa réduction $X_{red}$ (Lemme \ref{lemma-reduction-formal-algebraic-space})
est affine.
\end{lemma}

\begin{proof}
Par les Lemmes \ref{lemma-characterize-quasi-separated} et
\ref{lemma-characterize-quasi-compact} ainsi que par les
Définitions \ref{definition-separated} et \ref{definition-quasi-compact},
nous voyons que $X$ est quasi-compact et quasi-séparé.
D'après le lemme de Yasuda (Lemme \ref{lemma-structure-quasi-compact-quasi-separated}),
nous pouvons écrire $X = \colim X_\lambda$ comme une limite inductive filtrante
d'épaississements d'espaces algébriques. Cependant, chaque $X_\lambda$
est affine d'après Limits of Spaces, Lemma
\ref{spaces-limits-lemma-reduction-scheme}
car $(X_\lambda)_{red} = X_{red}$.
Ainsi, $X$ est un espace algébrique formel affine par définition.
\end{proof}






\section{Morphismes représentables par des espaces algébriques}
\label{section-representable}

\noindent
Soit $f : X \to Y$ un morphisme d'espaces algébriques formels qui
est représentable par des espaces algébriques. Pour ce type de morphismes,
nous disposons d'une théorie abondante, grâce au travail accompli
dans les chapitres consacrés aux espaces algébriques.

\begin{lemma}
\label{lemma-composition-representable}
Le composé de morphismes représentables par des espaces algébriques est
représentable par des espaces algébriques. Il en va de même pour les morphismes représentables
(par des schémas).
\end{lemma}

\begin{proof}
Voir Bootstrap, Lemma \ref{bootstrap-lemma-composition-transformation}.
\end{proof}

\begin{lemma}
\label{lemma-base-change-representable}
Tout changement de base d'un morphisme représentable par des espaces algébriques est
représentable par des espaces algébriques. Il en va de même pour les morphismes représentables
(par des schémas).
\end{lemma}

\begin{proof}
Voir Bootstrap, Lemma \ref{bootstrap-lemma-base-change-transformation}.
\end{proof}

\begin{lemma}
\label{lemma-permanence-representable}
Soit $S$ un schéma. Soient $f : X \to Y$ et $g : Y \to Z$ des morphismes
d'espaces algébriques formels sur $S$. Si $g \circ f : X \to Z$ est représentable
par des espaces algébriques, alors $f : X \to Y$ est représentable par des espaces algébriques.
\end{lemma}

\begin{proof}
Notons que la diagonale de $Y \to Z$ est représentable d'après le
Lemme \ref{lemma-diagonal-morphism-formal-algebraic-spaces}.
Ainsi, $X \to Y$ est représentable par des espaces algébriques d'après
Bootstrap, Lemma \ref{bootstrap-lemma-representable-by-spaces-permanence}.
\end{proof}

\noindent
La propriété d'être représentable par des espaces algébriques est locale sur la
source et sur le but.

\begin{lemma}
\label{lemma-representable-by-algebraic-spaces-local}
Soit $S$ un schéma. Soit $f : X \to Y$ un morphisme d'espaces algébriques
formels sur $S$. Les conditions suivantes sont équivalentes :
\begin{enumerate}
\item le morphisme $f$ est représentable par des espaces algébriques ;
\item il existe un diagramme commutatif
$$
\xymatrix{
U \ar[d] \ar[r] & V \ar[d] \\
X \ar[r] & Y
}
$$
où $U$, $V$ sont des espaces algébriques formels, les flèches verticales sont
représentables par des espaces algébriques, $U \to X$
est étale surjectif et $U \to V$ est représentable par des espaces algébriques ;
\item pour tout diagramme commutatif
$$
\xymatrix{
U \ar[d] \ar[r] & V \ar[d] \\
X \ar[r] & Y
}
$$
où $U$, $V$ sont des espaces algébriques formels et les flèches verticales sont
représentables par des espaces algébriques, le morphisme $U \to V$ est
représentable par des espaces algébriques ;
\item il existe un recouvrement $\{Y_j \to Y\}$ comme dans la
Définition \ref{definition-formal-algebraic-space}
et, pour chaque $j$, un recouvrement $\{X_{ji} \to Y_j \times_Y X\}$ comme dans la
Définition \ref{definition-formal-algebraic-space}, tels que
$X_{ji} \to Y_j$ soit représentable par des espaces algébriques pour tous $j$ et $i$ ;
\item il existe un recouvrement $\{X_i \to X\}$ comme dans la
Définition \ref{definition-formal-algebraic-space}
et, pour chaque $i$, une factorisation $X_i \to Y_i \to Y$, où $Y_i$
est un espace algébrique formel affine, où $Y_i \to Y$ est représentable
par des espaces algébriques, et telle que $X_i \to Y_i$ soit représentable par des espaces
algébriques ;
\item en ajouter ici.
\end{enumerate}
\end{lemma}

\begin{proof}
Il est clair que (1) implique (2), car on peut prendre $U = X$ et $V = Y$.
Réciproquement, (2) implique (1) d'après
Bootstrap, Lemma \ref{bootstrap-lemma-representable-by-spaces-cover}
appliqué à $U \to X \to Y$.

\medskip\noindent
Supposons (1) vraie et considérons un diagramme comme en (3).
Alors $U \to Y$ est représentable par des espaces algébriques
(en tant que composé $U \to X \to Y$ ; voir
Bootstrap, Lemma \ref{bootstrap-lemma-composition-transformation})
et se factorise par $V$. Ainsi, $U \to V$ est représentable par des
espaces algébriques d'après le Lemme \ref{lemma-permanence-representable}.

\medskip\noindent
Il est clair que (3) implique (2). Les conditions (1) -- (3) sont donc équivalentes.

\medskip\noindent
Observons que la condition (4) a un sens, car le produit fibré
$Y_j \times_Y X$ est un espace algébrique formel d'après le
Lemme \ref{lemma-fibre-products}.
Il est clair que (4) implique (5).

\medskip\noindent
Supposons $X_i \to Y_i \to Y$ comme en (5). Posons alors
$V = \coprod Y_i$ et $U = \coprod X_i$ ; on voit que
(5) implique (2).

\medskip\noindent
Enfin, supposons (1) -- (3) vraies.
Nous pouvons donc choisir un recouvrement quelconque $\{Y_j \to Y\}$ comme dans la
Définition \ref{definition-formal-algebraic-space}
et, pour chaque $j$, un recouvrement quelconque $\{X_{ji} \to Y_j \times_Y X\}$ comme dans la
Définition \ref{definition-formal-algebraic-space}.
Alors $X_{ij} \to Y_j$ est représentable par des espaces algébriques d'après (3),
et nous voyons que (4) est vraie. Ceci achève la démonstration.
\end{proof}

\begin{lemma}
\label{lemma-algebraic-space-over-affine-formal}
Soit $S$ un schéma. Soit $Y$ un espace algébrique formel affine sur $S$.
Soit $f : X \to Y$ une application de faisceaux sur $(\Sch/S)_{fppf}$ qui est
représentable par des espaces algébriques. Alors $X$ est un espace
algébrique formel.
\end{lemma}

\begin{proof}
Écrivons $Y = \colim Y_\lambda$ comme dans la
Définition \ref{definition-affine-formal-algebraic-space}.
Pour chaque $\lambda$, le produit fibré
$X \times_Y Y_\lambda$ est un espace algébrique.
Ainsi, $X = \colim X \times_Y Y_\lambda$ est un espace
algébrique formel d'après le Lemme \ref{lemma-colimit-is-formal}.
\end{proof}

\begin{lemma}
\label{lemma-representable-by-algebraic-spaces}
Soit $S$ un schéma. Soit $Y$ un espace algébrique formel sur $S$.
Soit $f : X \to Y$ une application de faisceaux sur $(\Sch/S)_{fppf}$ qui est
représentable par des espaces algébriques. Alors $X$ est un espace
algébrique formel.
\end{lemma}

\begin{proof}
Soit $\{Y_i \to Y\}$ comme dans la
Définition \ref{definition-formal-algebraic-space}.
Alors $X \times_Y Y_i \to X$ est une famille de morphismes
représentables par des espaces algébriques, étales et formant une famille
surjective. Il suffit donc de montrer que
$X \times_Y Y_i$ est un espace algébrique formel ; voir le
Lemme \ref{lemma-etale-covering-by-formal-algebraic-spaces}.
Cela résulte du Lemme \ref{lemma-algebraic-space-over-affine-formal}.
\end{proof}

\begin{lemma}
\label{lemma-affine-representable-by-algebraic-spaces}
Soit $S$ un schéma. Soit $f : X \to Y$ un morphisme
d'espaces algébriques formels affines qui est représentable par des
espaces algébriques. Alors $f$ est représentable (par des schémas) et affine.
\end{lemma}

\begin{proof}
Nous allons montrer que $f$ est affine ; il s'ensuivra alors que
$f$ est représentable et affine d'après
Morphisms of Spaces, Lemma \ref{spaces-morphisms-lemma-affine-local}.
Écrivons $Y = \colim Y_\mu$ et $X = \colim X_\lambda$ comme dans la
Définition \ref{definition-affine-formal-algebraic-space}.
Soit $T \to Y$ un morphisme, où $T$ est un schéma
sur $S$. Nous devons montrer que $X \times_Y T \to T$ est affine ; voir
Bootstrap, Definition
\ref{bootstrap-definition-property-transformation}.
Pour ce faire, nous pouvons supposer $T$ affine et devons démontrer
que $X \times_Y T$ est affine. Dans ce cas, $T \to Y$ se factorise
par $Y_\mu \to Y$ pour un certain $\mu$ ; voir le
Lemme \ref{lemma-factor-through-thickening}.
Puisque $f$ est quasi-compact, nous voyons que $X \times_Y T$ est
quasi-compact (Lemme \ref{lemma-characterize-quasi-compact-morphism}).
Ainsi, $X \times_Y T \to X$ se factorise par $X_\lambda$ pour un certain
$\lambda$. De même, $X_\lambda \to Y$ se factorise par $Y_\mu$
après avoir augmenté $\mu$. Alors
$X \times_Y T = X_\lambda \times_{Y_\mu} T$.
Nous concluons, car les produits fibrés de schémas affines sont affines.
\end{proof}

\begin{lemma}
\label{lemma-representable-affine}
Soit $S$ un schéma. Soit $\varphi : A \to B$ une application continue
d'anneaux topologiques faiblement admissibles sur $S$. Les conditions suivantes
sont équivalentes :
\begin{enumerate}
\item $\text{Spf}(\varphi) : \text{Spf}(B) \to \text{Spf}(A)$
est représentable par des espaces algébriques ;
\item $\text{Spf}(\varphi) : \text{Spf}(B) \to \text{Spf}(A)$
est représentable (par des schémas) ;
\item $\varphi$ est tendu ; voir la Définition \ref{definition-taut}.
\end{enumerate}
\end{lemma}

\begin{proof}
Les assertions (1) et (2) sont équivalentes d'après le
Lemma \ref{lemma-affine-representable-by-algebraic-spaces}.

\medskip\noindent
Supposons que les conditions équivalentes (1) et (2) soient vérifiées.
Si $I \subset A$ est un idéal faible de définition, alors
$\Spec(A/I) \to \text{Spf}(A)$ est représentable et est un épaississement
(cela ressort de la construction du spectre formel,
mais résulte aussi du
Lemma \ref{lemma-mcquillan-affine-formal-algebraic-space}).
Alors $\Spec(A/I) \times_{\text{Spf}(A)} \text{Spf}(B) \to \text{Spf}(B)$
est représentable et est un épaississement en tant que changement de base.
Ainsi, d'après le
Lemma \ref{lemma-mcquillan-affine-formal-algebraic-space}
il existe un idéal faible de définition $J(I) \subset B$ tel que
$\Spec(A/I) \times_{\text{Spf}(A)} \text{Spf}(B) = \Spec(B/J(I))$
comme sous-foncteurs de $\text{Spf}(B)$. Nous obtenons un diagramme cartésien
$$
\xymatrix{
\Spec(B/J(I)) \ar[d] \ar[r] & \Spec(A/I) \ar[d] \\
\text{Spf}(B) \ar[r] & \text{Spf}(A)
}
$$
Le Lemme \ref{lemma-fibre-product-affines-over-separated}
montre que $B/J(I) = B \widehat{\otimes}_A A/I$.
Il s'ensuit que $J(I)$ est l'adhérence de l'idéal $\varphi(I)B$ ; voir le
Lemma \ref{lemma-closure-image-ideal}.
Puisque $\text{Spf}(A) = \colim \Spec(A/I)$, où $I$ est comme ci-dessus,
nous obtenons $\text{Spf}(B) = \colim \Spec(B/J(I))$.
Ainsi, les idéaux $J(I)$ forment un système fondamental d'idéaux faibles
de définition (voir le
Lemma \ref{lemma-mcquillan-affine-formal-algebraic-space}).
La condition (3) est donc vérifiée.

\medskip\noindent
Supposons (3) vérifiée. Nous allons essentiellement reprendre en sens inverse les
arguments du paragraphe précédent.
Soit $I \subset A$ un idéal faible de définition.
D'après le Lemme \ref{lemma-fibre-product-affines-over-separated},
nous obtenons un diagramme cartésien
$$
\xymatrix{
\text{Spf}(B \widehat{\otimes}_A A/I) \ar[d] \ar[r] & \Spec(A/I) \ar[d] \\
\text{Spf}(B) \ar[r] & \text{Spf}(A)
}
$$
Si $J(I)$ est l'adhérence de $IB$, alors $J(I)$ est ouvert dans $B$
par le caractère tendu de $\varphi$. Ainsi, si $J$ est ouvert dans $B$ et $J \subset J(B)$,
alors $B/J \otimes_A A/I = B/(IB + J) = B/J(I)$, car
$J(I) = \bigcap_{J \subset B\text{ ouvert}} (IB + J)$ d'après le Lemme \ref{lemma-closed}.
Par conséquent, la limite définissant le produit tensoriel complété se stabilise et donne
$B \widehat{\otimes}_A A/I = B/J(I)$.
Ainsi, $\text{Spf}(B \widehat{\otimes}_A A/I) = \Spec(B/J(I))$.
Ceci prouve que $\text{Spf}(B) \times_{\text{Spf}(A)} \Spec(A/I)$
est représentable pour tout idéal faible de définition $I \subset A$.
Puisque tout morphisme $T \to \text{Spf}(A)$ avec $T$ quasi-compact
se factorise par $\Spec(A/I)$ pour un certain idéal faible de définition $I$
(Lemma \ref{lemma-factor-through-thickening})
nous concluons que $\text{Spf}(\varphi)$ est représentable, c'est-à-dire que
(2) est vérifiée. Ceci achève la démonstration.
\end{proof}

\begin{lemma}
\label{lemma-property-goes-up-affine-morphism}
Soit $S$ un schéma. Soit $Y$ un espace algébrique formel affine.
Soit $f : X \to Y$ une application de faisceaux sur $(\Sch/S)_{fppf}$ qui
est représentable et affine. Alors
\begin{enumerate}
\item $X$ est un espace algébrique formel affine ;
\item si $Y$ est à indexation dénombrable, alors $X$ l'est aussi ;
\item si $Y$ est à indexation dénombrable et classique, alors $X$ est
à indexation dénombrable et classique ;
\item si $Y$ est faiblement adique, alors $X$ est faiblement adique ;
\item si $Y$ est adique*, alors $X$ est adique* ;
\item si $Y$ est noethérien et si $f$ est (localement) de type fini, alors
$X$ est noethérien.
\end{enumerate}
\end{lemma}

\begin{proof}
Preuve de (1). Écrivons $Y = \colim_{\lambda \in \Lambda} Y_\lambda$ comme dans la
Définition \ref{definition-affine-formal-algebraic-space}.
Puisque $f$ est représentable et affine, les produits fibrés
$X_\lambda = Y_\lambda \times_Y X$ sont affines. De plus,
$X = \colim Y_\lambda \times_Y X$.
Ainsi, $X$ est un espace algébrique formel affine.

\medskip\noindent
Preuve de (2). Si $Y$ est à indexation dénombrable, alors, dans l'argument précédent,
nous pouvons supposer $\Lambda$ dénombrable.
Nous voyons aussitôt que $X$ est également à indexation dénombrable.

\medskip\noindent
Preuve de (3), (4) et (5). Dans chacun de ces cas, les hypothèses impliquent
que $Y$ est un espace algébrique formel affine à indexation dénombrable
(Lemma \ref{lemma-implications-between-types})
et donc que $X$ l'est aussi d'après (2). Nous pouvons ainsi écrire $X = \text{Spf}(A)$
et $Y = \text{Spf}(B)$ pour des $S$-algèbres topologiques faiblement
admissibles $A$ et $B$ ; voir le Lemme \ref{lemma-countably-indexed}.
D'après le Lemme \ref{lemma-morphism-between-formal-spectra},
le morphisme $f$ correspond à un homomorphisme continu de $S$-algèbres
$\varphi : B \to A$. Le
Lemme \ref{lemma-representable-affine} montre que $\varphi$ est tendu. Nous concluons que
(3) résulte du Lemme \ref{lemma-taut-ascent-admissible},
(4) résulte du Lemme \ref{lemma-taut-ascent-weakly-adic}, et
(5) résulte du Lemme \ref{lemma-taut-is-adic}.

\medskip\noindent
Preuve de (6). En combinant (3) avec le Lemme \ref{lemma-implications-between-types},
nous voyons que $X$ est adique*. Nous pouvons donc utiliser le critère du
Lemme \ref{lemma-characterize-noetherian-affine}.
Il nous dit d'abord que les schémas affines $Y_\lambda$ sont noethériens.
Puis $X_\lambda \to Y_\lambda$ est de type fini, donc $X_\lambda$
est également noethérien (Morphisms, Lemma
\ref{morphisms-lemma-finite-type-noetherian}).
Le critère nous dit alors que $X$ est noethérien, ce qui achève la
démonstration.
\end{proof}

\begin{lemma}
\label{lemma-property-goes-up-affine}
Soit $S$ un schéma. Soit $f : X \to Y$ un morphisme d'espaces algébriques
formels affines, représentable par des espaces algébriques. Alors
\begin{enumerate}
\item si $Y$ est à indexation dénombrable, alors $X$ l'est aussi ;
\item si $Y$ est à indexation dénombrable et classique, alors $X$ est à indexation
dénombrable et classique ;
\item si $Y$ est faiblement adique, alors $X$ est faiblement adique ;
\item si $Y$ est adique*, alors $X$ est adique* ;
\item si $Y$ est noethérien et si $f$ est (localement) de type fini, alors
$X$ est noethérien.
\end{enumerate}
\end{lemma}

\begin{proof}
Combiner les Lemmes \ref{lemma-affine-representable-by-algebraic-spaces} et
\ref{lemma-property-goes-up-affine-morphism}.
\end{proof}

\begin{example}
\label{example-representable-morphism-from-completion}
Soit $B$ un anneau topologique faiblement admissible. Soit $B \to A$
un homomorphisme d'anneaux (sans topologie). Nous pouvons alors considérer
$$
A^\wedge = \lim A/JA
$$
où la limite porte sur tous les idéaux faibles de définition $J$ de $B$.
Alors $A^\wedge$ (muni de la topologie limite) est un anneau
complet linéairement topologisé. Le noyau (ouvert) $I$
de la surjection $A^\wedge \to A/JA$ est l'adhérence de $JA^\wedge$ ; voir le
Lemme \ref{lemma-closed}. D'après le
Lemma \ref{lemma-topologically-nilpotent}
nous voyons que $I$ est formé d'éléments topologiquement nilpotents.
Ainsi, $I$ est un idéal faible de définition de $A^\wedge$, et nous concluons que
$A^\wedge$ est un anneau topologique faiblement admissible. Par conséquent,
$\varphi : B \to A^\wedge$ est une application tendue d'anneaux topologiques
faiblement admissibles et
$$
\text{Spf}(A^\wedge) \longrightarrow \text{Spf}(B)
$$
est un cas particulier du phénomène étudié dans le
Lemme \ref{lemma-representable-affine}.
\end{example}

\begin{remark}[Avertissement]
\label{remark-warning}
La discussion des Lemmes \ref{lemma-representable-affine},
\ref{lemma-property-goes-up-affine-morphism} et
\ref{lemma-property-goes-up-affine}
est optimale aux deux sens suivants :
\begin{enumerate}
\item Si $A$ et $B$ sont des anneaux faiblement admissibles et si $\varphi : A \to B$
est une application continue, alors
$\text{Spf}(\varphi) : \text{Spf}(B) \to \text{Spf}(A)$ n'est en général
pas représentable ;
\item si $f : Y \to X$ est un morphisme représentable d'espaces algébriques
formels affines et si $X = \text{Spf}(A)$ est de McQuillan,
il n'en résulte pas que $Y$ soit de McQuillan.
\end{enumerate}
Un exemple pour (1) consiste à prendre $A = k$, un corps (muni de la topologie discrète),
et $B = k[[t]]$ muni de la topologie $t$-adique.
Un exemple pour (2) est donné dans
Examples, Section \ref{examples-section-affine-formal-algebraic-space}.
\end{remark}

\noindent
Malgré l'avertissement précédent, nous avons le résultat suivant.

\begin{lemma}
\label{lemma-etale}
Soit $S$ un schéma. Soit $Y$ un espace algébrique formel affine de McQuillan
sur $S$, c'est-à-dire $Y = \text{Spf}(B)$ pour une $S$-algèbre topologique
faiblement admissible $B$. Alors il existe une équivalence de catégories entre
\begin{enumerate}
\item la catégorie des morphismes $f : X \to Y$
d'espaces algébriques formels affines qui sont représentables
par des espaces algébriques et étales ;
\item la catégorie des $B$-algèbres topologiques de la forme
$A^\wedge$, où $A$ est une $B$-algèbre étale et où
$A^\wedge = \lim A/JA$, avec $J \subset B$ parcourant les
idéaux faibles de définition de $B$.
\end{enumerate}
L'équivalence envoie $A^\wedge$ sur $X = \text{Spf}(A^\wedge)$.
En particulier, tout $X$ comme en (1) est de McQuillan.
\end{lemma}

\begin{proof}
Soit $A$ une $B$-algèbre étale. Alors $B/J \to A/JA$ est
étale pour tout idéal ouvert $J \subset B$. Par conséquent,
le morphisme $\text{Spf}(A^\wedge) \to Y$ est représentable et étale.
Le foncteur $\text{Spf}$ est pleinement fidèle d'après le
Lemma \ref{lemma-morphism-between-formal-spectra}.
Pour achever la démonstration, nous allons montrer au paragraphe suivant
que tout $X \to Y$ comme en (1) appartient à l'image essentielle.

\medskip\noindent
Choisissons un idéal faible de définition $J_0 \subset B$. Posons
$Y_0 = \Spec(B/J_0)$ et $X_0 = Y_0 \times_Y X$. Alors $X_0 \to Y_0$
est un morphisme étale de schémas affines (voir le
Lemma \ref{lemma-affine-representable-by-algebraic-spaces}).
Écrivons $X_0 = \Spec(A_0)$. D'après Algebra, Lemma \ref{algebra-lemma-lift-etale},
nous pouvons trouver un homomorphisme étale d'algèbres $B \to A$ tel que
$A_0 \cong A/J_0A$. Considérons un idéal de définition $J \subset J_0$.
Comme ci-dessus, nous pouvons écrire $\Spec(B/J) \times_Y X = \Spec(\bar A)$
pour un homomorphisme étale d'anneaux $B/J \to \bar A$. Alors
$B/J \to \bar A$ et $B/J \to A/JA$ sont tous deux des homomorphismes étales d'anneaux
relevant l'homomorphisme étale d'anneaux $B/J_0 \to A_0$. D'après
More on Algebra, Lemma \ref{more-algebra-lemma-locally-nilpotent-henselian}
il existe un unique isomorphisme de $B/J$-algèbres
$\varphi_J : A/JA \to \bar A$ relevant l'identification modulo $J_0$.
Comme les applications $\varphi_J$ sont uniques, elles sont compatibles lorsque $J$ varie.
Ainsi,
$$
X = \colim \Spec(B/J) \times_Y X = \colim \Spec(A/JA) = \text{Spf}(A)
$$
et nous voyons que le lemme est démontré.
\end{proof}

\begin{lemma}
\label{lemma-etale-surjective}
Avec les notations et hypothèses du Lemme \ref{lemma-etale}, soit
$f : X \to Y$ correspondant à $B \to A^\wedge$. Les conditions suivantes sont équivalentes :
\begin{enumerate}
\item $f : X \to Y$ est surjectif ;
\item $B \to A$ est fidèlement plat ;
\item pour tout idéal faible de définition $J \subset B$,
l'homomorphisme d'anneaux $B/J \to A/JA$ est fidèlement plat ;
\item pour un certain idéal faible de définition $J \subset B$,
l'homomorphisme d'anneaux $B/J \to A/JA$ est fidèlement plat.
\end{enumerate}
\end{lemma}

\begin{proof}
Soit $J \subset B$ un idéal faible de définition. Comme tout élément de $J$
est topologiquement nilpotent, nous voyons que tout élément de $1 + J$ est
inversible. Il s'ensuit que $J$ est contenu dans le radical de Jacobson de $B$
(Algebra, Lemma \ref{algebra-lemma-contained-in-radical}).
Ainsi, un homomorphisme plat d'anneaux $B \to A$ est fidèlement plat si et seulement si
$B/J \to A/JA$ est fidèlement plat
(Algebra, Lemma \ref{algebra-lemma-ff-rings}).
Nous voyons ainsi que (2) -- (4) sont équivalentes.
Si (1) est vérifiée, alors, pour tout idéal faible de définition $J \subset B$,
le morphisme
$\Spec(A/JA) = \Spec(B/J) \times_Y X \to \Spec(B/J)$ est surjectif,
ce qui implique (3). Réciproquement, supposons (3).
Un morphisme $T \to Y$ avec $T$ quasi-compact
se factorise par $\Spec(B/J)$ pour un certain idéal de définition $J$ de $B$
(Lemma \ref{lemma-factor-through-thickening}).
Ainsi, $X \times_Y T = \Spec(A/JA) \times_{\Spec(B/J)} T \to T$
est surjectif en tant que changement de base du morphisme surjectif
$\Spec(A/JA) \to \Spec(B/J)$. La condition (1) est donc vérifiée.
\end{proof}





\section{Types d'espaces algébriques formels}
\label{section-types}

\noindent
Dans cette section, nous définissons
les espaces algébriques formels
« localement noethériens »,
« localement adiques* »,
« localement faiblement adiques »,
« localement à indexation dénombrable et classiques », et
« localement à indexation dénombrable ».
Les types « localement adiques »,
« localement classiques » et
« localement de McQuillan »
manquent, car nous ne savons pas démontrer l'analogue des lemmes suivants
dans ces cas (il suffirait de démontrer l'analogue de ces
lemmes pour les recouvrements étales entre espaces algébriques formels affines).

\begin{lemma}
\label{lemma-iff-countably-indexed}
Soit $S$ un schéma. Soit $X \to Y$ un morphisme d'espaces algébriques
formels affines qui est représentable par des espaces algébriques,
surjectif et plat. Alors $X$ est à indexation dénombrable si et seulement
si $Y$ est à indexation dénombrable.
\end{lemma}

\begin{proof}
Supposons $X$ à indexation dénombrable. Écrivons $X = \colim X_n$ comme dans le
Lemme \ref{lemma-countable-affine-formal-algebraic-space}.
Écrivons $Y = \colim Y_\lambda$ comme dans la
Définition \ref{definition-affine-formal-algebraic-space}.
Pour tout $n$, nous pouvons choisir un $\lambda_n$ tel que
$X_n \to Y$ se factorise par $Y_{\lambda_n}$ ; voir le
Lemma \ref{lemma-factor-through-thickening}.
D'autre part, pour tout $\lambda$, le schéma
$Y_\lambda \times_Y X$ est affine
(Lemme \ref{lemma-affine-representable-by-algebraic-spaces})
et, par conséquent, $Y_\lambda \times_Y X \to X$ se factorise par
$X_n$ pour un certain $n$ (Lemme \ref{lemma-factor-through-thickening}).
Considérons le diagramme
$$
\xymatrix{
Y_\lambda \times_Y X \ar[r] \ar[d] & X_n \ar[r] \ar[d] & X \ar[d] \\
Y_\lambda \ar@{..>}[r] \ar@/_1pc/[rr] & Y_{\lambda_n} \ar[r] & Y
}
$$
Si nous montrons que la flèche pointillée existe, nous concluons que
$Y = \colim Y_{\lambda_n}$ et que $Y$ est à indexation dénombrable. Pour cela, choisissons
un $\mu$ tel que $\mu \geq \lambda$ et $\mu \geq \lambda_n$.
Ainsi, $Y_\lambda \to Y$ et $Y_{\lambda_n} \to Y$ se factorisent tous deux
par $Y_\mu \to Y$.
Écrivons $Y_\mu = \Spec(B_\mu)$ ; le sous-schéma fermé $Y_\lambda$ correspond à
$J \subset B_\mu$, et le sous-schéma fermé $Y_{\lambda_n}$ correspond à
$J' \subset B_\mu$. Nous voulons montrer que $J' \subset J$.
Le diagramme précédent donne $J'A_\mu \subset JA_\mu$,
où $Y_\mu \times_Y X = \Spec(A_\mu)$.
Puisque $X \to Y$ est surjectif et plat, le morphisme
$Y_\lambda \times_Y X \to Y_\lambda$ est un morphisme fidèlement plat
de schémas affines ; ainsi, $B_\mu \to A_\mu$ est
fidèlement plat. Donc $J' \subset J$, comme voulu.

\medskip\noindent
Supposons $Y$ à indexation dénombrable. Alors $X$ est à indexation dénombrable
d'après le Lemme \ref{lemma-property-goes-up-affine}.
\end{proof}

\begin{lemma}
\label{lemma-iff-cic}
Soit $S$ un schéma. Soit $X \to Y$ un morphisme d'espaces algébriques
formels affines qui est représentable par des espaces algébriques,
surjectif et plat. Alors $X$ est à indexation dénombrable et classique
si et seulement si $Y$ est à indexation dénombrable et classique.
\end{lemma}

\begin{proof}
Nous avons déjà vu l'implication dans un sens au
Lemme \ref{lemma-property-goes-up-affine}. Dans l'autre sens,
notons que, d'après le Lemme \ref{lemma-iff-countably-indexed}, nous pouvons supposer
$X$ et $Y$ tous deux à indexation dénombrable. Ainsi, $X = \text{Spf}(A)$
et $Y = \text{Spf}(B)$ pour des $S$-algèbres topologiques faiblement
admissibles $A$ et $B$ ; voir le Lemme \ref{lemma-countably-indexed}.
D'après le Lemme \ref{lemma-morphism-between-formal-spectra},
le morphisme $X \to Y$ correspond à un homomorphisme continu de $S$-algèbres
$\varphi : B \to A$. Le
Lemme \ref{lemma-representable-affine} montre que $\varphi$ est tendu.
Soit $J \subset B$ un idéal ouvert et soit $I \subset A$
l'adhérence de $JA$. D'après les
Lemmes \ref{lemma-fibre-product-affines-over-separated} et
\ref{lemma-closure-image-ideal}
nous voyons que $\Spec(B/J) \times_Y X = \Spec(A/I)$.
Ainsi, $B/J \to A/I$ est fidèlement plat (puisque $X \to Y$ est surjectif
et plat). Cela signifie que $\varphi : B \to A$ est comme dans la
Section \ref{section-taut-descent} (les rôles de $A$ et $B$ étant échangés).
Le lemme résulte donc du
Lemme \ref{lemma-taut-descent-weakly-admissible}.
\end{proof}

\begin{lemma}
\label{lemma-iff-weakly-adic}
Soit $S$ un schéma. Soit $X \to Y$ un morphisme d'espaces algébriques
formels affines qui est représentable par des espaces algébriques,
surjectif et plat. Alors $X$ est faiblement adique
si et seulement si $Y$ est faiblement adique.
\end{lemma}

\begin{proof}
La démonstration est exactement la même que celle du
Lemme \ref{lemma-iff-cic}, sauf qu'à la fin on utilise le
Lemme \ref{lemma-taut-descent-weakly-adic}.
\end{proof}

\begin{lemma}
\label{lemma-iff-adic-star}
Soit $S$ un schéma. Soit $X \to Y$ un morphisme d'espaces algébriques
formels affines qui est représentable par des espaces algébriques,
surjectif et plat. Alors $X$ est adique* si et seulement si $Y$ est adique*.
\end{lemma}

\begin{proof}
La démonstration est exactement la même que celle du
Lemme \ref{lemma-iff-cic}, sauf qu'à la fin on utilise le
Lemme \ref{lemma-taut-descent-adic-star}.
\end{proof}

\begin{lemma}
\label{lemma-iff-noetherian}
Soit $S$ un schéma. Soit $X \to Y$ un morphisme d'espaces algébriques
formels affines qui est représentable par des espaces algébriques,
surjectif, plat et (localement) de type fini. Alors $X$ est noethérien
si et seulement si $Y$ est noethérien.
\end{lemma}

\begin{proof}
Observons qu'un espace algébrique formel affine noethérien est adique* ; voir le
Lemme \ref{lemma-implications-between-types}. Ainsi, d'après le
Lemme \ref{lemma-iff-adic-star}, nous pouvons supposer que $X$ et $Y$
sont tous deux adiques*. Nous utiliserons le critère du
Lemme \ref{lemma-characterize-noetherian-affine}
pour établir le lemme. Écrivons $Y = \colim Y_n$
comme dans le Lemme \ref{lemma-countable-affine-formal-algebraic-space}.
Pour tout $n$, posons $X_n = Y_n \times_Y X$. Alors $X_n$ est un
schéma affine (Lemme \ref{lemma-affine-representable-by-algebraic-spaces})
et $X = \colim X_n$. Chacun des morphismes $X_n \to Y_n$ est
fidèlement plat et de type fini. Le lemme résulte donc du
fait que, dans cette situation, $X_n$ est noethérien si et seulement si $Y_n$
est noethérien ; voir
Algebra, Lemma \ref{algebra-lemma-descent-Noetherian} (pour la descente)
et
Algebra, Lemma \ref{algebra-lemma-Noetherian-permanence} (pour la montée).
\end{proof}

\begin{lemma}
\label{lemma-type-local}
Soit $S$ un schéma. Soit
$$
P \in
\left\{
\begin{matrix}
countably\ indexed,\\
countably\ indexed\ and\ classical,\\
weakly\ adic,\ adic*,\ Noetherian
\end{matrix}
\right\}
$$
Soit $X$ un espace algébrique formel sur $S$.
Les conditions suivantes sont équivalentes :
\begin{enumerate}
\item si $Y$ est un espace algébrique formel affine et si
$f : Y \to X$ est représentable par des espaces algébriques et étale,
alors $Y$ possède la propriété $P$ ;
\item pour un certain $\{X_i \to X\}_{i \in I}$ comme dans la
Définition \ref{definition-formal-algebraic-space},
chaque $X_i$ possède la propriété $P$.
\end{enumerate}
\end{lemma}

\begin{proof}
Il est clair que (1) implique (2). Supposons (2) et soit
$Y \to X$ comme en (1). Puisque les produits fibrés $X_i \times_X Y$
sont des espaces algébriques formels (Lemme \ref{lemma-fibre-products-general}),
nous pouvons choisir des recouvrements $\{X_{ij} \to X_i \times_X Y\}$ comme dans la
Définition \ref{definition-formal-algebraic-space}.
Puisque $Y$ est quasi-compact, il existe
$(i_1, j_1), \ldots, (i_n, j_n)$ tels que
$$
X_{i_1 j_1} \amalg \ldots \amalg X_{i_n j_n} \longrightarrow Y
$$
soit surjectif et étale. Alors $X_{i_kj_k} \to X_{i_k}$ est représentable
par des espaces algébriques et étale ; ainsi, $X_{i_kj_k}$ possède la propriété $P$ d'après le
Lemme \ref{lemma-property-goes-up-affine}.
Alors $X_{i_1 j_1} \amalg \ldots \amalg X_{i_n j_n}$ est un
espace algébrique formel affine possédant la propriété $P$ (nous omettons un petit détail
concernant les réunions disjointes finies d'espaces algébriques formels affines).
Nous concluons donc en appliquant l'un des
Lemmas \ref{lemma-iff-countably-indexed},
\ref{lemma-iff-cic},
\ref{lemma-iff-weakly-adic},
\ref{lemma-iff-adic-star}, et
\ref{lemma-iff-noetherian}.
\end{proof}

\noindent
Le lemme précédent permet de poser la définition suivante.

\begin{definition}
\label{definition-types-formal-algebraic-spaces}
Soit $S$ un schéma. Soit $X$ un espace algébrique formel sur $S$.
On dit que $X$ est
{\it localement à indexation dénombrable},
{\it localement à indexation dénombrable et classique},
{\it localement faiblement adique},
{\it localement adique*}, ou
{\it localement noethérien}
si les conditions équivalentes du Lemme \ref{lemma-type-local}
sont satisfaites pour la propriété correspondante.
\end{definition}

\noindent
Le complété formel d'un espace algébrique localement noethérien
le long d'une partie fermée est un espace algébrique formel localement noethérien.

\begin{lemma}
\label{lemma-formal-completion-types}
Soit $S$ un schéma. Soit $X$ un espace algébrique sur $S$.
Soit $T \subset |X|$ une partie fermée. Soit $X_{/T}$ le
complété formel de $X$ le long de $T$.
\begin{enumerate}
\item Si $X \setminus T \to X$ est quasi-compact,
alors $X_{/T}$ est localement adique* ;
\item si $X$ est localement noethérien, alors $X_{/T}$ est localement
noethérien.
\end{enumerate}
\end{lemma}

\begin{proof}
Choisissons un morphisme étale surjectif $U \to X$ avec $U = \coprod U_i$,
réunion disjointe de schémas affines ; voir Properties of Spaces, Lemma
\ref{spaces-properties-lemma-cover-by-union-affines}.
Soit $T_i \subset U_i$ l'image réciproque de $T$.
Nous avons $X_{/T} \times_X U_i = (U_i)_{/T_i}$
(Lemme \ref{lemma-map-completions-representable}).
Ainsi, $\{(U_i)_{/T_i} \to X_{/T}\}$ est un recouvrement comme dans la
Définition \ref{definition-formal-algebraic-space}.
De plus, si $X \setminus T \to X$ est quasi-compact, alors
$U_i \setminus T_i \to U_i$ l'est aussi, et si $X$ est localement noethérien,
$U_i$ l'est aussi. Le lemme résulte donc du cas affine, qui est le
Lemme \ref{lemma-affine-formal-completion-types}.
\end{proof}

\begin{remark}[Avertissement]
\label{remark-warning-completion}
Supposons $X = \Spec(A)$ et que $T \subset X$ soit le lieu des zéros d'un
idéal de type fini $I \subset A$. Soit $J = \sqrt{I}$ le
radical de $I$. Les définitions montrent alors que
$X_{/T} = \text{Spf}(A^\wedge)$, où $A^\wedge = \lim A/I^n$ est
le complété $I$-adique de $A$. D'autre part, l'application
$A^\wedge \to \lim A/J^n$ du complété $I$-adique vers
le complété $J$-adique peut ne pas être un isomorphisme d'anneaux.
À titre d'exemple, posons
$$
A = \bigcup\nolimits_{n \geq 1} \mathbf{C}[t^{1/n}]
$$
et $I = (t)$. Alors $J = \mathfrak m$ est l'idéal maximal
de l'anneau de valuation $A$ et $J^2 = J$. Ainsi, le complété $J$-adique
de $A$ est $\mathbf{C}$, tandis que le complété $I$-adique
est l'anneau de valuation décrit dans l'Exemple \ref{example-david-hansen}
(mais il est notamment facile de voir que $A \subset A^\wedge$).
\end{remark}

\begin{lemma}
\label{lemma-types-fibre-products}
Soit $S$ un schéma. Soient $X \to Y$ et $Z \to Y$ des
morphismes d'espaces algébriques formels sur $S$. Alors
\begin{enumerate}
\item Si $X$ et $Z$ sont localement à indexation dénombrable, alors $X \times_Y Z$
l'est aussi ;
\item si $X$ et $Z$ sont localement à indexation dénombrable et classiques,
alors $X \times_Y Z$ l'est aussi et est classique ;
\item si $X$ et $Z$ sont faiblement adiques, alors $X \times_Y Z$
est faiblement adique ;
\item si $X$ et $Z$ sont localement adiques*, alors $X \times_Y Z$ est
localement adique* ;
\item si $X$ et $Z$ sont localement noethériens et si $X_{red} \to Y_{red}$
est localement de type fini, alors $X \times_Y Z$ est localement noethérien.
\end{enumerate}
\end{lemma}

\begin{proof}
Choisissons un recouvrement $\{Y_j \to Y\}$ comme dans la
Définition \ref{definition-formal-algebraic-space}.
Pour chaque $j$, choisissons un recouvrement $\{X_{ji} \to Y_j \times_Y X\}$
comme dans la Définition \ref{definition-formal-algebraic-space}.
Pour chaque $j$, choisissons un recouvrement $\{Z_{jk} \to Y_j \times_Y Z\}$
comme dans la Définition \ref{definition-formal-algebraic-space}.
Observons que $X_{ji} \times_{Y_j} Z_{jk}$ est un
espace algébrique formel affine d'après le
Lemme \ref{lemma-fibre-product-affines-over-separated}.
Ainsi,
$$
\{X_{ji} \times_{Y_j} Z_{jk} \to X \times_Y Z\}
$$
est un recouvrement comme dans la Définition \ref{definition-formal-algebraic-space}.
Il suffit donc de prouver (1), (2), (3) et (4) lorsque $X$, $Y$ et $Z$
sont des espaces algébriques formels affines.

\medskip\noindent
Supposons $X$ et $Z$ à indexation dénombrable. Écrivons $X = \colim X_n$ et
$Z = \colim Z_m$ comme dans le
Lemme \ref{lemma-countable-affine-formal-algebraic-space}.
Écrivons $Y = \colim_{\lambda \in \Lambda} Y_\lambda$ comme dans la
Définition \ref{definition-affine-formal-algebraic-space}.
Pour tous $n$ et $m$, nous pouvons trouver $\lambda_{n, m} \in \Lambda$
tel que $X_n \to Y$ et $Z_m \to Y$ se factorisent par $Y_{\lambda_{n, m}}$
(voir par exemple le Lemme \ref{lemma-factor-through-thickening}).
Choisissons $\lambda_0 \in \Lambda$. Par récurrence sur $t \geq 1$,
choisissons un élément $\lambda_t \in \Lambda$ tel que
$\lambda_t \geq \lambda_{n, m}$ pour tous $1 \leq n, m \leq t$
et $\lambda_t \geq \lambda_{t - 1}$. Posons $Y' = \colim Y_{\lambda_t}$.
Alors $Y' \to Y$ est un monomorphisme tel que $X \to Y$ et $Z \to Y$
se factorisent par $Y'$. Nous pouvons donc remplacer $Y$ par $Y'$, c'est-à-dire
supposer $Y$ à indexation dénombrable.

\medskip\noindent
Supposons $X$, $Y$ et $Z$ à indexation dénombrable. D'après le
Lemme \ref{lemma-countably-indexed}, nous pouvons écrire
$X = \text{Spf}(A)$, $Y = \text{Spf}(B)$, $Z = \text{Spf}(C)$
pour des anneaux topologiques faiblement admissibles $A$, $B$ et $C$.
Les morphismes $X \to Y$ et $Z \to Y$ sont donnés par des
homomorphismes continus d'anneaux $B \to A$ et $B \to C$ ; voir le
Lemme \ref{lemma-morphism-between-formal-spectra}.
D'après le Lemme \ref{lemma-fibre-product-affines-over-separated},
nous voyons que $X \times_Y Z = \text{Spf}(A \widehat{\otimes}_B C)$
et que $A \widehat{\otimes}_B C$ est un anneau topologique faiblement admissible.
En particulier, nous voyons que $X \times_Y Z$ est à indexation dénombrable
d'après la partie (3) du Lemme \ref{lemma-completed-tensor-product}.
Ceci démontre (1).

\medskip\noindent
Preuve de (2). Dans ce cas, $X$ et $Z$ sont à indexation dénombrable ;
les arguments précédents montrent donc que $X \times_Y Z$ est le spectre formel
de $A \widehat{\otimes}_B C$, où $A$ et $C$ sont admissibles.
Alors $A \widehat{\otimes}_B C$ est admissible d'après la
partie (2) du Lemme \ref{lemma-completed-tensor-product}.

\medskip\noindent
Preuve de (3). Comme précédemment, nous concluons que $X \times_Y Z$ est le spectre formel
de $A \widehat{\otimes}_B C$, où $A$ et $C$ sont faiblement adiques.
Alors $A \widehat{\otimes}_B C$ est faiblement adique d'après le
Lemme \ref{lemma-completed-tensor-product-weakly-adic}.

\medskip\noindent
Preuve de (4). En raisonnant comme ci-dessus, cela résulte de la
partie (4) du Lemme \ref{lemma-completed-tensor-product}.

\medskip\noindent
Preuve de (5). Pour déduire le cas (5) de la
partie (5) du Lemme \ref{lemma-completed-tensor-product},
nous devons vérifier que les hypothèses coïncident.
Avec les notations du premier paragraphe de la démonstration,
si $X_{red} \to Y_{red}$ est localement de type fini, alors
$(X_{ji})_{red} \to (Y_j)_{red}$ est localement de type fini. Cela résulte de
Morphisms of Spaces, Lemma \ref{spaces-morphisms-lemma-finite-type-local}
et du fait que, dans le diagramme commutatif
$$
\xymatrix{
(X_{ji})_{red} \ar[d] \ar[r] & (Y_j)_{red} \ar[d] \\
X_{red} \ar[r] & Y_{red}
}
$$
les morphismes verticaux sont étales. En effet, nous avons
$(X_{ji})_{red} = X_{ij} \times_X X_{red}$ et
$(Y_j)_{red} = Y_j \times_Y Y_{red}$
par le Lemme \ref{lemma-reduction-smooth}. Comme ci-dessus, nous nous ramenons donc
au cas où $X$, $Y$, $Z$ sont des espaces algébriques formels affines,
où $X$, $Z$ sont noethériens et où $X_{red} \to Y_{red}$ est de type fini.
Ensuite, au deuxième paragraphe de la démonstration, nous avons remplacé $Y$ par $Y'$ ;
mais, par construction, $Y_{red} = Y'_{red}$, donc l'hypothèse de type fini
est préservée par ce remplacement. Nous voyons alors que
$X, Y, Z$ correspondent à $A, B, C$, et $X \times_Y Z$ à
$A \widehat{\otimes}_B C$, avec $A$, $C$ noethériens adiques.
Enfin, prendre la réduction revient à quotienter par
l'idéal des éléments topologiquement nilpotents
(Example \ref{example-reduction-affine-formal-spectrum})
et le fait que $X_{red} \to Y_{red}$ soit de type fini
signifie bien que $B/\mathfrak b \to A/\mathfrak a$ est de type fini.
Ceci achève la démonstration.
\end{proof}

\begin{lemma}
\label{lemma-structure-locally-noetherian}
Soit $S$ un schéma. Soit $X$ un espace algébrique formel localement noethérien
sur $S$. Alors $X = \colim X_n$ pour un système $X_1 \to X_2 \to X_3 \to \ldots$
d'épaississements d'ordre fini d'espaces algébriques localement noethériens sur $S$,
où $X_1 = X_{red}$ et où $X_n$ est le $n$-ième voisinage infinitésimal de
$X_1$ dans $X_m$ pour tout $m \geq n$.
\end{lemma}

\begin{proof}
Nous ne faisons qu'esquisser la démonstration et omettons certains détails.
Posons $X_1 = X_{red}$. Définissons $X_n \subset X$ comme le sous-foncteur
caractérisé par la règle suivante : un morphisme $f : T \to X$, où $T$ est un schéma, se factorise
par $X_n$ si et seulement si la $n$-ième puissance du faisceau d'idéaux
de l'immersion fermée $X_1 \times_X T \to T$ est nulle. Alors $X_n \subset X$
est un sous-faisceau, car l'annulation des modules quasi-cohérents se vérifie
localement pour la topologie fppf. Nous affirmons que $X_n \to X$ est représentable par des schémas,
est une immersion fermée et que $X = \colim X_n$ (comme faisceaux fppf).
Pour le vérifier, nous pouvons travailler localement pour la topologie étale sur $X$. Nous pouvons donc supposer que
$X = \text{Spf}(A)$ est un espace algébrique formel affine noethérien.
Alors $X_1 = \Spec(A/\mathfrak a)$, où $\mathfrak a \subset A$
est l'idéal des éléments topologiquement nilpotents de l'anneau topologique
adique noethérien $A$. Alors $X_n = \Spec(A/\mathfrak a^n)$,
ce qui donne le résultat voulu.
\end{proof}









\section{Morphismes et homomorphismes continus d'anneaux}
\label{section-morphisms-rings}

\noindent
Dans cette section, nous notons $\textit{WAdm}$ la catégorie des
anneaux topologiques faiblement admissibles et des homomorphismes continus d'anneaux.
Nous définissons les sous-catégories pleines
$$
\textit{WAdm} \supset
\textit{WAdm}^{count} \supset
\textit{WAdm}^{cic} \supset
\textit{WAdm}^{weakly\ adic} \supset
\textit{WAdm}^{adic*} \supset
\textit{WAdm}^{Noeth}
$$
dont les objets sont les suivants :
\begin{enumerate}
\item $\textit{WAdm}^{count}$ : les anneaux topologiques faiblement admissibles
$A$ qui possèdent un système fondamental dénombrable d'idéaux ouverts ;
\item $\textit{WAdm}^{cic}$ : les anneaux topologiques admissibles
$A$ qui possèdent un système fondamental dénombrable d'idéaux ouverts ;
\item $\textit{WAdm}^{weakly\ adic}$ : les anneaux topologiques faiblement adiques
(Section \ref{section-weakly-adic}) ;
\item $\textit{WAdm}^{adic*}$ : les anneaux topologiques adiques qui possèdent
un idéal de définition de type fini ;
\item $\textit{WAdm}^{Noeth}$ : les anneaux topologiques adiques qui
sont noethériens.
\end{enumerate}
Manifestement, les spectres formels de ces types d'anneaux sont les briques
élémentaires des espaces algébriques formels
localement à indexation dénombrable,
localement à indexation dénombrable et classiques,
localement faiblement adiques,
localement adiques* et
localement noethériens.

\medskip\noindent
Rappelons brièvement la relation entre les morphismes d'espaces algébriques
formels affines à indexation dénombrable et les
morphismes de $\textit{WAdm}^{count}$.
Soit $S$ un schéma. Soient $X$ et $Y$ des espaces algébriques formels
affines à indexation dénombrable. Écrivons $X = \text{Spf}(A)$
et $Y = \text{Spf}(B)$ pour des $S$-algèbres topologiques
$A$ et $B$ de $\textit{WAdm}^{count}$ ; voir le
Lemme \ref{lemma-countably-indexed}.
D'après le Lemme \ref{lemma-morphism-between-formal-spectra},
il existe une correspondance biunivoque entre les morphismes
$f : X \to Y$ et les homomorphismes continus
$$
\varphi : B \longrightarrow A
$$
de $S$-algèbres topologiques. Cette correspondance est donnée par
$f \mapsto f^\sharp$ et $\varphi \mapsto \text{Spf}(\varphi)$.

\medskip\noindent
Soit $S$ un schéma. Soit $f : X \to Y$ un morphisme d'espaces algébriques
formels localement à indexation dénombrable. Considérons un
diagramme commutatif
$$
\xymatrix{
U \ar[d] \ar[r] & V \ar[d] \\
X \ar[r] & Y
}
$$
où $U$ et $V$ sont des espaces algébriques formels affines et où $U \to X$ et $V \to Y$
sont représentables par des espaces algébriques et étales. D'après la
Définition \ref{definition-types-formal-algebraic-spaces} (donc aussi le
Lemme \ref{lemma-type-local}), $U$ et $V$ sont des espaces algébriques formels
affines à indexation dénombrable. La discussion du paragraphe
précédent montre que $U \to V$ est isomorphe à $\text{Spf}(\varphi)$
pour un certain homomorphisme continu
$$
\varphi : B \longrightarrow A
$$
de $S$-algèbres topologiques dans $\textit{WAdm}^{count}$.

\begin{lemma}
\label{lemma-completion-in-sub}
Soit $A \in \Ob(\textit{WAdm})$. Soit $A \to A'$ un homomorphisme
d'anneaux (sans topologie). Soit $(A')^\wedge = \lim_{I \subset A\text{ w.i.d}} A'/IA'$
l'objet de $\textit{WAdm}$ construit dans
l'Exemple \ref{example-representable-morphism-from-completion}.
\begin{enumerate}
\item Si $A$ appartient à $\textit{WAdm}^{count}$, il en est de même de $(A')^\wedge$.
\item Si $A$ appartient à $\textit{WAdm}^{cic}$, il en est de même de $(A')^\wedge$.
\item Si $A$ appartient à $\textit{WAdm}^{weakly\ adic}$, il en est de même de $(A')^\wedge$.
\item Si $A$ appartient à $\textit{WAdm}^{adic*}$, il en est de même de $(A')^\wedge$.
\item Si $A$ appartient à $\textit{WAdm}^{Noeth}$ et si $A'$ est noethérien, alors
$(A')^\wedge$ appartient à $\textit{WAdm}^{Noeth}$.
\end{enumerate}
\end{lemma}

\begin{proof}
Rappelons que $A \to (A')^\wedge$ est tendu ; voir la discussion de
l'Exemple \ref{example-representable-morphism-from-completion}.
Les assertions (1), (2), (3) et (4) résultent donc des Lemmes
\ref{lemma-taut-ascent-countable},
\ref{lemma-taut-ascent-admissible},
\ref{lemma-taut-ascent-weakly-adic}, et
\ref{lemma-taut-is-adic}.
Enfin, supposons $A$ noethérien et adique.
D'après (4), $(A')^\wedge$ est adique.
Algebra, Lemma \ref{algebra-lemma-completion-Noetherian-Noetherian}
montre que $(A')^\wedge$ est noethérien. Ainsi, (5) est vérifié.
\end{proof}

\begin{situation}
\label{situation-local-property}
Soit $P$ une propriété des morphismes de $\textit{WAdm}^{count}$.
Considérons les diagrammes commutatifs
\begin{equation}
\label{equation-localize}
\vcenter{
\xymatrix{
A \ar[r] & (A')^\wedge \\
B \ar[r] \ar[u]^\varphi & (B')^\wedge \ar[u]_{\varphi'}
}
}
\end{equation}
qui satisfont aux conditions suivantes :
\begin{enumerate}
\item $A$ et $B$ sont des objets de $\textit{WAdm}^{count}$ ;
\item $A \to A'$ et $B \to B'$ sont des homomorphismes étales d'anneaux ;
\item $(A')^\wedge = \lim A'/IA'$, resp.\  $(B')^\wedge = \lim B'/JB'$
où $I \subset A$, resp.\ $J \subset B$,
parcourt les idéaux faibles de définition de $A$, resp.\ de $B$ ;
\item $\varphi : B \to A$ et $\varphi' : (B')^\wedge \to (A')^\wedge$
sont continus.
\end{enumerate}
D'après le Lemme \ref{lemma-completion-in-sub}, les anneaux topologiques
$(A')^\wedge$ et $(B')^\wedge$ sont des objets de $\textit{WAdm}^{count}$.
Nous disons que $P$ est une {\it propriété locale} si les axiomes suivants sont vérifiés :
\begin{enumerate}
\item
\label{item-axiom-1}
pour tout diagramme (\ref{equation-localize}), on a
$P(\varphi) \Rightarrow P(\varphi')$,
\item
\label{item-axiom-2}
pour tout diagramme (\ref{equation-localize}) où $A \to A'$
est fidèlement plat, on a
$P(\varphi') \Rightarrow P(\varphi)$,
\item
\label{item-axiom-3}
si $P(B \to A_i)$ pour $i = 1, \ldots, n$, alors
$P(B \to \prod_{i = 1, \ldots, n} A_i)$.
\end{enumerate}
L'axiome (\ref{item-axiom-3})
a un sens puisque $\textit{WAdm}^{count}$ possède les produits finis.
\end{situation}

\begin{lemma}
\label{lemma-property-defines-property-morphisms}
Soit $S$ un schéma. Soit $f : X \to Y$ un morphisme d'espaces algébriques formels
localement à indexation dénombrable sur $S$.
Soit $P$ une propriété locale des morphismes de $\textit{WAdm}^{count}$.
Les conditions suivantes sont équivalentes :
\begin{enumerate}
\item pour tout diagramme commutatif
$$
\xymatrix{
U \ar[d] \ar[r] & V \ar[d] \\
X \ar[r] & Y
}
$$
où $U$ et $V$ sont des espaces algébriques formels affines, où $U \to X$ et $V \to Y$
sont représentables par des espaces algébriques et étales, le morphisme $U \to V$
correspond à un morphisme de $\textit{WAdm}^{count}$ qui possède la propriété $P$ ;
\item il existe un recouvrement $\{Y_j \to Y\}$ comme dans la
Définition \ref{definition-formal-algebraic-space} et, pour tout $j$,
un recouvrement $\{X_{ji} \to Y_j \times_Y X\}$ comme dans la
Définition \ref{definition-formal-algebraic-space},
tel que chaque $X_{ji} \to Y_j$ corresponde
à un morphisme de $\textit{WAdm}^{count}$ qui possède la propriété $P$ ;
\item il existe un recouvrement $\{X_i \to X\}$ comme dans la
Définition \ref{definition-formal-algebraic-space}
et, pour tout $i$, une factorisation $X_i \to Y_i \to Y$ où $Y_i$
est un espace algébrique formel affine, où $Y_i \to Y$ est représentable
par des espaces algébriques et étale, et où $X_i \to Y_i$ correspond
à un morphisme de $\textit{WAdm}^{count}$ qui possède la propriété $P$.
\end{enumerate}
\end{lemma}

\begin{proof}
Il est clair que (1) implique (2) et que (2) implique (3).
Supposons $\{X_i \to X\}$ et $X_i \to Y_i \to Y$ comme en (3),
et donnons-nous un diagramme comme en (1).
Puisque $Y_i \times_Y V$ est un espace algébrique formel
(Lemme \ref{lemma-fibre-products-general}), nous pouvons choisir
des recouvrements $\{Y_{ij} \to Y_i \times_Y V\}$ comme dans la
Définition \ref{definition-formal-algebraic-space}.
Pour tout $(i, j)$, nous pouvons de même choisir des recouvrements
$\{X_{ijk} \to Y_{ij} \times_{Y_i} X_i \times_X U\}$
comme dans la Définition \ref{definition-formal-algebraic-space}.
Puisque $U$ est quasi-compact, nous pouvons choisir
$(i_1, j_1, k_1), \ldots, (i_n, j_n, k_n)$ tels que
$$
X_{i_1 j_1 k_1} \amalg \ldots \amalg X_{i_n j_n k_n} \longrightarrow U
$$
soit surjectif. Pour $s = 1, \ldots, n$, considérons le diagramme commutatif
$$
\xymatrix{
& & & X_{i_s j_s k_s} \ar[ld] \ar[d] \ar[rd] \\
X \ar[d] & X_{i_s} \ar[l] \ar[d] &
X_{i_s} \times_X U \ar[l] \ar[d] & Y_{i_s j_s} \ar[ld] \ar[rd] &
X_{i_s} \times_X U \ar[d] \ar[r] &
U \ar[d] \ar[r] & X \ar[d] \\
Y & Y_{i_s} \ar[l] &
Y_{i_s} \times_Y V \ar[l] & &
Y_{i_s} \times_Y V \ar[r] &
V \ar[r] & Y
}
$$
Convenons que $P$ est vérifiée pour un morphisme d'espaces algébriques formels
affines à indexation dénombrable si elle est vérifiée pour le
morphisme correspondant de $\textit{WAdm}^{count}$. Remarquons que les morphismes
$X_{i_s j_s k_s} \to X_{i_s}$, $Y_{i_s j_s} \to Y_{i_s}$
sont donnés par des complétés d'homomorphismes étales d'anneaux ; voir le Lemme \ref{lemma-etale}.
Nous voyons donc que $P(X_{i_s} \to Y_{i_s})$ implique
$P(X_{i_s j_s k_s} \to Y_{i_s j_s})$ par l'axiome (\ref{item-axiom-1}).
Remarquons que les morphismes $Y_{i_s j_s} \to V$ sont donnés par des complétés
d'homomorphismes étales d'anneaux (même lemme que précédemment).
En appliquant l'axiome (\ref{item-axiom-2}) au diagramme
$$
\xymatrix{
X_{i_s j_s k_s} \ar@{=}[r] \ar[d] & X_{i_s j_s k_s} \ar[d] \\
Y_{i_s j_s} \ar[r] & V
}
$$
(cela est permis puisque les identités sont des homomorphismes fidèlement plats d'anneaux),
nous concluons que $P(X_{i_s j_s k_s} \to V)$ est vérifiée.
L'axiome (\ref{item-axiom-3}) montre alors que
$P(\coprod_{s = 1, \ldots, n} X_{i_s j_s k_s} \to V)$ est vérifiée.
Comme le morphisme $\coprod X_{i_s j_s k_s} \to U$ est surjectif
par construction, le morphisme correspondant de $\textit{WAdm}^{count}$
est le complété d'un homomorphisme étale et fidèlement plat d'anneaux ; voir le
Lemme \ref{lemma-etale-surjective}.
Une dernière application de l'axiome (\ref{item-axiom-2})
(avec $B' = B$) implique que $P(U \to V)$ est vraie, comme voulu.
\end{proof}

\begin{remark}[Variante adique*]
\label{remark-variant-adic-star}
Soit $P$ une propriété des morphismes de $\textit{WAdm}^{adic*}$.
Nous disons que $P$ est une {\it propriété locale} si les axiomes
(\ref{item-axiom-1}), (\ref{item-axiom-2}), (\ref{item-axiom-3})
de la Situation \ref{situation-local-property}
sont vérifiés pour les morphismes de $\textit{WAdm}^{adic*}$. De la même manière,
nous obtenons une variante du Lemme \ref{lemma-property-defines-property-morphisms}
pour les morphismes entre espaces algébriques formels localement adiques* sur $S$.
\end{remark}

\begin{remark}[Variante noethérienne]
\label{remark-variant-Noetherian}
Soit $P$ une propriété des morphismes de $\textit{WAdm}^{Noeth}$.
Nous disons que $P$ est une {\it propriété locale} si les axiomes
(\ref{item-axiom-1}), (\ref{item-axiom-2}), (\ref{item-axiom-3}),
de la Situation \ref{situation-local-property}
sont vérifiés pour les morphismes de $\textit{WAdm}^{Noeth}$. De la même manière,
nous obtenons une variante du Lemme \ref{lemma-property-defines-property-morphisms}
pour les morphismes entre espaces algébriques formels localement noethériens sur $S$.
\end{remark}

\begin{situation}
\label{situation-base-change-local-property}
Soit $P$ une propriété locale des morphismes de $\textit{WAdm}^{count}$ ; voir la
Situation \ref{situation-local-property}. Nous disons que $P$ est {\it stable par
changement de base} si, étant donnés $B \to A$ et $B \to C$ dans $\textit{WAdm}^{count}$,
on a $P(B \to A) \Rightarrow P(C \to A \widehat{\otimes}_B C)$.
Cela a un sens puisque $A \widehat{\otimes}_B C$ est un objet de
$\textit{WAdm}^{count}$ d'après le Lemme \ref{lemma-completed-tensor-product}.
\end{situation}

\begin{lemma}
\label{lemma-base-change-property-morphisms}
Soit $S$ un schéma. Soit $P$ une propriété locale des morphismes de
$\textit{WAdm}^{count}$ qui est stable par changement de base.
Soient $f : X \to Y$ et $g : Z \to Y$ des morphismes d'espaces algébriques formels
localement à indexation dénombrable sur $S$. Si $f$ satisfait aux conditions équivalentes du
Lemme \ref{lemma-property-defines-property-morphisms},
alors il en est de même de $\text{pr}_2 : X \times_Y Z \to Z$.
\end{lemma}

\begin{proof}
Choisissons un recouvrement $\{Y_j \to Y\}$ comme dans la
Définition \ref{definition-formal-algebraic-space}.
Pour tout $j$, choisissons un recouvrement $\{X_{ji} \to Y_j \times_Y X\}$
comme dans la Définition \ref{definition-formal-algebraic-space}.
Pour tout $j$, choisissons un recouvrement $\{Z_{jk} \to Y_j \times_Y Z\}$
comme dans la Définition \ref{definition-formal-algebraic-space}.
Remarquons que $X_{ji} \times_{Y_j} Z_{jk}$ est un
espace algébrique formel affine à indexation dénombrable ; voir le
Lemme \ref{lemma-types-fibre-products}.
Nous voyons alors que
$$
\{X_{ji} \times_{Y_j} Z_{jk} \to X \times_Y Z\}
$$
est un recouvrement comme dans la Définition \ref{definition-formal-algebraic-space}.
De plus, les morphismes $X_{ji} \times_{Y_j} Z_{jk} \to Z$
se factorisent par $Z_{jk}$. Par hypothèse,
$X_{ji} \to Y_j$ correspond à un morphisme $B_j \to A_{ji}$ de
$\text{WAdm}^{count}$ qui possède la propriété $P$.
Les morphismes $Z_{jk} \to Y_j$ correspondent à des morphismes $B_j \to C_{jk}$ dans
$\text{WAdm}^{count}$. Puisque
$X_{ji} \times_{Y_j} Z_{jk} =
\text{Spf}(A_{ji} \widehat{\otimes}_{B_j} C_{jk})$
par le Lemme \ref{lemma-fibre-product-affines-over-separated},
il suffit de montrer que
$C_{jk} \to A_{ji} \widehat{\otimes}_{B_j} C_{jk}$
possède la propriété $P$, ce que garantit précisément
la stabilité de $P$ par changement de base.
\end{proof}

\begin{remark}[Variante adique*]
\label{remark-base-change-variant-adic-star}
Soit $P$ une propriété locale des morphismes de $\textit{WAdm}^{adic*}$ ; voir la
Remarque \ref{remark-variant-adic-star}. Nous disons que $P$ est {\it stable par
changement de base} si, étant donnés $B \to A$ et $B \to C$ dans $\textit{WAdm}^{adic*}$,
on a $P(B \to A) \Rightarrow P(C \to A \widehat{\otimes}_B C)$.
Cela a un sens puisque $A \widehat{\otimes}_B C$ est un objet de
$\textit{WAdm}^{adic*}$ d'après le Lemme \ref{lemma-completed-tensor-product}.
De la même manière, nous obtenons une variante du
Lemme \ref{lemma-base-change-property-morphisms}
pour les morphismes entre espaces algébriques formels localement adiques* sur $S$.
\end{remark}

\begin{remark}[Variante noethérienne]
\label{remark-base-change-variant-Noetherian}
Soit $P$ une propriété locale des morphismes de $\textit{WAdm}^{Noeth}$ ; voir la
Remarque \ref{remark-variant-Noetherian}. Nous disons que $P$ est
{\it stable par changement de base} si, étant donnés $B \to A$ et $B \to C$
dans $\textit{WAdm}^{Noeth}$, la propriété $P(B \to A)$
implique à la fois que $A \widehat{\otimes}_B C$ est adique
noethérien\footnote{Voir le Lemme \ref{lemma-completed-tensor-product}
pour un critère.} et que
$P(C \to A \widehat{\otimes}_B C)$.
De la même manière, nous obtenons une variante du
Lemme \ref{lemma-base-change-property-morphisms}
pour les morphismes entre espaces algébriques formels localement noethériens sur $S$.
\end{remark}

\begin{remark}[Autre variante noethérienne]
\label{remark-base-change-variant-variant-Noetherian}
Soient $P$ et $Q$ des propriétés locales des morphismes de
$\textit{WAdm}^{Noeth}$ ; voir la Remarque \ref{remark-variant-Noetherian}.
Nous disons que {\it $P$ est stable par changement de base par $Q$}
si, étant donnés $B \to A$ et $B \to C$ dans $\textit{WAdm}^{Noeth}$
qui satisfont à $P(B \to A)$ et $Q(B \to C)$, alors
$A \widehat{\otimes}_B C$ est adique noethérien et
$P(C \to A \widehat{\otimes}_B C)$ est vérifiée.
En raisonnant exactement comme dans la démonstration du
Lemme \ref{lemma-base-change-property-morphisms},
nous obtenons l'assertion suivante :
étant donnés des morphismes $f : X \to Y$ et $g : Y \to Z$ d'espaces algébriques
formels localement noethériens sur $S$
tels que
\begin{enumerate}
\item les conditions équivalentes du
Lemme \ref{lemma-property-defines-property-morphisms}
soient vérifiées pour $f$ et $P$ ;
\item les conditions équivalentes du
Lemme \ref{lemma-property-defines-property-morphisms}
soient vérifiées pour $g$ et $Q$ ;
\end{enumerate}
alors les conditions équivalentes du
Lemme \ref{lemma-property-defines-property-morphisms}
sont vérifiées pour $\text{pr}_2 : X \times_Y Z \to Z$ et $P$.
\end{remark}

\begin{situation}
\label{situation-composition-local-property}
Soit $P$ une propriété locale des morphismes de $\textit{WAdm}^{count}$ ; voir la
Situation \ref{situation-local-property}. Nous disons que $P$ est {\it stable par
composition} si, étant donnés $B \to A$ et $C \to B$ dans $\textit{WAdm}^{count}$,
on a $P(B \to A) \wedge P(C \to B) \Rightarrow P(C \to A)$.
\end{situation}

\begin{lemma}
\label{lemma-composition-property-morphisms}
Soit $S$ un schéma. Soit $P$ une propriété locale des morphismes de
$\textit{WAdm}^{count}$ qui est stable par composition.
Soient $f : X \to Y$ et $g : Y \to Z$ des morphismes d'espaces algébriques formels
localement à indexation dénombrable sur $S$. Si $f$ et $g$
satisfont aux conditions équivalentes du
Lemme \ref{lemma-property-defines-property-morphisms},
alors il en est de même de $g \circ f : X \to Z$.
\end{lemma}

\begin{proof}
Choisissons un recouvrement $\{Z_k \to Z\}$ comme dans la
Définition \ref{definition-formal-algebraic-space}.
Pour tout $k$, choisissons un recouvrement $\{Y_{kj} \to Z_k \times_Z Y\}$
comme dans la Définition \ref{definition-formal-algebraic-space}.
Pour tous $k$ et $j$, choisissons un recouvrement $\{X_{kji} \to Y_{kj} \times_Y X\}$
comme dans la Définition \ref{definition-formal-algebraic-space}.
Si $f$ et $g$
satisfont aux conditions équivalentes du
Lemme \ref{lemma-property-defines-property-morphisms},
alors $X_{kji} \to Y_{jk}$ et $Y_{jk} \to Z_k$
correspondent aux flèches
$B_{kj} \to A_{kji}$ et $C_k \to B_{kj}$ de
$\text{WAdm}^{count}$ qui possèdent la propriété $P$.
Il en est donc de même de leurs composées, ce qui conclut.
\end{proof}

\begin{remark}[Variante adique*]
\label{remark-composition-variant-adic-star}
Soit $P$ une propriété locale des morphismes de $\textit{WAdm}^{adic*}$ ; voir la
Remarque \ref{remark-variant-adic-star}. Nous disons que $P$ est {\it stable par
composition} si, étant donnés $B \to A$ et $C \to B$ dans $\textit{WAdm}^{adic*}$,
on a $P(B \to A) \wedge P(C \to B) \Rightarrow P(C \to A)$.
De la même manière, nous obtenons une variante du
Lemme \ref{lemma-composition-property-morphisms}
pour les morphismes entre espaces algébriques formels localement adiques* sur $S$.
\end{remark}

\begin{remark}[Variante noethérienne]
\label{remark-composition-variant-Noetherian}
Soit $P$ une propriété locale des morphismes de $\textit{WAdm}^{Noeth}$ ; voir la
Remarque \ref{remark-variant-Noetherian}. Nous disons que $P$ est
{\it stable par composition} si, étant donnés $B \to A$ et $C \to B$
dans $\textit{WAdm}^{Noeth}$, on a
$P(B \to A) \wedge P(C \to B) \Rightarrow P(C \to A)$.
De la même manière, nous obtenons une variante du
Lemme \ref{lemma-composition-property-morphisms}
pour les morphismes entre espaces algébriques formels localement noethériens sur $S$.
\end{remark}

\begin{situation}
\label{situation-permanence-local-property}
Soit $P$ une propriété locale des morphismes de $\textit{WAdm}^{count}$ ; voir la
Situation \ref{situation-local-property}. Nous disons que $P$
{\it possède la propriété de simplification}
si, étant donnés $B \to A$ et $C \to B$ dans $\textit{WAdm}^{count}$,
on a $P(C \to B) \wedge P(C \to A) \Rightarrow P(B \to A)$.
\end{situation}

\begin{lemma}
\label{lemma-permanence-property-morphisms}
Soit $S$ un schéma. Soit $P$ une propriété locale des morphismes de
$\textit{WAdm}^{count}$ qui possède la propriété de simplification.
Soient $f : X \to Y$ et $g : Y \to Z$ des morphismes d'espaces algébriques formels
localement à indexation dénombrable sur $S$. Si $g \circ f$ et $g$
satisfont aux conditions équivalentes du
Lemme \ref{lemma-property-defines-property-morphisms},
alors il en est de même de $f : X \to Y$.
\end{lemma}

\begin{proof}
Choisissons un recouvrement $\{Z_k \to Z\}$ comme dans la
Définition \ref{definition-formal-algebraic-space}.
Pour tout $k$, choisissons un recouvrement $\{Y_{kj} \to Z_k \times_Z Y\}$
comme dans la Définition \ref{definition-formal-algebraic-space}.
Pour tous $k$ et $j$, choisissons un recouvrement $\{X_{kji} \to Y_{kj} \times_Y X\}$
comme dans la Définition \ref{definition-formal-algebraic-space}.
Faisons correspondre à $X_{kji} \to Y_{jk}$ et $Y_{jk} \to Z_k$
les flèches
$B_{kj} \to A_{kji}$ et $C_k \to B_{kj}$ de
$\text{WAdm}^{count}$.
Si $g \circ f$ et $g$ satisfont aux conditions équivalentes du
Lemme \ref{lemma-property-defines-property-morphisms},
alors $C_k \to B_{kj}$ et $C_k \to A_{kji}$ satisfont à $P$.
Ainsi, $B_{kj} \to A_{kji}$ y satisfait aussi, ce qui conclut.
\end{proof}

\begin{remark}[Variante adique*]
\label{remark-permanence-variant-adic-star}
Soit $P$ une propriété locale des morphismes de $\textit{WAdm}^{adic*}$ ; voir la
Remarque \ref{remark-variant-adic-star}. Nous disons que $P$ {\it possède la propriété de
simplification} si, étant donnés $B \to A$ et $C \to B$ dans $\textit{WAdm}^{adic*}$,
on a $P(C \to A) \wedge P(C \to B) \Rightarrow P(B \to A)$.
De la même manière, nous obtenons une variante du
Lemme \ref{lemma-composition-property-morphisms}
pour les morphismes entre espaces algébriques formels localement adiques* sur $S$.
\end{remark}

\begin{remark}[Variante noethérienne]
\label{remark-permanence-variant-Noetherian}
Soit $P$ une propriété locale des morphismes de $\textit{WAdm}^{Noeth}$ ; voir la
Remarque \ref{remark-variant-Noetherian}. Nous disons que $P$
{\it possède la propriété de simplification} si, étant donnés $B \to A$ et $C \to B$
dans $\textit{WAdm}^{Noeth}$, on a
$P(C \to B) \wedge P(C \to A) \Rightarrow P(C \to B)$.
De la même manière, nous obtenons une variante du
Lemme \ref{lemma-composition-property-morphisms}
pour les morphismes entre espaces algébriques formels localement noethériens sur $S$.
\end{remark}














\section{Homomorphismes d'anneaux tendus et représentabilité par des espaces algébriques}
\label{section-taut}

\noindent
Dans cette section, nous montrons brièvement que les morphismes entre espaces algébriques
formels localement à indexation dénombrable correspondent, localement pour la topologie étale,
à des homomorphismes continus tendus entre anneaux topologiques faiblement
admissibles qui possèdent des systèmes fondamentaux dénombrables d'idéaux ouverts.
En fait, cela ressort assez clairement du
Lemme \ref{lemma-representable-affine},
et nous invitons le lecteur à sauter cette section.

\begin{lemma}
\label{lemma-representable-property-rings}
Soit $B \to A$ une flèche de $\textit{WAdm}^{count}$.
Les conditions suivantes sont équivalentes :
\begin{enumerate}
\item[(a)] $B \to A$ est tendu (Définition \ref{definition-taut}) ;
\item[(b)] pour $B \supset J_1 \supset J_2 \supset J_3 \supset \ldots$
étant un système fondamental d'idéaux faibles de définition, il existe
un diagramme commutatif
$$
\xymatrix{
A \ar[r] & \ldots \ar[r] & A_3 \ar[r] & A_2 \ar[r] & A_1 \\
B \ar[r] \ar[u] & \ldots \ar[r] & B/J_3 \ar[r] \ar[u] &
B/J_2 \ar[r] \ar[u] & B/J_1 \ar[u]
}
$$
tel que $A_{n + 1}/J_nA_{n + 1} = A_n$ et $A = \lim A_n$
comme anneau topologique.
\end{enumerate}
De plus, ces conditions équivalentes définissent une propriété locale,
c'est-à-dire qu'elles satisfont aux axiomes (\ref{item-axiom-1}), (\ref{item-axiom-2}),
(\ref{item-axiom-3}).
\end{lemma}

\begin{proof}
L'équivalence de (a) et (b) est immédiate. Nous donnerons ci-dessous une
démonstration algébrique des axiomes, mais il se trouve que nous les avons déjà
démontrés. En effet, d'après le Lemme \ref{lemma-representable-affine},
les conditions équivalentes (a) et (b) signifient que le
morphisme correspondant d'espaces algébriques formels affines est représentable
par des espaces algébriques.
Puisque cette condition est « locale sur la source et le but pour la topologie étale » d'après le
Lemme \ref{lemma-representable-by-algebraic-spaces-local},
nous obtenons immédiatement les axiomes (\ref{item-axiom-1}), (\ref{item-axiom-2}) et
(\ref{item-axiom-3}).

\medskip\noindent
Démonstration algébrique directe de (\ref{item-axiom-1}), (\ref{item-axiom-2}),
(\ref{item-axiom-3}). Donnons-nous un diagramme (\ref{equation-localize}) comme dans la
Situation \ref{situation-local-property}.
D'après l'Exemple \ref{example-representable-morphism-from-completion},
les homomorphismes $A \to (A')^\wedge$ et $B \to (B')^\wedge$
satisfont à (a) et (b).

\medskip\noindent
Supposons (a) et (b) vérifiées pour $\varphi$. Soit $J \subset B$ un idéal faible
de définition. Alors l'adhérence de $JA$, resp.\ de $J(B')^\wedge$,
est un idéal faible de définition $I \subset A$, resp.\ $J' \subset (B')^\wedge$.
L'adhérence de $I(A')^\wedge$ est alors un idéal faible de définition
$I' \subset (A')^\wedge$. Un argument topologique montre que $I'$ est aussi
l'adhérence de $J(A')^\wedge$ et de $J'(A')^\wedge$.
Enfin, lorsque $J$ parcourt un système fondamental d'idéaux faibles de définition
de $B$, les idéaux $I$ et $I'$ en font autant dans $A$ et $(A')^\wedge$.
Il s'ensuit que (a) est vérifiée pour $\varphi'$. Cela démontre (\ref{item-axiom-1}).

\medskip\noindent
Supposons $A \to A'$ fidèlement plat et (a) et (b) vérifiées pour $\varphi'$.
Soit $J \subset B$ un idéal faible de définition. En appliquant (a) et (b)
aux homomorphismes $B \to (B')^\wedge \to (A')^\wedge$, nous trouvons que
l'adhérence $I'$ de $J(A')^\wedge$ est un idéal faible de définition.
En particulier, $I'$ est ouvert, donc l'image réciproque de $I'$
dans $A$ est ouverte. Nous avons alors (explication ci-dessous)
\begin{align*}
A \cap I'
& =
A \cap \bigcap (J(A')^\wedge + \Ker((A')^\wedge \to A'/I_0A')) \\
& =
A \cap \bigcap \Ker((A')^\wedge \to A'/JA' + I_0 A') \\
& = \bigcap (JA + I_0)
\end{align*}
qui est l'adhérence de $JA$ d'après le Lemme \ref{lemma-closed}.
Les intersections portent sur les idéaux faibles de définition $I_0 \subset A$.
La première égalité vient du fait qu'un système fondamental de voisinages de
$0$ dans $(A')^\wedge$ est formé des noyaux des homomorphismes $(A')^\wedge \to A'/I_0A'$.
La deuxième égalité est triviale. La troisième vient du fait que $A \to A'$
est fidèlement plat ; voir
Algebra, Lemma \ref{algebra-lemma-faithfully-flat-universally-injective}.
Ainsi, l'adhérence de $JA$ est ouverte. D'après le Lemme \ref{lemma-topologically-nilpotent},
l'adhérence de $JA$
est un idéal faible de définition de $A$. Enfin, étant donné un idéal faible
de définition $I \subset A$, nous pouvons trouver $J$ tel que
$J(A')^\wedge$ soit contenu dans l'adhérence de $I(A')^\wedge$
d'après la propriété (a) pour $B \to (B')^\wedge$ et $\varphi'$.
Nous voyons donc que (a) est vérifiée pour $\varphi$. Cela démontre (\ref{item-axiom-2}).

\medskip\noindent
Nous omettons la démonstration de (\ref{item-axiom-3}).
\end{proof}

\begin{lemma}
\label{lemma-representable-local-property}
Soit $S$ un schéma. Soit $f : X \to Y$ un morphisme d'espaces algébriques formels
localement à indexation dénombrable sur $S$.
Les conditions suivantes sont équivalentes :
\begin{enumerate}
\item pour tout diagramme commutatif
$$
\xymatrix{
U \ar[d] \ar[r] & V \ar[d] \\
X \ar[r] & Y
}
$$
où $U$ et $V$ sont des espaces algébriques formels affines, où $U \to X$ et $V \to Y$
sont représentables par des espaces algébriques et étales, le morphisme $U \to V$
correspond à un homomorphisme tendu $B \to A$ de $\textit{WAdm}^{count}$ ;
\item il existe un recouvrement $\{Y_j \to Y\}$ comme dans la
Définition \ref{definition-formal-algebraic-space} et, pour tout $j$,
un recouvrement $\{X_{ji} \to Y_j \times_Y X\}$ comme dans la
Définition \ref{definition-formal-algebraic-space},
tel que chaque $X_{ji} \to Y_j$ corresponde
à un homomorphisme tendu d'anneaux dans $\textit{WAdm}^{count}$ ;
\item il existe un recouvrement $\{X_i \to X\}$ comme dans la
Définition \ref{definition-formal-algebraic-space}
et, pour tout $i$, une factorisation $X_i \to Y_i \to Y$ où $Y_i$
est un espace algébrique formel affine, où $Y_i \to Y$ est représentable
par des espaces algébriques et étale, et où $X_i \to Y_i$ correspond
à un homomorphisme tendu d'anneaux dans $\textit{WAdm}^{count}$ ;
\item $f$ est représentable par des espaces algébriques.
\end{enumerate}
\end{lemma}

\begin{proof}
Le fait qu'un homomorphisme de $\textit{WAdm}^{count}$ soit
« tendu » est une propriété locale d'après le
Lemme \ref{lemma-representable-property-rings}.
Le Lemme \ref{lemma-property-defines-property-morphisms}
nous dit donc exactement que (1), (2) et (3) sont équivalentes.
D'autre part, d'après le Lemme \ref{lemma-representable-affine},
le caractère « tendu » des homomorphismes de $\textit{WAdm}^{count}$ correspond exactement à la
« représentabilité par des espaces algébriques » des morphismes correspondants
d'espaces algébriques formels affines à indexation dénombrable.
Ainsi, l'implication (1) $\Rightarrow$ (2) du
Lemme \ref{lemma-representable-by-algebraic-spaces-local}
montre que (4) implique (1) du présent lemme.
De même, l'implication (4) $\Rightarrow$ (1) du
Lemme \ref{lemma-representable-by-algebraic-spaces-local}
montre que (2) implique (4) du présent lemme.
\end{proof}









\section{Morphismes adiques}
\label{section-adic}

\noindent
Cette section met en correspondance la notion parfois employée de
« morphisme adique » $f : X \to Y$ entre espaces algébriques formels localement adiques*
$X$ et $Y$, d'une part avec la représentabilité de $f$ par des
espaces algébriques et, d'autre part, avec notre notion d'
homomorphisme continu tendu d'anneaux. Rappelons d'abord que
le caractère tendu équivaut au caractère adique pour les anneaux adiques dont
l'idéal de définition est de type fini.

\begin{lemma}
\label{lemma-adic-homomorphism}
Soient $A$ et $B$ des anneaux topologiques préadiques. Soit
$\varphi : A \to B$ un homomorphisme continu d'anneaux.
\begin{enumerate}
\item Si $\varphi$ est adique, alors $\varphi$ est tendu.
\item Si $B$ est complet, si $A$ possède un idéal de définition
de type fini et si $\varphi$ est tendu, alors $\varphi$ est adique.
\end{enumerate}
En particulier, les conditions « $\varphi$ est adique » et « $\varphi$ est tendu »
sont équivalentes dans la catégorie $\textit{WAdm}^{adic*}$.
\end{lemma}

\begin{proof}
La partie (1) est le Lemme \ref{lemma-adic-taut}.
La partie (2) est le Lemme \ref{lemma-taut-is-adic}.
La dernière assertion résulte de (1) et (2).
\end{proof}

\noindent
Soit $S$ un schéma. Soit $f : X \to Y$ un morphisme d'espaces algébriques
formels localement adiques* sur $S$.
D'après le Lemme \ref{lemma-representable-local-property},
les conditions suivantes sont équivalentes :
\begin{enumerate}
\item $f$ est représentable par des espaces algébriques (autrement dit, les
conditions équivalentes du
Lemme \ref{lemma-representable-by-algebraic-spaces-local} sont vérifiées) ;
\item pour tout diagramme commutatif
$$
\xymatrix{
U \ar[d] \ar[r] & V \ar[d] \\
X \ar[r] & Y
}
$$
où $U$ et $V$ sont des espaces algébriques formels affines, où $U \to X$ et $V \to Y$
sont représentables par des espaces algébriques et étales, le morphisme $U \to V$
correspond à un homomorphisme adique\footnote{De manière équivalente, tendu d'après le
Lemme \ref{lemma-adic-homomorphism}.} dans $\textit{WAdm}^{adic*}$.
\end{enumerate}
Dans cette situation, nous dirons que {\it $f$ est un morphisme adique} (la définition
formelle figure ci-dessous). Cette notion ou terminologie ne sera définie et employée
que pour les morphismes entre espaces algébriques formels localement adiques*,
car, autrement, nous ne disposons pas de l'équivalence entre (1) et (2) ci-dessus.

\begin{definition}
\label{definition-adic-morphism}
Soit $S$ un schéma. Soit $f : X \to Y$ un morphisme d'espaces algébriques formels
sur $S$. Supposons $X$ et $Y$ localement adiques*. Nous disons que $f$ est
un {\it morphisme adique} si $f$ est représentable par des espaces algébriques.
Voir la discussion précédente.
\end{definition}











\section{Morphismes de type fini}
\label{section-finite-type}

\noindent
Compte tenu de l'organisation du projet Stacks, la définition suivante
est véritablement celle qui convient, et les notions plus fortes devraient
porter un autre nom.

\begin{definition}
\label{definition-finite-type}
Soit $S$ un schéma. Soit $f : Y \to X$ un morphisme d'espaces algébriques formels
sur $S$.
\begin{enumerate}
\item Nous disons que $f$ est {\it localement de type fini}
si $f$ est représentable par des espaces algébriques et est localement
de type fini au sens de
Bootstrap, Definition \ref{bootstrap-definition-property-transformation}.
\item Nous disons que $f$ est de {\it type fini} si $f$ est localement de type fini et
quasi-compact (Définition \ref{definition-quasi-compact-morphism}).
\end{enumerate}
\end{definition}

\noindent
Nous étudierons la relation entre les morphismes de type fini de
certains espaces algébriques formels et les
homomorphismes continus d'anneaux $A \to B$ qui sont topologiquement de type fini
à la Section \ref{section-tft}.

\begin{lemma}
\label{lemma-characterize-finite-type}
Soit $S$ un schéma. Soit $f : X \to Y$ un morphisme d'espaces algébriques formels
sur $S$. Les conditions suivantes sont équivalentes :
\begin{enumerate}
\item $f$ est de type fini ;
\item $f$ est représentable par des espaces algébriques et est de type fini au
sens de
Bootstrap, Definition \ref{bootstrap-definition-property-transformation}.
\end{enumerate}
\end{lemma}

\begin{proof}
Cela résulte de Bootstrap, Lemma
\ref{bootstrap-lemma-transformations-property-implication},
de l'implication « quasi-compact $+$ localement de type fini
$\Rightarrow$ de type fini » pour les morphismes d'espaces algébriques, et du
Lemme \ref{lemma-quasi-compact-representable}.
\end{proof}

\begin{lemma}
\label{lemma-composition-finite-type}
La composée de morphismes de type fini est de type fini.
Il en est de même de la propriété « localement de type fini ».
\end{lemma}

\begin{proof}
Voir Bootstrap, Lemma \ref{bootstrap-lemma-composition-transformation-property}
et utiliser Morphisms of Spaces, Lemma
\ref{spaces-morphisms-lemma-composition-finite-type}.
\end{proof}

\begin{lemma}
\label{lemma-base-change-finite-type}
Un changement de base d'un morphisme de type fini est de type fini.
Il en est de même de la propriété « localement de type fini ».
\end{lemma}

\begin{proof}
Voir Bootstrap, Lemma \ref{bootstrap-lemma-base-change-transformation-property}
et utiliser Morphisms of Spaces, Lemma
\ref{spaces-morphisms-lemma-base-change-finite-type}.
\end{proof}

\begin{lemma}
\label{lemma-permanence-finite-type}
Soit $S$ un schéma. Soient $f : X \to Y$ et $g : Y \to Z$ des morphismes
d'espaces algébriques formels sur $S$. Si $g \circ f : X \to Z$ est localement de
type fini, alors $f : X \to Y$ est localement de type fini.
\end{lemma}

\begin{proof}
Le Lemme \ref{lemma-permanence-representable} montre que $f$ est
représentable par des espaces algébriques. Soit $T$ un schéma et soit
$T \to Z$ un morphisme. Nous pouvons alors appliquer
Morphisms of Spaces, Lemma \ref{spaces-morphisms-lemma-permanence-finite-type}
aux morphismes
$T \times_Z X \to T \times_Z Y \to T$ d'espaces algébriques pour conclure.
\end{proof}

\noindent
La propriété d'être localement de type fini est locale sur la source et sur le but.

\begin{lemma}
\label{lemma-finite-type-local}
Soit $S$ un schéma. Soit $f : X \to Y$ un morphisme d'espaces algébriques formels
sur $S$. Les conditions suivantes sont équivalentes :
\begin{enumerate}
\item le morphisme $f$ est localement de type fini ;
\item il existe un diagramme commutatif
$$
\xymatrix{
U \ar[d] \ar[r] & V \ar[d] \\
X \ar[r] & Y
}
$$
où $U$, $V$ sont des espaces algébriques formels, où les flèches verticales sont
représentables par des espaces algébriques et étales, où $U \to X$
est surjectif et où $U \to V$ est localement de type fini ;
\item pour tout diagramme commutatif
$$
\xymatrix{
U \ar[d] \ar[r] & V \ar[d] \\
X \ar[r] & Y
}
$$
où $U$, $V$ sont des espaces algébriques formels et où les flèches verticales
sont représentables par des espaces algébriques et étales, le morphisme
$U \to V$ est localement de type fini ;
\item il existe un recouvrement $\{Y_j \to Y\}$ comme dans la
Définition \ref{definition-formal-algebraic-space}
et, pour tout $j$, un recouvrement $\{X_{ji} \to Y_j \times_Y X\}$ comme dans la
Définition \ref{definition-formal-algebraic-space}, tels que
$X_{ji} \to Y_j$ soit localement de type fini pour tous $j$ et $i$ ;
\item il existe un recouvrement $\{X_i \to X\}$ comme dans la
Définition \ref{definition-formal-algebraic-space}
et, pour tout $i$, une factorisation $X_i \to Y_i \to Y$ où $Y_i$
est un espace algébrique formel affine, où $Y_i \to Y$ est représentable
par des espaces algébriques et étale, telle que $X_i \to Y_i$ soit
localement de type fini ;
\item ajouter davantage ici.
\end{enumerate}
\end{lemma}

\begin{proof}
Dans chacun des cinq cas, le morphisme $f : X \to Y$ est représentable
par des espaces algébriques ; voir le
Lemme \ref{lemma-representable-by-algebraic-spaces-local}.
Nous utiliserons ce fait ci-dessous sans autre mention.

\medskip\noindent
Il est clair que (1) implique (2), car nous pouvons prendre $U = X$ et $V = Y$.
Réciproquement, donnons-nous un diagramme comme en (2). Soit $T$ un schéma et
soit $T \to Y$ un morphisme. Nous pouvons alors considérer
$$
\xymatrix{
U \times_Y T \ar[d] \ar[r] & V \times_Y T \ar[d] \\
X \times_Y T \ar[r] & T
}
$$
Les flèches verticales sont étales et la flèche horizontale supérieure
est localement de type fini comme changement de base d'un tel morphisme.
Par conséquent, Morphisms of Spaces, Lemma
\ref{spaces-morphisms-lemma-finite-type-local} montre que
$X \times_Y T \to T$ est localement de type fini. Autrement dit,
(1) est vérifiée.

\medskip\noindent
Supposons (1) vraie et considérons un diagramme comme en (3).
Alors $U \to Y$ est localement de type fini
(en tant que composée $U \to X \to Y$ ; voir
Bootstrap, Lemma \ref{bootstrap-lemma-composition-transformation-property}).
Soit $T$ un schéma et soit $T \to V$ un morphisme.
Alors la projection $T \times_V U \to T$ se factorise sous la forme
$$
T \times_V U = (T \times_Y U) \times_{(V \times_Y V)} V
\to T \times_Y U \to T
$$
La seconde flèche est localement de type fini (comme changement de base
de la composée $U \to X \to Y$), et
la première est le changement de base de la diagonale $V \to V \times_Y V$,
qui est localement de type fini d'après le
Lemme \ref{lemma-diagonal-morphism-formal-algebraic-spaces}.

\medskip\noindent
Il est clair que (3) implique (2). Ainsi, (1) à (3) sont équivalentes.

\medskip\noindent
Remarquons que la condition de (4) a un sens, puisque le produit fibré
$Y_j \times_Y X$ est un espace algébrique formel d'après le
Lemme \ref{lemma-fibre-products}.
Il est clair que (4) implique (5).

\medskip\noindent
Supposons $X_i \to Y_i \to Y$ comme en (5). Posons alors
$V = \coprod Y_i$ et $U = \coprod X_i$ pour voir que
(5) implique (2).

\medskip\noindent
Enfin, supposons (1) à (3) vraies.
Nous pouvons donc choisir un recouvrement quelconque $\{Y_j \to Y\}$ comme dans la
Définition \ref{definition-formal-algebraic-space}
et, pour tout $j$, un recouvrement quelconque $\{X_{ji} \to Y_j \times_Y X\}$ comme dans la
Définition \ref{definition-formal-algebraic-space}.
Alors $X_{ij} \to Y_j$ est localement de type fini d'après (3),
et nous voyons que (4) est vraie. Cela achève la démonstration.
\end{proof}

\begin{example}
\label{example-finite-type-from-finite-type-ring-map}
Soit $S$ un schéma. Soit $A$ un anneau topologique faiblement admissible sur $S$.
Soit $A \to A'$ un homomorphisme d'anneaux de type fini. Alors
$$
(A')^\wedge = \lim_{I \subset A\ w.i.d.} A'/IA'
$$
est un anneau faiblement admissible, et le morphisme correspondant
$\text{Spf}((A')^\wedge) \to \text{Spf}(A)$ est représentable ;
voir l'Exemple \ref{example-representable-morphism-from-completion}.
Si $T \to \text{Spf}(A)$ est un morphisme où $T$ est un schéma
quasi-compact, alors il se factorise par $\Spec(A/I)$ pour un certain idéal faible
de définition $I \subset A$ (Lemme \ref{lemma-factor-through-thickening}).
Alors $T \times_{\text{Spf}(A)} \text{Spf}((A')^\wedge)$
est égal à $T \times_{\Spec(A/I)} \Spec(A'/IA')$, et
nous voyons que $\text{Spf}((A')^\wedge) \to \text{Spf}(A)$ est
de type fini.
\end{example}

\begin{lemma}
\label{lemma-locally-finite-type-locally-noetherian}
Soit $S$ un schéma. Soit $f : X \to Y$ un morphisme d'espaces algébriques
formels sur $S$. Si $Y$ est localement noethérien et si
$f$ est localement de type fini, alors $X$ est localement noethérien.
\end{lemma}

\begin{proof}
Choisissons $\{Y_j \to Y\}$ et $\{X_{ij} \to Y_j \times_Y X\}$
comme dans le Lemme \ref{lemma-finite-type-local}. Il résulte alors
du Lemme \ref{lemma-property-goes-up-affine-morphism}
que chaque $X_{ij}$ est noethérien. Cela démontre le lemme.
\end{proof}

\begin{lemma}
\label{lemma-fibre-product-Noetherian}
Soit $S$ un schéma. Soient $f : X \to Y$ et $Z \to Y$ des morphismes
d'espaces algébriques formels sur $S$. Si $Z$ est localement
noethérien et si $f$ est localement de type fini, alors
$Z \times_Y X$ est localement noethérien.
\end{lemma}

\begin{proof}
Le morphisme $Z \times_Y X \to Z$ est localement de type fini d'après le
Lemme \ref{lemma-base-change-finite-type}. Le résultat découle donc
du Lemme \ref{lemma-locally-finite-type-locally-noetherian}.
\end{proof}








\section{Morphismes surjectifs}
\label{section-surjective}

\noindent
D'après le Lemme \ref{lemma-reduction-surjective}, la définition suivante n'entre
pas en conflit avec les définitions existantes pour les morphismes d'espaces
algébriques ou pour les morphismes d'espaces algébriques formels représentables
par des espaces algébriques.

\begin{definition}
\label{definition-surjective}
Soit $S$ un schéma. Un morphisme $f : X \to Y$ d'espaces algébriques formels
sur $S$ est dit {\it surjectif} s'il induit un morphisme surjectif
$X_{red} \to Y_{red}$ sur les espaces algébriques réduits sous-jacents.
\end{definition}

\begin{lemma}
\label{lemma-composition-surjective}
La composée de deux morphismes surjectifs est un morphisme surjectif.
\end{lemma}

\begin{proof}
Démonstration omise.
\end{proof}

\begin{lemma}
\label{lemma-base-change-surjective}
Un changement de base d'un morphisme surjectif est un morphisme surjectif.
\end{lemma}

\begin{proof}
Démonstration omise.
\end{proof}

\begin{lemma}
\label{lemma-characterize-surjective}
Soit $S$ un schéma. Soit $f : X \to Y$ un morphisme d'espaces algébriques formels
sur $S$. Les conditions suivantes sont équivalentes :
\begin{enumerate}
\item $f$ est surjectif ;
\item pour tout schéma $T$ et tout morphisme $T \to Y$,
la projection $X \times_Y T \to T$ est un morphisme surjectif
d'espaces algébriques formels ;
\item pour tout schéma affine $T$ et tout morphisme $T \to Y$,
la projection $X \times_Y T \to T$ est un morphisme surjectif d'espaces algébriques
formels ;
\item il existe un recouvrement $\{Y_j \to Y\}$ comme dans la
Définition \ref{definition-formal-algebraic-space},
tel que chaque $X \times_Y Y_j \to Y_j$ soit un morphisme surjectif
d'espaces algébriques formels ;
\item il existe un morphisme surjectif $Z \to Y$
d'espaces algébriques formels tel que $X \times_Y Z \to Z$ soit surjectif ;
\item ajouter davantage ici.
\end{enumerate}
\end{lemma}

\begin{proof}
Démonstration omise.
\end{proof}







\section{Monomorphismes}
\label{section-monomorphisms}

\noindent
Voici la définition.

\begin{definition}
\label{definition-monomorphism}
Soit $S$ un schéma.
Un morphisme d'espaces algébriques formels sur $S$ est appelé un
{\it monomorphisme} s'il est un morphisme injectif de faisceaux.
\end{definition}

\noindent
Voici un exemple. Soit $X$ un espace algébrique et soit
$T \subset |X|$ une partie fermée. Alors le morphisme
$X_{/T} \to X$ du complété formel de $X$ le long de $T$ vers $X$
est un monomorphisme. En particulier, les monomorphismes d'espaces algébriques
formels ne sont généralement pas représentables.

\begin{lemma}
\label{lemma-composition-monomorphism}
La composée de deux monomorphismes est un monomorphisme.
\end{lemma}

\begin{proof}
Démonstration omise.
\end{proof}

\begin{lemma}
\label{lemma-base-change-monomorphism}
Un changement de base d'un monomorphisme est un monomorphisme.
\end{lemma}

\begin{proof}
Démonstration omise.
\end{proof}

\begin{lemma}
\label{lemma-characterize-monomorphisms}
Soit $S$ un schéma. Soit $f : X \to Y$ un morphisme d'espaces algébriques formels
sur $S$. Les conditions suivantes sont équivalentes :
\begin{enumerate}
\item $f$ est un monomorphisme ;
\item pour tout schéma $T$ et tout morphisme $T \to Y$,
la projection $X \times_Y T \to T$ est un monomorphisme
d'espaces algébriques formels ;
\item pour tout schéma affine $T$ et tout morphisme $T \to Y$,
la projection $X \times_Y T \to T$ est un monomorphisme d'espaces algébriques
formels ;
\item il existe un recouvrement $\{Y_j \to Y\}$ comme dans la
Définition \ref{definition-formal-algebraic-space},
tel que chaque $X \times_Y Y_j \to Y_j$ soit un monomorphisme
d'espaces algébriques formels ;
\item il existe une famille de morphismes $\{Y_j \to Y\}$ telle
que $\coprod Y_j \to Y$ soit une surjection de faisceaux sur
$(\Sch/S)_{fppf}$ et que chaque $X \times_Y Y_j \to Y_j$ soit un
monomorphisme pour tout $j$ ;
\item il existe un morphisme $Z \to Y$ d'espaces algébriques formels
qui est représentable par des espaces algébriques, surjectif, plat et localement
de présentation finie, tel que $X \times_Y Z \to X$ soit un monomorphisme ;
\item ajouter davantage ici.
\end{enumerate}
\end{lemma}

\begin{proof}
Démonstration omise.
\end{proof}





\section{Immersions fermées}
\label{section-closed-immersions}

\noindent
Voici la définition.

\begin{definition}
\label{definition-closed-immersion}
Soit $S$ un schéma. Soit $f : Y \to X$ un morphisme d'espaces algébriques formels
sur $S$. Nous disons que $f$ est une {\it immersion fermée}
si $f$ est représentable par des espaces algébriques et est une immersion fermée
au sens de
Bootstrap, Definition \ref{bootstrap-definition-property-transformation}.
\end{definition}

\noindent
Lors de la lecture de cette section, veuillez sauter les lemmes préliminaires obligés.

\begin{lemma}
\label{lemma-composition-closed-immersion}
La composée de deux immersions fermées est une immersion fermée.
\end{lemma}

\begin{proof}
Démonstration omise.
\end{proof}

\begin{lemma}
\label{lemma-base-change-closed-immersion}
Un changement de base d'une immersion fermée est une immersion fermée.
\end{lemma}

\begin{proof}
Démonstration omise.
\end{proof}

\begin{lemma}
\label{lemma-characterize-closed-immersion}
Soit $S$ un schéma. Soit $f : X \to Y$ un morphisme d'espaces algébriques formels
sur $S$. Les conditions suivantes sont équivalentes :
\begin{enumerate}
\item $f$ est une immersion fermée ;
\item pour tout schéma $T$ et tout morphisme $T \to Y$, la projection
$X \times_Y T \to T$ est une immersion fermée ;
\item pour tout schéma affine $T$ et tout morphisme $T \to Y$,
la projection $X \times_Y T \to T$ est une immersion fermée ;
\item il existe un recouvrement $\{Y_j \to Y\}$ comme dans la
Définition \ref{definition-formal-algebraic-space},
tel que chaque $X \times_Y Y_j \to Y_j$ soit une immersion fermée ;
\item il existe un morphisme $Z \to Y$ d'espaces algébriques formels
qui est représentable par des espaces algébriques, surjectif, plat et localement
de présentation finie, tel que $X \times_Y Z \to X$ soit une
immersion fermée ;
\item ajouter davantage ici.
\end{enumerate}
\end{lemma}

\begin{proof}
Démonstration omise.
\end{proof}

\begin{lemma}
\label{lemma-closed-immersion-into-McQuillan}
Soit $S$ un schéma. Soit $X$ un espace algébrique formel affine de McQuillan
sur $S$. Soit $f : Y \to X$ une immersion fermée d'espaces algébriques formels
sur $S$. Alors $Y$ est un espace algébrique formel affine de McQuillan et
$f$ correspond à un homomorphisme continu $A \to B$ de $S$-algèbres topologiques
faiblement admissibles qui est tendu, dont le noyau est fermé et dont
l'image est dense.
\end{lemma}

\begin{proof}
Écrivons $X = \text{Spf}(A)$, où $A$ est un anneau topologique faiblement admissible.
Soit $I_\lambda$ un système fondamental d'idéaux faiblement admissibles
de définition de $A$. Alors $Y \times_X \Spec(A/I_\lambda)$ est
un sous-schéma fermé de $\Spec(A/I_\lambda)$,
donc affine (Définition \ref{definition-closed-immersion}).
Écrivons $Y \times_X \Spec(A/I_\lambda) = \Spec(B_\lambda)$.
L'homomorphisme d'anneaux $A/I_\lambda \to B_\lambda$
est surjectif. Par conséquent, les projections
$$
B = \lim B_\lambda \longrightarrow B_\lambda
$$
sont surjectives, puisque les composées $A \to B \to B_\lambda$ sont surjectives.
Il s'ensuit que $Y$ est de McQuillan d'après le
Lemme \ref{lemma-mcquillan-affine-formal-algebraic-space}.
L'homomorphisme d'anneaux $A \to B$ est tendu d'après le Lemme \ref{lemma-representable-affine}.
Le noyau est fermé parce que $B$ est complet et que $A \to B$ est
continu. Enfin, comme $A \to B_\lambda$ est surjectif pour tout $\lambda$,
nous voyons que l'image de $A$ dans $B$ est dense.
\end{proof}

\noindent
Malgré le résultat précédent, nous ne savons pas, en général, comment
se comportent les immersions fermées lorsque le but est un
espace algébrique formel affine de McQuillan ; voir la Remarque \ref{remark-questions}.

\begin{example}
\label{example-closed-immersion-from-quotient}
Soit $S$ un schéma. Soit $A$ un anneau topologique faiblement admissible sur $S$.
Soit $K \subset A$ un idéal fermé. En posant
$$
B = (A/K)^\wedge = \lim_{I \subset A\ w.i.d.} A/(I + K)
$$
le morphisme $\text{Spf}(B) \to \text{Spf}(A)$ est représentable ; voir
l'Exemple \ref{example-representable-morphism-from-completion}.
Si $T \to \text{Spf}(A)$ est un morphisme où $T$ est un schéma
quasi-compact, alors il se factorise par $\Spec(A/I)$ pour un certain idéal faible
de définition $I \subset A$ (Lemme \ref{lemma-factor-through-thickening}).
Alors $T \times_{\text{Spf}(A)} \text{Spf}(B)$
est égal à $T \times_{\Spec(A/I)} \Spec(A/(K + I))$, et
nous voyons que $\text{Spf}(B) \to \text{Spf}(A)$ est une immersion fermée.
Le noyau de $A \to B$ est $K$, puisque $K$ est fermé, mais
il faut prendre garde qu'en général l'homomorphisme d'anneaux $A \to B = (A/K)^\wedge$
n'est pas nécessairement surjectif.
\end{example}

\begin{lemma}
\label{lemma-monomorphism-iso-over-red}
Soit $S$ un schéma. Soit $f : X \to Y$ un morphisme d'espaces algébriques
formels. Supposons que
\begin{enumerate}
\item $f$ soit représentable par des espaces algébriques ;
\item $f$ soit un monomorphisme ;
\item l'inclusion $Y_{red} \to Y$ se factorise par $f$ ;
\item $f$ soit localement de type fini ou $Y$ soit localement noethérien.
\end{enumerate}
Alors $f$ est une immersion fermée.
\end{lemma}

\begin{proof}
Les hypothèses (2) et (3) impliquent que
$X_{red} = X \times_Y Y_{red} = Y_{red}$.
Nous utiliserons ce fait sans autre mention.

\medskip\noindent
Si $Y' \to Y$ est un morphisme étale d'espaces algébriques formels sur $S$,
alors le changement de base $f' : X \times_Y Y' \to Y'$ satisfait aux conditions
(1) à (4). D'après le Lemme \ref{lemma-characterize-closed-immersion},
nous pouvons donc supposer que $Y$ est un espace algébrique formel affine.

\medskip\noindent
Écrivons $Y = \colim_{\lambda \in \Lambda} Y_\lambda$ comme dans la
Définition \ref{definition-affine-formal-algebraic-space}.
Alors $X_\lambda = X \times_Y Y_\lambda$ est un espace algébrique
muni d'un monomorphisme $f_\lambda : X_\lambda \to Y_\lambda$
qui induit un isomorphisme $X_{\lambda, red} \to Y_{\lambda, red}$.
Ainsi, $X_\lambda$ est un schéma affine d'après
Limits of Spaces, Proposition \ref{spaces-limits-proposition-affine}
(puisque $X_{\lambda, red} \to X_\lambda$ est surjectif et entier).
Pour achever la démonstration, il suffit de montrer que
$X_\lambda \to Y_\lambda$ est une immersion fermée,
ce que nous ferons au paragraphe suivant.

\medskip\noindent
Soit $X \to Y$ un monomorphisme de schémas affines tel que
$X_{red} = X \times_Y Y_{red} = Y_{red}$. En général, cela
n'implique pas que $X \to Y$ soit une immersion fermée ; voir
Examples, Section \ref{examples-section-epimorphism-not-surjective}.
Cependant, sous notre hypothèse (4), nous savons qu'au paragraphe précédent
soit $X_\lambda \to Y_\lambda$ est de type fini, soit $Y_\lambda$
est noethérien. Cela signifie que
$X \to Y$ correspond à un homomorphisme d'anneaux $R \to A$ tel que
$R/I \to A/IA$ soit un isomorphisme, où $I \subset R$ est le
nilradical (c'est-à-dire l'idéal localement nilpotent maximal de $R$),
et que soit $R \to A$ est de type fini, soit $R$ est noethérien.
Dans le premier cas, $R \to A$ est surjectif d'après
Algebra, Lemma \ref{algebra-lemma-surjective-mod-locally-nilpotent},
et, dans le second, $I$ est de type fini, donc
nilpotent ; ainsi $R \to A$ est surjectif par le lemme de Nakayama ; voir
Algebra, Lemma \ref{algebra-lemma-NAK}, partie (11).
\end{proof}








\section{Séries restreintes}
\label{section-restricted-power-series}

\noindent
Soit $A$ un anneau topologique complet pour une
topologie linéaire (More on Algebra, Definition
\ref{more-algebra-definition-topological-ring}).
Soit $I_\lambda$ un système fondamental d'idéaux ouverts.
Soit $r \geq 0$ un entier. Dans ce contexte, on note souvent
par
$$
A\{x_1, \ldots, x_r\} =
\lim_\lambda A/I_\lambda[x_1, \ldots, x_r] =
\lim_\lambda (A[x_1, \ldots, x_r]/I_\lambda A[x_1, \ldots, x_r])
$$
muni de la topologie limite. Autrement dit, il s'agit
du complété de l'anneau de polynômes relativement aux
idéaux $I_\lambda$. On peut considérer les éléments de $A\{x_1, \ldots, x_r\}$ comme des
séries
$$
f = \sum\nolimits_{E = (e_1, \ldots, e_r)} a_E x_1^{e_1} \ldots x_r^{e_r}
$$
en $x_1, \ldots, x_r$ à coefficients $a_E \in A$ qui tendent
vers zéro dans la topologie de $A$. Autrement dit, pour tout $\lambda$,
tous les $a_E$, sauf un nombre fini, appartiennent à $I_\lambda$.
Pour cette raison, les éléments de $A\{x_1, \ldots, x_r\}$ sont parfois
appelés des {\it séries restreintes}.
On note parfois cet anneau $A\langle x_1, \ldots, x_r\rangle$ ;
nous nous abstiendrons d'employer cette notation.

\begin{remark}[Propriété universelle des séries restreintes]
\label{remark-universal-property}
\begin{reference}
\cite[Chapter 0, 7.5.3]{EGA}
\end{reference}
Soit $A \to C$ un homomorphisme continu d'anneaux complets munis d'une topologie linéaire.
Alors tout homomorphisme de $A$-algèbres $A[x_1, \ldots x_r] \to C$ se prolonge de manière unique en un
homomorphisme continu $A\{x_1, \ldots, x_r\} \to C$ sur les séries restreintes.
\end{remark}

\begin{remark}
\label{remark-I-adic-completion-and-restricted-power-series}
Soit $A$ un anneau et soit $I \subset A$ un idéal. Si $A$ est complet pour la topologie
$I$-adique, alors le complété $I$-adique $A[x_1, \ldots, x_r]^\wedge$ de
$A[x_1, \ldots, x_r]$ est, comme anneau, l'anneau des séries restreintes sur $A$.
Cependant, il n'est pas clair que $A[x_1, \ldots, x_r]^\wedge$ soit complet pour la topologie
$I$-adique. Nous considérons la topologie sur $A\{x_1, \ldots, x_r\}$
comme la topologie limite (qui est toujours complète), tandis que nous considérons souvent
la topologie sur $A[x_1, \ldots, x_r]^\wedge$ comme la topologie $I$-adique
(qui n'est pas toujours complète). Si $I$ est de type fini, alors
$A\{x_1, \ldots, x_r\} = A[x_1, \ldots, x_r]^\wedge$ comme anneaux topologiques ;
voir Algebra, Lemma \ref{algebra-lemma-hathat-finitely-generated}.
\end{remark}





\section{Algèbres topologiquement de type fini}
\label{section-tft}

\noindent
Voici notre définition. Celle-ci ne fait pas l'objet d'un consensus général.
De nombreux auteurs imposent des conditions supplémentaires, souvent parce qu'ils ne
s'intéressent qu'à certains types d'anneaux et non au cas le plus général.

\begin{definition}
\label{definition-topologically-finite-type}
Soit $A \to B$ un homomorphisme continu d'anneaux topologiques
(More on Algebra, Definition \ref{more-algebra-definition-topological-ring}).
Nous disons que $B$ est {\it topologiquement de type fini sur} $A$ s'il
existe un homomorphisme de $A$-algèbres $A[x_1, \ldots, x_n] \to B$ dont
l'image est dense dans $B$.
\end{definition}

\noindent
Si $A$ est un anneau complet muni d'une topologie linéaire, alors l'anneau des
séries restreintes $A\{x_1, \ldots, x_r\}$ est topologiquement de type
fini sur $A$. Si $k$ est un corps, alors l'anneau de séries formelles
$k[[x_1, \ldots, x_r]]$ est topologiquement de type fini sur $k$.

\medskip\noindent
Pour les homomorphismes continus tendus d'anneaux topologiques faiblement admissibles,
le fait d'être topologiquement de type fini correspond exactement aux morphismes
de type fini entre les espaces algébriques formels affines associés.

\begin{lemma}
\label{lemma-topologically-finite-type-finite-type}
Soit $S$ un schéma. Soit $\varphi : A \to B$ un homomorphisme continu d'anneaux
topologiques faiblement admissibles sur $S$. Les conditions suivantes
sont équivalentes :
\begin{enumerate}
\item $\text{Spf}(\varphi) : Y = \text{Spf}(B) \to \text{Spf}(A) = X$
est de type fini ;
\item $\varphi$ est tendu et $B$ est topologiquement de type fini sur $A$.
\end{enumerate}
\end{lemma}

\begin{proof}
Nous pouvons utiliser le Lemme \ref{lemma-representable-affine} pour relier le caractère tendu de
$\varphi$ à la représentabilité de $\text{Spf}(\varphi)$. Nous l'emploierons
ci-dessous sans autre mention. Il s'ensuit que $X = \colim \Spec(A/I)$
et $Y = \colim \Spec(B/J(I))$, où $I \subset A$ parcourt les idéaux faibles
de définition de $A$ et où $J(I)$ est l'adhérence de $IB$ dans $B$.

\medskip\noindent
Supposons (2).
Choisissons un homomorphisme d'anneaux $A[x_1, \ldots, x_r] \to B$ d'image dense.
Alors $A[x_1, \ldots, x_r] \to B \to B/J(I)$ a lui aussi une image dense,
ce qui signifie qu'il est surjectif. Par conséquent, $B/J(I)$ est de
type fini sur $A/I$. Soit $T \to X$ un morphisme où
$T$ est un schéma quasi-compact. Alors $T \to X$ se factorise par
$\Spec(A/I)$ pour un certain $I$ (Lemme \ref{lemma-factor-through-thickening}).
Alors $T \times_X Y = T \times_{\Spec(A/I)} \Spec(B/J(I))$ ; voir la démonstration du
Lemme \ref{lemma-representable-affine}.
Par conséquent, $T \times_Y X \to T$ est de type fini comme changement de base du
morphisme $\Spec(B/J(I)) \to \Spec(A/I)$ qui est de type
fini. Ainsi, (1) est vraie.

\medskip\noindent
Supposons (1). Choisissons un $I \subset A$ quelconque comme ci-dessus. Puisque
$\Spec(A/I) \times_X Y = \Spec(B/J(I))$, nous voyons que $A/I \to B/J(I)$
est de type fini. Choisissons $b_1, \ldots, b_r \in B$
dont les images engendrent $B/J(I)$ sur $A/I$. Nous affirmons que l'image
de l'homomorphisme d'anneaux $A[x_1, \ldots, x_r] \to B$ qui envoie $x_i$ sur $b_i$
est dense. Pour le montrer, soit $I' \subset I$ un second idéal faible
de définition. Alors nous avons
$$
B/(J(I') + IB) = B/J(I)
$$
parce que $J(I)$ est l'adhérence de $IB$ et que $J(I')$ est ouvert.
Nous pouvons donc appliquer Algebra, Lemma
\ref{algebra-lemma-surjective-mod-locally-nilpotent}
pour voir que $(A/I')[x_1, \ldots, x_r] \to B/J(I')$ est surjectif.
Ainsi, (2) est vraie, ce qui achève la démonstration.
\end{proof}

\noindent
Soit $A$ un anneau topologique complet pour une topologie
linéaire. Soit $(I_\lambda)$ un système fondamental d'idéaux ouverts.
Soit $\mathcal{C}$ la catégorie des systèmes projectifs $(B_\lambda)$ où
\begin{enumerate}
\item $B_\lambda$ est une $A/I_\lambda$-algèbre de type fini ;
\item $B_\mu \to B_\lambda$ est un homomorphisme de $A/I_\mu$-algèbres
qui induit un isomorphisme $B_\mu/I_\lambda B_\mu \to B_\lambda$.
\end{enumerate}
Les morphismes de $\mathcal{C}$ sont donnés par des systèmes compatibles d'homomorphismes.

\begin{lemma}
\label{lemma-category-affine-over}
Soit $S$ un schéma. Soit $X$ un espace algébrique formel affine sur $S$.
Supposons $X$ de McQuillan et soit $A$ l'anneau topologique faiblement
admissible associé à $X$. Il existe alors une antiéquivalence de catégories
entre
\begin{enumerate}
\item la catégorie $\mathcal{C}$ introduite ci-dessus ;
\item la catégorie des morphismes $Y \to X$ de type fini entre
espaces algébriques formels affines.
\end{enumerate}
\end{lemma}

\begin{proof}
Soit $(I_\lambda)$ un système fondamental d'idéaux faiblement admissibles
de définition de $A$. Considérons $Y$ comme en (2). Alors
$Y \times_X \Spec(A/I_\lambda)$ est affine
(Définition \ref{definition-finite-type}
et Lemme \ref{lemma-affine-representable-by-algebraic-spaces}).
Écrivons $Y \times_X \Spec(A/I_\lambda) = \Spec(B_\lambda)$.
L'homomorphisme d'anneaux $A/I_\lambda \to B_\lambda$ est de type fini
parce que $\Spec(B_\lambda) \to \Spec(A/I_\lambda)$ est de
type fini
(d'après la Définition \ref{definition-finite-type}).
Alors $(B_\lambda)$ est un objet de $\mathcal{C}$.

\medskip\noindent
Réciproquement, étant donné un objet $(B_\lambda)$ de $\mathcal{C}$, nous pouvons poser
$Y = \colim \Spec(B_\lambda)$. C'est un espace algébrique formel
affine. Nous affirmons que
$$
Y \times_X \Spec(A/I_\lambda) =
\left(\colim_\mu \Spec(B_\mu)\right) \times_X \Spec(A/I_\lambda) =
\Spec(B_\lambda)
$$
Pour le montrer, il suffit d'obtenir les mêmes valeurs en évaluant
sur un schéma quasi-compact $U$. Un morphisme
$U \to \left(\colim_\mu \Spec(B_\mu)\right) \times_X \Spec(A/I_\lambda)$
provient d'un morphisme
$U \to \Spec(B_\mu) \times_{\Spec(A/I_\mu)} \Spec(A/I_\lambda)$
pour un certain $\mu \geq \lambda$ (appliquer
deux fois le Lemme \ref{lemma-factor-through-thickening}). Puisque
$\Spec(B_\mu) \times_{\Spec(A/I_\mu)} \Spec(A/I_\lambda) = \Spec(B_\lambda)$
d'après notre seconde hypothèse sur les objets de $\mathcal{C}$,
cela démontre l'assertion. Nous pouvons alors montrer que le morphisme
$Y \to X$ est de type fini. En effet, remarquons que,
pour tout morphisme $U \to X$ où $U$ est un schéma quasi-compact, nous obtenons
une factorisation $U \to \Spec(A/I_\lambda) \to X$ pour un certain $\lambda$
(voir le lemme cité ci-dessus). Par conséquent,
$$
Y \times_X U =
Y \times_X \Spec(A/I_\lambda)) \times_{\Spec(A/I_\lambda)} U =
\Spec(B_\lambda) \times_{\Spec(A/I_\lambda)} U
$$
est un schéma de type fini sur $U$, comme voulu. Ainsi, la construction
$(B_\lambda) \mapsto \colim \Spec(B_\lambda)$ fournit bien un foncteur
de la catégorie (1) vers la catégorie (2).

\medskip\noindent
Pour achever la démonstration, montrons que les constructions précédentes
définissent des foncteurs quasi-inverses entre les catégories (1) et (2).
Dans un sens, il faut montrer que
$$
\left(\colim_\mu \Spec(B_\mu)\right) \times_X \Spec(A/I_\lambda) =
\Spec(B_\lambda)
$$
pour tout objet $(B_\lambda)$ de la catégorie $\mathcal{C}$.
Nous l'avons démontré ci-dessus. Dans
l'autre sens, il faut montrer que
$$
Y = \colim (Y \times_X \Spec(A/I_\lambda))
$$
pour $Y$ dans la catégorie (2). Là encore, cela résulte de l'évaluation sur des
objets tests quasi-compacts et du fait que $X = \colim \Spec(A/I_\lambda)$.
\end{proof}

\begin{remark}
\label{remark-questions}
Soit $A$ un anneau topologique faiblement admissible et soit $(I_\lambda)$
un système fondamental d'idéaux faibles de définition. Soit $X = \text{Spf}(A)$ ;
autrement dit, $X$ est un espace algébrique formel affine de McQuillan.
Soit $f : Y \to X$ un morphisme d'espaces algébriques formels affines.
En général, $Y$ n'est pas nécessairement de McQuillan. Plus précisément,
nous pouvons poser les questions suivantes :
\begin{enumerate}
\item Supposons que $f : Y \to X$ soit une immersion fermée. Alors
$Y$ est de McQuillan et $f$ correspond à un homomorphisme continu
$\varphi : A \to B$ d'anneaux topologiques faiblement admissibles
qui est tendu, dont le noyau $K \subset A$ est un idéal fermé et
dont l'image $\varphi(A)$ est dense dans $B$ ; voir le
Lemme \ref{lemma-closed-immersion-into-McQuillan}.
Quelles conditions sur $A$ garantissent que $B = (A/K)^\wedge$ comme dans
l'Exemple \ref{example-closed-immersion-from-quotient} ?
\item Quelles conditions sur $A$ garantissent que les immersions fermées
$f : Y \to X$ correspondent aux quotients $A/K$ de $A$ par des idéaux fermés,
autrement dit que l'homomorphisme continu correspondant $\varphi$ est surjectif
et ouvert ?
\item Supposons que $f : Y \to X$ soit de type fini. Nous obtenons alors
$Y = \colim \Spec(B_\lambda)$, où $(B_\lambda)$ est un objet de
$\mathcal{C}$ d'après le Lemme \ref{lemma-category-affine-over}.
Dans ce cas, il existe bien un entier fixe $r$ tel
que $B_\lambda$ soit engendré par $r$ éléments sur $A/I_\lambda$ pour
tout $\lambda$ (l'argument est pour l'essentiel déjà donné dans la démonstration de
(1) $\Rightarrow$ (2) du
Lemme \ref{lemma-topologically-finite-type-finite-type}).
Cependant, il n'est pas clair que les projections
$\lim B_\lambda \to B_\lambda$ soient surjectives ; autrement dit,
il n'est pas clair que $Y$ soit de McQuillan.
Existe-t-il un exemple où $Y$ n'est pas de McQuillan ?
\item Supposons que $f : Y \to X$ soit de type fini et que $Y$ soit de McQuillan.
Alors $f$ correspond à un homomorphisme continu $\varphi : A \to B$ d'anneaux topologiques
faiblement admissibles. En fait, $\varphi$ est tendu et
$B$ est topologiquement de type fini sur $A$ ; voir le
Lemme \ref{lemma-topologically-finite-type-finite-type}.
Autrement dit, $f$ se factorise sous la forme
$$
Y \longrightarrow \mathbf{A}^r_X \longrightarrow X
$$
où la première flèche est une immersion fermée d'espaces algébriques formels
affines de McQuillan. Cependant, les questions (1) et
(2) se posent alors pour $Y \to \mathbf{A}^r_X$.
\end{enumerate}
Nous répondrons ci-dessous à ces questions lorsque
$X$ est à indexation dénombrable, c'est-à-dire lorsque $A$ possède un système fondamental
dénombrable d'idéaux ouverts. Si vous connaissez des réponses à ces questions
dans une plus grande généralité, ou des contre-exemples, veuillez écrire à
\href{mailto:stacks.project@gmail.com}{stacks.project@gmail.com}.
\end{remark}

\begin{lemma}
\label{lemma-closed-immersion-into-countably-indexed}
Soit $S$ un schéma. Soit $X$ un espace algébrique formel affine à indexation
dénombrable sur $S$. Soit $f : Y \to X$ une immersion fermée d'espaces algébriques
formels sur $S$. Alors $Y$ est un espace algébrique formel affine à indexation dénombrable
et $f$ correspond à $A \to A/K$, où $A$ est un objet de
$\textit{WAdm}^{count}$
(Section \ref{section-morphisms-rings})
et où $K \subset A$ est un idéal fermé.
\end{lemma}

\begin{proof}
D'après le Lemme \ref{lemma-countably-indexed},
nous avons $X = \text{Spf}(A)$, où $A$ est un objet de
$\textit{WAdm}^{count}$. Comme une immersion fermée est représentable
et affine, le Lemme \ref{lemma-property-goes-up-affine-morphism} montre
que $Y$ est un espace algébrique formel affine à indexation dénombrable.
En appliquant de nouveau le Lemme \ref{lemma-countably-indexed},
nous voyons donc que $Y = \text{Spf}(B)$, avec $B$ objet de
$\textit{WAdm}^{count}$. D'après le Lemme \ref{lemma-closed-immersion-into-McQuillan},
nous concluons que $f$ est donné par un morphisme $A \to B$ de
$\textit{WAdm}^{count}$ qui est tendu et d'image dense.
Pour achever la démonstration, nous appliquons le
Lemme \ref{lemma-dense-image-surjective}.
\end{proof}

\begin{lemma}
\label{lemma-quotient-restricted-power-series}
Soit $B \to A$ une flèche de $\textit{WAdm}^{count}$ ; voir la
Section \ref{section-morphisms-rings}.
Les conditions suivantes sont équivalentes :
\begin{enumerate}
\item[(a)] $B \to A$ est tendu et $B/J \to A/I$ est de type fini pour
tout idéal faible de définition $J \subset B$, où l'idéal $I \subset A$ est
l'adhérence de $JA$ ;
\item[(b)] $B \to A$ est tendu et $B/J_\lambda \to A/I_\lambda$
est de type fini pour un système cofinal $(J_\lambda)$
d'idéaux faibles de définition de $B$, où
$I_\lambda \subset A$ est l'adhérence de $J_\lambda A$ ;
\item[(c)] $B \to A$ est tendu et $A$ est topologiquement de type
fini sur $B$ ;
\item[(d)] $A$ est isomorphe, comme $B$-algèbre topologique, à un quotient de
$B\{x_1, \ldots, x_n\}$ par un idéal fermé.
\end{enumerate}
De plus, ces conditions équivalentes définissent une propriété locale,
c'est-à-dire qu'elles satisfont aux
axiomes (\ref{item-axiom-1}), (\ref{item-axiom-2}), (\ref{item-axiom-3}).
\end{lemma}

\begin{proof}
Les implications (a) $\Rightarrow$ (b), (c) $\Rightarrow$ (a),
(d) $\Rightarrow$ (c) résultent immédiatement des définitions.
Supposons (b) vérifiée et soient $J \subset B$ et $I \subset A$ comme en (a).
Choisissons un diagramme commutatif
$$
\xymatrix{
A \ar[r] & \ldots \ar[r] & A_3 \ar[r] & A_2 \ar[r] & A_1 \\
B \ar[r] \ar[u] & \ldots \ar[r] & B/J_3 \ar[r] \ar[u] &
B/J_2 \ar[r] \ar[u] & B/J_1 \ar[u]
}
$$
tel que $A_{n + 1}/J_nA_{n + 1} = A_n$ et $A = \lim A_n$, comme dans le
Lemme \ref{lemma-representable-property-rings}.
Pour tout $m$, il existe un $\lambda$ tel que $J_\lambda \subset J_m$.
Comme $B/J_\lambda \to A/I_\lambda$ est de type fini, il s'ensuit
que $B/J_m \to A/I_m$ est de type fini.
Soient $\alpha_1, \ldots, \alpha_n \in A_1$ des générateurs de $A_1$ sur
$B/J_1$. Puisque $A$ est la limite dénombrable d'un système à morphismes de
transition surjectifs, nous pouvons trouver $a_1, \ldots, a_n \in A$ dont les images sont
$\alpha_1, \ldots, \alpha_n$ dans $A_1$. D'après la
Remarque \ref{remark-universal-property}, nous obtenons un homomorphisme continu
$B\{x_1, \ldots, x_n\} \to A$ qui envoie $x_i$ sur $a_i$. Cet homomorphisme
induit des surjections $(B/J_m)[x_1, \ldots, x_n] \to A_m$ d'après
Algebra, Lemma \ref{algebra-lemma-surjective-mod-locally-nilpotent}.
Pour $m \geq 1$, nous obtenons une suite exacte courte
$$
0 \to K_m \to (B/J_m)[x_1, \ldots, x_n] \to A_m \to 0
$$
Les morphismes de transition induits $K_{m + 1} \to K_m$ sont surjectifs parce que
$A_{m + 1}/J_mA_{m + 1} = A_m$. Par conséquent, la limite projective de ces
suites exactes courtes est exacte ; voir
Algebra, Lemma \ref{algebra-lemma-ML-exact-sequence}.
Puisque $B\{x_1, \ldots, x_n\} = \lim (B/J_m)[x_1, \ldots, x_n]$
et $A = \lim A_m$,
nous concluons que $B\{x_1, \ldots, x_n\} \to A$ est surjectif et ouvert.
Comme $A$ est complet, le noyau est un idéal fermé. Nous voyons ainsi que
(a), (b), (c) et (d) sont équivalentes.

\medskip\noindent
Donnons-nous un diagramme (\ref{equation-localize}) comme dans la
Situation \ref{situation-local-property}.
D'après l'Exemple \ref{example-finite-type-from-finite-type-ring-map},
les homomorphismes $A \to (A')^\wedge$ et $B \to (B')^\wedge$
satisfont à (a), (b), (c) et (d). De plus, d'après le
Lemme \ref{lemma-representable-property-rings},
pour démontrer les axiomes (\ref{item-axiom-1}) et (\ref{item-axiom-2}),
nous pouvons supposer que $B \to A$ et $(B')^\wedge \to (A')^\wedge$
sont tous deux tendus. Choisissons un idéal faible de définition $J \subset B$. Soient
$J' \subset (B')^\wedge$, $I \subset A$, $I' \subset (A')^\wedge$
les adhérences de $J(B')^\wedge$, $JA$, $J(A')^\wedge$.
D'après ce qui précède, il suffit de considérer le
diagramme commutatif
$$
\xymatrix{
A/I \ar[r] & (A')^\wedge/I' \\
B/J \ar[r] \ar[u]^{\overline{\varphi}} &
(B')^\wedge/J' \ar[u]_{\overline{\varphi}'}
}
$$
et de montrer : (1) $\overline{\varphi}$ est de type fini
$\Rightarrow \overline{\varphi}'$ est de type fini ; (2) si
$A \to A'$ est fidèlement plat, alors
$\overline{\varphi}'$ est de type fini $\Rightarrow \overline{\varphi}$
est de type fini. Remarquons que $(B')^\wedge/J' = B'/JB'$ et
$(A')^\wedge/I' = A'/IA'$ par construction des topologies sur
$(B')^\wedge$ et $(A')^\wedge$. En particulier, les homomorphismes horizontaux
du diagramme sont étales. La partie (1) résulte alors de
Algebra, Lemma \ref{algebra-lemma-compose-finite-type},
et la partie (2) de
Descent, Lemma \ref{descent-lemma-finite-type-local-source-fppf-algebra},
puisque l'homomorphisme d'anneaux $A/I \to (A')^\wedge/I' = A'/IA'$ est fidèlement plat
et étale.

\medskip\noindent
Nous omettons la démonstration de l'axiome (\ref{item-axiom-3}).
\end{proof}

\begin{lemma}
\label{lemma-quotient-restricted-power-series-admissible}
Dans le Lemme \ref{lemma-quotient-restricted-power-series},
si $B$ est admissible (par exemple adique), alors les conditions équivalentes
(a) à (d) sont aussi équivalentes à la condition suivante :
\begin{enumerate}
\item[(e)] $B \to A$ est tendu et $B/J \to A/I$ est de type fini pour
un certain idéal de définition $J \subset B$, où $I \subset A$ est
l'adhérence de $JA$.
\end{enumerate}
\end{lemma}

\begin{proof}
Il suffit de montrer que (e) implique (a). Soit $J' \subset B$ un idéal faible
de définition et soit $I' \subset A$ l'adhérence de $J'A$. Nous devons
montrer que $B/J' \to A/I'$ est de type fini. Si l'assertion correspondante
est vraie pour le plus petit idéal faible de définition $J'' = J' \cap J$, alors elle
est vraie pour $J'$. Nous pouvons donc supposer $J' \subset J$. Comme $J$ est un idéal
de définition (et non seulement un idéal faible de définition), nous avons
$J^n \subset J'$ pour un certain $n \geq 1$. Nous pouvons donc considérer le
diagramme
$$
\xymatrix{
0 \ar[r] & I/I' \ar[r] & A/I' \ar[r] & A/I \ar[r] & 0 \\
0 \ar[r] & J/J' \ar[r] \ar[u] & B/J' \ar[r] \ar[u] & B/J \ar[r] \ar[u] & 0
}
$$
à lignes exactes. Puisque $I' \subset A$ est ouvert et que
$I$ est l'adhérence de $J A$, nous voyons que $I/I' = (J/J') \cdot A/I'$.
Comme $J/J'$ est un idéal nilpotent et que $B/J \to A/I$ est de type fini,
Algebra, Lemma \ref{algebra-lemma-finite-type-mod-nilpotent} montre
que $A/I'$ est de type fini sur $B/J'$, comme voulu.
\end{proof}

\begin{lemma}
\label{lemma-representable-affine-finite-type}
Soit $S$ un schéma. Soit $f : X \to Y$ un morphisme d'espaces algébriques
formels affines. Supposons $Y$ à indexation dénombrable.
Les conditions suivantes sont équivalentes :
\begin{enumerate}
\item $f$ est localement de type fini ;
\item $f$ est de type fini ;
\item $f$ correspond à un morphisme $B \to A$ de $\textit{WAdm}^{count}$
(Section \ref{section-morphisms-rings})
qui satisfait aux conditions équivalentes du
Lemme \ref{lemma-quotient-restricted-power-series}.
\end{enumerate}
\end{lemma}

\begin{proof}
Puisque $X$ et $Y$ sont affines, il est clair que les conditions (1)
et (2) sont équivalentes. Dans les cas (1) et (2), le morphisme $f$
est représentable par des espaces algébriques par définition, donc
affine d'après le Lemme \ref{lemma-affine-representable-by-algebraic-spaces}.
Ainsi, si (1) ou (2) est vérifiée, nous voyons que
$X$ est à indexation dénombrable d'après le
Lemme \ref{lemma-property-goes-up-affine-morphism}.
Écrivons $X = \text{Spf}(A)$ et $Y = \text{Spf}(B)$
pour des $S$-algèbres topologiques $A$ et $B$ de $\textit{WAdm}^{count}$ ; voir le
Lemme \ref{lemma-countably-indexed}. D'après le
Lemme \ref{lemma-morphism-between-formal-spectra},
nous voyons que $f$ correspond à un homomorphisme continu $B \to A$.
Le résultat découle alors du
Lemme \ref{lemma-topologically-finite-type-finite-type}.
\end{proof}

\begin{lemma}
\label{lemma-finite-type-local-property}
Soit $S$ un schéma. Soit $f : X \to Y$ un morphisme d'espaces algébriques formels
localement à indexation dénombrable sur $S$.
Les conditions suivantes sont équivalentes :
\begin{enumerate}
\item pour tout diagramme commutatif
$$
\xymatrix{
U \ar[d] \ar[r] & V \ar[d] \\
X \ar[r] & Y
}
$$
où $U$ et $V$ sont des espaces algébriques formels affines, où $U \to X$ et $V \to Y$
sont représentables par des espaces algébriques et étales, le morphisme $U \to V$
correspond à un morphisme de $\textit{WAdm}^{count}$ qui est
tendu et topologiquement de type fini ;
\item il existe un recouvrement $\{Y_j \to Y\}$ comme dans la
Définition \ref{definition-formal-algebraic-space} et, pour tout $j$,
un recouvrement $\{X_{ji} \to Y_j \times_Y X\}$ comme dans la
Définition \ref{definition-formal-algebraic-space},
tel que chaque $X_{ji} \to Y_j$ corresponde
à un morphisme de $\textit{WAdm}^{count}$ qui est
tendu et topologiquement de type fini ;
\item il existe un recouvrement $\{X_i \to X\}$ comme dans la
Définition \ref{definition-formal-algebraic-space}
et, pour tout $i$, une factorisation $X_i \to Y_i \to Y$ où $Y_i$
est un espace algébrique formel affine, où $Y_i \to Y$ est représentable
par des espaces algébriques et étale, et où $X_i \to Y_i$ correspond
à un morphisme de $\textit{WAdm}^{count}$ qui est
tendu et topologiquement de type fini ;
\item $f$ est localement de type fini.
\end{enumerate}
\end{lemma}

\begin{proof}
D'après le Lemme \ref{lemma-quotient-restricted-power-series},
la propriété
$P(\varphi)=$ « $\varphi$ est tendu et topologiquement de type fini »
est locale sur $\text{WAdm}^{count}$. Par conséquent, le
Lemme \ref{lemma-property-defines-property-morphisms}
montre que les conditions (1), (2) et (3) sont équivalentes.
D'autre part, d'après le Lemme \ref{lemma-representable-affine-finite-type},
la condition $P$ sur les morphismes de $\textit{WAdm}^{count}$
correspond exactement au fait que les morphismes d'espaces algébriques formels
affines à indexation dénombrable soient localement de type fini.
Ainsi, l'implication (1) $\Rightarrow$ (3) du
Lemme \ref{lemma-finite-type-local}
montre que (4) implique (1) du présent lemme.
De même, l'implication (4) $\Rightarrow$ (1) du
Lemme \ref{lemma-finite-type-local}
montre que (2) implique (4) du présent lemme.
\end{proof}







\section{Axiomes de séparation pour les morphismes}
\label{section-separation-axioms}

\noindent
Cette section est l'analogue de
Morphisms of Spaces, Section \ref{spaces-morphisms-section-separation-axioms}
pour les morphismes d'espaces algébriques formels.

\begin{definition}
\label{definition-separated-morphism}
Soit $S$ un schéma.
Soit $f : X \to Y$ un morphisme d'espaces algébriques formels sur $S$.
Soit $\Delta_{X/Y} : X \to X \times_Y X$ le morphisme diagonal.
\begin{enumerate}
\item Nous disons que $f$ est {\it séparé} si $\Delta_{X/Y}$ est une immersion fermée.
\item Nous disons que $f$ est {\it quasi-séparé} si $\Delta_{X/Y}$ est quasi-compact.
\end{enumerate}
\end{definition}

\noindent
Puisque $\Delta_{X/Y}$ est représentable (par des schémas) d'après le
Lemme \ref{lemma-diagonal-morphism-formal-algebraic-spaces},
nous pouvons tester ces propriétés en considérant les morphismes $T \to X \times_Y X$
issus de schémas affines $T$ et en vérifiant si
$$
E = T \times_{X \times_Y X} X \longrightarrow T
$$
est quasi-compact ou une immersion fermée ; voir le
Lemme \ref{lemma-quasi-compact-representable} ou la
Définition \ref{definition-closed-immersion}.
Remarquons que le schéma $E$ est l'égalisateur de deux morphismes
$a, b : T \to X$ qui coïncident comme morphismes vers $Y$,
et que $E \to T$ est un monomorphisme localement de type fini.

\begin{lemma}
\label{lemma-base-change-separated}
Tous les axiomes de séparation énumérés dans la
Définition \ref{definition-separated-morphism}
sont stables par changement de base.
\end{lemma}

\begin{proof}
Soient $f : X \to Y$ et $Y' \to Y$ des morphismes d'espaces algébriques formels.
Soit $f' : X' \to Y'$ le changement de base de $f$ par $Y' \to Y$. Alors
$\Delta_{X'/Y'}$ est le changement de base de $\Delta_{X/Y}$ par
le morphisme $X' \times_{Y'} X' \to X \times_Y X$. Chacune des propriétés
de la diagonale employées dans la Définition \ref{definition-separated-morphism}
est stable par changement de base. Le lemme est donc vrai.
\end{proof}

\begin{lemma}
\label{lemma-fibre-product-after-map}
Soit $S$ un schéma. Soient $f : X \to Z$, $g : Y \to Z$ et $Z \to T$
des morphismes d'espaces algébriques formels sur $S$. Considérons le
morphisme induit $i : X \times_Z Y \to X \times_T Y$. Alors
\begin{enumerate}
\item $i$ est représentable (par des schémas), localement de type fini,
localement quasi-fini, séparé et est un monomorphisme ;
\item si $Z \to T$ est séparé, alors $i$ est une immersion fermée ;
\item si $Z \to T$ est quasi-séparé, alors $i$ est quasi-compact.
\end{enumerate}
\end{lemma}

\begin{proof}
Par la théorie générale des catégories, le diagramme suivant
$$
\xymatrix{
X \times_Z Y \ar[r]_i \ar[d] & X \times_T Y \ar[d] \\
Z \ar[r]^-{\Delta_{Z/T}} \ar[r] & Z \times_T Z
}
$$
est un diagramme de produit fibré. Ainsi, $i$ est le changement de base du
morphisme diagonal $\Delta_{Z/T}$. Le lemme résulte donc
du Lemme \ref{lemma-diagonal-morphism-formal-algebraic-spaces}.
\end{proof}

\begin{lemma}
\label{lemma-composition-separated}
Tous les axiomes de séparation énumérés dans la
Définition \ref{definition-separated-morphism}
sont stables par composition des morphismes.
\end{lemma}

\begin{proof}
Soient $f : X \to Y$ et $g : Y \to Z$ des morphismes d'espaces algébriques formels
auxquels s'applique l'axiome considéré.
La diagonale $\Delta_{X/Z}$ est la composée
$$
X \longrightarrow X \times_Y X \longrightarrow X \times_Z X.
$$
Notre axiome de séparation est défini en exigeant que la diagonale
possède une certaine propriété $\mathcal{P}$. D'après le
Lemme \ref{lemma-fibre-product-after-map} ci-dessus, nous voyons que
la seconde flèche possède aussi cette propriété. Le lemme en résulte,
puisque la composée de morphismes (représentables) possédant la propriété
$\mathcal{P}$ possède encore la propriété $\mathcal{P}$.
\end{proof}

\begin{lemma}
\label{lemma-separated-local}
Soit $S$ un schéma. Soit $f : X \to Y$ un morphisme d'espaces algébriques formels
sur $S$. Soit $\mathcal{P}$ l'un quelconque des axiomes de séparation de la
Définition \ref{definition-separated-morphism}.
Les conditions suivantes sont équivalentes :
\begin{enumerate}
\item $f$ possède la propriété $\mathcal{P}$ ;
\item pour tout schéma $Z$ et tout morphisme $Z \to Y$, le
changement de base $Z \times_Y X \to Z$ de $f$ possède la propriété $\mathcal{P}$ ;
\item pour tout schéma affine $Z$ et tout morphisme $Z \to Y$, le
changement de base $Z \times_Y X \to Z$ de $f$ possède la propriété $\mathcal{P}$ ;
\item pour tout schéma affine $Z$ et tout morphisme $Z \to Y$, le
produit $Z \times_Y X$ possède la propriété $\mathcal{P}$ comme espace algébrique formel (voir la
Définition \ref{definition-separated}) ;
\item il existe un recouvrement $\{Y_j \to Y\}$ comme dans la
Définition \ref{definition-formal-algebraic-space},
tel que le changement de base $Y_j \times_Y X \to Y_j$ possède
la propriété $\mathcal{P}$ pour tout $j$.
\end{enumerate}
\end{lemma}

\begin{proof}
Nous utiliserons à plusieurs reprises le
Lemme \ref{lemma-base-change-separated}
sans autre mention. En particulier, il est clair que
(1) implique (2) et que (2) implique (3).

\medskip\noindent
Supposons (3) et soit $Z \to Y$ un morphisme où $Z$ est un schéma affine.
Soient $U$, $V$ des schémas affines et soient $a : U \to Z \times_Y X$
et $b : V \to Z \times_Y X$ des morphismes. Alors
$$
U \times_{Z \times_Y X} V =
(Z \times_Y X) \times_{\Delta, (Z \times_Y X) \times_Z (Z \times_Y X)}
(U \times_Z V)
$$
et nous voyons qu'il est quasi-compact si $\mathcal{P} =$ « quasi-séparé »,
ou qu'il s'agit d'un schéma affine muni d'une immersion fermée dans
$U \times_Z V$ si $\mathcal{P} =$ « séparé ». Ainsi, (4) est vérifiée.

\medskip\noindent
Supposons (4) et soit $Z \to Y$ un morphisme où $Z$ est un schéma affine.
Soient $U$, $V$ des schémas affines et soient $a : U \to Z \times_Y X$
et $b : V \to Z \times_Y X$ des morphismes. En lisant l'argument précédent
à rebours, nous voyons que $U \times_{Z \times_Y X} V \to U \times_Z V$
est quasi-compact si $\mathcal{P} =$ « quasi-séparé », ou une immersion
fermée si $\mathcal{P} =$ « séparé ». Puisque nous pouvons choisir $U$ et
$V$ comme ci-dessus de façon que $U$ parcoure un
recouvrement étale de $Z \times_Y X$, nous trouvons
que les morphismes correspondants
$$
U \times_Z V \to (Z \times_Y X) \times_Z (Z \times_Y X)
$$
forment un recouvrement étale par des affines. Nous concluons donc que
$\Delta : (Z \times_Y X) \to (Z \times_Y X) \times_Z (Z \times_Y X)$
est quasi-compact, resp.\ une immersion fermée. Ainsi, (3) est vérifiée.

\medskip\noindent
Montrons que (3) implique (5). Supposons (3) et soit
$\{Y_j \to Y\}$ comme dans la
Définition \ref{definition-formal-algebraic-space}.
Nous devons montrer que les morphismes
$$
\Delta_j :
Y_j \times_Y X
\longrightarrow
(Y_j \times_Y X) \times_{Y_j} (Y_j \times_Y X) =
Y_j \times_Y X \times_Y X
$$
possèdent la propriété correspondante (c'est-à-dire sont quasi-compacts ou des immersions fermées).
Écrivons $Y_j = \colim Y_{j, \lambda}$ comme dans la
Définition \ref{definition-affine-formal-algebraic-space}.
En remplaçant $Y_j$ par $Y_{j, \lambda}$ dans la formule ci-dessus, nous obtenons la
propriété par notre hypothèse (3). Puisque la flèche affichée
est la colimite des flèches $\Delta_{j, \lambda}$ et que nous
pouvons vérifier si $\Delta_j$ possède la propriété correspondante
après changement de base par des schémas affines qui s'envoient dans
$Y_j \times_Y X \times_Y X$, nous concluons par le
Lemme \ref{lemma-factor-through-thickening}.

\medskip\noindent
Montrons que (5) implique (1). Soit $\{Y_j \to Y\}$ comme en (5).
Nous avons alors le diagramme de produit fibré
$$
\xymatrix{
\coprod Y_j \times_Y X \ar[r] \ar[d] &
X \ar[d] \\
\coprod Y_j \times_Y X \times_Y X \ar[r] &
X \times_Y X
}
$$
Par hypothèse, la flèche verticale de gauche est quasi-compacte ou une immersion fermée.
Il résulte de
Spaces, Lemma \ref{spaces-lemma-descent-representable-transformations-property}
que la flèche verticale de droite est, elle aussi, quasi-compacte ou une
immersion fermée.
\end{proof}







\section{Morphismes propres}
\label{section-proper}

\noindent
Voici la définition que nous utiliserons.

\begin{definition}
\label{definition-proper}
Soit $S$ un schéma. Soit $f : Y \to X$ un morphisme d'espaces algébriques formels
sur $S$. Nous disons que $f$ est {\it propre}
si $f$ est représentable par des espaces algébriques et est propre au sens de
Bootstrap, Definition \ref{bootstrap-definition-property-transformation}.
\end{definition}

\noindent
Il résulte des définitions qu'un morphisme propre est de type fini.

\begin{lemma}
\label{lemma-proper-local}
Soit $S$ un schéma. Soit $f : X \to Y$ un morphisme d'espaces algébriques formels
sur $S$. Les conditions suivantes sont équivalentes :
\begin{enumerate}
\item $f$ est propre ;
\item pour tout schéma $Z$ et tout morphisme $Z \to Y$, le
changement de base $Z \times_Y X \to Z$ de $f$ est propre ;
\item pour tout schéma affine $Z$ et tout morphisme $Z \to Y$, le
changement de base $Z \times_Y X \to Z$ de $f$ est propre ;
\item pour tout schéma affine $Z$ et tout morphisme $Z \to Y$, le
produit $Z \times_Y X$ est un espace algébrique propre sur $Z$ ;
\item il existe un recouvrement $\{Y_j \to Y\}$ comme dans la
Définition \ref{definition-formal-algebraic-space},
tel que le changement de base $Y_j \times_Y X \to Y_j$ soit propre pour tout $j$.
\end{enumerate}
\end{lemma}

\begin{proof}
Démonstration omise.
\end{proof}

\begin{lemma}
\label{lemma-base-change-proper}
Les morphismes propres d'espaces algébriques formels sont préservés par changement de base.
\end{lemma}

\begin{proof}
C'est une conséquence immédiate du Lemme \ref{lemma-proper-local}
et de la transitivité du changement de base.
\end{proof}







\section{Espaces algébriques formels et recouvrements fpqc}
\label{section-fpqc}

\noindent
Cette section est l'analogue de Properties of Spaces, Section
\ref{spaces-properties-section-fpqc}. Veuillez lire cette section
au préalable.

\begin{lemma}
\label{lemma-sheaf-fpqc}
\begin{slogan}
Les espaces algébriques formels sont des faisceaux fpqc
\end{slogan}
Soit $S$ un schéma. Soit $X$ un espace algébrique formel sur $S$. Alors
$X$ vérifie la propriété de faisceau pour la topologie fpqc.
\end{lemma}

\begin{proof}
La démonstration est {\bf identique} à celle de
Properties of Spaces, Proposition
\ref{spaces-properties-proposition-sheaf-fpqc}.
Puisque $X$ est un faisceau pour la topologie de Zariski, il
suffit de montrer ce qui suit. Étant donné un morphisme plat surjectif
entre schémas affines $f : T' \to T$, on a :
$X(T)$ est l'égalisateur des deux applications $X(T') \to X(T' \times_T T')$.
Voir Topologies, Lemma \ref{topologies-lemma-sheaf-property-fpqc}.

\medskip\noindent
Soient $a, b : T \to X$ deux morphismes tels que $a \circ f = b \circ f$.
Nous devons montrer que $a = b$. Considérons le produit fibré
$$
E = X \times_{\Delta_{X/S}, X \times_S X, (a, b)} T.
$$
Par le Lemme \ref{lemma-diagonal-formal-algebraic-space},
le morphisme $\Delta_{X/S}$ est un monomorphisme représentable. Ainsi,
$E \to T$ est un monomorphisme de schémas. Notre hypothèse selon laquelle
$a \circ f = b \circ f$ implique que $T' \to T$ se factorise (uniquement) par $E$.
Considérons le diagramme commutatif
$$
\xymatrix{
T' \times_T E \ar[r] \ar[d] & E \ar[d] \\
T' \ar[r] \ar@/^5ex/[u] \ar[ru] & T
}
$$
Puisque la projection $T' \times_T E \to T'$ est un monomorphisme
muni d'une section, nous en déduisons qu'elle est un isomorphisme. Ainsi,
$E \to T$ est un isomorphisme d'après
Descent, Lemma \ref{descent-lemma-descending-property-isomorphism}.
Cela signifie que $a = b$, comme souhaité.

\medskip\noindent
Ensuite, soit $c : T' \to X$ un morphisme tel que les deux composées
$T' \times_T T' \to T' \to X$ coïncident. Nous devons trouver un morphisme
$a : T \to X$ dont la composée avec $T' \to T$ est $c$. Choisissons un
schéma formel affine $U$ et un morphisme \'etale $U \to X$ tels que l'image
de $|U| \to |X_{red}|$ contienne l'image de $|c| : |T'| \to |X_{red}|$.
C'est possible d'après la
Définition \ref{definition-formal-algebraic-space},
Properties of Spaces, Lemma \ref{spaces-properties-lemma-topology-points},
le fait qu'une réunion finie d'espaces algébriques formels affines est un
espace algébrique formel affine, et le fait que $|T'|$ est quasi-compact
(petit argument omis). Le morphisme $U \to X$ est représentable
par des schémas (Lemme \ref{lemma-presentation-representable}) et
séparé (Lemme \ref{lemma-separated-from-separated}). Ainsi,
$$
V = U \times_{X, c} T' \longrightarrow T'
$$
est un morphisme \'etale et séparé de schémas. Il est également surjectif
par notre choix de $U \to X$ (si l'on ne souhaite pas le démontrer, on peut
remplacer $U$ par une réunion disjointe d'espaces algébriques formels affines afin que
$U \to X$ soit surjectif ; tout le reste fonctionne également). Le fait que
$c \circ \text{pr}_0 = c \circ \text{pr}_1$ signifie que nous obtenons une
donnée de descente sur $V/T'/T$
(Descent, Definition \ref{descent-definition-descent-datum})
car
\begin{align*}
V \times_{T'} (T' \times_T T')
& =
U \times_{X, c \circ \text{pr}_0} (T' \times_T T') \\
& =
(T' \times_T T') \times_{c \circ \text{pr}_1, X} U \\
& =
(T' \times_T T') \times_{T'} V
\end{align*}
Le morphisme $V \to T'$ est ind-quasi-affine d'après
More on Morphisms, Lemma
\ref{more-morphisms-lemma-etale-separated-ind-quasi-affine}
(car les morphismes \'etales sont localement quasi-finis, voir
Morphisms, Lemma \ref{morphisms-lemma-etale-locally-quasi-finite}).
D'après More on Groupoids, Lemma \ref{more-groupoids-lemma-ind-quasi-affine},
la donnée de descente est effective. Soit $W \to T$ un morphisme
tel qu'il existe un isomorphisme $\alpha : T' \times_T W \to V$
compatible avec la donnée de descente sur $V$ et la donnée de descente canonique
sur $T' \times_T W$. Alors $W \to T$ est surjectif et \'etale
(Descent, Lemmas \ref{descent-lemma-descending-property-surjective} et
\ref{descent-lemma-descending-property-etale}).
Considérons la composée
$$
b' : T' \times_T W \longrightarrow V = U \times_{X, c} T' \longrightarrow U
$$
Les deux composées
$b' \circ (\text{pr}_0, 1), 
b' \circ (\text{pr}_1, 1) :
(T' \times_T T') \times_T W \to T' \times_T W \to U$
coïncident par notre choix de $\alpha$ et la propriété correspondante de $c$
(calcul omis). Ainsi, $b'$ descend en un morphisme $b : W \to U$ d'après
Descent, Lemma \ref{descent-lemma-fpqc-universal-effective-epimorphisms}.
Le diagramme
$$
\xymatrix{
T' \times_T W \ar[r] \ar[d] & W \ar[r]_b & U \ar[d] \\
T' \ar[rr]^c &  & X
}
$$
est commutatif. Cela signifie que nous avons démontré l'existence
de $a$ localement pour la topologie \'etale sur $T$, i.e., nous avons un $a' : W \to X$.
Toutefois, puisque nous avons démontré l'unicité
dans le premier paragraphe, cette solution locale pour la topologie \'etale
satisfait la condition de recollement, i.e., nous avons
$\text{pr}_0^*a' = \text{pr}_1^*a'$ dans $X(W \times_T W)$.
Puisque $X$ est un faisceau \'etale, il existe un unique $a \in X(T)$ qui se restreint
à $a'$ sur $W$.
\end{proof}





\section{Morphismes issus de schémas formels affines}
\label{section-map-out-of}

\noindent
Nous démontrons quelques résultats qui seront utiles plus tard.
Dans l'article \cite{Bhatt-Algebraize}, le lecteur trouvera
des résultats très généraux de même nature.

\begin{lemma}
\label{lemma-map-into-affine}
Soit $S$ un schéma. Soit $A$ une algèbre topologique faiblement admissible
sur $S$. Soit $X$ un schéma affine sur $S$. Alors
l'application naturelle
$$
\Mor_S(\Spec(A), X)
\longrightarrow
\Mor_S(\text{Spf}(A), X)
$$
est bijective.
\end{lemma}

\begin{proof}
Si $X$ est affine, disons $X = \Spec(B)$, il résulte du
Lemme \ref{lemma-morphism-between-formal-spectra}
que les morphismes $\text{Spf}(A) \to \Spec(B)$ correspondent aux homomorphismes continus
d'algèbres sur $S$, $B \to A$, où $B$ porte la topologie discrète.
Ce sont simplement les homomorphismes d'algèbres sur $S$, qui correspondent aux morphismes
$\Spec(A) \to \Spec(B)$.
\end{proof}

\begin{lemma}
\label{lemma-map-into-scheme}
Soit $S$ un schéma. Soit $A$ une algèbre topologique faiblement admissible
sur $S$ telle que $A/I$ soit un anneau local pour un certain idéal faible
de définition $I \subset A$. Soit $X$ un schéma sur $S$. Alors
l'application naturelle
$$
\Mor_S(\Spec(A), X)
\longrightarrow
\Mor_S(\text{Spf}(A), X)
$$
est bijective.
\end{lemma}

\begin{proof}
Soit $\varphi : \text{Spf}(A) \to X$ un morphisme. Puisque $\Spec(A/I)$
est local, nous voyons que $\varphi$ envoie $\Spec(A/I)$ dans un
ouvert affine $U \subset X$. Cela implique alors que $\Spec(A/J)$
s'envoie dans $U$ pour tout idéal de définition $J$. Nous pouvons donc
appliquer le Lemme \ref{lemma-map-into-affine} pour voir que $\varphi$ provient
d'un morphisme $\Spec(A) \to X$. Cela démontre la surjectivité de l'application.
Nous omettons la démonstration de l'injectivité.
\end{proof}

\begin{lemma}
\label{lemma-map-into-algebraic-space}
Soit $S$ un schéma. Soit $R$ une algèbre locale noethérienne complète sur $S$.
Soit $X$ un espace algébrique sur $S$. Alors l'application naturelle
$$
\Mor_S(\Spec(R), X)
\longrightarrow
\Mor_S(\text{Spf}(R), X)
$$
est bijective.
\end{lemma}

\begin{proof}
Soit $\mathfrak m$ l'idéal maximal de $R$. Nous devons montrer que
$$
\Mor_S(\Spec(R), X) \longrightarrow \lim \Mor_S(\Spec(R/\mathfrak m^n), X)
$$
est bijective pour $R$ comme ci-dessus.

\medskip\noindent
Injectivité : soient $x, x' : \Spec(R) \to X$
deux morphismes ayant la même image dans le membre de droite.
Considérons le produit fibré
$$
T = \Spec(R) \times_{(x, x'), X \times_S X, \Delta} X
$$
Alors $T$ est un schéma et $T \to \Spec(R)$ est localement de type fini,
est un monomorphisme, est séparé et est localement quasi-fini ; voir
Morphisms of Spaces, Lemma \ref{spaces-morphisms-lemma-properties-diagonal}.
En particulier, $T$ est localement noethérien ; voir
Morphisms, Lemma \ref{morphisms-lemma-finite-type-noetherian}.
Soit $t \in T$ l'unique point s'envoyant sur le point fermé de $\Spec(R)$,
qui existe puisque $x$ et $x'$ coïncident sur $R/\mathfrak m$. Alors
$R \to \mathcal{O}_{T, t}$ est un homomorphisme local d'anneaux noethériens tel que
$R/\mathfrak m^n \to \mathcal{O}_{T, t}/\mathfrak m^n\mathcal{O}_{T, t}$
est un isomorphisme pour tout $n$ (car $x$ et $x'$ coïncident sur
$\Spec(R/\mathfrak m^n)$ pour tout $n$). Puisque $\mathcal{O}_{T, t}$
s'injecte dans son complété (voir
Algebra, Lemma \ref{algebra-lemma-intersect-powers-ideal-module-zero})
nous concluons que $R = \mathcal{O}_{T, t}$. Ainsi, $x$ et $x'$ coïncident
sur $R$.

\medskip\noindent
Surjectivité : soit $(x_n)$ un élément du membre de droite.
Choisissons un schéma $U$ et un morphisme \'etale surjectif $U \to X$. 
Notons $x_0 : \Spec(k) \to X$ le morphisme induit sur le corps résiduel
$k = R/\mathfrak m$. Le morphisme de schémas
$U \times_{X, x_0} \Spec(k) \to \Spec(k)$ est \'etale et surjectif.
Ainsi, $U \times_{X, x_0} \Spec(k)$ est une réunion disjointe non vide de spectres
d'extensions finies séparables de $k$ ; voir
Morphisms, Lemma \ref{morphisms-lemma-etale-over-field}.
Nous pouvons donc trouver une extension finie séparable $k'/k$
et un $k'$-point $u_0 : \Spec(k') \to U$ tels que
$$
\xymatrix{
\Spec(k') \ar[d] \ar[r]_-{u_0} & U \ar[d] \\
\Spec(k) \ar[r]^-{x_0} & X
}
$$
soit commutatif. Soit $R \subset R'$ une extension \'etale finie d'anneaux locaux
noethériens complets qui induit $k'/k$ sur les corps résiduels
(voir Algebra, Lemmas \ref{algebra-lemma-henselian-cat-finite-etale} et
\ref{algebra-lemma-complete-henselian}). Notons $x'_n$ la restriction
de $x_n$ à $\Spec(R'/\mathfrak m^nR')$. D'après
More on Morphisms of Spaces, Lemma
\ref{spaces-more-morphisms-lemma-etale-formally-etale}
on peut trouver un élément
$(u'_n) \in \lim \Mor_S(\Spec(R'/\mathfrak m^nR'), U)$
qui s'envoie sur $(x'_n)$. D'après le Lemme \ref{lemma-map-into-scheme},
la famille $(u'_n)$ provient d'un unique
morphisme $u' : \Spec(R') \to U$. Notons $x' : \Spec(R') \to X$ la
composée. Remarquons que $R' \otimes_R R'$ est un produit fini de spectres d'anneaux
locaux noethériens complets auxquels notre discussion s'applique.
Ainsi, le diagramme
$$
\xymatrix{
\Spec(R' \otimes_R R') \ar[r] \ar[d] & \Spec(R') \ar[d]^{x'} \\
\Spec(R') \ar[r]^{x'} & X
}
$$
est commutatif par l'injectivité démontrée ci-dessus et le fait que
$x'_n$ est la restriction de $x_n$, qui est défini sur $R/\mathfrak m^n$.
Puisque $\{\Spec(R') \to \Spec(R)\}$ est un recouvrement fppf, nous concluons
que $x'$ descend en un morphisme $x : \Spec(R) \to X$.
Nous omettons la démonstration du fait que $x_n$ est la restriction de $x$ à
$\Spec(R/\mathfrak m^n)$.
\end{proof}

\begin{lemma}
\label{lemma-adic-into-completion}
Soit $S$ un schéma. Soit $X$ un espace algébrique sur $S$.
Soit $T \subset |X|$ un sous-ensemble fermé tel que
$X \setminus T \to X$ soit quasi-compact. Soit $R$ une algèbre locale
noethérienne complète sur $S$. Alors un morphisme adique $p : \text{Spf}(R) \to X_{/T}$
correspond à un unique morphisme $g : \Spec(R) \to X$ tel
que $g^{-1}(T) = \{\mathfrak m_R\}$. 
\end{lemma}

\begin{proof}
L'énoncé a un sens parce que $X_{/T}$ est adique* d'après le
Lemme \ref{lemma-formal-completion-types} (et que nous sommes donc
autorisés à employer le terme adique pour les morphismes ; voir la
Définition \ref{definition-adic-morphism}).
Soit $p$ donné. D'après le Lemme \ref{lemma-map-into-algebraic-space},
nous obtenons un unique morphisme $g : \Spec(R) \to X$ correspondant à
la composée $\text{Spf}(R) \to X_{/T} \to X$.
Soit $Z \subset X$ le sous-espace fermé induit réduit de support
$T$. Le morphisme d'inclusion $Z \to X$ correspond à un morphisme
$Z \to X_{/T}$. Puisque $p$ est adique, il est représentable par des espaces
algébriques et nous obtenons
$$
\text{Spf}(R) \times_{X_{/T}} Z = \text{Spf}(R) \times_X Z
$$
est un espace algébrique muni d'une immersion fermée dans $\text{Spf}(R)$.
(L'égalité vaut parce que $X_{/T} \to X$ est un monomorphisme.)
Ainsi, ce produit fibré est égal à $\Spec(R/J)$ pour un certain
idéal $J \subset R$ qui contient $\mathfrak m_R^{n_0}$ pour un certain
$n_0 \geq 1$. Cela implique que $\Spec(R) \times_X Z$
est un sous-schéma fermé de $\Spec(R)$,
disons $\Spec(R) \times_X Z = \Spec(R/I)$, dont l'intersection avec
$\Spec(R/\mathfrak m_R^n)$ pour $n \geq n_0$ est égale à $\Spec(R/J)$.
En termes algébriques, cela signifie que
$I + \mathfrak m_R^n = J + \mathfrak m_R^n = J$ pour tout $n \geq n_0$.
D'après le théorème d'intersection de Krull,
cela implique que $I = J$, ce qui conclut.
\end{proof}








\section{Le petit site \'etale d'un espace algébrique formel}
\label{section-etale-sites}

\noindent
La motivation de la définition suivante vient des schémas
formels classiques : l'espace topologique sous-jacent d'un schéma formel
$(\mathfrak X, \mathcal{O}_\mathfrak X)$
est l'espace topologique sous-jacent de la réduction $\mathfrak X_{red}$.

\medskip\noindent
Une remarque importante est la suivante. Supposons que $X$ soit un espace algébrique
de réduction $X_{red}$ (Properties of Spaces, Definition
\ref{spaces-properties-definition-reduced-induced-space}).
Alors nous avons
$$
X_{spaces, \etale} = X_{red, spaces, \etale},\quad
X_\etale = X_{red, \etale},\quad
X_{affine, \etale} = X_{red, affine, \etale}
$$
d'après More on Morphisms of Spaces, Theorem
\ref{spaces-more-morphisms-theorem-topological-invariance} et
Lemma \ref{spaces-more-morphisms-lemma-topological-invariance}.
Par conséquent, la définition suivante n'entre pas en conflit avec la notion
existante lorsque notre espace algébrique formel se trouve être
un espace algébrique.

\begin{definition}
\label{definition-etale-sites}
Soit $S$ un schéma. Soit $X$ un espace algébrique formel de
réduction $X_{red}$ (Lemme \ref{lemma-reduction-formal-algebraic-space}).
\begin{enumerate}
\item Le {\it petit site \'etale} $X_\etale$ de $X$ est
le site $X_{red, \etale}$ de Properties of Spaces, Definition
\ref{spaces-properties-definition-etale-site}.
\item Le site $X_{spaces, \etale}$ est le site
$X_{red, spaces, \etale}$ de Properties of Spaces, Definition
\ref{spaces-properties-definition-spaces-etale-site}.
\item Le site $X_{affine, \etale}$ est le site
$X_{red, affine, \etale}$ de Properties of Spaces, Lemma
\ref{spaces-properties-lemma-alternative}.
\end{enumerate}
\end{definition}

\noindent
Dans le Lemme \ref{lemma-identify-spaces-etale}, nous verrons que
$X_{spaces, \etale}$ peut se décrire au moyen de
morphismes d'espaces algébriques formels qui sont représentables
par des espaces algébriques et \'etales.
D'après Properties of Spaces, Lemmas
\ref{spaces-properties-lemma-compare-etale-sites} et
\ref{spaces-properties-lemma-alternative}
nous avons les identifications
\begin{equation}
\label{equation-etale-topos}
\Sh(X_\etale) = \Sh(X_{spaces, \etale}) = \Sh(X_{affine, \etale})
\end{equation}
Nous appellerons cela le {\it (petit) topos \'etale} de $X$.

\begin{lemma}
\label{lemma-functoriality-etale-site}
Soit $S$ un schéma.
Soit $f : X \to Y$ un morphisme d'espaces algébriques formels sur $S$.
\begin{enumerate}
\item Il existe un foncteur continu
$Y_{spaces, \etale} \to X_{spaces, \etale}$
qui induit un morphisme de sites
$$
f_{spaces, \etale} : X_{spaces, \etale} \to Y_{spaces, \etale}.
$$
\item La règle $f \mapsto f_{spaces, \etale}$ est compatible avec les
compositions, autrement dit $(f \circ g)_{spaces, \etale}
= f_{spaces, \etale} \circ g_{spaces, \etale}$ (voir
Sites, Definition \ref{sites-definition-composition-morphisms-sites}).
\item Le morphisme de topos associé à $f_{spaces, \etale}$
induit, via (\ref{equation-etale-topos}), un morphisme de topos
$f_{small} : \Sh(X_\etale) \to \Sh(Y_\etale)$
dont la construction est compatible aux compositions.
\end{enumerate}
\end{lemma}

\begin{proof}
Le seul point est que $f$ induit un morphisme des réductions
$X_{red} \to Y_{red}$ d'après le Lemme \ref{lemma-reduction-formal-algebraic-space}.
Ce lemme résulte donc immédiatement du lemme correspondant pour les
morphismes d'espaces algébriques (Properties of Spaces,
Lemma \ref{spaces-properties-lemma-functoriality-etale-site}).
\end{proof}

\noindent
Si le morphisme d'espaces algébriques formels $X \to Y$ est \'etale,
alors le morphisme de topos $\Sh(X_\etale) \to \Sh(Y_\etale)$
est une localisation. Voici un énoncé.

\begin{lemma}
\label{lemma-etale-morphism-topoi}
Soient $S$ un schéma et $f : X \to Y$ un morphisme
d'espaces algébriques formels sur $S$. Supposons $f$ représentable
par des espaces algébriques et \'etale. Dans ce cas, il existe un
foncteur cocontinu $j : X_\etale \to Y_\etale$.
Le morphisme de topos $f_{small}$ est le
morphisme de topos associé à $j$ ; voir
Sites, Lemma \ref{sites-lemma-cocontinuous-morphism-topoi}.
De plus, $j$ est également continu ; ainsi,
Sites, Lemma \ref{sites-lemma-when-shriek} s'applique.
\end{lemma}

\begin{proof}
Cela résultera immédiatement du cas des espaces algébriques
(Properties of Spaces, Lemma
\ref{spaces-properties-lemma-etale-morphism-topoi})
si nous pouvons montrer que le morphisme induit $X_{red} \to Y_{red}$
est \'etale. Remarquons que $X \times_Y Y_{red}$ est un
espace algébrique, \'etale sur l'espace algébrique réduit $Y_{red}$,
et donc lui-même réduit (d'après notre définition des espaces algébriques
réduits dans Properties of Spaces, Section
\ref{spaces-properties-section-types-properties}).
Ainsi, $X_{red} = X \times_Y Y_{red}$, comme souhaité.
\end{proof}

\begin{lemma}
\label{lemma-affine-identify-affine-etale}
Soit $S$ un schéma. Soit $X$ un espace algébrique formel affine sur $S$.
Alors $X_{affine, \etale}$ est équivalent à la catégorie dont les objets
sont les morphismes $\varphi : U \to X$ d'espaces algébriques formels tels que
\begin{enumerate}
\item $U$ est un espace algébrique formel affine,
\item $\varphi$ est représentable par des espaces algébriques et \'etale.
\end{enumerate}
\end{lemma}

\begin{proof}
Notons $\mathcal{C}$ la catégorie introduite dans le lemme.
Remarquons que, pour $\varphi : U \to X$ dans $\mathcal{C}$, le
morphisme $\varphi$ est représentable (par des schémas) et affine ; voir le
Lemme \ref{lemma-affine-representable-by-algebraic-spaces}.
Rappelons que $X_{affine, \etale} = X_{red, affine, \etale}$.
Nous pouvons donc définir un foncteur
$$
\mathcal{C} \longrightarrow X_{affine, \etale},\quad
(U \to X) \longmapsto U \times_X X_{red}
$$
car $U \times_X X_{red}$ est un schéma affine.

\medskip\noindent
Pour achever la démonstration, nous allons construire un foncteur quasi-inverse.
À savoir, écrivons $X = \colim X_\lambda$ comme dans la
Définition \ref{definition-affine-formal-algebraic-space}.
Pour chaque $\lambda$, l'inclusion $X_{red} \subset X_\lambda$
est un épaississement. Ainsi, pour tout $\lambda$, nous avons une
équivalence
$$
X_{red, affine, \etale} = X_{\lambda, affine, \etale}
$$
par exemple d'après
More on Algebra, Lemma \ref{more-algebra-lemma-locally-nilpotent-henselian}.
Ainsi, si $U_{red} \to X_{red}$ est un morphisme \'etale avec
$U_{red}$ affine, nous obtenons un système de morphismes \'etales
$U_\lambda \to X_\lambda$ de schémas affines compatibles aux
morphismes de transition du système qui définit $X$. Nous pouvons donc prendre
$$
U = \colim U_\lambda
$$
comme espace algébrique formel affine sur $X$. La construction donne
$U \times_X X_\lambda = U_\lambda$. Cela montre que $U \to X$ est
représentable et \'etale. Nous omettons la vérification du fait que les constructions
sont mutuellement inverses.
\end{proof}

\begin{lemma}
\label{lemma-affine-etale-mcquillan}
Soit $S$ un schéma. Soit $X$ un espace algébrique
formel affine sur $S$. Supposons que $X$ soit de McQuillan, i.e.,
égal à $\text{Spf}(A)$ pour une certaine
algèbre topologique faiblement admissible sur $S$, notée $A$.
Alors $(X_{affine, \etale})^{opp}$ est équivalent à
la catégorie suivante :
\begin{enumerate}
\item ses objets sont les $A$-algèbres de la forme
$B^\wedge = \lim B/JB$, où $A \to B$ est un homomorphisme d'anneaux \'etale
et $J$ parcourt les idéaux faibles de définition de $A$, et
\item ses morphismes sont les homomorphismes continus de $A$-algèbres.
\end{enumerate}
\end{lemma}

\begin{proof}
Combiner les Lemmes \ref{lemma-affine-identify-affine-etale} et \ref{lemma-etale}.
\end{proof}

\begin{lemma}
\label{lemma-identify-spaces-etale}
Soit $S$ un schéma. Soit $X$ un espace algébrique formel sur $S$.
Alors $X_{spaces, \etale}$ est équivalent à la catégorie dont les objets
sont les morphismes $\varphi : U \to X$ d'espaces algébriques formels tels que
$\varphi$ soit représentable par des espaces algébriques et \'etale.
\end{lemma}

\begin{proof}
Notons $\mathcal{C}$ la catégorie introduite dans le lemme.
Rappelons que $X_{spaces, \etale} = X_{red, spaces, \etale}$.
Nous pouvons donc définir un foncteur
$$
\mathcal{C} \longrightarrow X_{spaces, \etale},\quad
(U \to X) \longmapsto U \times_X X_{red}
$$
car $U \times_X X_{red}$ est un espace algébrique \'etale sur $X_{red}$.

\medskip\noindent
Pour achever la démonstration, nous allons construire un foncteur quasi-inverse.
Choisissons un objet $\psi : V \to X_{red}$ de $X_{red, spaces, \etale}$.
Considérons le foncteur
$U_{V, \psi} : (\Sch/S)_{fppf} \to \textit{Sets}$ donné par
$$
U_{V, \psi}(T) = \{(a, b) \mid
a : T \to X,
\ b : T \times_{a, X} X_{red} \to V,
\ \psi \circ b = a|_{T \times_{a, X} X_{red}}\}
$$
Nous affirmons que la transformation $U_{V, \psi} \to X$, $(a, b) \mapsto a$
définit un objet de la catégorie $\mathcal{C}$.
Montrons d'abord que $U_{V, \psi}$ est un espace algébrique formel.
Remarquons que $U_{V, \psi}$ est un faisceau pour la topologie fppf (certains détails
sont omis). Supposons ensuite que $X_i \to X$ soit un recouvrement \'etale par des
espaces algébriques formels affines comme dans la
Définition \ref{definition-formal-algebraic-space}.
Posons $V_i = V \times_{X_{red}} X_{i, red}$ et notons
$\psi_i : V_i \to X_{i, red}$ la projection. Alors 
nous avons
$$
U_{V, \psi} \times_X X_i = U_{V_i, \psi_i}
$$
par un argument formel, car $X_{i, red} = X_i \times_X X_{red}$
(puisque $X_i \to X$ est représentable par des espaces algébriques et \'etale).
Il suffit donc de montrer que $U_{V_i, \psi_i}$ est un
espace algébrique formel affine, car nous disposerons alors
d'un recouvrement $U_{V_i, \psi_i} \to U_{V, \psi}$ comme dans la
Définition \ref{definition-formal-algebraic-space}.
D'autre part, nous avons vu dans la démonstration du
Lemme \ref{lemma-etale-morphism-topoi}
que $\psi_i : V_i \to X_i$ est le changement de
base d'un morphisme représentable et \'etale
$U_i \to X_i$ d'espaces algébriques formels affines.
Il n'est alors pas difficile de voir que $U_i = U_{V_i, \psi_i}$,
comme souhaité.

\medskip\noindent
Nous omettons la vérification du fait que $U_{V, \psi} \to X$
est représentable par des espaces algébriques et \'etale.
Nous obtenons ainsi notre foncteur $(V, \psi) \mapsto (U_{V, \psi} \to X)$
dans l'autre sens.
Nous omettons la vérification du fait que les constructions
sont mutuellement inverses.
\end{proof}

\begin{lemma}
\label{lemma-identify-affine-etale}
Soit $S$ un schéma. Soit $X$ un espace algébrique formel sur $S$.
Alors $X_{affine, \etale}$ est équivalent à la catégorie dont les objets
sont les morphismes $\varphi : U \to X$ d'espaces algébriques formels tels que
\begin{enumerate}
\item $U$ est un espace algébrique formel affine,
\item $\varphi$ est représentable par des espaces algébriques et \'etale.
\end{enumerate}
\end{lemma}

\begin{proof}
Cela résulte de la combinaison des Lemmes \ref{lemma-identify-spaces-etale} et
\ref{lemma-characterize-affine}.
\end{proof}






\section{Le faisceau structural}
\label{section-structure-sheaf}


\noindent
Comme dans la section \ref{section-etale-sites}, la discussion de la présente
section est motivée par les schémas formels classiques. Ceux-ci sont définis comme des couples
$(\mathfrak X, \mathcal{O}_\mathfrak X)$ où $\mathcal{O}_\mathfrak X$
est un faisceau d'anneaux topologiques.

\medskip\noindent
Soit $X$ un espace algébrique formel. Un {\it faisceau structural de $X$}
est un faisceau d'anneaux topologiques $\mathcal{O}_X$ sur le site \'etale
$X_\etale$ (défini dans la section \ref{section-etale-sites}) tel que
$$
\mathcal{O}_X(U_{red}) = \lim \Gamma(U_\lambda, \mathcal{O}_{U_\lambda})
$$
comme anneaux topologiques chaque fois que
\begin{enumerate}
\item $\varphi : U \to X$ est un morphisme d'espaces algébriques formels,
\item $U$ est un espace algébrique formel affine,
\item $\varphi$ est représentable par des espaces algébriques et \'etale,
\item $U_{red} \to X_{red}$ est l'objet affine correspondant de
$X_\etale$, voir le
lemme \ref{lemma-identify-affine-etale},
\item $U = \colim U_\lambda$ est une présentation de $U$ comme limite inductive au sens de la
définition \ref{definition-affine-formal-algebraic-space}.
\end{enumerate}
Les faisceaux structuraux existent, mais peuvent se comporter de façon inattendue.

\begin{lemma}
\label{lemma-structure-sheaf}
Tout espace algébrique formel possède un faisceau structural.
\end{lemma}

\begin{proof}
Soient $S$ un schéma et $X$ un espace algébrique formel sur $S$.
D'après (\ref{equation-etale-topos}), il suffit de construire $\mathcal{O}_X$
comme faisceau d'anneaux topologiques sur $X_{affine, \etale}$.
Notons $\mathcal{C}$ la catégorie dont les objets
sont les morphismes $\varphi : U \to X$ d'espaces algébriques formels tels que
$U$ soit un espace algébrique formel affine et que
$\varphi$ soit représentable par des espaces algébriques et \'etale.
D'après le lemme \ref{lemma-identify-affine-etale},
le foncteur $U \mapsto U_{red}$ est une équivalence de catégories
$\mathcal{C} \to X_{affine, \etale}$. La règle donnée
avant le lemme fournit donc déjà $\mathcal{O}_X$ comme préfaisceau
d'anneaux topologiques sur $X_{affine, \etale}$. Il suffit ainsi de vérifier
la condition de faisceau.

\medskip\noindent
Par définition de $X_{affine, \etale}$, un recouvrement correspond à une famille finie
$\{g_i : U_i \to U\}_{i = 1, \ldots, n}$ de morphismes de $\mathcal{C}$
telle que $\{U_{i, red} \to U_{red}\}$ soit un recouvrement \'etale.
Les morphismes $g_i$ sont représentables par des espaces algébriques
(lemme \ref{lemma-permanence-representable}), donc affines
(lemme \ref{lemma-affine-representable-by-algebraic-spaces}).
Ensuite, $g_i$ est \'etale (cela résulte formellement du
lemme \ref{spaces-properties-lemma-etale-permanence} de Properties of Spaces,
puisque $U_i$ et $U$ sont \'etales sur $X$ au sens de la
section \ref{bootstrap-section-representable-by-spaces-properties} de Bootstrap).
Enfin, écrivons $U = \colim U_\lambda$ comme dans la
définition \ref{definition-affine-formal-algebraic-space}.

\medskip\noindent
Ces préparatifs étant faits, démontrons la propriété de
faisceau comme suit. Pour chaque $\lambda$, posons
$U_{i, \lambda} = U_i \times_U U_\lambda$ et
$U_{ij, \lambda} = (U_i \times_U U_j) \times_U U_\lambda$.
Selon ce qui précède, ce sont des schémas affines, $\{U_{i, \lambda} \to U_\lambda\}$
est un recouvrement \'etale, et
$U_{ij, \lambda} = U_{i, \lambda} \times_{U_\lambda} U_{j, \lambda}$.
De plus, $U_i = \colim U_{i, \lambda}$ et
$U_i \times_U U_j = \colim U_{ij, \lambda}$.
Pour chaque $\lambda$, on a une suite exacte
$$
0 \to
\Gamma(U_\lambda, \mathcal{O}_{U_\lambda}) \to
\prod\nolimits_i \Gamma(U_{i, \lambda}, \mathcal{O}_{U_{i, \lambda}}) \to
\prod\nolimits_{i, j} \Gamma(U_{ij, \lambda}, \mathcal{O}_{U_{ij, \lambda}})
$$
car le faisceau structural sur $U_\lambda$ vérifie la condition de faisceau
pour la topologie \'etale
(voir la proposition de \'Etale Cohomology
\ref{etale-cohomology-proposition-quasi-coherent-sheaf-fpqc}).
Comme les limites commutent aux limites, la limite projective de ces
suites exactes est une suite exacte
$$
0 \to
\lim \Gamma(U_\lambda, \mathcal{O}_{U_\lambda}) \to
\prod\nolimits_i \lim \Gamma(U_{i, \lambda}, \mathcal{O}_{U_{i, \lambda}}) \to
\prod\nolimits_{i, j} \lim
\Gamma(U_{ij, \lambda}, \mathcal{O}_{U_{ij, \lambda}})
$$
Cela signifie que
$$
0 \to
\mathcal{O}_X(U_{red}) \to
\prod\nolimits_i \mathcal{O}_X(U_{i, red}) \to
\prod\nolimits_{i, j} \mathcal{O}_X((U_i \times_U U_j)_{red})
$$
est exacte en tant que suite de groupes abéliens. Il reste à
montrer que la topologie limite sur $\mathcal{O}_X(U_{red})$ est induite
par la topologie limite sur
$\prod_{i = 1, \ldots, n} \mathcal{O}_X(U_{i, red})$.
Les noyaux des applications
$\mathcal{O}_X(U_{red}) \to \mathcal{O}_{U_\lambda}(U_\lambda)$
forment un système fondamental de sous-modules ouverts. Comme ils sont les images
réciproques des noyaux des applications
$\prod_i \mathcal{O}_X(U_{i, red}) \to
\prod_i \mathcal{O}_{U_{i, \lambda}}(U_{i, \lambda})$
et que ces noyaux forment un système fondamental de sous-modules ouverts
pour $\mathcal{O}_X(U_{red}) \to \mathcal{O}_{U_\lambda}(U_\lambda)$,
le résultat est acquis.
\end{proof}

\begin{remark}
\label{remark-not-enough-sections}
Le faisceau structural ne possède pas toujours « assez de sections ».
Dans la section \ref{examples-section-affine-formal-algebraic-space} de Examples,
nous avons vu qu'il existe des espaces algébriques formels affines qui
ne sont pas de McQuillan, et même des exemples dont les points ne sont pas
séparés par les fonctions régulières.
\end{remark}

\noindent
Dans le lemme suivant, nous montrons que le faisceau structural d'un schéma formel
affine à indexation dénombrable a une cohomologie supérieure nulle.
Sans l'hypothèse d'indexation dénombrable, ce résultat est sans doute faux en général.

\begin{lemma}
\label{lemma-higher-vanishing-structure-sheaf}
Si $X$ est un espace algébrique formel affine à indexation dénombrable, alors
$H^n(X_\etale, \mathcal{O}_X) = 0$ pour $n > 0$.
\end{lemma}

\begin{proof}
Nous pouvons travailler avec $X_{affine, \etale}$, puisque celui-ci donne le même topos.
Nous allons appliquer le lemme
\ref{sites-cohomology-lemma-cech-vanish-collection}
de Cohomology on Sites pour établir cette annulation. Comme $X_{affine, \etale}$
possède les sommes disjointes finies, on est ramené au complexe de {\v C}ech
d'un recouvrement donné par une seule flèche $\{U_{red} \to V_{red}\}$
dans $X_{affine, \etale} = X_{red, affine, \etale}$
(voir \'Etale Cohomology, lemme \ref{etale-cohomology-lemma-cech-complex}).
Il faut donc montrer que
$$
0 \to \mathcal{O}_X(V_{red}) \to \mathcal{O}_X(U_{red}) \to
\mathcal{O}_X(U_{red} \times_{V_{red}} U_{red}) \to \ldots
$$
est exacte. Nous le ferons ci-dessous dans le cas $V_{red} = X_{red}$.
Le cas général se démontre exactement de la même manière.

\medskip\noindent
Rappelons que $X = \text{Spf}(A)$, où $A$ est un
anneau topologique faiblement admissible possédant un système
fondamental dénombrable d'idéaux faibles de définition.
Nous avons vu dans les lemmes
\ref{lemma-affine-identify-affine-etale} et
\ref{lemma-affine-etale-mcquillan}
que l'objet $U_{red}$ de $X_{affine, \etale}$
correspond à un morphisme $U \to X$ d'espaces algébriques formels
affines, représentable par des espaces algébriques et \'etale,
et que $U = \text{Spf}(B^\wedge)$, où $B$ est une $A$-algèbre \'etale.
La règle définissant le faisceau structural donne
$$
\mathcal{O}_X(U_{red}) = B^\wedge
$$
Rappelons que $B^\wedge = \lim B/JB$, où la limite porte sur les
idéaux faibles de définition $J \subset A$.
En développant les définitions, on obtient
$$
\mathcal{O}_X(U_{red} \times_{X_{red}} U_{red}) = (B \otimes_A B)^\wedge
$$
et ainsi de suite. Puisque $U \to X$ est un recouvrement, l'application $A \to B$
est fidèlement plate, voir le lemme \ref{lemma-etale-surjective}.
Le complexe
$$
0 \to A \to B \to B \otimes_A B \to B \otimes_A B \otimes_A B \to \ldots
$$
est donc {\bf universellement} exact, voir le lemme \ref{descent-lemma-ff-exact} de Descent.
Notre but est de montrer que
$$
H^n(0 \to A^\wedge \to B^\wedge \to (B \otimes_A B)^\wedge
\to (B \otimes_A B \otimes_A B)^\wedge \to \ldots)
$$
est nul pour $n > 0$. Pour comprendre la situation, décomposons notre
complexe exact (avant passage au complété) en suites exactes courtes
$$
0 \to A \to B \to M_1 \to 0,\quad
0 \to M_i \to B^{\otimes_A i + 1} \to M_{i + 1} \to 0
$$
D'après ce qui précède, ce sont des suites exactes courtes universellement
exactes. Ainsi $JM_i = M_i \cap J(B^{\otimes_A i + 1})$ pour
tout idéal $J$ de $A$. En particulier, la
topologie de $M_i$ comme sous-module de $B^{\otimes_A i + 1}$
est identique à celle de $M_i$ comme module quotient de
$B^{\otimes_A i}$. Comme il existe dans $A$ un système fondamental dénombrable
d'idéaux faibles de définition, les suites
$$
0 \to A^\wedge \to B^\wedge \to M_1^\wedge \to 0,\quad
0 \to M_i^\wedge \to (B^{\otimes_A i + 1})^\wedge \to M_{i + 1}^\wedge \to 0
$$
restent exactes d'après le lemme \ref{lemma-ses}. Cela démontre le lemme.
\end{proof}

\begin{remark}
\label{remark-bad-quasi-coherent}
Même si le faisceau structural possède de bonnes propriétés, cela ne
signifie pas qu'il existe une bonne théorie des modules quasi-cohérents. Par exemple,
dans la section \ref{examples-section-nonabelian-QCoh} de Examples,
nous avons vu que, pour presque tout espace algébrique formel affine noethérien,
la notion la plus naturelle de module quasi-cohérent conduit à une
catégorie de modules qui n'est pas abélienne.
\end{remark}





\section{Limites inductives d'espaces algébriques formels}
\label{section-colimits-of-formal}

\noindent
Dans cette section, nous généralisons le résultat de la
section \ref{section-global-colimits}
au cas des systèmes de morphismes d'espaces algébriques formels.
Remarquons que, dans les lemmes ci-dessous, la condition
« $f_{\lambda \mu} : X_\lambda \to X_\mu$ est une immersion fermée
induisant un isomorphisme $X_{\lambda, red} \to X_{\mu, red}$ »
peut se reformuler ainsi :
« $f_{\lambda \mu}$ est représentable et est un épaississement ».

\begin{lemma}
\label{lemma-colimit-affine-formal}
Soit $S$ un schéma. Donnons-nous un ensemble ordonné filtrant
$\Lambda$ et un système d'espaces algébriques formels affines
$(X_\lambda, f_{\lambda \mu})$ indexé par $\Lambda$, où chaque
$f_{\lambda \mu} : X_\lambda \to X_\mu$ est une immersion fermée
induisant un isomorphisme $X_{\lambda, red} \to X_{\mu, red}$.
Alors $X = \colim_{\lambda \in \Lambda} X_\lambda$
est un espace algébrique formel affine sur $S$.
\end{lemma}

\begin{proof}
On peut écrire
$X_\lambda = \colim_{\omega \in \Omega_\lambda} X_{\lambda, \omega}$
comme limite inductive de schémas affines indexée par un ensemble ordonné filtrant
$\Omega_\lambda$, de telle sorte que les morphismes de transition
$X_{\lambda, \omega} \to X_{\lambda, \omega'}$
soient des épaississements. Pour tous $\lambda, \mu \in \Lambda$
et $\omega \in \Omega_\lambda$, avec $\mu \geq \lambda$,
il existe un $\omega' \in \Omega_\mu$ tel que le morphisme
$X_{\lambda, \omega} \to X_\mu$ se factorise par $X_{\mu, \omega'}$, voir le
lemme \ref{lemma-factor-through-thickening}. Le
morphisme $X_{\lambda, \omega} \to X_{\mu, \omega'}$
est alors une immersion fermée induisant un isomorphisme sur les réductions,
donc un épaississement. Posons
$\Omega = \coprod_{\lambda \in \Lambda} \Omega_\lambda$
et convenons que $(\lambda, \omega) \leq (\mu, \omega')$ si et seulement si
$\lambda \leq \mu$ et si $X_{\lambda, \omega} \to X_\mu$ se factorise par
$X_{\mu, \omega'}$. Il résulte de ce qui précède que
$\Omega$ est un ensemble ordonné filtrant et que
$X = \colim_{\lambda \in \Lambda} X_\lambda =
\colim_{(\lambda, \omega) \in \Omega} X_{\lambda, \omega}$.
Ceci achève la démonstration.
\end{proof}

\begin{lemma}
\label{lemma-colimit-formal-spaces-is-formal-space}
Soit $S$ un schéma. Donnons-nous un ensemble ordonné filtrant
$\Lambda$ et un système d'espaces algébriques formels
$(X_\lambda, f_{\lambda \mu})$ indexé par $\Lambda$, où chaque
$f_{\lambda \mu} : X_\lambda \to X_\mu$ est une immersion fermée
induisant un isomorphisme $X_{\lambda, red} \to X_{\mu, red}$.
Alors $X = \colim_{\lambda \in \Lambda} X_\lambda$
est un espace algébrique formel sur $S$.
\end{lemma}

\begin{proof}
Comme la limite inductive est prise dans la catégorie des faisceaux fppf,
$X$ est un faisceau. Choisissons et fixons $\lambda \in \Lambda$. Choisissons
un recouvrement $\{X_{i, \lambda} \to X_\lambda\}$ comme dans la définition
\ref{definition-formal-algebraic-space}. En particulier,
En particulier, $\{X_{i, \lambda, red} \to X_{\lambda, red}\}$ est un
recouvrement \'etale par des schémas affines.
Pour chaque $\mu \geq \lambda$, il existe un diagramme cartésien
$$
\xymatrix{
X_{i, \lambda} \ar[r] \ar[d] & X_{i, \mu} \ar[d] \\
X_\lambda \ar[r] & X_\mu
}
$$
dont les flèches verticales sont \'etales. En effet, le morphisme \'etale
$X_{i, \lambda, red} \to X_{\lambda, red} = X_{\mu, red}$
correspond à un morphisme \'etale $X_{i, \mu} \to X_\mu$
d'espaces algébriques formels, où $X_{i, \mu}$ est un espace
algébrique formel affine, voir le
lemme \ref{lemma-affine-identify-affine-etale}.
Le même lemme montre que le changement de base de $X_{i, \mu}$ à $X_\lambda$
coïncide avec $X_{i, \lambda}$. Il en résulte aussi que
$X_{i, \mu} = X_\mu \times_{X_{\mu'}} X_{i, \mu'}$ pour
$\mu' \geq \mu \geq \lambda$. Posons $X_i = \colim X_{i, \mu}$.
Alors $X_{i, \mu} = X_i \times_X X_\mu$ (comme foncteurs).
Comme tout morphisme $T \to X = \colim X_\mu$
issu d'un schéma affine (ou quasi-compact) $T$ se factorise par $X_\mu$
pour un certain $\mu$, on conclut que
$\colim X_{i, \mu} \to \colim X_\mu$ est \'etale.
Ainsi, s'il est possible de montrer que $\colim X_{i, \mu}$ est un espace
algébrique formel affine, le lemme en résultera.
Notons que les morphismes $X_{i, \mu} \to X_{i, \mu'}$
sont des immersions fermées, car ils s'obtiennent par changement de base de l'immersion fermée
$X_\mu \to X_{\mu'}$. Enfin, le morphisme
$X_{i, \mu, red} \to X_{i, \mu', red}$ est un isomorphisme,
puisque $X_{\mu, red} \to X_{\mu', red}$ en est un.
On est donc ramené au cas traité dans le
lemme \ref{lemma-colimit-affine-formal}.
\end{proof}








\section{Recomplétion}
\label{section-recompletion}

\noindent
Dans cette section, nous définissons le complété d'un espace algébrique
formel le long d'une partie fermée de sa réduction. C'est la généralisation
naturelle de la section \ref{section-completion}.

\begin{lemma}
\label{lemma-completion-affine-formal-is-affine-formal}
Soit $S$ un schéma. Soit $X$ un espace algébrique formel affine sur $S$.
Soit $T \subset |X_{red}|$ une partie fermée. Alors le foncteur
$$
X_{/T} : (\Sch/S)_{fppf} \longrightarrow \textit{Sets},\quad
U \longmapsto \{f : U \to X : f(|U|) \subset T\}
$$
est un espace algébrique formel affine.
\end{lemma}

\begin{proof}
Écrivons $X = \colim X_\lambda$ comme dans la
définition \ref{definition-affine-formal-algebraic-space}.
Alors $X_{\lambda, red} = X_{red}$ et nous considérons
$T$ comme une partie fermée de $|X_\lambda| = |X_{\lambda, red}|$.
D'après le lemme \ref{lemma-completion-affine-is-affine-formal-algebraic-space},
pour chaque $\lambda$, le complété
$(X_\lambda)_{/T}$ est un espace algébrique formel affine.
Les morphismes de transition $(X_\lambda)_{/T} \to (X_\mu)_{/T}$ sont des
immersions fermées, car ils s'obtiennent par changement de base des morphismes de transition
$X_\lambda \to X_\mu$, voir le lemme \ref{lemma-map-completions-representable}.
De plus, les morphismes $((X_\lambda)_{/T})_{red} \to ((X_\mu)_/T)_{red}$
sont des isomorphismes d'après le lemme \ref{lemma-reduction-completion}.
Comme $X_{/T} = \colim (X_\lambda)_{/T}$, on conclut
par le lemme \ref{lemma-colimit-affine-formal}.
\end{proof}

\begin{lemma}
\label{lemma-completion-fas-is-fas}
Soit $S$ un schéma. Soit $X$ un espace algébrique formel sur $S$.
Soit $T \subset |X_{red}|$ une partie fermée. Alors le foncteur
$$
X_{/T} : (\Sch/S)_{fppf} \longrightarrow \textit{Sets},\quad
U \longmapsto \{f : U \to X \mid f(|U|) \subset T\}
$$
est un espace algébrique formel.
\end{lemma}

\begin{proof}
Le foncteur $X_{/T}$ est un faisceau fppf, car si
$\{U_i \to U\}$ est un recouvrement fppf, alors
$\coprod |U_i| \to |U|$ est surjectif.

\medskip\noindent
Choisissons un recouvrement $\{g_i : X_i \to X\}_{i \in I}$ comme dans la définition
\ref{definition-formal-algebraic-space}.
Les morphismes $X_i \times_X X_{/T} \to X_{/T}$ sont \'etales
(voir le lemme de Spaces
\ref{spaces-lemma-base-change-representable-transformations-property})
et l'application $\coprod X_i \times_X X_{/T} \to X_{/T}$ est une surjection de
faisceaux. Il suffit donc de montrer que $X_{/T} \times_X X_i$
est un espace algébrique formel affine. Un point à valeurs dans $U$ de
$X_i \times_X X_{/T}$ est un morphisme $U \to X_i$ dont l'image est
contenue dans la partie fermée
$|g_{i, red}|^{-1}(T) \subset |X_{i, red}|$. Cela résulte donc du
lemme \ref{lemma-completion-affine-formal-is-affine-formal}.
\end{proof}

\begin{definition}
\label{definition-completion-formal-algebraic-space}
Soit $S$ un schéma. Soit $X$ un espace algébrique formel sur $S$.
Soit $T \subset |X_{red}|$ une partie fermée. L'espace algébrique formel
$X_{/T}$ du lemme \ref{lemma-completion-is-formal-algebraic-space}
est appelé le {\it complété de $X$ le long de $T$}.
\end{definition}

\noindent
Soit $f : X \to X'$ un morphisme d'espaces algébriques formels sur un schéma $S$.
Supposons que $T \subset |X_{red}|$ et $T' \subset |X'_{red}|$
soient des parties fermées telles que $|f_{red}|(T) \subset T'$. Il est alors clair que
$f$ définit un morphisme d'espaces algébriques formels
$$
X_{/T} \longrightarrow X'_{/T'}
$$
entre les complétés.

\begin{lemma}
\label{lemma-map-recompletions}
Soit $S$ un schéma. Soit $f : X' \to X$ un morphisme
d'espaces algébriques formels sur $S$. Soit $T \subset |X_{red}|$
une partie fermée et soit $T' = |f_{red}|^{-1}(T) \subset |X'_{red}|$.
Alors
$$
\xymatrix{
X'_{/T'} \ar[r] \ar[d] & X' \ar[d]^f \\
X_{/T} \ar[r] & X
}
$$
est un diagramme cartésien d'espaces algébriques formels sur $S$.
\end{lemma}

\begin{proof}
Observons en effet que les flèches horizontales sont des monomorphismes
par construction. Il suffit donc de montrer qu'un morphisme
$g : U \to X'$ issu d'un schéma $U$ définit un point de $X'_{/T}$
si et seulement si $f \circ g$ définit un point de $X_{/T}$.
Autrement dit, il faut montrer que
$g(U)$ est contenu dans $T' \subset |X'_{red}|$
si et seulement si $(f \circ g)(U)$ est contenu dans $T \subset |X_{red}|$.
Cela résulte immédiatement du choix de $T'$ comme
image réciproque de $T$.
\end{proof}

\begin{lemma}
\label{lemma-reduction-recompletion}
Soit $S$ un schéma. Soit $X$ un espace algébrique formel sur $S$.
Soit $T \subset |X_{red}|$ une partie fermée. La réduction
$(X_{/T})_{red}$ du complété $X_{/T}$ de $X$ le long de $T$ est
le sous-espace fermé induit réduit $Z$ de $X_{red}$ correspondant à $T$.
\end{lemma}

\begin{proof}
Il résulte du lemme \ref{lemma-reduction-formal-algebraic-space},
de la définition de Properties of Spaces
\ref{spaces-properties-definition-reduced-induced-space}
(qui utilise le lemme de Properties of Spaces
\ref{spaces-properties-lemma-reduced-closed-subspace} pour construire $Z$),
et de la définition de $X_{/T}$ que
$Z$ et $(X_{/T})_{red}$ sont des espaces algébriques réduits
caractérisés par la même propriété universelle :
un morphisme $g : Y \to X$ dont la source est un espace algébrique réduit
se factorise par eux si et seulement si $|Y|$ s'envoie dans $T \subset |X|$.
\end{proof}

\begin{lemma}
\label{lemma-recompletion-affine-types}
Soit $S$ un schéma. Soit $X$ un espace algébrique formel affine sur $S$.
Soit $T \subset X_{red}$ une partie fermée et soit $X_{/T}$
le complété formel de $X$ le long de $T$. Alors
\begin{enumerate}
\item $X_{/T}$ est un espace algébrique formel affine,
\item si $X$ est de McQuillan, alors $X_{/T}$ est de McQuillan,
\item si $|X_{red}| \setminus T$ est quasi-compact et si $X$
est à indexation dénombrable, alors $X_{/T}$ est à indexation dénombrable,
\item si $|X_{red}| \setminus T$ est quasi-compact et si $X$
est adique*, alors $X_{/T}$ est adique*,
\item si $X$ est noethérien, alors $X_{/T}$ est noethérien.
\end{enumerate}
\end{lemma}

\begin{proof}
L'assertion (1) est le lemme \ref{lemma-completion-affine-formal-is-affine-formal}.
Si $X$ est de McQuillan, alors $X = \text{Spf}(A)$ pour un anneau topologique
faiblement admissible $A$. Alors $X_{/T} \to X \to \Spec(A)$ vérifie
la propriété (2) du lemme \ref{lemma-mcquillan-affine-formal-algebraic-space} ;
par conséquent, $X_{/T}$ est de McQuillan, voir la
définition \ref{definition-types-affine-formal-algebraic-space}.

\medskip\noindent
Supposons $X$ et $T$ comme dans (3).
Alors $X = \text{Spf}(A)$, où $A$ possède un système fondamental
$A \supset I_1 \supset I_2 \supset I_3 \supset \ldots$
constitué d'idéaux faibles de définition, voir le lemme \ref{lemma-countably-indexed}.
D'après le lemme \ref{algebra-lemma-qc-open} de Algebra,
on peut trouver un idéal de type fini
$\overline{J} = (\overline{f}_1, \ldots, \overline{f}_r) \subset A/I_1$
tel que $T$ soit défini par $\overline{J}$ dans $\Spec(A/I_1) = |X_{red}|$.
Choisissons des relèvements $f_i \in A$ de $\overline{f}_i$.
Si $Z = \Spec(B)$ est un schéma affine et si $g : Z \to X$ est
un morphisme tel que $g(Z) \subset T$ (au sens ensembliste), alors
$g^\sharp : A \to B$ se factorise par $A/I_n$ pour un certain $n$
et $g^\sharp(f_i)$ est nilpotent dans $B$ pour tout $i$. Ainsi,
$J_{m, n} = (f_1, \ldots, f_r)^m + I_n$ s'envoie sur zéro dans $B$
pour certains $n, m \geq 1$. Il en résulte que $X_{/T}$ est le spectre formel de
$\lim_{n, m} A/J_{m, n}$ et est donc à indexation dénombrable.
Cela démontre (3).

\medskip\noindent
Démonstration de (4). L'argument est le même que dans (3).
Toutefois, on peut ici choisir $I_n = I^n$ pour un idéal de type fini
$I \subset A$. Il est alors clair que $X_{/T}$ est le spectre formel de
$\lim A/J^n$, où $J = (f_1, \ldots, f_r) + I$. Quelques détails sont omis.

\medskip\noindent
Démonstration de (5). Dans ce cas, $X_{red}$ est le spectre d'un anneau noethérien
et l'hypothèse que $|X_{red}| \setminus T$ est quasi-compact
est donc satisfaite. Comme dans la démonstration de (4), on voit alors que
$X_{/T}$ est le spectre de $\lim A/J^n$, qui est un anneau topologique
adique noethérien, voir le
lemme \ref{algebra-lemma-completion-Noetherian-Noetherian} de Algebra.
\end{proof}

\begin{lemma}
\label{lemma-recompletion-types}
Soit $S$ un schéma. Soit $X$ un espace algébrique formel sur $S$.
Soit $T \subset X_{red}$ une partie fermée et soit $X_{/T}$
le complété formel de $X$ le long de $T$. Alors
\begin{enumerate}
\item si $X_{red} \setminus T \to X_{red}$ est quasi-compact et si $X$
est localement à indexation dénombrable, alors $X_{/T}$ l'est aussi,
\item si $X_{red} \setminus T \to X_{red}$ est quasi-compact et si $X$
est localement adique*, alors $X_{/T}$ l'est aussi, et
\item si $X$ est localement noethérien, alors $X_{/T}$ l'est aussi.
\end{enumerate}
\end{lemma}

\begin{proof}
Choisissons un recouvrement $\{X_i \to X\}$ comme dans la
définition \ref{definition-formal-algebraic-space}.
Soit $T_i \subset X_{i, red}$ l'image réciproque de $T$.
On a $X_i \times_X X_{/T} = (X_i)_{/T_i}$
(lemme \ref{lemma-map-recompletions}).
Ainsi, $\{(X_i)_{/T_i} \to X_{/T}\}$ est un recouvrement comme dans la
définition \ref{definition-formal-algebraic-space}.
De plus, si $X_{red} \setminus T \to X_{red}$ est quasi-compact, alors
$X_{i, red} \setminus T_i \to X_{i, red}$ l'est aussi ;
et si $X$ est localement à indexation dénombrable, localement adique* ou
localement noethérien, alors $X_i$ est respectivement à indexation dénombrable, adique*
ou noethérien. Le lemme résulte donc du cas affine, à savoir le
lemme \ref{lemma-recompletion-affine-types}.
\end{proof}








\section{Complété le long d'un sous-espace fermé}
\label{section-completion-subspace}

\noindent
Cette section est l'analogue de la section \ref{section-completion}
pour les complétés le long d'un sous-espace fermé.

\begin{definition}
\label{definition-completion-subspace}
Soit $S$ un schéma. Soit $X$ un espace algébrique sur $S$.
Soit $Z \subset X$ un sous-espace fermé et notons $Z_n \subset X$
le $n$-ième voisinage infinitésimal. L'espace algébrique formel
$$
X^\wedge_Z = \colim Z_n
$$
(voir le lemme \ref{lemma-colimit-formal-spaces-is-formal-space})
est appelé le {\it complété de $X$ le long de $Z$}.
\end{definition}

\noindent
Par exemple, si $X = \Spec(A)$ et si $Z$ est défini par l'idéal $I \subset A$,
alors
$$
X^\wedge_Z = \colim \Spec(A/I^n) = \text{Spf}(B)
$$
où $B = \lim A/I^n$ est muni de la topologie limite. Notons que $B$ est un
anneau topologique faiblement adique, mais n'est pas adique en général.

\medskip\noindent
Revenons au cas général. Si $T = |Z|$, il existe un morphisme canonique
$X^\wedge_Z \to X_{/T}$ qui compare les complétés le long de $Z$
et de $T$ (section \ref{section-completion}) et qui n'est pas nécessairement un isomorphisme
(voir le lemme \ref{lemma-when-isomorphism}).

\medskip\noindent
Soit $f : X \to X'$ un morphisme d'espaces algébriques sur un schéma $S$.
Supposons que $Z \subset X$ et $Z' \subset X'$ soient des sous-espaces fermés
tels que $f|_Z$ envoie $Z$ dans $Z'$ et induise un morphisme $Z \to Z'$.
Il est alors clair que $f$ définit un morphisme d'espaces algébriques formels
$$
X^\wedge_Z \longrightarrow (X')^\wedge_{Z'}
$$
entre les complétés.

\begin{lemma}
\label{lemma-map-completions-subspaces-representable}
Soit $S$ un schéma. Soit $f : X' \to X$ un morphisme
d'espaces algébriques sur $S$. Soit $Z \subset X$
un sous-espace fermé et soit $Z' = f^{-1}(Z) = X' \times_X Z$.
Alors
$$
\xymatrix{
(X')^\wedge_{Z'} \ar[r] \ar[d] & X' \ar[d]^f \\
X^\wedge_Z \ar[r] & X
}
$$
est un diagramme cartésien de faisceaux. En particulier, le morphisme
$(X')^\wedge_{Z'} \to X^\wedge_Z$ est représentable par des espaces algébriques.
\end{lemma}

\begin{proof}
Supposons en effet que $Y \to X$ soit un morphisme d'un schéma dans $X$ tel
que $Y \to X$ se factorise par $Z$. Alors $Y \times_X X' \to X$
est un morphisme d'espaces algébriques tel que $Y \times_X X' \to X'$
se factorise par $Z'$. Puisque $Z'_n = X' \times_X Z_n$ pour tout $n \geq 1$,
il en va de même pour les voisinages infinitésimaux.
Le carré cartésien de foncteurs résulte donc
des formules $X^\wedge_Z = \colim Z_n$ et
$(X')^\wedge_{Z'} = \colim Z'_n$.
\end{proof}

\begin{lemma}
\label{lemma-reduction-completion-subspace}
Soit $S$ un schéma. Soit $X$ un espace algébrique sur $S$.
Soit $Z \subset X$ un sous-espace fermé. La réduction $(X^\wedge_Z)_{red}$
du complété $X^\wedge_Z$ de $X$ le long de $Z$ est $Z_{red}$.
\end{lemma}

\begin{proof}
La démonstration est omise.
\end{proof}

\begin{lemma}
\label{lemma-affine-formal-completion-subspace-types}
Soit $S$ un schéma. Soit $X = \Spec(A)$ un schéma affine sur $S$.
Soit $Z \subset X$ le sous-schéma fermé correspondant à l'idéal
$I \subset A$. Alors
\begin{enumerate}
\item L'espace algébrique formel affine $X^\wedge_Z$ est faiblement adique.
\item Si $I$ est de type fini, alors
$X^\wedge_Z = \text{Spf}(A^\wedge)$, où $A^\wedge$ est le
complété $I$-adique de $A$.
\item Si $Z \to X$ est de présentation finie, c'est-à-dire si $I$ est
de type fini, alors $X^\wedge_Z$ est adique*.
\item Si $X$ est noethérien, alors $X^\wedge_Z$ est noethérien.
\end{enumerate}
\end{lemma}

\begin{proof}
La démonstration est omise.
\end{proof}

\begin{lemma}
\label{lemma-formal-completion-subspace-types}
Soit $S$ un schéma. Soit $X$ un espace algébrique sur $S$.
Soit $Z \subset X$ un sous-espace fermé. Soit $X^\wedge_Z$ le
complété formel de $X$ le long de $Z$.
\begin{enumerate}
\item L'espace algébrique formel $X^\wedge_Z$ est localement faiblement adique.
\item Si $Z \to X$ est de présentation finie,
alors $X^\wedge_Z$ est localement adique*.
\item Si $X$ est localement noethérien, alors $X_Z$ est localement
noethérien.
\end{enumerate}
\end{lemma}

\begin{proof}
La démonstration est omise.
\end{proof}

\begin{lemma}
\label{lemma-when-isomorphism}
Soit $S$ un schéma. Soit $X$ un espace algébrique sur $S$.
Soit $Z \subset X$ un sous-espace fermé et posons $T = |Z|$. Le
morphisme canonique $c : X^\wedge_Z \to X_{/T}$ est un isomorphisme si $Z \to X$
est de présentation finie, mais pas en général.
\end{lemma}

\begin{proof}
Les deux constructions commutent à la localisation \'etale ; il suffit donc
de le démontrer lorsque $X$ est affine. Écrivons $X = \Spec(A)$ et supposons que $Z$ corresponde
à l'idéal $I \subset A$. Si $I$ est de type fini, alors
$X^\wedge_Z$ et $X_{/T}$ sont tous deux égaux à $\text{Spf}(A^\wedge)$, où $A^\wedge$
est le complété $I$-adique de $A$, voir les
lemmes \ref{lemma-affine-formal-completion-types} et
\ref{lemma-affine-formal-completion-subspace-types}.
Si $A = \mathbf{Z}[x_1, x_2, \ldots]$ et $I = (x_1, x_2, \ldots)$,
alors $\Spec(A/(x_n^n, n \geq 1)) \to X$ se factorise par $X_{/T}$,
mais pas par $X^\wedge_Z$ ; par conséquent, $c$ n'est pas un isomorphisme.
\end{proof}






\begin{multicols}{2}[\section{Autres chapitres}]
\noindent
Préliminaires
\begin{enumerate}
\item \hyperref[introduction-section-phantom]{Introduction}
\item \hyperref[conventions-section-phantom]{Conventions}
\item \hyperref[sets-section-phantom]{Théorie des ensembles}
\item \hyperref[categories-section-phantom]{Catégories}
\item \hyperref[topology-section-phantom]{Topologie}
\item \hyperref[sheaves-section-phantom]{Faisceaux sur les espaces}
\item \hyperref[sites-section-phantom]{Sites et faisceaux}
\item \hyperref[stacks-section-phantom]{Champs}
\item \hyperref[fields-section-phantom]{Corps}
\item \hyperref[algebra-section-phantom]{Algèbre commutative}
\item \hyperref[brauer-section-phantom]{Groupes de Brauer}
\item \hyperref[homology-section-phantom]{Algèbre homologique}
\item \hyperref[derived-section-phantom]{Catégories dérivées}
\item \hyperref[simplicial-section-phantom]{Méthodes simpliciales}
\item \hyperref[more-algebra-section-phantom]{Compléments d'algèbre}
\item \hyperref[smoothing-section-phantom]{Lissage des morphismes d'anneaux}
\item \hyperref[modules-section-phantom]{Faisceaux de modules}
\item \hyperref[sites-modules-section-phantom]{Modules sur les sites}
\item \hyperref[injectives-section-phantom]{Objets injectifs}
\item \hyperref[cohomology-section-phantom]{Cohomologie des faisceaux}
\item \hyperref[sites-cohomology-section-phantom]{Cohomologie sur les sites}
\item \hyperref[dga-section-phantom]{Algèbre différentielle graduée}
\item \hyperref[dpa-section-phantom]{Algèbre à puissances divisées}
\item \hyperref[sdga-section-phantom]{Faisceaux différentiels gradués}
\item \hyperref[hypercovering-section-phantom]{Hyperrecouvrements}
\end{enumerate}
Schémas
\begin{enumerate}
\setcounter{enumi}{25}
\item \hyperref[schemes-section-phantom]{Schémas}
\item \hyperref[constructions-section-phantom]{Constructions de schémas}
\item \hyperref[properties-section-phantom]{Propriétés des schémas}
\item \hyperref[morphisms-section-phantom]{Morphismes de schémas}
\item \hyperref[coherent-section-phantom]{Cohomologie des schémas}
\item \hyperref[divisors-section-phantom]{Diviseurs}
\item \hyperref[limits-section-phantom]{Limites de schémas}
\item \hyperref[varieties-section-phantom]{Variétés}
\item \hyperref[topologies-section-phantom]{Topologies sur les schémas}
\item \hyperref[descent-section-phantom]{Descente}
\item \hyperref[perfect-section-phantom]{Catégories dérivées de schémas}
\item \hyperref[more-morphisms-section-phantom]{Compléments sur les morphismes}
\item \hyperref[flat-section-phantom]{Compléments sur la platitude}
\item \hyperref[groupoids-section-phantom]{Schémas en groupoïdes}
\item \hyperref[more-groupoids-section-phantom]{Compléments sur les schémas en groupoïdes}
\item \hyperref[etale-section-phantom]{Morphismes étales de schémas}
\end{enumerate}
Thèmes de théorie des schémas
\begin{enumerate}
\setcounter{enumi}{41}
\item \hyperref[chow-section-phantom]{Homologie de Chow}
\item \hyperref[intersection-section-phantom]{Théorie de l'intersection}
\item \hyperref[pic-section-phantom]{Schémas de Picard des courbes}
\item \hyperref[weil-section-phantom]{Théories cohomologiques de Weil}
\item \hyperref[adequate-section-phantom]{Modules adéquats}
\item \hyperref[dualizing-section-phantom]{Complexes dualisants}
\item \hyperref[duality-section-phantom]{Dualité pour les schémas}
\item \hyperref[discriminant-section-phantom]{Discriminants et différentes}
\item \hyperref[derham-section-phantom]{Cohomologie de de Rham}
\item \hyperref[local-cohomology-section-phantom]{Cohomologie locale}
\item \hyperref[algebraization-section-phantom]{Géométrie algébrique et formelle}
\item \hyperref[curves-section-phantom]{Courbes algébriques}
\item \hyperref[resolve-section-phantom]{Résolution des surfaces}
\item \hyperref[models-section-phantom]{Réduction semi-stable}
\item \hyperref[functors-section-phantom]{Foncteurs et morphismes}
\item \hyperref[equiv-section-phantom]{Catégories dérivées de variétés}
\item \hyperref[pione-section-phantom]{Groupes fondamentaux de schémas}
\item \hyperref[etale-cohomology-section-phantom]{Cohomologie étale}
\item \hyperref[crystalline-section-phantom]{Cohomologie cristalline}
\item \hyperref[proetale-section-phantom]{Cohomologie pro-étale}
\item \hyperref[relative-cycles-section-phantom]{Cycles relatifs}
\item \hyperref[more-etale-section-phantom]{Compléments de cohomologie étale}
\item \hyperref[trace-section-phantom]{La formule des traces}
\end{enumerate}
Espaces algébriques
\begin{enumerate}
\setcounter{enumi}{64}
\item \hyperref[spaces-section-phantom]{Espaces algébriques}
\item \hyperref[spaces-properties-section-phantom]{Propriétés des espaces algébriques}
\item \hyperref[spaces-morphisms-section-phantom]{Morphismes d'espaces algébriques}
\item \hyperref[decent-spaces-section-phantom]{Espaces algébriques décents}
\item \hyperref[spaces-cohomology-section-phantom]{Cohomologie des espaces algébriques}
\item \hyperref[spaces-limits-section-phantom]{Limites d'espaces algébriques}
\item \hyperref[spaces-divisors-section-phantom]{Diviseurs sur les espaces algébriques}
\item \hyperref[spaces-over-fields-section-phantom]{Espaces algébriques sur des corps}
\item \hyperref[spaces-topologies-section-phantom]{Topologies sur les espaces algébriques}
\item \hyperref[spaces-descent-section-phantom]{Descente et espaces algébriques}
\item \hyperref[spaces-perfect-section-phantom]{Catégories dérivées d'espaces}
\item \hyperref[spaces-more-morphisms-section-phantom]{Compléments sur les morphismes d'espaces}
\item \hyperref[spaces-flat-section-phantom]{Platitude sur les espaces algébriques}
\item \hyperref[spaces-groupoids-section-phantom]{Groupoïdes dans les espaces algébriques}
\item \hyperref[spaces-more-groupoids-section-phantom]{Compléments sur les groupoïdes dans les espaces}
\item \hyperref[bootstrap-section-phantom]{Amorçage}
\item \hyperref[spaces-pushouts-section-phantom]{Sommes amalgamées d'espaces algébriques}
\end{enumerate}
Thèmes de géométrie
\begin{enumerate}
\setcounter{enumi}{81}
\item \hyperref[spaces-chow-section-phantom]{Groupes de Chow des espaces}
\item \hyperref[groupoids-quotients-section-phantom]{Quotients de groupoïdes}
\item \hyperref[spaces-more-cohomology-section-phantom]{Compléments sur la cohomologie des espaces}
\item \hyperref[spaces-simplicial-section-phantom]{Espaces simpliciaux}
\item \hyperref[spaces-duality-section-phantom]{Dualité pour les espaces}
\item \hyperref[formal-spaces-section-phantom]{Espaces algébriques formels}
\item \hyperref[restricted-section-phantom]{Algébrisation des espaces formels}
\item \hyperref[spaces-resolve-section-phantom]{Résolution des surfaces revisitée}
\end{enumerate}
Théorie des déformations
\begin{enumerate}
\setcounter{enumi}{89}
\item \hyperref[formal-defos-section-phantom]{Théorie formelle des déformations}
\item \hyperref[defos-section-phantom]{Théorie des déformations}
\item \hyperref[cotangent-section-phantom]{Le complexe cotangent}
\item \hyperref[examples-defos-section-phantom]{Problèmes de déformation}
\end{enumerate}
Champs algébriques
\begin{enumerate}
\setcounter{enumi}{93}
\item \hyperref[algebraic-section-phantom]{Champs algébriques}
\item \hyperref[examples-stacks-section-phantom]{Exemples de champs}
\item \hyperref[stacks-sheaves-section-phantom]{Faisceaux sur les champs algébriques}
\item \hyperref[criteria-section-phantom]{Critères de représentabilité}
\item \hyperref[artin-section-phantom]{Axiomes d'Artin}
\item \hyperref[quot-section-phantom]{Espaces de Quot et de Hilbert}
\item \hyperref[stacks-properties-section-phantom]{Propriétés des champs algébriques}
\item \hyperref[stacks-morphisms-section-phantom]{Morphismes de champs algébriques}
\item \hyperref[stacks-limits-section-phantom]{Limites de champs algébriques}
\item \hyperref[stacks-cohomology-section-phantom]{Cohomologie des champs algébriques}
\item \hyperref[stacks-perfect-section-phantom]{Catégories dérivées de champs}
\item \hyperref[stacks-introduction-section-phantom]{Introduction aux champs algébriques}
\item \hyperref[stacks-more-morphisms-section-phantom]{Compléments sur les morphismes de champs}
\item \hyperref[stacks-geometry-section-phantom]{La géométrie des champs}
\end{enumerate}
Thèmes de théorie des espaces de modules
\begin{enumerate}
\setcounter{enumi}{107}
\item \hyperref[moduli-section-phantom]{Champs de modules}
\item \hyperref[moduli-curves-section-phantom]{Espaces de modules de courbes}
\end{enumerate}
Divers
\begin{enumerate}
\setcounter{enumi}{109}
\item \hyperref[examples-section-phantom]{Exemples}
\item \hyperref[exercises-section-phantom]{Exercices}
\item \hyperref[guide-section-phantom]{Guide bibliographique}
\item \hyperref[desirables-section-phantom]{Desiderata}
\item \hyperref[coding-section-phantom]{Style de programmation}
\item \hyperref[obsolete-section-phantom]{Obsolète}
\item \hyperref[fdl-section-phantom]{Licence de documentation libre GNU}
\item \hyperref[index-section-phantom]{Index généré automatiquement}
\end{enumerate}
\end{multicols}


\bibliography{my}
\bibliographystyle{amsalpha}

\end{document}

